% !TEX program = xelatex
\documentclass[12pt,a4paper,fontset=fandol]{ctexbook}

% ==================== 页面 ====================
\usepackage[top=2.5cm,bottom=2.5cm,left=2.5cm,right=2.5cm]{geometry}

% ==================== 字体 ====================
\usepackage{fontspec}
\usepackage{iffont}
% CJK 钉 fandol（两侧都取自 TeX 树里的同一份 FandolSong-Regular.otf）；
% mono 钉 DejaVu Sans Mono：CI 由 fontconfig 命中同名，本机没有该系统字体，
% 就用 MiKTeX dejavu 包里的同一版本 TTF（kpathsea 按文件名命中）。
\IfFontExistsTF{DejaVu Sans Mono}
  {\setmonofont{DejaVu Sans Mono}[Scale=MatchLowercase]}
  {\setmonofont[Extension=.ttf,
      UprightFont=DejaVuSansMono,
      BoldFont=DejaVuSansMono-Bold,
      ItalicFont=DejaVuSansMono-Oblique,
      BoldItalicFont=DejaVuSansMono-BoldOblique]
     {DejaVuSansMono}[Scale=MatchLowercase]}

% ==================== 颜色 ====================
\usepackage[dvipsnames,svgnames,x11names]{xcolor}
\definecolor{codebg}{HTML}{F5F5F5}
\definecolor{codeframe}{HTML}{E0E0E0}
\definecolor{cppkeyword}{HTML}{1A1AFF}
\definecolor{cppcomment}{HTML}{2E8B2E}
\definecolor{cppstring}{HTML}{B22222}
\definecolor{linenumgray}{HTML}{999999}
\definecolor{ubheadbg}{HTML}{5B1F26}
\definecolor{okgreen}{HTML}{1A7F37}
\definecolor{errred}{HTML}{B62324}
\definecolor{stdpurple}{HTML}{5B3E9D}

% ==================== listings ====================
\usepackage{listings}

% ---- 所有代码块共用的排版（先定义，各语言 style 组合它）----
\lstdefinestyle{codebase}{
  basicstyle=\small\ttfamily,
  keywordstyle=\color{cppkeyword}\bfseries,
  commentstyle=\color{cppcomment}\itshape,
  stringstyle=\color{cppstring},
  numberstyle=\tiny\color{linenumgray},
  numbers=left,
  frame=single,
  rulecolor=\color{codeframe},
  framesep=6pt,
  breaklines=true,
  tabsize=4,
  columns=flexible,
  keepspaces=true,
  backgroundcolor=\color{codebg},
  showstringspaces=false,
  escapeinside={(*@}{@*)},
}

\lstdefinestyle{cppstyle}{
  style=codebase,
  language=C++,
  morekeywords={constexpr,consteval,constinit,decltype,auto,
    concept,requires,co_await,co_yield,co_return,
    noexcept,override,final,nullptr,static_assert,
    template,typename,namespace,using,mutable,volatile,
    explicit,operator,friend,inline,virtual,
    thread_local,alignas,alignof},
  deletekeywords={int,char,bool,float,double,long,short,unsigned,signed,void,wchar_t},
  emph={size_t,int32_t,uint32_t,int64_t,uint64_t,ssize_t,ptrdiff_t,
    string,vector,map,set,unique_ptr,shared_ptr,optional,variant,any,
    span,string_view,atomic,memory_order,bit_cast,start_lifetime_as},
  emphstyle=\color{teal!80!black},
}
\lstset{style=cppstyle}

% ---- C： libc、属性与内建函数 ----
\lstdefinestyle{cstyle}{
  style=codebase,
  language=C,
  morekeywords={bool,FILE,restrict,inline,thread_local},
  emph={size_t,int32_t,uint32_t,int64_t,uint64_t,ssize_t,ptrdiff_t,
    uintptr_t,intptr_t,char16_t,wchar_t,va_list},
  emphstyle=\color{teal!80!black},
}

% ---- 本册自建 shell 方言：编译器、消毒器与取证命令行 ----
% 注意：邻册的 ShellPM / ShellLD 只在各自文件内定义，不能借用。
\lstdefinelanguage{ShellUB}{
  morekeywords={if,then,else,elif,fi,for,in,do,done,while,case,esac,
    function,return,local,export,source,echo,cd,set,exec,trap,read,
    printf,test,exit,shift,eval,alias},
  sensitive=true,
  morecomment=[l]{\#},
  morestring=[b]",
  morestring=[d]',
}
\lstdefinestyle{shellstyle}{
  style=codebase,
  language=ShellUB,
  keywordstyle=\color{cppkeyword}\bfseries,
  emph={gcc,g++,clang,clang++,cl,link,cmake,ninja,make,
    grep,sed,awk,xargs,jq,python3,pip,pwsh,bash,git,tar,mkdir,touch,
    ls,cat,wc,diff,find,sort,uniq,tee,rm,cp,mv,objdump,nm,readelf,
    strings,valgrind,gdb,lldb,llvm, opt, llc, taskset, ulimit, timeout,
    ccache,sanitizers,ubsan,asan,lsan,tsan,msan},
  emphstyle=\color{teal!80!black}\bfseries,
}

% ---- 本册自建 CMake 方言：门禁拒收裸 language=CMake ----
\lstdefinelanguage{CMakeUB}{
  morekeywords={cmake_minimum_required,project,add_executable,
    add_library,target_link_libraries,target_include_directories,
    target_compile_definitions,target_compile_options,set,if,elseif,
    else,endif,find_package,option,message,install,file,
    target_sources,add_subdirectory,enable_language,set_target_properties,
    check_cxx_compiler_flag,list,string},
  sensitive=false,
  morecomment=[l]{\#},
  morestring=[b]",
}
\lstdefinestyle{cmakestyle}{
  style=codebase,
  language=CMakeUB,
  keywordstyle=\color{cppkeyword}\bfseries,
}

% ---- 汇编与机器面：AT&T 与 ARM64 的读法（本册自建方言）----
% listings 的关键词默认只含字母，.text 这类伪指令必须以 alsoletter={.} 放行。
\lstdefinelanguage{UBAsm}{
  alsoletter={.},
  morekeywords={.text,.data,.bss,.section,.globl,.global,.local,.type,
    .size,.align,.p2align,.byte,.word,.long,.quad,.string,.asciz,.file,
    .ident,.weak,.hidden,.set,.equ,.macro,.endm,.if,.else,.endif,
    .cfi_startproc,.cfi_endproc,.cfi_def_cfa_offset,.save,.rept,.endr,
    mov,movq,movl,movzbl,movsx,lea,add,sub,imul,idiv,and,or,xor,not,
    shl,shr,sar,cmp,test,jmp,je,jne,jl,jg,jle,jge,ja,jb,call,ret,push,
    pop,setne,setl,cmovne,cs,rep,bswap,cqto,cltd,endbr64,
    ldp,stp,bl,br,cbz,cbnz,adrp,adr,mul,sdiv,ldr,str,
    sxtw,uxtb,rev,rev32,clz},
  sensitive=true,
  morecomment=[l]{\#},
  morecomment=[n]{/*}{*/},
  morestring=[b]",
}
\lstdefinestyle{asmstyle}{
  style=codebase,
  language=UBAsm,
  keywordstyle=\color{stdpurple}\bfseries,
}

% ---- 工具输出：编译器诊断、消毒器报告、反汇编清单（无语言，纯文本）----
\lstdefinestyle{outstyle}{
  style=codebase,
  language={},
  numbers=none,
  frame=single,
  morecomment=[l][\color{stdpurple}\bfseries]{\#},
  emph={runtime,error,warning,UBSAN,SignedIntegerOverflow,
    UndefinedBehavior,LeakSanitizer,ThreadSanitizer,data_race,
    misaligned,unaligned,overflow,trap,poison,volatile},
  emphstyle=\color{errred}\bfseries,
}

% ==================== tcolorbox ====================
\usepackage[most]{tcolorbox}

\newtcolorbox{guideline}[1][]{
  enhanced,breakable,
  colback=blue!5!white,colframe=blue!75!black,
  fonttitle=\bfseries,title=要点,#1,
  boxed title style={colback=blue!75!black},
  attach boxed title to top left={yshift=-2mm,xshift=2mm},
  arc=2mm,boxrule=0.8mm,
}
\newtcolorbox{compare}[1][]{
  enhanced,breakable,
  colback=green!3!white,colframe=green!50!black,
  fonttitle=\bfseries,title=对比,#1,
  boxed title style={colback=green!50!black},
  attach boxed title to top left={yshift=-2mm,xshift=2mm},
  arc=2mm,boxrule=0.8mm,
}
\newtcolorbox{warning}[1][]{
  enhanced,breakable,
  colback=red!5!white,colframe=red!75!black,
  fonttitle=\bfseries,title=常见陷阱,#1,
  boxed title style={colback=red!75!black},
  attach boxed title to top left={yshift=-2mm,xshift=2mm},
  arc=2mm,boxrule=0.8mm,
}
\newtcolorbox{note}[1][]{
  enhanced,breakable,
  colback=orange!5!white,colframe=orange!80!black,
  fonttitle=\bfseries,title=注意,#1,
  boxed title style={colback=orange!80!black},
  attach boxed title to top left={yshift=-2mm,xshift=2mm},
  arc=2mm,boxrule=0.8mm,
}
\newtcolorbox{bestpractice}[1][]{
  enhanced,breakable,
  colback=purple!5!white,colframe=purple!60!black,
  fonttitle=\bfseries,title=最佳实践,#1,
  boxed title style={colback=purple!60!black},
  attach boxed title to top left={yshift=-2mm,xshift=2mm},
  arc=2mm,boxrule=0.8mm,
}
% 本册专用：标准原文条款框与可复现案例框
\newtcolorbox{clause}[1][]{
  enhanced,breakable,
  colback=stdpurple!4!white,colframe=stdpurple,
  fonttitle=\bfseries,title=标准怎么说,#1,
  boxed title style={colback=stdpurple},
  attach boxed title to top left={yshift=-2mm,xshift=2mm},
  arc=2mm,boxrule=0.8mm,
}
\newtcolorbox{casebox}[1][]{
  enhanced,breakable,
  colback=errred!4!white,colframe=errred,
  fonttitle=\bfseries,title=可复现案例,#1,
  boxed title style={colback=errred},
  attach boxed title to top left={yshift=-2mm,xshift=2mm},
  arc=2mm,boxrule=0.8mm,
}

% ==================== 章节标题 ====================
\usepackage{titlesec}
\usepackage{fancyhdr}
\usepackage{graphicx}
\usepackage{amssymb}
\usepackage{amsmath}
\usepackage{tabularx}
\usepackage{booktabs}
\usepackage{longtable}
\usepackage{array}
\usepackage{multirow}
\usepackage{enumitem}
\usepackage{seqsplit}
\usepackage{float}

\titleformat{\chapter}[display]
  {\normalfont\huge\bfseries\color{ubheadbg}}
  {\CTEXthechapter}{20pt}{\Huge}
\titleformat{\section}
  {\normalfont\Large\bfseries\color{blue!70!black}}
  {\thesection}{1em}{}
\titleformat{\subsection}
  {\normalfont\large\bfseries\color{blue!60!black}}
  {\thesubsection}{1em}{}

\pagestyle{fancy}
\fancyhf{}
\fancyhead[LE,RO]{\thepage}
\fancyhead[RE]{\leftmark}
\fancyhead[LO]{\rightmark}
\setlength{\headheight}{14.5pt}

\setlist[itemize]{leftmargin=2em,itemsep=2pt}
\setlist[enumerate]{leftmargin=2em,itemsep=2pt}
\emergencystretch=6em

% ==================== 超链接 ====================
\usepackage{hyperref}
\hypersetup{
  colorlinks=true,
  linkcolor=blue!70!black,
  urlcolor=blue!60!black,
  bookmarks=true,
  pdfauthor={C++ Reference},
  pdftitle={C/C++ 未定义行为与标准一致性参考手册：UB、实现定义、未指定与前向兼容},
}

% ==================== 自定义命令 ====================
\newcommand{\cpp}[1]{C++#1}
\newcommand{\Fn}[1]{\texttt{#1()}}
\newcommand{\Tp}[1]{\texttt{#1}}
\newcommand{\Opt}[1]{\texttt{-\/-#1}}
\newcommand{\Bk}{\discretionary{}{}{}}
\newcommand{\Lang}[1]{\textsf{#1}}
\newcommand{\CL}{\Lang{C}}
\newcommand{\CPP}{\Lang{C++}}
% 五种行为类别：本册全篇统一用这五个宏，避免同义多形
\newcommand{\UB}{\textbf{\textsf{UB}}}
\newcommand{\IDB}{\textbf{\textsf{ID}}}
\newcommand{\USB}{\textbf{\textsf{US}}}
\newcommand{\NDR}{\textbf{\textsf{NDR}}}
\newcommand{\ILLF}{\textbf{\textsf{IF}}}
% 标准引用与实现物件
\newcommand{\Std}[1]{\textsc{[#1]}}
\newcommand{\Clause}[1]{\textsc{\small #1}}
\newcommand{\Hdr}[1]{\texttt{<#1>}}
\newcommand{\Feat}[1]{\texttt{#1}}
\newcommand{\Warn}[1]{\texttt{-\/-#1}}
\newcommand{\Mx}[1]{\textsf{#1}}

% ====================================================================
% book 双页模式默认 \flushbottom，会靠拉伸行距把每页撑满；本仓库框体与长表多，
% 撑页产生的 Underfull \vbox 是排版噪声而非缺陷，这里统一改用小底部对齐。
\raggedbottom

\begin{document}

\frontmatter

\begin{titlepage}
\centering
\vspace*{4cm}
{\Huge\bfseries C/C++ \\ 未定义行为与标准一致性参考手册\par}
\vspace{1cm}
{\LARGE 未定义、实现定义、未指定与前向兼容\par}
\vspace{0.4cm}
{\Large 对象生命周期、指针与别名、求值顺序、\\ 并发内存模型与标准库约定\par}
\vfill
{\large \today\par}
\end{titlepage}

\tableofcontents

\mainmatter

\cleardoublepage

% ====================================================================
% 骨架冒烟章：把本册全部 listing 方言、框体与自定义宏各用一次，
% 保证 preamble 里的每一项都在编译路径上。正式章节合并时整章替换掉。
% ====================================================================% 章节正文由 frag_chNN.tex 合并插入到下面这行之前（merge_frags.py）
\part{契约与分类}
\chapter{五种行为类别与 as-if 规则}\label{ch:ub01-taxonomy}

这一章只做一件事：把标准用来描述程序含义的那几个词分清，并给全册定下读法。
后面十四章谈到的每一类事故，几乎都能先落到本章的某个格子里；落不进格子的，
多半是术语混用而不是代码疑难。

\section{标准对程序作出的三层承诺}\label{sec:ub01-three-layers}

标准并不是一口气承诺“程序会按你写的样子执行”。它分三层给出承诺，
每一层的违约后果由不同的人承担，这也是五种行为类别的由来。

\subsection{第一层：语法、链接与合法性}\label{sec:ub01-layer-grammar}

最外层是可编译性：词法与语法要合法，声明要能被找到，链接期要有唯一定义，
同一程序里的同名实体不能给出互相矛盾的定义。这一层的承诺是硬的，
因为它在程序执行之前就能判定。违反它的后果是程序不合法，标准用
\ILLF{} 标记；绝大多数情况下实现会给出诊断，少数情况下标准明确写了
“不需要诊断”（\NDR{}）。

这一层有两个容易被忽略的细节。第一，链接期的失败不是运行期的事故：
调用一个只有声明没有定义的函数，程序根本不会产生可执行文件，
它连“执行”都没有，谈不上任何行为类别。第二，\NDR{} 的常见来源正是跨
翻译单元的重复定义与一单元内的两个不同 \Tp{constexpr} 定义一类问题：
判定它需要实现付出与优化相当的代价，所以标准放弃了诊断义务。

\subsection{第二层：语义要求与抽象机}\label{sec:ub01-layer-semantics}

第二层是执行语义：标准先描述一台抽象机（\Std{intro.execution}
\Clause{执行阶段}），再规定什么样的执行序列算作符合要求的实现。
这一层的承诺是有条件的：它只对“不含未定义行为的执行”给出要求。
换句话说，抽象机不是对每一行源码的承诺，而是对一整条执行轨迹的承诺；
轨迹里只要有一段落在 \UB{} 上，标准对这段之后的所有事情都不再负责。
本章第六节会把这句话讲透，因为它是全册最重要的一句话。

\subsection{第三层：实现提供的行为}\label{sec:ub01-layer-impl}

第三层是标准故意不定的部分：它列出一组候选答案，把选择权交给实现。
这一层里有两个词必须区分干净——\IDB{} 要求实现作出选择并且写进文档，
\USB{} 只要求实现作出选择，不要求它把选择告诉任何人。此外，翻译上限
（单个文件能容纳的宏数量、嵌套深度一类）、语言环境的行为、由实现定义的
输入输出细节，也都落在这一层。标准还区分 \Clause{hosted} 与
\Clause{freestanding} 两类一致实现：前者必须提供完整标准库，后者只承诺
一个很小的头文件集合，同一份代码在两类实现里的合法边界并不相同。
标准为什么要把这一层故意空出来，第~\ref{ch:ub02-inevitable} 章给的答案与
\CL{}/\CPP{} 的设计无关：那是“标准”这一文体自身的性质
（\S\ref{sec:ub02-standard-blanks}）。

\section{五类行为与五个记号}\label{sec:ub01-five-kinds}

全册统一使用五个记号：良定义行为（不特别标记，正文里直接说“良好定义”）、
\UB{}、\IDB{}、\USB{}、\ILLF{}，再加一个作为 \ILLF{} 限定词使用的 \NDR{}。
这张表是全册的坐标系，后面每一章的每个案例都要能在表里指出自己那一格。

\begin{longtable}{@{}>{\raggedright\arraybackslash}p{0.11\textwidth} >{\raggedright\arraybackslash}p{0.20\textwidth} >{\raggedright\arraybackslash}p{0.13\textwidth} >{\raggedright\arraybackslash}p{0.17\textwidth} >{\raggedright\arraybackslash}p{0.15\textwidth}@{}}
\caption{五种行为类别：谁承担后果、有无文档义务、是否要求诊断}\label{tab:ub01-kinds}\\
\toprule
类别 & 标准给出的要求 & 后果承担方 & 实现的文档义务 & 诊断要求 \\
\midrule
\endfirsthead
\toprule
类别 & 标准给出的要求 & 后果承担方 & 实现的文档义务 & 诊断要求 \\
\midrule
\endhead
\bottomrule
\endfoot
良好定义 & 给出结果，或给出不互相矛盾的要求 & 无人违约 & 无 & 不需要 \\
\UB{} & 对这段执行序列不作任何要求 & 程序员 & 无 & 不要求 \\
\IDB{} & 候选集合固定，实现必须选定一个 & 实现（选定后） & 有，必须文档化 & 不要求 \\
\USB{} & 候选集合内任选，怎么选都合规 & 程序员 & 无强制 & 不要求 \\
\ILLF{} & 程序不合法，实现没有义务接受 & 程序员 & 无 & 通常要求 \\
\NDR{} & 程序不合法，但明确免除诊断 & 程序员 & 无 & 明确不要求 \\
\end{longtable}

\subsection{UB：整条执行失去要求}\label{sec:ub01-ub}

\UB{} 的定义式说法是：程序使用某语言构造时，后果在此标准中不作要求。
关键在于“不作要求”指向的是\emph{执行}，不是源码。同一行代码在一条执行
轨迹里越界，在另一条里没有越界，只有前者失去保证。这就是为什么
“这段代码有问题”这种说法本身就不准确：有问题的是某个输入下的某条路径。

\subsection{ID：有文档义务的选项}\label{sec:ub01-idb}

\IDB{} 是标准里唯一把义务压在实现身上的类别。典型条目是 \Tp{char} 的
符号性、\Tp{wchar\_t} 的宽度与编码、移位超过宽度时的结果、
\Tp{size\_t} 的位宽。这类事情的正确处理姿势是查实现文档，而不是猜：
在文档化的范围内，代码是\emph{可移植到该实现}的，只是不可移植到全部实现。
第~\ref{ch:ub13-implementation} 章专门处理这一类。

\subsection{US：没有文档义务的选项}\label{sec:ub01-usb}

\USB{} 与 \IDB{} 的唯一区别是文档义务。函数参数的求值顺序、
\Tp{std::numeric\_limits} 里某些由环境决定的成员取哪一侧，都是 \USB{}：
实现可以换选择而不通知任何人。因此依赖 \USB{} 的代码比依赖 \IDB{} 的代码
更脆——后者至少有一份可以援引的文档。第~\ref{ch:ub08-evaluation} 章与第~\ref{ch:ub14-conventions} 章分别处理求值顺序
与探测手段。

\subsection{IF 与 NDR：合法性问题不要混进行为问题}\label{sec:ub01-if-ndr}

\ILLF{} 是程序层面而不是执行层面的判定：它意味着实现没有义务让你的
程序跑起来。大多数 \ILLF{} 会立刻变成编译器诊断，这是最容易处理的一类。
\NDR{} 则是诊断义务被明确免除的 \ILLF{}，它的危险不在后果恶劣，
而在“编译器不说话”这一事实被误读为“编译器认可”。两者的详细判据见
附录~\ref{ap:ubA-kinds}。

\begin{clause}[title={标准怎么说：五个词的分工}]
标准的写法是：先用术语条目给每个类别下定义（\Std{defns.undefined}
\Clause{未定义行为}、\Std{defns.unspecified} \Clause{未指定行为} 一类），
再在正文的每个具体规则处声明它属于哪一类：写 “is undefined” 的是 \UB{}，
写 “is implementation-defined” 的是 \IDB{}，写 “is unspecified” 的是
\USB{}，写 “The program is ill-formed” 的是 \ILLF{}，后面追加
“no diagnostic required” 的就是 \NDR{}。术语条目位于标准前部的定义章，
而行为类别的使用规则集中在符合性一章（\Std{intro.implementation}
\Clause{实现定义行为}）。本册只给小节标签，不给编号：编号随标准版本
漂移，请以你手上那一版草案或正式文本为准。
\end{clause}

\section{术语条目应该这样查}\label{sec:ub01-defns}

\Std{defns.parameter.pointable} 一类的条目说明了术语章的用法：
它规定一个词在这部标准里\emph{只}有这一个含义。查阅顺序有三个建议。
第一，先看术语条目，确定类别词的准确含义；第二，回到具体规则的正文，
看它用的是哪一个措辞；第三，只有在确认这是 \IDB{} 之后，才去查实现文档，
而不是反过来。许多“标准到底怎么说”的争论，其实是把第二、三步的顺序弄反了。

\CL 的组织方式略有不同：它的术语集中在前面的术语章（例如行为、
未定义行为、未指定行为、实现定义行为各占一节），并在资料性附录 Annex J
里逐条汇总所有未定义与未指定行为。想通读一类风险，Annex J 是比正文
更快的入口。\CPP 没有等价的全量汇总表，只能按小节检索。

\subsection{条件支持：C 里的第六种说法}\label{sec:ub01-cond-supported}

\CL 还用了 \Clause{conditionally supported} 这个词：某些构造标准允许实现
完全不支持，但一旦支持就必须按规定的语义支持。它既不是 \UB{} 也不是
\IDB{}，而是“可选实现特性”。跨编译器写扩展语法时，最容易踩到的就是它——
踩到时你会同时看到诊断和“其实也能编过”的矛盾现象，因为不同实现
对同一构造的支持度不同。

\section{as-if 规则：可观测行为的边界}\label{sec:ub01-asif}

as-if 规则是编译器权力来源的正式表述。它的内容可以压缩成一句话：
只要最终的可观测行为与抽象机的某次执行一致，实现可以对程序作任何改写，
包括删掉你写的代码。因此“可观测行为”的外延就是优化自由度的外延。

\subsection{可观测行为包含什么}\label{sec:ub01-observable}

标准的口径大致包含四件事。第一，程序正常终止时写入最终文件状态的
输出序列；第二，对 \Tp{volatile} 对象的读写，其次数、顺序与取值都有要求；
第三，由实现定义的输入与输出侧效应；第四，在并发程序里，
同步操作与原子操作的相对顺序也属于可观测面（\Std{intro.multithread}
\Clause{多线程执行与同步}，\Std{intro.races} \Clause{数据竞争}）。
异常抛出与捕获在 \CPP 里同样被算进可观测行为，这限制了库函数调用的删除。
除此之外的一切——中间变量存不存、循环展开几次、检查留不留——
标准都不关心。

\begin{lstlisting}[style=asmstyle, caption={示意：边界检查留在指令流里，越界走定义良好的中止路径（说明性写法，非本册实测反汇编）}]
    .text
    .globl  at
at:
    movl    (%rdi), %eax        /* 先取长度：边界信息在对象里，检查有依据 */
    cmpl    %esi, %eax
    jle     trap                /* 越界：进入定义良好的中止路径 */
    movl    4(%rdi,%rsi,4), %eax
    ret
trap:
    call    abort
    ret
\end{lstlisting}

\subsection{推论：你认为必要的检查可以被删掉}\label{sec:ub01-elide}

从 as-if 规则直接推出本章最实用的一条结论：如果某个检查只会在
未定义行为的执行中触发，那么删掉它不违反任何承诺，优化器就会删掉它。
两个经典形状如下。

第一个是判空之后赋值。如果代码在前面已经无条件解引用过同一个指针，
那么“指针为空”这条路只可能属于 \UB{} 的执行，于是后判的空检查可被推断为
恒真而被消除。第二个是有符号溢出参与循环条件。因为标准不承诺回绕，
优化器可以假定溢出不会发生，进而把“溢出后变小所以退出循环”这类
写法判成恒假，把循环改成无限循环或者整体删除。

这两条推论都不是某个编译器的怪癖，而是规则本身的后果。另一册专门讲编译器
如何把这些假设进一步用于优化，本册只负责确认：假设的存在是合法的，
写代码的人不能依赖它不会兑现。

\section{UB 是承诺的撤回，而不是错误}\label{sec:ub01-retraction}

“错误”这个词暗含三件事：有可判定的状态、有可捕获的机制、有可依赖的记忆。
\UB{} 三件都不给。它不是程序的一种状态，而是标准对该次执行撤回了全部要求。
因此“这段代码有 UB，所以会崩”也是不准确的推论：崩是一种可能的结局，
不崩同样是一种可能的结局，而“这次没崩”不包含任何保证。

\begin{lstlisting}[style=cppstyle, language=C++, caption={两段各含一处未定义行为的代码，用来观察档位差异}]
constexpr int scale(int v) { return 2 * v + 1; }   // 有符号溢出：不承诺回绕
int put_at(int *p, int v) {
    if (p) { *p = v; }          // 判空：可能因为先前已解引用而被消除
    return v;
}
int worst(int v) {
    if (scale(v) < v) { return -1; }   // 依赖回绕的分支
    return 0;
}
\end{lstlisting}

\subsection{空指针解引用的两种结局}\label{sec:ub01-null}

在没有优化的档位上，空指针解引用通常在硬件层得到陷阱（缺页或访问违例），
进程被终止，看起来“像是一次可诊断的错误”。但请注意这件事的机制：
终止你的不是标准，也不是编译器，而是操作系统对某段地址的保护策略。
在另一个环境里结局会不同——\Clause{freestanding} 实现、内核态代码或者
把第 0 页映射为可读写的调试环境里，写入 \Tp{nullptr} 可以安静地成功。
而在高优化档位上，判空检查本身可能已经按上一节的推论消失，
于是这段代码从“检查再写”变成“直接写”，程序连崩都不会崩在你想的地方。

\begin{casebox}[title={同一份代码的两种命运：空指针与有符号溢出}]
上一节的 \Fn{put\_at} 与 \Fn{worst} 在两个档位下的定性差异：

\begin{itemize}
  \item 低优化档位：\Fn{put\_at} 的判空真实存在，空指针进入即被跳过或
        在别处崩；\Fn{worst} 走的是回绕后的分支，返回一个稳定但错误的值。
        现象可复现，容易误判为“已验证”。
  \item 高优化档位：判空可能在指令层面消失；\Fn{worst} 里的比较被判定为
        永不成立，整个分支连同可能的返回路径一起被折叠，函数只剩一条出口。
        此时“错误值”都不再保证，甚至可能变成不返回。
\end{itemize}

两种结局的差别不是编译器坏了，而是标准从未承诺过这两个表达式的含义：
档位只是选择了撤回承诺后的哪一种可能。
\end{casebox}

\subsection{有符号溢出的两种结局}\label{sec:ub01-overflow}

有符号整数溢出是 \UB{} 里流传最广、也最容易被误传的一条。事实的骨架是：
\Std{expr} \Clause{表达式}一类的规则明确说，运算结果不能在结果类型的
可表示范围内时行为未定义；而 \Tp{unsigned} 的回绕是\emph{良好定义}的，
它不属于 \UB{}，也不属于 \USB{}。因此“无符号减到负数是 UB”是常见的
错误说法，正确的方向恰好相反。

\Tp{-fwrapv} 一类的开关可以改变这件事：它让有符号溢出按回绕定义，
代价是取消一批以“不会溢出”为前提的优化。这是典型的用优化空间换
行为确定性的交易，属于工具问题而不是标准问题，第~\ref{ch:ub03-sanitizers} 章与附录~\ref{ap:ubB-switches} 会
把开关面集中整理。

\begin{lstlisting}[style=shellstyle, caption={命令行样例：在两个档位上观察同一份源码的诊断面}]
# 同一份源码，分别用消毒档与优化档编译
g++ -std=c++23 -O0 -fsanitize=undefined t.cpp -o t0
g++ -std=c++23 -O2 -fsanitize=undefined t.cpp -o t2
clang++ -std=c++23 -O2 -fno-sanitize-recover=undefined t.cpp -o t2c
# 需要保留回绕语义时改用有副作用的开关，并确认优化空间已经变小
g++ -std=c++23 -O2 -fwrapv t.cpp -o t2w
\end{lstlisting}

\begin{lstlisting}[style=outstyle, caption={工具输出样例：消毒器报告形状示意，非本仓库实测值}]
# 报告形状示意：具体措辞、是否带表达式原文、是否继续执行都以实现文档为准
runtime error: signed integer overflow: 1073741823 + 1 cannot be
               represented in type 'int'
\end{lstlisting}

\begin{warning}[title={把 UB 当可捕获的错误是范畴错误}]
最常见的三个误解都要在这一章清掉：

\begin{enumerate}
  \item \CPP{} 的 \Tp{try}/\Tp{catch} 抓不到 \UB{}。异常是良好定义的
        控制转移机制，它只能捕获程序仍然遵守标准的那部分执行里
        抛出的东西。越界写坏栈之后再指望某个 \Tp{catch(...)} 兜住，
        是把两件事拼成了幻觉。
  \item “加了检查就安全”只在检查本身仍在指令流里时成立。优化档位可以
        让检查消失，此时你依赖的是编译器的当前决策而不是标准。
  \item “以前能跑”不是证据。同一份代码在相同输入上重复运行给出相同结果，
        只说明这一条执行轨迹的偶然形状，不说明任何承诺的存在。
\end{enumerate}

操作系统层的信号或结构化异常处理确实能拦住某些内存事故，但那是平台
行为而不是语言行为：它只在“事故恰好表现为一次可 traps 的访问”时有效，
而且不同平台的等价性以平台文档为准。
\end{warning}

\section{给每一类行为配一套手段}\label{sec:ub01-means}

分类的价值在于手段不再混用。五类行为各有各的检查面，把它们对调是
工程上最常见的浪费：用断言去管 \IDB{} 毫无意义，用文档去管 \UB{} 也
不会让越界变得合法。

\begin{compare}[title={类别与可用检查面的对应}]
\begin{itemize}
  \item \UB{}：有动态工具可查（消毒器与内存工具，见
        第~\ref{ch:ub03-sanitizers} 章）；断言只能覆盖你预先想到的谓词。
  \item \IDB{}：只能查文档与探测宏；任何运行期检查都在实现选定之后，
        检查不出“另一个实现会怎么选”。
  \item \USB{}：文档也不保证有；能用特性宏或运行期探测确认的就探测，
        不能确认的就把它当随机数处理。
  \item \ILLF{}：编译器诊断覆盖绝大部分，剩下的用严格警告档在 CI 里挡。
  \item \NDR{}：没有诊断可以依赖，只能靠类型系统与编译期门禁把不合法
        重新变成不合法。\Tp{static\_assert} 一类的写法正是这一格的常规
        工具，它把运行期的意外提前到编译期，让编译器替你说那句话。
\end{itemize}
\end{compare}

\begin{bestpractice}[title={按类别选择处理手段的四步法}]
\begin{enumerate}
  \item 先分类：写出这段代码依赖的那条保证，然后查它是五类中的哪一类。
        分不出来就不要开始写防御代码。
  \item 再决定要不要依赖：\IDB{} 的依赖必须附实现名与文档出处；
        \USB{} 的依赖默认拒绝，除非有探测手段。
  \item 给 \UB{} 建检查面：消毒档进 CI，覆盖真实输入分布；断言写成
        谓词而不是写成注释，并且清楚断言在 \Tp{NDEBUG} 下会消失。
  \item 留证据：每个跨平台分支旁边写明它依赖的类别、验证过的两个实现，
        以及一旦换编译器时该看哪一段文档。
\end{enumerate}
\end{bestpractice}

\begin{note}[title={分类的边际收益}]
实际排障里，把事故的类别说对，往往比立刻找到根因更省时间：
它决定了你应该去查标准、查实现文档、还是去开消毒器重跑。
三者混用的典型症状是同一个 bug 被三个人用三种方法各自“修好”了一次。
\end{note}

\section{本册路线图}\label{sec:ub01-roadmap}

剩下十三章与三个附录各自回答一个问题，顺序按“先能测、再懂模型、
最后谈约定”安排：

\begin{itemize}
  \item 第~\ref{ch:ub02-inevitable} 章：为什么任何语言都免不了这一类问题，
        别的语言把它换了什么名字；为什么留白是“标准”这一文体的性质，
        与起草方式和学科都无关（\S\ref{sec:ub02-standard-blanks}）；
        以及“触发 \UB{} 就是错误”这句话的三条限定
        （\S\ref{sec:ub02-three-limits}）。
  \item 第~\ref{ch:ub03-sanitizers} 章：消毒器能看见哪几类，看不见哪几类。
  \item 第~\ref{ch:ub04-lifetime} 章：对象生命周期与不确定值是一切内存事故的起点。
  \item 第~\ref{ch:ub05-pointer} 章：指针算术、尾后指针与比较的合法边界。
  \item 第~\ref{ch:ub06-aliasing} 章：严格别名规则为什么能改变生成代码。
  \item 第~\ref{ch:ub07-alignment} 章：对齐不足与类型双关的可行范围。
  \item 第~\ref{ch:ub08-evaluation} 章：求值顺序、副作用与序列点。
  \item 第~\ref{ch:ub09-integers} 章：整数表示、溢出、移位与位操作。
  \item 第~\ref{ch:ub10-memory-model} 章：内存模型、原子操作与数据竞争。
  \item 第~\ref{ch:ub11-terminate} 章：异常路径上的终止与不可恢复状态。
  \item 第~\ref{ch:ub12-library-pre} 章：标准库前置条件被违反时的表现。
  \item 第~\ref{ch:ub13-implementation} 章：\IDB{} 清单与跨实现差异。
  \item 第~\ref{ch:ub14-conventions} 章：\USB{}、特性探测与库的约定。
  \item 第~\ref{ch:ub15-forward} 章：前向兼容——新版本如何把旧代码变成 \ILLF{}。
  \item 附录~\ref{ap:ubA-kinds}：按类别汇总的判据总表。
  \item 附录~\ref{ap:ubB-switches}：开关、诊断选项与工具面总表。
  \item 附录~\ref{ap:ubC-cases}：可复现案例的完整复现步骤。
\end{itemize}

\begin{guideline}[title={读完全册应该建立的一个习惯}]
拿到任何一段可疑代码，先做一件事：把它依赖的每条保证归到
良好定义、\UB{}、\IDB{}、\USB{}、\ILLF{}（含 \NDR{}）五格里的一格。
归得进去，就知道该查标准、查文档还是开工具；归不进去，说明你还没
真正读懂这段代码在依赖什么。分类不是术语练习，它是排障路线的起点。
\end{guideline}

\chapter{为什么任何语言都免不了未定义行为}\label{ch:ub02-inevitable}

本章是全册的立论章。它要回答的不是“\UB{} 有哪些”，而是一个更前置的问题：
为什么这门语言必须有这个概念，别的语言又是如何绕开或者改名绕开的。
答案决定了一件事——把 \UB{} 当成 \CL/\CPP 的设计缺陷来批评，
或者当成可以靠纪律彻底消除的污点来治理，都是错的。

\section{核心论断：四个口子与三条出路}\label{sec:ub02-trilemma}

只要一门语言允许下面四类操作，它就绕不开一次取舍。这四类操作是语言
表达能力的代价，不是某门语言的偏好：

\begin{enumerate}
  \item 越出可表示域的算术：加法结果超出类型能表示的范围。
  \item 以类型不匹配的方式访问存储：把一个地址当作另一种布局来读写。
  \item 读取未初始化内容：读一块从未被赋予值的存储。
  \item 不受约束的共享可变状态：两个执行序列同时改同一块存储。
\end{enumerate}

这四类操作不是某门语言的偏好，而是语言提供底层控制力时一并给出的能力。
先看它们为什么关不掉，再看语言能选的三条出路。

\subsection{四类口子为什么关不掉}\label{sec:ub02-why-open}

定宽整数之所以是定宽，因为它对应机器里一个寄存器与一条指令；一旦允许
溢出，就越出了可表示域，而硬件给出的那个结果并不由语言负责。
把地址当作另一种布局来读，是分配器、序列化、网络报文、内存映射文件
与向量指令的共同需求；禁止它等于禁止这一整类程序。读取未初始化内容
之所以可能发生，是因为把每一块新内存清零本身就是一项持续付出的代价。
不受约束的共享可变状态，则是多线程与共享地址空间换并行度的直接后果。

这四件事没有一件是“写法不小心”造成的，它们是语言承诺提供底层控制力时
一并提供的能力。设计者能做的是决定它们在哪里可用：全场景可用（\CL/\CPP）、
只在显式声明的区域可用（\Lang{Rust}、\Lang{C\#}、\Lang{Swift} 的逃生舱）、
或者原则上不提供（\Lang{Go} 的零初始化与语言层无裸指针）。
“不提供”这一档最接近消灭问题，但它同时也是最贵的一档，
因为它把代价摊到每一个程序头上，而不是摊到需要那份控制力的程序头上。

\subsection{三条出路与其代价}\label{sec:ub02-three-ways}

对这四类操作，语言设计只有三条出路，没有第四条。

第一条：为它们定义明确行为。要么每次算术都做溢出判断，要么每次访问
都带上类型与边界信息，要么读之前保证已初始化。代价是运行期检查、
对象元数据，或者放弃某些硬件指令。

第二条：放弃一部分优化。一旦允许溢出、跨布局访问与竞争，编译器就
不能假定它们不发生，许多以“不发生”为前提的改写全部作废。
这一条对第~\ref{ch:ub01-taxonomy} 章的 as-if 规则是直接的减税。

第三条：每次操作都加运行时检查。这一条其实是第一条的实现方式，
但它把成本变成显性：检查必须留在指令流里，元数据必须随时可取。

所谓“未定义行为”只是第一条路选了一种极端形式：既不定义行为，
也不加检查，而是撤回承诺，把优化空间留给实现。它是定价方式，
不是道德缺陷。一门语言如果拒绝这第三种定价，它就必须在另外两条里
付账——问题从来不是“要不要付”，而是“在哪一层付”。

\begin{guideline}[title={一条可以拿去检验的判据}]
问任何一门语言：越界读写、溢出、读未初始化、并发竞争这四件事，
它在哪一层付账——编译期规则、运行期检查、还是撤回报错并交回程序员？
答不出来的一定是还在依赖宿主实现而没有写完契约的语言，
而不是“没有这个问题”的语言。
\end{guideline}

\section{成本对照：九类语言把账记在哪里}\label{sec:ub02-cost}

下表只写机制，不写性能数字：任何“慢几倍”“体积增加百分之几”之类的
说法都依赖实现、档位与工作负载，本册一律不作断言。

\begin{longtable}{@{}>{\raggedright\arraybackslash}p{0.13\textwidth} >{\raggedright\arraybackslash}p{0.24\textwidth} >{\raggedright\arraybackslash}p{0.24\textwidth} >{\raggedright\arraybackslash}p{0.15\textwidth}@{}}
\caption{四类越界操作在不同语言里的处置方式与代价落点（只述机制）}\label{tab:ub02-cost}\\
\toprule
语言 & 越界访问的处置 & 未初始化读取的处置 & 代价落在哪一层 \\
\midrule
\endfirsthead
\toprule
语言 & 越界访问的处置 & 未初始化读取的处置 & 代价落在哪一层 \\
\midrule
\endhead
\bottomrule
\endfoot
\CL 与 \CPP & 解引用越界指针或写下标越界的元素是 \UB{}；尾后指针可以构成但不保证可解引用 & 不确定值可被读，读到的是 \UB{} 的一个来源；填充位有专门豁免 & 运行期不付账，代价以优化空间的形式收回 \\
\Lang{Java} & 数组访问带边界信息，越界抛出定义好的异常；虚拟机可消除证明得了的检查 & 局部变量有明确赋值规则，未赋值的用不到；字段总有默认值 & 编译期规则加运行期元数据，检查留在指令流 \\
\Lang{C\#} & 默认检查越界并抛出异常；\Tp{unsafe} 加指针算术绕过检查 & 值类型必须显式赋值后才可用；引用类型默认为空引用 & 默认在运行期，选择逃生舱后转给程序员 \\
\Lang{Rust} & 安全子集里索引越界即中止（定义良好）；\Tp{unsafe} 块内指针越界仍按同一套 \UB{} 判据衡量 & 未初始化的内存要通过 \Tp{MaybeUninit} 这类显式包装才能出现 & 编译期借用检查付账，\Tp{unsafe} 那一层重新变成运行期风险 \\
\Lang{Swift} & 索引越界与整数溢出都触发定义好的运行时中止；可选值解包为 \Tp{nil} 也有规定行为 & 变量必须先初始化才能读；类型系统挡住裸内存 & 默认在运行期；\Tp{Unsafe} 前缀的类型是显式逃生舱 \\
\Lang{Go} & 切片与数组越界产生定义好的运行时 panic & 所有变量零初始化，把“读未初始化”这一整类问题变成定义良好的读零 & 运行期检查；竞争另由内存模型与检测器处理 \\
\Lang{Python} 与 \Lang{Lua} & 语言层索引越界抛出各自的定义好的异常 & 语言层没有裸露的未初始化概念 & 风险被转移到宿主 C API 与扩展层，由扩展作者付账 \\
\Lang{JavaScript} & 越界读写数组或类型化数组返回 \Tp{undefined} 或被忽略，都是定义好的结果 & 变量声明后即为 \Tp{undefined}，不存在裸内存 & 定义良好但精度与数值语义由实现与规范细节决定 \\
\end{longtable}

\subsection{读这张表的三种姿势}\label{sec:ub02-cost-read}

第一，检查不是免费的。Java、C\#、Swift、Go 把边界判断写进指令流，
就必须让对象随身携带长度或类型信息；这与 \CL/\CPP “数组名退化为地址、
长度由程序员保管”是两种不同的存储契约，不是两种不同的责任心。

第二，全部这些语言都留了逃生舱。C\# 有 \Tp{unsafe}，Rust 有
\Tp{unsafe}，Swift 有 \Tp{Unsafe} 系列类型，Java 有历史上的
\Tp{sun.misc.Unsafe} 与后来的 \Tp{VarHandle}，Python 与 Lua 有各自的
C API。逃生舱的作用不是提供新功能，而是让某些必须贴近机器的代码
能够存在。一进逃生舱，语言就把检查责任交回程序员，随之交回的还有
\UB{} 这一整套定价方式。

第三，代价落点决定治理手段。落在编译期的问题只能在类型系统与
借用规则里修；落在运行期的问题只能靠检查与 profiling；
落在程序员手里的部分只能靠工具与纪律。第~\ref{ch:ub03-sanitizers} 章的工具面之所以对
\CL/\CPP 特别重要，正是因为这一门语言
把账记在了程序员手上。

\subsection{第四个口子：共享可变状态怎么定价}\label{sec:ub02-races}

表格的三列讲的是访问与初始化，并发那一格值得单独说，因为它是四类口子里
唯一一个“单线程永远遇不到”的。各家的处置仍然是同一条光谱：

\begin{itemize}
  \item \Lang{Rust} 在编译期用借用检查加上 \Tp{Send} 与 \Tp{Sync}
        两个标记特质，把可跨线程共享的东西收窄到不会互相冲突的那一类。
        代价换成了“有些写法写不出来”，收益是竞态在编译期消失。
  \item \Lang{Java} 与 \Lang{C\#} 给出定义良好的内存模型：未同步的并发读写
        得到的值由模型规定其可见性范围，而不是让实现随意处置；
        长期未同步的写入允许被缓存，因此程序仍可能读到旧值。
        这不是 \UB{}，但它同样需要程序员理解规则才能写对。
  \item \Lang{Go} 的内存模型文档明确表态：带竞争的程序行为不可预测。
        这一句承担的职责与 \UB{} 完全一致，只是它靠 \Tp{-race}
        这类显式开关的检测器提供反馈。
  \item \Lang{Swift} 在调试构建里提供排他性检查，并把跨任务可变状态
        收进 actor 一类隔离机制；裸指针与手动并发仍落在逃生舱里。
  \item \CL/\CPP 的态度与前面几章一致：数据竞争本身被标准判为
        \UB{}（\Std{intro.races} \Clause{数据竞争}），同时给出
        \Tp{std::atomic} 与内存序作为“想买回确定性请按这里”的入口。
        第~\ref{ch:ub10-memory-model} 章处理这条线。
\end{itemize}

\begin{lstlisting}[style=shellstyle, caption={同一份 C++ 源码在“定义良好的邻居语言”与“撤回承诺的语言”里的不同处置：命令行侧的对照}]
# 有符号溢出：交给回绕定义，代价是取消一批优化前提
g++ -std=c++23 -O2 -fwrapv t.cpp -o t_wrap
# 交给动态工具报告，代价是运行期插桩
g++ -std=c++23 -O1 -fsanitize=undefined t.cpp -o t_check
# Go 一侧不需要选择：越界与整数溢出都由语言规定处置
go vet ./...
\end{lstlisting}

\section{UB 的变体：没有这个词的语言怎么承担同一职责}\label{sec:ub02-variants}

\UB{} 在别处不叫 \UB{}，但它承担的职责一件没少。三种形态值得单独认识。

\subsection{Rust 的 unsafe 判据}\label{sec:ub02-rust}

Rust 把语言切成两半：安全子集里越界与竞争被编译期规则挡在门外；
\Tp{unsafe} 块里则存在与 \CPP 同构的一套未定义行为判据——构造悬垂引用、
用错误布局解读内存、破坏 \Tp{unsafe} 函数自己声明的前置条件。
这套判据在很长时间内写在开发手册而非正式语义规范里，
并且由 \Mx{Miri} 一类工具在实践中界定：工具报错的那条规则，
往往先于文字规范成文。也就是说，Rust 不是“消灭了 UB”，
而是把 UB 的边界收进了一小块显式声明的区域，并且用工具替它执法。

\subsection{JVM 的 final 字段与 happens-before}\label{sec:ub02-jvm}

Java 的内存模型经历过一次可称为“承认现实”的修订：为了让即时编译器
有权把 final 字段的读取上提或缓存，规范必须允许某些未加同步的观察
看到尚未冻结完成的对象。后来引入的 \Clause{happens-before} 关系把这件事
重写成有边界的规则：在没有建立同步的读线上，字段可见性并不保证。
这不是 \UB{} 这个词，但它解决的是同一个问题——语言的承诺必须给
优化留出可表述的空间，否则承诺无法兑现。

\subsection{宿主 API 里的未定义契约}\label{sec:ub02-host}

解释器的 C API 是最干净的对照。语言层把内存管理收进运行时之后，
扩展层又把契约原样交回程序员：引用计数失衡、栈索引越界、
没有按宿主规定持锁就触碰解释器状态——这些行为的后果在宿主文档里
被写成“未定义”或“必须自行保证”，而不是抛出异常。换句话说，
一门没有 \UB{} 的语言可以在它的扩展边界上重新发明 \UB{}，
因为它需要的从来不是“有 UB”，而是“有一层不再承诺的边界”。

\subsection{同一职责在不同语言里的登记表}\label{sec:ub02-registry}

把上面三种形态连同表格里的机制收成一张对照表，可以直接看出：
名字不同，职责相同，执法者也各不一样。

\begin{table}
\centering
\caption{同一层边界在各语言里的称呼与执法方式（机制描述，不含性能断言）}
\label{tab:ub02-registry}
\begin{tabularx}{0.76\textwidth}{@{}>{\raggedright\arraybackslash}p{0.16\textwidth} >{\raggedright\arraybackslash}p{0.26\textwidth} X@{}}
\toprule
语言 & 这一层的称呼 & 由谁执法 \\
\midrule
\CL 与 \CPP & \UB{}，写进标准的术语条目 & 标准撤回要求，后果由测试与工具承担 \\
\Lang{Rust} & \Tp{unsafe} 块内仍然存在的 \UB{} 判据 & 编译期规则加 \Mx{Miri} 一类动态语义工具 \\
\Lang{Java} & 同步之外的可见性缺口，规范用关系式描述 & 虚拟机实现与规范条文，异常兜住越界 \\
\Lang{C\#} & 不安全上下文里的裸指针操作 & 编译器要求显式开关，运行时不再提供检查 \\
\Lang{Swift} & 前置条件失败即中止，\Tp{Unsafe} 系列类型例外 & 运行时断言；逃生舱交回程序员 \\
\Lang{Go} & 带竞争的程序行为不可预测 & 语言层检查兜住越界，竞态靠显式检测器 \\
\Lang{Python} 与 \Lang{Lua} & 宿主 API 文档里的“必须自行保证” & 扩展作者，运行时只在语言层负责 \\
\Lang{JavaScript} & 越界读写有规定结果，数值精度另议 & 规范定义语言层，引擎只在优化内自由 \\
\bottomrule
\end{tabularx}
\end{table}

\begin{clause}[title={为什么这一层不可省略}]
只要实现允许越出可表示域的算术、允许以另一种布局访问存储、
允许读取从未赋值的存储、并且不禁止不受约束的共享可变状态，
它就必然要在三条出路里选一条。选择“撤回承诺”的语言用 \UB{} 这个词；
选择“定义并检查”的语言用异常、panic 或中止来称呼同一层的边界。
标准的措辞可以在 \Std{defns.undefined} \Clause{未定义行为} 一类的
术语条目里核对，但这条论断并不依赖任何一版标准的编号。
\end{clause}

\begin{compare}[title={C/C++ 特殊在哪里：不是有 UB，而是谁说了算}]
\begin{itemize}
  \item 判定权的位置：\CL/\CPP 把“这个操作是否越出可表示域、是否越界、
        是否有竞争”的判定权交给程序员，标准要求实现假定程序员不会越界。
        带检查的语言把判定权留在运行时，程序员不需要预先证明。
  \item 优化的定价基础：正因为判定权外移，编译器可以按“不会发生”
        来定价每一次改写。这是同一枚硬币的另一面，而不是额外的攻击。
  \item 违约的可见性：交给运行时的语言，违约会立刻变成异常或 panic；
        交给程序员的语言，违约可能一次都不显现。前者更贵，后者更隐蔽，
        这是效率与反馈速度之间的交换，不是对错的交换。
\end{itemize}
\end{compare}

\section{两组常见混淆}\label{sec:ub02-confusions}

\subsection{UB 不等于未初始化内存}\label{sec:ub02-not-uninit}

这两件事在全册必须分开使用。未初始化内存是 \UB{} 的\emph{一个来源}：
读到未初始化对象持有的不确定值，标准不保证该值的含义，填充位与
若干特殊情形另有豁免。但 \UB{} 的外延远大于它——有符号溢出、越界解引用、
违反严格别名、破坏前置条件、并发数据竞争，全都不涉及“未初始化”三个字。

反过来也不成立：读到未初始化内存并不总是 \UB{}。在某些语言里它被
定义为固定模式（调试构建里的填充字节），在 \Lang{Go} 里根本不存在
（零初始化把它取消了），而在 \CL/\CPP 的极少数情形下标准明确放行。
把两者当同义词，会直接导致工具选型错误：内存工具与消毒器覆盖的
不是同一件事，第~\ref{ch:ub04-lifetime} 章与第~\ref{ch:ub03-sanitizers} 章分头处理。

\begin{warning}[title={混淆的第一个症状：把 msan 当 ubsan 用}]
“这个函数里有未初始化读取，所以它一定有 UB”与“这里有 UB，
所以内存工具一定能报”都推不出来。未初始化读取只在你\emph{用到}那个值
参与决策或输出时才显现；越界写可能一行报告都不产生，直到很久以后
在别处崩掉。工具覆盖面与被测类别的对应关系，是第~\ref{ch:ub03-sanitizers} 章的正题。
\end{warning}

\subsection{ID 与 US 不是弱一等的 UB}\label{sec:ub02-not-weak-ub}

第二组混淆是把 \IDB{} 与 \USB{} 当成 \UB{} 的轻量版。它们是三类不同的
承诺：\UB{} 意味着没有任何实现义务；\IDB{} 意味着实现必须选定一个答案
并且\emph{写进文档}；\USB{} 意味着实现必须选定一个答案但不必告诉你。
把三者混为一谈的后果很实际：

\begin{itemize}
  \item 对 \IDB{} 写防御代码是浪费。它已经在文档里给出确定答案，
        你要做的是查文档，或者把它归一化。
  \item 对 \USB{} 查文档常常查不到。它只能被探测或者被回避。
  \item 对 \UB{} 查文档是徒劳。没有答案可查，只能改代码。
\end{itemize}

探测的思路在本册只点到为止：数值特征问 \Tp{std::numeric{\Bk}\_limits}，
库特征问 \Feat{\_\_cpp\_lib\_*} 一类的特性测试宏，语言特征问
\Feat{\_\_cplusplus} 与实现自定义宏。细节留给第~\ref{ch:ub13-implementation} 与
第~\ref{ch:ub14-conventions} 章。

\begin{lstlisting}[style=cppstyle, language=C++, caption={可移植的饱和加：先判后算，不依赖溢出行为（上下界来自实现提供的 limits 头文件）}]
constexpr int add_sat(int a, int b) {
    if (b > 0 and a > INT_MAX - b) { return INT_MAX; }
    if (b < 0 and a < INT_MIN - b) { return INT_MIN; }
    return a + b;
}
\end{lstlisting}

\begin{lstlisting}[style=shellstyle, caption={探测实现是否已提供标准饱和算术：查特性测试宏，而不是猜版本}]
# 让预处理器把全部宏吐出来，再在其中找目标特性宏
g++ -std=c++23 -dM -E -x c++ /dev/null | grep -F __cpp_lib_saturate_cast
# 查不到就按本册的写法自带饱和实现，查到就切换到库版本
\end{lstlisting}

\begin{casebox}[title={在别的语言里良定义、在 C++ 里未定义的等价写法}]
\textbf{场景：取容器最后一个之后的位置并参与比较。}

\begin{itemize}
  \item \Lang{Java}、\Lang{C\#}、\Lang{Go}：索引超出长度即抛出或 panic，
        越界这件事本身有确定的、可依赖的后果；把它写进控制流是合法的。
  \item \Lang{JavaScript}：越界读返回 \Tp{undefined}，越界写被忽略，
        同样是定义好的结果。
  \item \CL/\CPP：构成尾后指针本身是允许的，\emph{解引用}它是 \UB{}；
        继续向后算出“尾后加一”这类位置，比较与相减的合法性都不再有保证。
        把“越界后返回某个值”当作分支策略，等于在依赖撤回后的承诺。
\end{itemize}

\textbf{可移植写法}：把判界放在访问之前，而不是把访问放在判界之前。
用长度或 \Tp{std::span} 一类带边界的视图先做谓词判断，越界路径显式返回；
算术侧同理——如上例的 \Fn{add\_sat} 先判后算，或探测
\Feat{\_\_cpp\_lib\_saturate\_cast} 后改用库提供的 \Tp{std::saturate\_cast}，
查不到就保留手写版本。两种写法都不要求实现回答“溢出之后是什么”。
\end{casebox}

\section{标准这一文体的性质：留白不是疏漏}\label{sec:ub02-standard-blanks}

前面几节把 \UB{} 解释成语言设计在表达力与优化空间之间作出的定价。
这一节要把视野抬高一层，因为本书认为这才是问题的根：\UB{} 不是
\CL{}/\CPP{} 的设计缺陷，它是“标准”这一文体的性质。任何一份标准，
不分起草方式、不分它管的是不是计算机语言，都不追求把世界穷举规定完，
它必然留下一层空白。所谓“未定义行为”是这层空白在 \CL{}/\CPP{} 里的
一种\emph{处置方式}，不是空白本身。

\subsection{标准为什么必然留白}\label{sec:ub02-blank-why}

四件事同时成立，写规范的人没有一件躲得开。

\begin{enumerate}
  \item 组合不可枚举。语言构造两两组合出的程序空间是无穷的，规范只能按规则
        局部约定，无法逐个情形给出结论。规则与规则之间没覆盖到的交界就是空白；
        它不是没想到，而是穷举在原则上不可能完成。
  \item 判定不可完成。若给每一处空白补上“必须诊断”或“必须如何”，常常要求实现
        先判定一个关于程序运行期性质的一般问题——这块存储是否会被越出边界访问，
        这两个执行序列是否真的可能并发。这类问题的一般形式没有算法可解，
        所以“先判定再据此处理”这一类补法从一开始就不成立，写得出措辞的
        只剩“一旦发生就算违约”。
  \item 共识有成本。标准是多方签下的协议。把每一处空白都填成硬要求，等于要求
        每个签署方都为别人需要的能力付实现成本；谈判里真正省事的写法，是把无法
        共同负担的部分显式交出去。这正是 \IDB{} 与 \USB{} 两类留白存在的理由，
        也是它们与 \UB{} 的差别所在。
  \item 前向要留余地。规定得越完整，未来硬件与实现策略的可选项越少。以整型溢出
        为例：钉成“按补码回绕”会排除某几类机器，钉成“陷入或饱和”会排除另几类；
        写成不承诺，两边都还能留在这份标准里。第~\ref{ch:ub15-forward} 章讲的是
        这笔余地的另一面——今天的空白会被后续版本填上，填上之后旧写法就可能
        变成 \ILLF{}。
\end{enumerate}

\subsection{与起草方式无关}\label{sec:ub02-blank-not-drafting}

最常见的解释是“委员会流程不够好，才留下这么多空白”。这个解释经不起分类检验。
把语言的规范化路径排一遍，每一条都带同一层空白，只是位置不同。

\begin{itemize}
  \item 自下而上的委员会标准：多轮提案、逐条表决。它的措辞里专门设有
        \Clause{未定义行为}、\Clause{实现定义行为} 这一类术语来登记空白
        （\Std{defns.undefined}、\Std{intro.implementation} 一类），
        说明留白是流程的有意产物，不是笔误。
  \item 实现先行、事后成文：先有一批事实行为，再由人把它们写下来。这条路看似
        最不容易漏，实际只是把空白挪到“哪一份实现算权威”这个问题上，
        版本分叉与方言差异就是它那份空白的表现形式。
  \item 单一实现即事实规范：规范与代码同生命期，看上去最接近穷举。但实现自己
        也无法穷举输入组合，未覆盖之处只能被描述成“目前跑出来是这样”，
        而这不是承诺。
  \item 组织内部的编码规范：常有人以为把语言子集收窄到可静态分析就能消灭空白。
        面向安全关键领域的 \CL{} 子集规范是这一路的代表，它确实删掉了大量
        未定义构造，代价是把表达能力的边界一并删掉；而它自己没覆盖的库与平台
        行为，仍然要引用外部文档补齐。
\end{itemize}

四条路径都留下了同一层东西。这正说明空白不是流程质量的产物，改进流程不会让它消失。

\subsection{也不是计算机语言标准的专利}\label{sec:ub02-blank-not-cs}

换一个学科看得更清楚：“规定”这一文体在别处同样不追求穷举。性能化的建筑规范
给出必须满足的性能目标，把具体构造做法留给当地条款与工程师判断；成文法无法
预见所有案情，缺口由解释与判例填补；诊疗指南给出推荐，不覆盖每一个个体差异；
计量标准定义单位与溯源方法，不规定每一次测量的操作细节。这些文本里都能找到
“本文件不作规定”一类的句子，它承担的功能与 \Std{defns.undefined} 那一格完全对应。

区别只在后果由谁承担。上面那些学科把空白处的判断责任交给执行者——工程师、
法官、医生、实验者；\CL{}/\CPP{} 把同样的责任交给程序员与实现者，只是它额外
提供了一整套工具来检查自己的判断。本册后十三章的全部内容就是这套工具。

\begin{guideline}[title={看到“标准没写”，先分辨是哪一种留白}]
\begin{itemize}
  \item 标准故意把选择交给实现，并要求写进文档：\IDB{}。去处是实现文档，
        查得到就能用，但要附实现名与版本。
  \item 标准故意把选择交给实现，不要求公开：\USB{}。去处是探测与约定，
        第~\ref{ch:ub14-conventions} 章展开。
  \item 标准拒绝作出承诺：\UB{}。去处是工具与纪律，不是文档。
  \item 标准本就不管这一层：库的契约、平台 ABI、加载器与线程库的行为。
        去处是另一份文档，别在语言标准里找它。
  \item 标准还没来得及写：版本推进中尚未冻结的空白。去处是
        第~\ref{ch:ub15-forward} 章。
\end{itemize}
\end{guideline}

\begin{note}[title={本节的一句话}]
\UB{} 的存在不需要用“语言设计粗糙”来解释。它是标准这种文体在不可枚举、
不可判定、共识有成本、前向要留余地这四件事面前的必然产物。一门语言能选的
从来不是“要不要留空白”，而是“这处空白用什么方式定价”：\IDB{} 写进文档，
\USB{} 不写，\ILLF{} 直接拒收，\NDR{} 放弃诊断义务，\UB{} 撤回承诺。
\end{note}

\section{三条限定：编译期计算、边界触发与固化行为}\label{sec:ub02-three-limits}

前两节把 \UB{} 说成标准的留白与语言的定价方式。这一节给三条限定，它们针对的是
同一句常被顺口说出来的话——“触发了 \UB{} 就是错误”。这句话作为一般判断是对的，
但在下面三种场合需要改写，否则结论会走偏。

\subsection{编译期计算能排除绝大多数，但绝大多数问题不可编译期计算}\label{sec:ub02-limit-constexpr}

这一条只属于 \CPP{}。\CL{} 没有把任意函数在编译期求值这套机制，它的常量表达式
只覆盖有限的折叠。\CPP{} 的常量求值（\Tp{constexpr} 函数被常量实参调用、
\Tp{consteval}、\Tp{static\_assert}、模板实参一类上下文）有一条硬规则：
求值过程中只要走到 \UB{}，这条表达式就\emph{不是常量表达式}，于是整个程序不合法。
换句话说，标准在那一格写下的“不承诺”，在常量上下文里被实现换成了“编译期拒收”。
在这个覆盖面内，绝大多数 \UB{} 确实被排除了——不是被修好了，而是被提前到了
一个不可能上线的时刻。

\begin{lstlisting}[style=cppstyle, language=C++, caption={同一个函数在两种上下文里的判定强度完全不同}]
constexpr int scale(int v) { return 2 * v + 1; }   // 有符号溢出：标准不承诺回绕

// 常量上下文：v 在编译期已知，溢出被真的算出来，整条表达式不是常量表达式。
// 下面这一行会直接编译失败，诊断信息指向 scale 的调用点。
// static_assert(scale(1 << 30) > 0, "check width");

int run(int argc) {
    return scale(argc);   // 运行期实参：同一个溢出，编译器一句话也不说
}
\end{lstlisting}

限制也就在这里，而且有两层。第一层是输入：绝大多数真实事故触发条件依赖
编译期拿不到的东西——外部输入的长度、来自共享库的返回值、分配器给出的地址、
并发时序、文件内容。这些东西无法常量化，于是同一段代码在测试里能被
\Tp{static\_assert} 覆盖，在线上就没人核查。第二层是常量上下文自身的边界：
它禁掉了一批操作（输入输出、线程与原子、超出限制的动态存储），
可量化的范围还随版本在扩大，具体到哪一档能用要以实现为准。
因此“能过编译期”与“没有 \UB{}”是两个不同的轴：常量求值器只对
\emph{能算}的部分负责，它不是消毒器，也不证明剩下的部分安全。

正确的用法是把编译期计算当成\emph{收窄触发面}的手段：能常量化的那部分
一律常量化，剩下的仍然按第~\ref{ch:ub03-sanitizers} 章的工具面去覆盖。
把它当成判据，就是把“实现替我算了一遍”误读成“标准替我保证了”。

\subsection{触发者常常不在库内，甚至不在这个进程或这台机器上}\label{sec:ub02-limit-crossborder}

“这次 \UB{} 是这段库的代码造成的”是另一个默认说法，它在一类常见结构里不成立：
库通过外部 ABI 对外提供服务——也就是以 C ABI 暴露的不透明句柄加一族函数，
创建、读、写、销毁各自一个入口，由调用方决定顺序与时机。这种边界上，
违约的动作通常发生在\emph{调用方}：多销毁一次、销毁后继续用、
把两侧对布局的理解不一致的指针传来传去。

更麻烦的是，触发者可以完全不在同一个执行环境里。对象经由共享内存交给另一个进程；
结构体被序列化到另一端，由另一台机器上的另一份代码按同一份头文件解读；
句柄被写进持久格式，下次启动时才被恢复。此时“谁在什么时刻越出承诺”
与“谁的代码里写着那行越界的访问”分属两个程序、两个地址空间，
甚至两个设备。本册第 1 章那句“标准对整条执行轨迹负责”在这里要反过来读：
轨迹可以跨过程序边界。

由此得到一个反直觉的结论：在这种结构里给库补一层生命周期管理
（引用计数、有效性表、延迟释放、销毁后置空）通常\emph{修不好}这件事，
还会带来新的问题。它对使用者侧的可观察行为做了改写：原本会崩的用例从此静默成功，
原本失败的集成测试变绿，释放时机被推迟从而改变了内存曲线、析构顺序与
外部资源的归还时刻；而在跨进程、跨设备的那一类场景里，那份数据仍然按错误的布局
被解读——库内的检查根本看不到远端。现象被压住了，违约还在。

处理这类 \UB{} 的正确落点不在库内部，而在边界上，本册第~\ref{sec:ub11-boundary} 节、
第~\ref{ch:ub12-library-pre} 章与第~\ref{ch:ub14-conventions} 章分别展开了这些手段：

\begin{itemize}
  \item 把所有权写进接口文档：谁创建、谁销毁、能否跨线程转移、句柄在销毁后
        是否仍被别的模块持有，每一条都是前置条件，不是注释。
  \item 边界上把违约映射成\emph{显式可返回的失败}（错误码、\Tp{std::optional}），
        而不是在库内部兜住。兜住就是改写语义。
  \item 检查放在\emph{使用者侧}的构建里跑：消毒器在调用方的二进制上才能看见
        调用方的调用顺序，库自己那份单元测试覆盖不到集成时序。
  \item 跨进程与跨设备不用“按结构体直接读字节”，改用带版本号与显式布局的格式
        （第~\ref{sec:ub04-case-wire-union} 节与第~\ref{ch:ub07-alignment} 章的
        布局、填充讨论在这里是同一件事的另一面）。
\end{itemize}

\begin{guideline}[title={问三个问题再决定“谁来修”}]
\begin{enumerate}
  \item 违约动作与越界代码在同一个地址空间里吗？不在，就没有任何库内检查能看见它。
  \item 修完之后，调用方观察得到的行为变了吗？变了，就要走接口版本而不是内部兜底。
  \item 使用者侧的集成测试能否复现？不能复现的检查，只是把失败推迟到线上。
\end{enumerate}
三个问题里有一个答“否”，就不要用“加生命周期管理”来回答 \UB{}。
\end{guideline}

\subsection{被主流实现固化的行为，可以在部分场合被接受}\label{sec:ub02-limit-frozen}

第三条限定最容易被滥用，所以先划清界线：这里说的\emph{不是} \IDB{} 与 \USB{}——
那两类是标准明确交出去的选择，前者还有文档义务（第~\ref{ch:ub13-implementation} 章）。
这里说的是标准\emph{撤回承诺}、而所有主流实现却给出同一个结果的情形。常见形状有：
补码回绕（\cpp{20} 起标准已强制补码\emph{表示}，溢出依然是 \UB{}，但机器上的结果是
确定的）、经由联合成员做类型双关（某些实现族把它的语义写进了自己的文档）、
关闭严格别名检查后的别名行为、以及把指针连同整数一起搬过边界再还原
（provenance，第~\ref{sec:ub05-prov} 节）。

这类写法在部分场合可以被接受——被接受的意思是\emph{风险被有记录地承担下来}，
不是它变成了定义行为。可接受的前提是四条，缺一就不要写：

\begin{enumerate}
  \item 写下边界：实现族、架构、开关档位与版本范围，越出这个范围就是未验证区。
  \item 用断言或探测把假设钉住（第~\ref{sec:ub13-detect-pin} 节），让换实现时
        第一次编译就失败，而不是第一次上线才失败。
  \item 在第二个实现上验证过，或明确记录放弃可移植性并写清放弃的代价。
  \item 附出处：把“哪份文档的哪一节支持这个假设”写在代码旁边。
\end{enumerate}

代价也要写下来。固化只是当下这份实现的选择：编译器升级、优化档位或整程序优化
一开，同一条路径的结果可能变；消毒器与静态分析会照常报告这类写法，需要显式抑制，
而抑制本身要在抑制点留下理由（附录~\ref{ap:ubB-switches} 收集了开关面）。
一条底线：如果这段代码需要跨两个实现族，或需要接受外部不可信输入，
那么它就不属于“可以被接受”的那一部分。

\begin{note}[title={三条限定共同不改的那句话}]
“触发 \UB{} 就是错误”作为一般判断仍然成立。三条限定改的是这句话的归属，不是它的对错：
编译期计算改变的是\emph{谁在什么时候说出这句话}（实现替你把不承诺换成拒收）；
跨边界触发改变的是这句话该\emph{记在谁的账上}（未必是库，未必是同一个进程）；
固化行为改变的是说完这句话之后\emph{跟什么动作}（不是修复违约，而是把风险钉住并留档）。
\end{note}

\section{把结论压成可操作的形式}\label{sec:ub02-practice}

\begin{lstlisting}[style=outstyle, caption={工具输出样例：同一类问题在不同层的可见形式（示意形状，非实测）}]
# 编译期规则的语言：违规被拒收，信息指向声明与借用规则
# 运行期检查的语言：违规表现为抛出、panic 或 trap，可被上层看见
# 撤回承诺的语言：编译器不说话，运行一次也不报错，问题留在源码里
\end{lstlisting}

三段输出对应三条出路，它们真正的区别是\emph{反馈时机}：
编译期规则最早给出信号，运行期检查在触发时给出信号，
撤回承诺则可能一直不给信号。反馈越晚，成本越倾向于落在测试与
事故处理上；反馈越早，成本越倾向于落在语言能表达的范围内。
\CL/\CPP 选择把表达能力的上限抬高，同时把反馈时机推后，
这就是它需要第~\ref{ch:ub03-sanitizers} 章整套工具的原因，也是本册存在的理由。

代价的讨论只到定性为止。可以说“加了检查会留下数量级的改写空间损失”，
也可以说“取消某些优化前提会缩小向量化与范围推断的可用范围”；
不能说“慢几倍”或“体积涨多少”——这些数字取决于实现、档位与工作负载，
本册不给。

\subsection{三条出路在工程上的对应物}\label{sec:ub02-engineering}

把三条出路翻成日常动作，会得到一份很直的对应表，它可以当作
任何一次代码评审的开场白。

\begin{itemize}
  \item 选了“编译期规则”的语言，工程动作是设计类型：让非法状态
        无法表示，让越界与竞争在写出来的那一刻就被拒收。
        \CL/\CPP 里等价的动作是类型包装、\Tp{std::span} 一类带边界的视图
        与 \Tp{not\_null} 风格的约束（在第~\ref{ch:ub12-library-pre} 与
        第~\ref{ch:ub14-conventions} 章展开）。
  \item 选了“运行期检查”的语言，工程动作是配置与观测：把检查留在
        构建里，把触发次数交给 profiling，把异常处理写成控制流的一部分。
  \item 选了“撤回承诺”的语言，工程动作只剩两件：把可能的违约变成
        可看见的报告（消毒器、静态分析、警告档），以及把必须依赖
        实现选择的地方显式写下来并附文档出处。
\end{itemize}

本册十四章到十五章做的正是第三件事：前十二章把各类违约逐个摊开，
第~\ref{ch:ub13-implementation} 章整理 \IDB{} 的清单，第~\ref{ch:ub14-conventions} 章给探测与约定的写法，
第~\ref{ch:ub15-forward} 章讨论版本推进时旧写法如何变成 \ILLF{}。这一章的结论到此为止：\UB{} 不是意外的产物，
而是一门语言在“表达力、优化空间、反馈时机”三者之间作出选择之后
必然留下的那一层边界。

\begin{note}[title={一条可检验的推论}]
判断一段代码是否安全，等价于判断它的每一个操作分别落在这五类的
哪一格（第~\ref{ch:ub01-taxonomy} 章）：全部落在良好定义里，代码在任何一致
实现上都成立；落在 \IDB{} 或 \USB{} 里，代码成立到实现文档或探测结果
划出的那条线为止；任何一格落到 \UB{} 或 \NDR{} 上，代码的正确性就不再由
标准保证，只能由测试覆盖到的那些输入暂时保证。这条推论可以直接拿来
做代码审查：说不出类别的行，先当作待分类处理。
\end{note}

\chapter{消毒器与静态分析的口径}\label{ch:ub03-sanitizers}

判定一段代码有没有踩到未定义行为，靠的不是运气，而是三套互补的手段：
运行时插桩的消毒器、编译期推断的静态分析器，以及程序自己写下的断言与库检查模式。
三者的能力边界完全不同，把它们的结论混在一起读，就会得到误导性的“全绿”。
本章给出三套手段的分工、常用开关的确切含义、报告的读法，
以及一份能直接落进持续集成的分档策略。行为类别本身的划分见 第~\ref{ch:ub01-taxonomy} 章；
本册后面每一章列出的具体 UB，几乎都能在下面某一种手段的覆盖面里找到位置。

\section{三种手段各自的口径}\label{sec:ub03-division}

\textbf{消毒器}（sanitizer）是编译期插桩、运行期检查的机制：编译器在每个可疑操作前
插入一段检查代码，命中时打印报告，程序按配置继续或中止。它的前提是这段代码被执行过，
且命中的数据恰好落在被插桩的那个点上。它不需要模型，也不需要人补注解，
代价是运行时开销与覆盖率依赖测试用例。

\textbf{静态分析器}不看真实输入，而是在编译期对程序做抽象解释、数据流或路径推断，
把“可能到达的错误状态”作为警告报出来。它可以覆盖没被执行到的分支，
但必须对无法判定的情形做保守假设：要么漏报，要么误报，两者只能取舍，不能同时消掉。

\textbf{断言与库检查模式}是程序自己写下的契约：\Fn{assert}、\Tp{static\_assert}、
标准库的调试模式、容器与分配器的自校验。它们检查的是程序层面的不变式，
而不是标准层面的行为类别：越界的下标在契约层被抓住，通常比在 \UB{} 层被抓住便宜得多。

\begin{compare}[title={三类手段能覆盖哪一类行为}]
运行时插桩擅长\UB{}：有符号溢出、越界下标、未对齐访问、空指针解引用、
读取生命周期已结束的对象，这些都有明确的“动作点”，插桩即可判定。
编译期推断擅长\IDB{} 的可疑用法与\ILLF{} 片段：未初始化读取、所有权错配、
返回局部量地址、约定被违反，这些不需要执行也能看出结构问题。
断言与库检查模式擅长\USB{} 与契约层：当标准允许多个答案而程序只要其中一个时，
写进契约比写进类型系统更省事。三类手段都覆盖不了\NDR{}：
标准连诊断都不要求，工具自然没有义务发现。
\end{compare}

\begin{clause}[title={工具不是合规裁判}]
标准对\UB{} 只说“不做任何要求”（\Std{defns.undefined}），对\NDR{} 说“不要求诊断”。
消毒器报出错误并不改变标准结论，没报也不改变。因此“消毒器不报错”是一个测试证据，
不是合规证明；反过来，编译器有义务给出的诊断（语法错误、\ILLF{} 的必需诊断项）
也不能靠关掉消毒器来回避。
\end{clause}

\section{UBSan：把语义层的未定义变成一条报告}\label{sec:ub03-ubsan}

\subsection{开关族与子检查的展开}\label{sec:ub03-ubsan-expansion}

\Opt{fsanitize=undefined} 不是一项检查，而是一组子检查的简称；
它展开成哪些子集、叫什么名字，GCC 与 Clang 略有差异，个别检查名在版本间还改过名。
下面的写法是把展开写清楚、再按项目需要关掉其中几项的做法。

\begin{lstlisting}[style=shellstyle, language=ShellUB, caption={UBSan 开关族：展开、关项与定位增强}]
# 展开写法：把关心的子检查逐条列出，避免依赖简称的版本含义
gcc -std=c17 -g -fno-omit-frame-pointer \
    -fsanitize=signed-integer-overflow \
    -fsanitize=integer-divide-by-zero,div-by-zero \
    -fsanitize=bounds,bounds-strict,alignment \
    -fsanitize=unreachable,null,return \
    -fsanitize=object-size,float-cast-overflow \
    -fsanitize=float-divide-by-zero \
    -fsanitize=bool,enum,vptr,builtin -c widget.c

# 简称写法（GCC 与 Clang 的展开集合略有差异，个别项需要显式补上）
clang++ -std=c++23 -g -fsanitize=undefined t.cpp
clang++ -std=c++23 -g -fsanitize=undefined \
        -fno-sanitize=vptr -fno-sanitize=alignment mixed.cpp

# shift 类：移位方向、负数移位与移位宽度越界，需单列
gcc -O1 -g -fsanitize=undefined -fsanitize=shift \
    -fsanitize=shift-base -fsanitize=shift-exponent bits.c

# 只跑语义层、不做插桩对照：把某个子检查单独列出，确认它能否复现
clang++ -std=c++23 -g -fsanitize=signed-integer-overflow narrow.c
\end{lstlisting}

\begin{table}[htbp]
\centering
\caption{UBSan 子检查名与其对应的行为类别}\label{tab:ub03-ubsan-checks}
\begin{tabularx}{\textwidth}{@{}>{\raggedright\arraybackslash}p{0.24\textwidth} >{\raggedright\arraybackslash}X >{\raggedright\arraybackslash}p{0.16\textwidth}@{}}
\toprule
子检查名 & 它抓的是哪一类行为 & 本册去处 \\
\midrule
signed-integer-overflow & 有符号运算结果溢出，\UB{}，也是优化反推前提的常见来源 & 第~\ref{ch:ub09-integers} 章 \\
integer-divide-by-zero、div-by-zero & 除以零，整数侧是 \UB{} & 第~\ref{ch:ub09-integers} 章 \\
bounds、bounds-strict & 数组下标越界；前者只在有大小信息时下检查，后者更激进 & 第~\ref{ch:ub05-pointer} 章 \\
alignment & 未对齐访问，含对齐后的指针被强转回更宽类型 & 第~\ref{ch:ub07-alignment} 章 \\
null、nonnull-attribute、returns-nonnull-attribute & 空指针解引用、非空约定被违反 & 第~\ref{ch:ub05-pointer} 章 \\
unreachable & 控制流到达声明为不可达的位置 & 第~\ref{ch:ub08-evaluation} 章 \\
vptr & 虚表指针损坏，通常是析构后经由引用使用对象 & 第~\ref{ch:ub11-terminate} 章 \\
bool、enum & 对象表示落在类型值域之外 & 第~\ref{ch:ub04-lifetime} 章 \\
float-cast-overflow、float-divide-by-zero & 浮点转换越出整数域，浮点除零按\IDB{} 处理 & 第~\ref{ch:ub09-integers} 章 \\
object-size、function & 访问超出对象声明大小、以错误类型调用函数指针 & 第~\ref{ch:ub05-pointer} 章 \\
builtin & 内建函数参数不满足前提 & 第~\ref{ch:ub12-library-pre} 章 \\
\bottomrule
\end{tabularx}
\end{table}

\subsection{报告之后是继续还是中止}\label{sec:ub03-recover-trap}

UBSan 的报告默认是“可恢复”的：打印一条 runtime error 之后继续执行。
这对定位有用，但对测试门禁是灾难——错误之后的程序状态已经不可信，
后面的断言失败多半是噪声。两种收紧方式：

\begin{itemize}
\item \Opt{fno-sanitize-recover=all}：命中的检查一律中止进程，
      并让退出码非零。可以只对个别项收紧：把
      \Opt{fno{\Bk}-sanitize{\Bk}-recover} 后面接上子检查名，
      与 \Opt{fsanitize{\Bk}-recover=alignment} 混用，用于逐步收敛。
\item \Opt{fsanitize-trap}：把检查点编译成非法指令，不再调用运行时库。
      体积小、无外部依赖，代价是没有源码位置、也没有可读的文本报告，
      现场只剩一条信号与栈帧，需要事后符号化。
\end{itemize}

两个开关都只影响“报出来之后做什么”，不改变判定集合本身。
运行期还有环境变量可调，例如 \Tp{UBSAN\_OPTIONS} 的
\Tp{print\_stacktrace}（把回溯一并打出来）与 \Tp{halt\_on\_error}，
可用取值以所用运行时版本的文档为准。

\begin{lstlisting}[style=outstyle, caption={三种报告形状的差别（示意，非本仓库实测）}]
# 默认（可恢复）：有位置，但程序继续跑完
runtime error: signed integer overflow:
    -2147483648 * -1 cannot be represented in type int
    widget.c:42:9

# -fno-sanitize-recover=all：同一条报告，进程随即中止，退出码非零
runtime error: index 8 out of bounds for type int [4]
    table.h:17:5

# -fsanitize-trap：只有地址与信号，位置信息消失
Illegal instruction (core dumped)
\end{lstlisting}

\subsection{位置缺失与栈不展开}\label{sec:ub03-missing-location}

实践中最常见的抱怨是“报告没有源码位置”或“回溯只有地址”。
两个直接原因：

\begin{enumerate}
\item \textbf{没带 \Opt{g}}。消毒器的定位信息依赖调试信息把插桩点与文件行号绑起来；
      没有调试信息时，位置退化成 \Tp{unknown}，行号列显示为 0，
      栈里的地址也无法符号化。
\item \textbf{帧指针被省}。优化档位默认省帧指针，栈展开器只能靠调试信息里的
      展开记录走栈；一旦报告来自被内联或被尾调用的位置，回溯就断在半路。
      加上 \Opt{fno{\Bk}-omit{\Bk}-frame{\Bk}-pointer} 通常是最省事的补救。
\end{enumerate}

\begin{note}[title={位置缺失时先读这三行}]
如果 \Tp{unknown file} 出现在最前，先怀疑调试信息；如果位置有、调用者是空，
先怀疑帧指针与链接选项；如果两者都正常但行号看着不像出事的地方，
多半是内联把位置归到了调用点——用 -O0 或 -O1 复现一次再判定。
\end{note}

\section{ASan 与 UBSan 的边界}\label{sec:ub03-asan-ubsan}

\subsection{内存层与语义层}\label{sec:ub03-layers}

ASan 处理的是\emph{内存访问本身}是否越界：它在堆块周围布红区（redzone）、
在栈与全局对象周围布影子内存，把每次分配与释放登记下来，
于是能抓 heap-buffer-overflow、stack-buffer-overflow、global-buffer-overflow、
use-after-return、use-after-free、double-free，以及配套的泄漏检查。
UBSan 处理的是\emph{操作在语义上}是否被标准许可：下标越界、值域越界、
未对齐、溢出、非空约定被违反。

两者的交集比多数人以为的小。UBSan 的 \Tp{bounds} 只在编译器能确定数组大小时插入检查，
对通过指针传递的缓冲区无能为力；ASan 只会检查这次访问有没有碰到红区，
对“值恰好写坏了但不越界”“下标合法但类型错”完全沉默。

\begin{warning}[title={两个开关不能互相顶替}]
\Opt{fsanitize=undefined} 通过不代表内存访问合法：越界写如果落在同一个对象的
相邻成员上，\Tp{bounds} 看不到，ASan 才看得到；反之 \Opt{fsanitize=address} 全绿
也不代表没有 \UB{}：有符号溢出、未对齐访问、读取非活跃 union 成员都不会有内存越界。
另外，TSan 与 ASan 通常不能同时启用，需要按轮次分开跑；
MSan 依赖额外的初始化追踪，覆盖面以所用实现文档为准。
\end{warning}

\subsection{容器注记与硬件标签}\label{sec:ub03-container-hw}

\textbf{容器注记}。标准库容器用内部指针表示区间，裸的越界检查插不进去。
为此 Clang 提供 \Opt{fsanitize=container-overflow}（部分文档把这一族写作
container-annotate，名称以所用版本的 \Opt{fsanitize} 列表为准），
它依赖标准库头文件里的容器注记来知道区间边界；
没有注记时该检查会静默降级，看起来像“全绿”。

\textbf{硬件标签}。\Opt{fsanitize=hwaddress} 是 ASan 的硬件辅助形态，
面向带内存标签扩展的 ARM 架构与配套内核：把可访问性检查交给硬件，
软件只负责报告与红区管理。定位与软件 ASan 不同：覆盖面按标签粒度受限，
开销明显低于软件形态，但只能在具备相应硬件与系统支持的机器上跑。
具体支持矩阵以工具链文档为准。

\section{MSVC 与 Windows 侧的口径}\label{sec:ub03-msvc}

\textbf{Clang-based ASan}：\Opt{fsanitize=address}。MSVC 的这项支持基于
LLVM 的 ASan 实现，需要链接随工具链分发的运行时库；为了拿到可读的回溯，
通常要同时关掉帧指针省略（\Opt{Oy-}）。不同版本覆盖的检查项数量不一样，
栈上缓冲区与部分检测能力是逐步补齐的，覆盖面请以所用 Visual Studio 文档为准。

\textbf{Code Analysis}：\Opt{analyze}。它报的是\emph{规则族}，编号集中在
C264xx 与 C268xx 两段。与未定义行为相关的常见族是：
指针与下标算术越界（C26481、C26482、C26483 一类）、
窄类型算术溢出（C26451 一族）、未初始化变量与指针（C26495、C26494 一类）、
返回局部量地址与生命周期错用（C26487 一类）、裸 \Tp{new} 与 \Tp{delete} 配对（C26409）。
单个编号在不同版本里会迁移、合并或改名，把它当规则族读，
不要把它当成稳定的接口；权威列表以 MSVC 的警告目录为准。

\textbf{\Opt{sdl}} 是安全开发生命周期检查集：启用一组缓冲区与调用约定相关的
编译期检查，并把若干安全相关警告提升为错误，同时打开运行时的安全检查开关。
它的覆盖面小于 \Opt{analyze}，但可以整仓库一次性打开，适合做基线。
三项的推荐组合是：\Opt{sdl} 常开，\Opt{analyze} 在改动目录上开，
\Opt{fsanitize=address} 在 Debug 配置下按测试轮次开。

\section{静态分析三件套对照}\label{sec:ub03-static-tools}

\begin{longtable}{@{}>{\raggedright\arraybackslash}p{0.15\textwidth} >{\raggedright\arraybackslash}p{0.21\textwidth} >{\raggedright\arraybackslash}p{0.25\textwidth} >{\raggedright\arraybackslash}p{0.18\textwidth}@{}}
\caption{静态分析手段：实现、开关族、能报的 UB 类别与误报代价}\label{tab:ub03-static-tools}\\
\toprule
实现 & 开关族 & 能报的行为类别 & 误报代价 \\
\midrule
\endfirsthead
\toprule
实现 & 开关族 & 能报的行为类别 & 误报代价 \\
\midrule
\endhead
\bottomrule
\endfoot
GCC 警告集 & \Opt{Wall}、\Opt{Wextra} 加本节末尾列出的专项警告 &
未初始化读取、能静态算出的越界与悬垂、严格别名违反、
移位宽度越界、fallthrough、求值顺序问题的一部分 &
绝大多数检查依赖优化档位，低档不报；未初始化与严格别名两支是误报大户，
需要逐条判后关闭或降级 \\
GCC 分析器 & \Opt{fanalyzer} 及其 \Opt{Wanalyzer-}\,* 子项族 &
内存类为主：use-after-free、double-free、分配未释放、越界写、
未初始化读、文件与分配器的配对错误 &
路径敏感，代价最高；大函数会超时或被内存上限截断，
需要按目录开并调 \Opt{param} 上限 \\
Clang 静态分析器 & \Opt{analyze} 或 scan-build，检查器族 core、cplusplus、
security、unix & 空指针解引用、未初始化值、new 与 delete 配对、
越界读写、API 约定误用 & 只看单函数与模型化的调用，误报中等；
跨翻译单元的部分看不到 \\
clang-tidy & clang-tidy 的 \Tp{bugprone-}\,* 命名族（细则见本节末尾） &
use-after-move、悬垂句柄、裸指针算术、可疑内存操作、
整数除法与窄化转换、约定层误用 & 走语法树不做数值推断，误报较低但会漏；
模板密集的代码明显变慢，需要按检查名白名单开 \\
MSVC & \Opt{analyze} 的 C264xx、C268xx 规则族 & 未初始化变量与指针、
生命周期错用、下标与指针算术越界、窄类型算术溢出 &
与标准库版本绑定，规则号会迁移；把它当规则族而非稳定接口 \\
\end{longtable}

表里两格的细则展开如下。GCC 的警告集在 \Opt{Wall}、\Opt{Wextra} 之外，
与未定义行为直接相关的专项是：\Opt{Wuninitialized} 与
\Opt{Wmaybe{\Bk}-uninitialized}（未初始化读取）、\Opt{Warray{\Bk}-bounds=2}
（静态越界，级别越高越激进）、\Opt{Wstringop{\Bk}-overflow}（写入越过对象边界）、
\Opt{Wdangling{\Bk}-pointer} 与 \Opt{Wfree{\Bk}-nonheap{\Bk}-object}
（悬垂指针与错误释放）、\Opt{Wuse{\Bk}-after{\Bk}-free}（释放后使用）、
\Opt{Wstrict{\Bk}-aliasing=2}（严格别名，级别要调到 2 才有实用价值）、
\Opt{Wshift{\Bk}-count-overflow}、\Opt{Wimplicit{\Bk}-fallthrough}、
\Opt{Wnonnull{\Bk}-compare}、\Opt{Wreturn-type} 与 \Opt{Wsequence-point}。
clang-tidy 侧除了整族打开 \Tp{bugprone-}\,*，
其中 \Tp{use-after-move}、\Tp{dangling-handle}、\Tp{integer-division}、
\Tp{undefined-memory-manipulation}、\Tp{signed-char-misuse}、
\Tp{too-small-loop-variable} 这几支与本册话题重合度最高；
再按名字补上约定族 \Tp{cppcoreguidelines{\Bk}-pro{\Bk}-type{\Bk}-}\,* 与
\Tp{cppcoreguidelines{\Bk}-pro{\Bk}-bounds{\Bk}-}\,*，
它们报的是强制转换与指针算术这两类最常写进接口层的用法。
以上名字的可用集合随版本变动，逐条确认以所用工具的检查列表为准。

\begin{bestpractice}[title={分析器的结果要进基线，不要进噪声}]
第一次全仓扫描的输出不要直接当门禁：逐条分类成\emph{真问题}、\emph{需要更多信息}、
\emph{不打算管}三类，把第三类的检查名写进配置文件，并把这份配置连同判定理由一起提交。
之后只在改动目录上开全量检查，让门禁只对新出现的编号说话。
没有这一步，第二次扫描的结果会被第一次的存量淹没，检查也就废了。
\end{bestpractice}

\section{接进构建与持续集成}\label{sec:ub03-ci}

检查档位要写进构建配置，而不是写进某个人的命令行习惯里。下面这段 \Mx{CMake} 是最小形状：
两个开关决定检查集合，默认全部关闭，本地构建与流水线跑同一份源码，差别只来自传进去的选项。
接进配置阶段之后，门禁的开关面就有了唯一出处，评审时只需看那一个缓存变量。

\begin{lstlisting}[style=cmakestyle, caption={把消毒器做成可开关的检查档位}]
# 由 CI 传 -DUB_CHECKS=ON -DUB_ASAN=ON 打开；默认关闭，本地构建不受影响
if(UB_CHECKS)
  target_compile_options(core PRIVATE -g -fno-omit-frame-pointer)
  target_link_options(core PRIVATE -g -fno-omit-frame-pointer)
  if(UB_UBSAN)
    target_compile_options(core PRIVATE -fsanitize=undefined)
    target_compile_options(core PRIVATE -fno-sanitize-recover=all)
    target_link_options(core PRIVATE -fsanitize=undefined)
  endif()
  if(UB_ASAN)
    target_compile_options(core PRIVATE -fsanitize=address)
    target_link_options(core PRIVATE -fsanitize=address)
  endif()
  target_compile_definitions(core PRIVATE UB_CONTRACT_CHECKS=1)
endif()
\end{lstlisting}

\begin{guideline}[title={分档策略：按档位决定检查集合}]
\textbf{第一档（每个 PR）}：Debug 或 \Tp{-O0}/\Tp{-Og} 加完整调试信息，
\Opt{fsanitize=undefined} 加 \Opt{fsanitize=address}，
\Opt{fno-sanitize-recover=all} 让任何一次命中直接让任务失败，
输入用固定的语料库，保证每次跑的是同一份数据。

\textbf{第二档（每日或每次合并）}：同一份源码换 \Tp{-O1}/\Tp{-O2} 再跑一遍，
检查集合不变但优化档位变化；同时补一轮静态分析，只针对改动的目录。

\textbf{第三档（发布构建）}：不挂消毒器（运行时开销与部署要求不划算），
只保留断言与库检查模式，且契约检查用独立宏控制，
保证“检查代码本身不会因为 \Tp{NDEBUG} 变成死代码后再被优化掉”。
\end{guideline}

\section{为什么“消毒器全绿”不等于没有 UB}\label{sec:ub03-green}

四个必须写进评估结论的理由：

\begin{enumerate}
\item \textbf{未覆盖的路径}。消毒器只在执行到的地方检查。
      没跑到的分支、没触发的错误处理、没进入的循环轮次，都是零信息。
      覆盖率报告能说明“跑过哪些行”，不能说明“跑过哪些值”。
\item \textbf{依赖具体数据的 UB}。有符号溢出、下标越界、未对齐访问、
      读取未初始化成员，都取决于运行时的实际数值。
      换一个输入形态就可能从合法变越界；测试语料的规模不能替代数值空间的覆盖。
\item \textbf{优化档位差异}。消毒器跑在低优化档位，发布构建跑在高优化档位，
      两者对“可达性”的推断不同：高档位会把标准认为不可能的分支删掉，
      于是低档位有检查的位置在发布二进制里根本不存在（见 第~\ref{sec:ub03-case-checkdrop} 节）。
\item \textbf{插桩本身会改变内存布局}。ASan 的红区让紧邻的越界很快暴露，
      也同时把“只有靠布局才会撞到别的对象”的那类错误盖住；
      没有红区时，同一越界可能只是安静地写坏相邻成员。
\end{enumerate}

还有一条纪律常被忽略：消毒器轮的结论只有和发布轮的\emph{输入}对齐时才有可比性。
随机 fuzz 跑出来的“零报告”不能当成发布构建安全的证据，
但同一份固定回归语料在两个档位上的行为差异，几乎总能定位到一个具体的优化推断——
这正是下一节案例的取证入口。

\section{案例：release 构建里消失的越界检查}\label{sec:ub03-case-checkdrop}

\begin{casebox}[title={检查写在那儿，却一次也没触发}]
一个环形缓冲区的写入函数把下标检查写成“溢出即越界”的形式，
Debug 加消毒器时一切正常；发布构建里越界计数从未增加，
即使测试明确喂了会越界的下标。下面是简化后的复现骨架：
\end{casebox}

\begin{lstlisting}[style=cppstyle, caption={被当成不可能发生而删掉的前提}]
struct Ring {
    int   slot[8];
    int   count;
    bool  overflowed;
};

void push(Ring* r, int i) {
    // 意图：i 为负或过大时记录一次越界，然后什么都不写
    if (i + 8 < 8) {              // 等价于 i < 0，但写成了有符号加法
        r->overflowed = true;
        return;
    }
    if (i >= 8) {                 // 上界检查
        r->overflowed = true;
        return;
    }
    r->slot[i] = r->count;        // 只有这里被消毒器的 bounds 覆盖
    r->count = r->count + 1;      // 计数自身也可能有符号溢出
}
\end{lstlisting}

问题出在第一行的判断。标准不许有符号溢出“回绕”，
于是编译器有权认为 \Tp{i + 8} 溢出这件事不可能发生；
从“不可能发生”反推出前提 \Tp{i >= INT\_MIN - 8}，
再进一步：在优化器眼里，这条加法不会改变符号，
于是 \Tp{i + 8 < 8} 被判成恒假，整段记录越界的分支被删。
Debug 档位保留了这条分支（插桩与低优化让推断变钝），
所以消毒器看到的是“检查在，而且没触发”——它没触发不是因为下标合法，
而是因为在发布档位上它压根不存在。

读法上的两条纪律：

\begin{itemize}
\item 下标一律用无符号或先做一次显式的有符号到无符号的转换，
      把比较写成标准承诺回绕的形式，让检查的可达性不依赖溢出不可能存在。
\item 越界检查必须能被\emph{同一份输入}在 Debug 加消毒器档位上打中；
      打不中的检查要当成不存在的检查。
\end{itemize}

怀疑检查被删时，不要靠读汇编猜：把计数器写进日志或返回错误码，
用 Debug 与发布两档各跑一次同一输入，比较行为差异；
再用 \Opt{fno-sanitize-recover=all} 让越界立刻中止，确认命中点在源码的哪一行。
差异存在，就是标准前提被用来反推的证据；差异不存在，则要去查另一层
（对齐、别名或生命周期，见 第~\ref{ch:ub04-lifetime} 章 与 第~\ref{ch:ub06-aliasing} 章）。

\clearpage

\part{内存与对象}
\chapter{生命周期、对象表示与 trap 表示}\label{ch:ub04-lifetime}

内存类问题在 \UB{} 谱系里占的位置最特殊：它既不像有符号溢出那样只落在一个表达式上，
也不像数据竞争那样要扯上并发模型，而是把“这块存储此刻归谁、里面装的字节算不算一个值”
这两件事压到了每一条读写上。本章把这条概念链摆直：存储期、对象生命周期、关联的存储、
放置 \Fn{new} 如何结束一个生命周期并开始另一个；然后是对象表示、值表示与 padding，
以及哪些类型的表示里允许出现 trap。最后给出标准库为这三件事准备的三个出口，
以及按字节读写这条豁免通道到底宽在什么地方。

\section{存储期、生命周期与关联的存储}\label{sec:ub04-concepts}

\subsection{四个概念各管什么}\label{sec:ub04-four-concepts}

\textbf{存储期}说的是存储本身活多久：静态存储期、线程存储期、自动存储期、
动态分配得到的存储期。它是存储的属性，不是对象的属性。

\textbf{对象生命周期}说的是一个对象在它所占的那段存储里“活着”的时间区间。
对没有自定义析构的类型，生命周期从初始化完成开始；
对有用户提供的析构的类型，从构造完成开始、到析构调用开始结束。
自动存储期到作用域末尾结束，但生命周期可以更早：显式析构，或同一块存储被交给新对象。

\textbf{关联的存储}把两者接起来：对象在整个生命周期内都有关联的存储；
指针类型除了自身占据的存储，还额外关联它所指向对象的存储。
这个概念存在的原因很实际——判空、解引用与生命周期传递讨论的都是“关联的存储”，
而不是“对象所在的字节”。

\textbf{重用的时机}：存储被释放、被交给新对象、或线程结束时，
上面这段区间就到此为止。生命周期结束到存储消失之间的那一段，
是可以合法操作的，但只能按底层字节操作。

\begin{lstlisting}[style=cppstyle, caption={生命周期结束、存储仍在：三种写法三种命运}]
struct Widget {
    int id;
    Widget(int v) : id(v) {}
    ~Widget() { id = -1; }
};

void demo(unsigned char* storage) {
    Widget* first = new (storage) Widget(1);   // 在已有存储上开始生命周期
    int keep = first->id;                      // 合法：对象此刻活着
    first->~Widget();                          // 生命周期结束，storage 仍然可用

    int again = first->id;                     // 读取不在生命周期内的对象：UB
    Widget* second = new (storage) Widget(2);  // 同一存储上开始新的生命周期

    int fresh = second->id;                    // 合法，但要经 second 这个指针读
    (void)keep; (void)again;
}
\end{lstlisting}

\subsection{放置 \Fn{new} 结束谁、开始谁}\label{sec:ub04-placement-new}

放置形式的 \Fn{new} 不分配存储，它做两件事：结束原先占据这段存储的那个对象的生命周期，
在同一段存储上开始一个新对象的生命周期。因此它的安全前提全是\emph{布局}层面的：
这段存储要足够大、对齐要满足新类型、原先的对象要么没有自定义析构、
要么已经被显式析构。三条里任何一条不成立，都是 \UB{}，
而不是“看运气的实现差异”。

新对象与旧对象的关系决定旧指针能不能继续用。同类型、非 const、
不含 const 或引用成员这类条件下，指向旧对象的指针会自动指向新对象；
类型不同、或者新对象在 const 方面变了，就必须走 \Fn{std::launder}（见 第~\ref{sec:ub04-library-escapes} 节）。

\subsection{读取不在生命周期内的对象属于哪一类}\label{sec:ub04-dead-read}

结论要说得干脆：读取一个生命周期已经结束、又还没开始的对象，属于 \UB{}，
不是 \IDB{}，也不是 \USB{}——标准没有为它准备任何答案。
UBSan 的 \Tp{vptr} 与 ASan 的 use-after-free、use-after-return
是这类问题最常见的两个观测面（口径见 第~\ref{ch:ub03-sanitizers} 章）。

两条常被引用的例外：
\begin{itemize}
\item 平凡可复制的类类型对象在生命周期结束后、存储被释放或重用之前，
      可以把这段存储当作无符号字符数组来操作，读到什么字节就是什么字节。
\item 反向同理：把一段已分配的存储当作原始存储使用时，
      只要没有以“该类型的对象”的身份去访问它，本身不构成违规。
\end{itemize}
这两条例外只赦免\emph{字节层}的访问，不赦免任何把字节当值来读的尝试。

\begin{clause}[title={标准怎么说：生命周期与对象表示}]
标准的定义式条款给的是区间而不是保证：对象在生命周期内有关联的存储，
存储期内它占据这段存储；对象表示是它所占据的存储中的字节序列，
值表示则决定这些字节如何对应到该类型的值上。
凡标准说“不在生命周期内访问”“使用了不确定的值”，
给出的都是 \UB{}；凡标准说“由实现决定”，才轮到 \IDB{} 与实现文档（\Std{basic.life}）。
\end{clause}

\section{对象表示、值表示与 padding}\label{sec:ub04-representation}

\subsection{三个词各指什么}\label{sec:ub04-three-words}

\textbf{对象表示}（object representation）：一个对象占据的那段存储里的字节序列，
也就是 \Tp{unsigned char} 数组能看到的原始内容。
\textbf{值表示}（value representation）：对象表示中\emph{用来表示该类型的值}的那部分位。
\textbf{padding bits}：对象表示里既不属于值表示、也不属于对齐要求的那部分位——
更准确地说，是为了满足对齐而被塞进来的位，标准不规定其内容。

三者的关系解释了本章后半部的所有现象：
比较运算看的是值表示，按字节复制搬运的是对象表示，
而 padding 位是对象表示中标准\emph{没有承诺}的部分。

\subsection{哪些类型允许有 trap representation}\label{sec:ub04-trap-reps}

“对象表示里的某些字节序列不对应任何值”，这样的位模式叫 trap representation。
标准不规定它的存在与否：类型的值表示、padding 位、是否有 trap，
全都落在 \IDB{} 的范围内，由实现文档说明（C 侧的表述见 \Std{6.2.6.1}）。
于是：

\begin{itemize}
\item \Tp{bool} 与 \Tp{\_Bool}：标准只保证 false 与 true 的表示是 0 与 1，
      其余字节序列是不是 trap、读出来算不算合法值，由实现决定。
      把一个非零字节塞进去再当 \Tp{bool} 读，
      轻则得到不可预期的真值，重则触发硬件陷阱，
      在多数现代实现上则表现为“同一个对象两次读出不同结果”。
\item 浮点类型：NaN 的位模式在 IEEE 754 上是合法的，
      但标准并不保证任意位序列都对应一个浮点值；把随机字节写进 \Tp{double}
      再比较大小，结果不受标准保护。把全 1 当 0 来清零是常见习惯，不是承诺。
\item 指针类型：非法地址、带标签的低比特、非零的空表示，
      都可能在解引用之前就已经越出了该类型的值域。
      这也是 UBSan 把 \Tp{null} 与 \Tp{object-size} 单列为检查的原因。
\end{itemize}

\begin{lstlisting}[style=cstyle, caption={把未经承诺的位模式写进值域很窄的类型}]
#include <string.h>

union Probe {
    unsigned char raw[sizeof(float)];
    float         value;
};

int main(void) {
    union Probe p;
    for (unsigned i = 0; i < sizeof p.raw; i = i + 1) {
        p.raw[i] = 0xFF;             /* 全 1：不保证对应任何浮点值 */
    }
    return p.value > 0.0f ? 1 : 0;   /* 非活跃成员，加上可疑位模式，两层问题 */
}
\end{lstlisting}

\begin{warning}[title={memset 清零不等于造出一个合法值}]
用 \Fn{memset} 把一块内存清成 0 再当 \Tp{double} 或 \Tp{bool} 用，
在多数平台上恰好工作；但标准没有承诺全零的字节序列一定是 0.0 的表示，
也没有承诺它一定是 \Tp{bool} 的合法值。可移植的做法是：
先按\emph{声明的类型}赋一个真实的字面量，再让这块内存参与后续读写；
需要跨越类型搬运字节时，走 \Tp{std::bit\_cast}（第~\ref{sec:ub04-bit-cast} 节）。
\end{warning}

\section{三个标准库出口}\label{sec:ub04-library-escapes}

\subsection{\Tp{std::bit\_cast}：同尺寸之间的按位重解释}\label{sec:ub04-bit-cast}

\Feat{std::bit\_cast}（\cpp{20}，\Hdr{bit}）在两个类型之间复制对象表示：
要求两侧大小相同、都是平凡可复制类型，且源是一个\emph{已经构造好}的对象（按值传入）。
它是 constexpr 的，可以在编译期做 IEEE 位模式的操作。

它解决的是“不要用 \Fn{reinterpret\_cast} 去读一个不是该类型的对象”这条弯路；
它不解决三件事：不开始任何生命周期（结果只是个值，不是对象）、
不保证目标类型的值表示合法（字节序列落在 trap 上就照原样带过去）、
也不填充 padding（源对象的 padding 位内容一并复制）。

\subsection{\Tp{std::start\_lifetime\_as}：在既有存储上开始生命周期}\label{sec:ub04-start-lifetime}

\Tp{std::start{\Bk}\_lifetime{\Bk}\_as} 与
\Tp{std::start{\Bk}\_lifetime{\Bk}\_as{\Bk}\_array}（\cpp{23}，\Hdr{memory}）
做的是隐式对象创建：给定一段已分配且足够对齐的存储，
在其中创建一个隐式生命周期的对象，把\emph{存储里现有的字节}作为它的对象表示。

它解决的是反序列化那一类场景：收到一段字节，想要一个指向结构体的合法指针，
而不必先造一个临时对象再 \Fn{memcpy} 过去。它不解决的是“值”的问题：
它只保证生命周期的开始，不跑构造函数，不初始化成员，
于是所有标量成员仍是它们当前字节所对应的那个值——含 padding，也可能含不确定值；
它也不适用于没有隐式生命周期的类型。

\subsection{\Fn{std::launder}：把指针挪回新对象}\label{sec:ub04-launder}

\Fn{std::launder}（\cpp{17}，\Hdr{utility}）的用法很窄：
放置 \Fn{new} 在同一段存储上开始的新对象，
如果不满足“旧指针自动指向新对象”的条件（类型不同、const 属性变了等），
旧指针继续用就是 \UB{}。\Fn{std::launder} 返回一个指向新对象的指针，
并把“这个地址上的对象是新的”这件事告诉编译器。

它不解决的事更值得记住：它不能让一个生命周期已结束的对象复活，
不能替代正确的对齐与大小前提，也不能把越界读变成合法读；
它要求新对象与旧对象在标准允许的关系之内，否则结果仍不被承诺。

\begin{lstlisting}[style=cppstyle, caption={三个出口各自的典型用法}]
#include <bit>
#include <memory>
#include <utility>
#include <cstdint>

struct Head {
    std::uint32_t seq;
    std::uint16_t port;
    std::uint16_t flags;
};

float bits_to_float(std::uint32_t v) {
    return std::bit_cast<float>(v);          // 同尺寸、两侧平凡可复制
}

Head* view(unsigned char* bytes) {
    return std::start_lifetime_as<Head>(bytes);  // 生命周期开始，值照旧
}

struct Node { const int tag; };
void remake(unsigned char* buf) {
    Node* old = new (buf) Node{7};
    Node* now = std::launder(new (buf) Node{9}); // const 成员变了：旧指针不可再用
    (void)old; (void)now;
}
\end{lstlisting}

\begin{compare}[title={三个出口能做什么、不能做什么}]
\Tp{std::bit\_cast} 造的是\emph{值}，不碰生命周期，也不管字节合不合法；
\Tp{std::start{\Bk}\_lifetime{\Bk}\_as} 造的是\emph{对象}，但不赋值、不构造，
现存的字节就是它的对象表示；
\Fn{std::launder} 修的是\emph{指针的身份}，前提是新对象已经在那块存储上开始了生命周期。
三者都不能把越界访问、未对齐访问或已结束的声明周期变成合法操作，
也不能替你把 padding 位填成一个确定的值。
\end{compare}

\section{按字节读写的豁免}\label{sec:ub04-byte-access}

\subsection{合法的那条路}\label{sec:ub04-byte-path}

标准给字节层留了一条明确的通道：任何对象都可以被当作一串 \Tp{unsigned char}、
\Tp{signed char} 或 \Tp{std::byte} 来访问（\Tp{std::byte} 见 \Hdr{cstddef}，\cpp{17}），
这条通道不受严格别名规则的约束，也是 \Fn{memcpy}、\Fn{memmove}
能作用于任意平凡可复制对象的原因。方向反过来时同样宽松：
把一串字节写回一个\emph{已经存在}的对象，是允许的。

要区分“允许的操作”与“保证的意义”。字节层搬运不保证任何值语义：
搬过去的字节可能落在 padding 上，可能是 trap 表示，
也可能与目标类型的对齐要求无关。可移植的序列化写法只有两类：
\begin{enumerate}
\item 逐字段写出、逐字段读回，宽度与字节序在协议里写死，
      读回时用赋值把每个成员显式初始化；
\item 固定宽度的整数组与目标类型之间用 \Tp{std::bit\_cast}，
      或先 \Fn{memcpy} 进一个已零初始化的对象。
\end{enumerate}

两种写法的差别在发送端最直观：把整个结构体交给 \Fn{memcpy} 是\emph{被赦免的操作}，
它不会坏在别名规则上，但会把 padding 位与本地的对齐、字节序一起发出去；
逐字段编码多写几行，换来的是协议里只有一个由宽度与位移定义的格式。
第~\ref{sec:ub04-case-wire-union} 节 的接收端样例给的是同一条纪律的另一半。

\subsection{padding 参与比较与哈希}\label{sec:ub04-padding-compare}

含 padding 的对象上，三类操作最容易出问题：
\begin{itemize}
\item \textbf{整对象比较}。用 \Fn{memcmp} 或按字节实现的相等判定比较两个结构体时，
      padding 位进入结果。逻辑上相等的两个对象可以判不等。
      \cpp{20} 的合成员 \Fn{operator==} 逐成员比较，绕开了 padding，
      但绕不开“成员本身未初始化”。
\item \textbf{按字节算哈希}。把对象表示直接喂给哈希函数，
      padding 内容会让同一逻辑值散落到不同桶里，缓存与去重全部失效。
\item \textbf{作为键或写进流}。未初始化的 padding 参与上述任一路径，
      读到的就是不确定值；非平凡可复制类型上这类访问是 \UB{}，
      平凡类型上则是“值不确定”的标准说法。
\end{itemize}

\begin{bestpractice}[title={按字节读写守则}]
一、要发出去的一律逐字段编码，别把结构体当协议；
二、收字节的一侧先把目标对象零初始化（\Tp{Rec r{};}），再逐字段赋值，
或者用 \Tp{std::start{\Bk}\_lifetime{\Bk}\_as} 之前先把存储清零；
三、参与比较与哈希的对象必须保证 padding 已被写干净，
做法是零初始化后赋值，或给类型加显式的序列化函数；
四、要用 \Tp{std::bit\_cast} 就限定在固定宽度整数与同尺寸类型之间，
并在旁边注明两端的目标布局。
\end{bestpractice}

\section{union 的活跃成员规则}\label{sec:ub04-union-active-member}

联合体同一时刻只有一个成员是活跃的：向某个成员写入使它变活跃，
其余成员的生命周期结束。读取非活跃成员，标准层面读作“访问一个生命周期已结束的对象”，
落在 \UB{} 一侧。C++ 与 C 在这一点上的\emph{工程口径}不同：

\begin{itemize}
\item \CPP{}：读取非活跃成员历来不被标准正面许可；
      对成员都是平凡类型的联合体，实现普遍支持类型双关这一写法，
      但它是实现扩展而非标准承诺。匿名联合体与含非平凡成员的联合体另有构造与析构约束。
\item \CL{}：结构体共享公共前缀（common sequence）有明确条款，
      联合体侧的类型双关同样依赖实现扩展；
      严格别名规则对 \Tp{unsigned char} 与联合体本身的例外要分开读。
\end{itemize}

关于 \cpp{11} 前后的口径变化：措辞从“最后一次写入者胜”的说明性段落，
逐步改成用生命周期与活跃成员来定义，近年草案还在动；
具体到某个写法是否被当前标准许可，以所用语言的当前草案与各实现文档为准。
需要跨类型看字节时，本章给出的三个出口比联合体更靠得住。

\begin{note}[title={联合体什么时候才是对的工具}]
联合体适合表示“同一块存储在不同时刻按不同解释使用”的\emph{ tagged} 语义
（配一个判别标签、手工管理活跃成员），或者在实现文档明确支持的前提下做定长视图。
如果目的只是把 \Tp{std::uint32\_t} 看成 \Tp{float}，
\Tp{std::bit\_cast} 是无前提、无生命周期负担的写法。
\end{note}

\section{初始化三档与未初始化读取}\label{sec:ub04-initialization}

未初始化读取之所以高发，多数时候是因为选错了初始化形式。三档的保证强度不一样：

\begin{longtable}{@{}>{\raggedright\arraybackslash}p{0.15\textwidth} >{\raggedright\arraybackslash}p{0.23\textwidth} >{\raggedright\arraybackslash}p{0.24\textwidth} >{\raggedright\arraybackslash}p{0.16\textwidth}@{}}
\caption{三档初始化给出的保证，与未初始化读取的行为归属}\label{tab:ub04-init-kinds}\\
\toprule
档位 & 触发写法与内置类型效果 & 类类型与成员 & 读取后果的归属 \\
\midrule
\endfirsthead
\toprule
档位 & 触发写法与内置类型效果 & 类类型与成员 & 读取后果的归属 \\
\midrule
\endhead
\bottomrule
\endfoot
零初始化 & 每个字节按数值 0 写；指针得空指针，
浮点得 +0.0，\Tp{bool} 得 false；静态与线程存储期对象必然先经过这一档 &
递归到每个成员与基类子对象，同样置零 &
无 trap 风险；读到的是确定的 0 \\
值初始化 & \Tp{T v{};}、\Tp{T v;}（\cpp{11} 起带默认成员初始化者）
与 \Fn{new} \Tp{T()}；标量落到零初始化 &
无用户默认构造时先零初始化再按需默认初始化；
有用户默认构造时调用它；
聚合按成员递归（\cpp{20} 起措辞与早先有别，细节以当前文档为准） &
成员若被构造函数接管则确定；
被漏掉的成员退回默认初始化的问题 \\
默认初始化 & \Tp{T v;} 对内置类型\emph{什么都不做}，值是 indeterminate；
自动存储期、堆上裸 \Fn{new} \Tp{int} 都走这一档 &
调用默认构造；未写进初始化列表与默认成员初始化的成员保持不确定 &
读取不确定值：非平凡可复制类型上是 \UB{}；
平凡类型上按“值不确定”处理，多数运算随之失去意义 \\
\end{longtable}

\begin{guideline}[title={三档选择}]
拿不准就写花括号：局部变量一律 \Tp{T v{};}，
结构体一律聚合初始化把每个成员点名；
只有在“性能关键且能证明 padding 无害”的场合才允许默认初始化，
并且要在旁边写清谁负责填这个洞。
未初始化读取的代价不在调试档，而在发布档：
不确定值会被优化器当成“任何值都可能”来反推前提，
于是分支被删、边界检查消失（第~\ref{ch:ub03-sanitizers} 章 的案例是同一机制）。
\end{guideline}

\section{案例：把网络字节解释成本地结构体}\label{sec:ub04-case-wire-union}

\begin{casebox}[title={同一份报文，两次运行给出不同结果}]
服务收到 12 字节报文，代码把字节灌进联合体，再按本地结构体打印与比较：
\end{casebox}

\begin{lstlisting}[style=cppstyle, caption={踩坑写法与其可移植替身}]
#include <cstdint>
#include <cstring>

struct Msg {                 // 本地布局：seq 前有 4 字节 padding
    std::uint32_t seq;
    std::uint16_t port;
    std::uint16_t kind;
    std::uint64_t stamp;
};

union Wire {                 // 有问题的写法
    unsigned char raw[sizeof(Msg)];
    Msg           msg;
};

bool same_a(const unsigned char* in, const Msg* expect) {
    Wire w;                                  // 默认初始化：raw 内容不确定
    std::memcpy(w.raw, in, sizeof w.raw);    // 灌字节，msg 从未被构造
    if (w.msg.stamp != expect->stamp) {      // 读非活跃成员，padding 参与比较
        return false;
    }
    return w.msg.seq == expect->seq;
}

bool same_b(const unsigned char* in, const Msg* expect) {  // 替身：逐字段解析
    Msg cur{};                               // 零初始化，padding 被写干净
    cur.seq = (std::uint32_t(in[0]) << 24) | (std::uint32_t(in[1]) << 16)
            | (std::uint32_t(in[2]) << 8) | std::uint32_t(in[3]);
    cur.port = (std::uint16_t(in[4]) << 8) | std::uint16_t(in[5]);
    cur.kind = (std::uint16_t(in[6]) << 8) | std::uint16_t(in[7]);
    cur.stamp = 0;
    for (int i = 0; i < 4; i = i + 1) {
        cur.stamp = (cur.stamp << 8) | std::uint64_t(in[8 + i]);
    }
    if (cur.seq != expect->seq) { return false; }
    if (cur.port != expect->port) { return false; }
    if (cur.kind != expect->kind) { return false; }
    return cur.stamp == expect->stamp;
}
\end{lstlisting}

\textbf{为什么两次运行不一样}。三处叠加：
其一，\Tp{Msg} 为了对齐 \Tp{stamp} 在 \Tp{port}、\Tp{kind} 之后留下 padding 位，
这些位在对象表示里存在、在值表示里没有承诺，报文里对应的字节是几全看发送端布局；
其二，\Tp{Wire w;} 是默认初始化，\Tp{raw} 里除了随后被 \Fn{memcpy} 覆盖的那几个字节，
其余（含填充区）保持不确定，第二次运行栈上残留不同，比较结果就不同；
其三，本地 \Tp{Msg} 的字段宽度、对齐与字节序和线上格式没有任何对应关系，
同一份字节在 32 位与 64 位目标上会解析成不同的字段。

\textbf{可移植写法}。两条路，选一条并坚持：
（一）像 \Tp{same\_b} 那样逐字段解析，宽度与字节序写进协议，
先零初始化再赋值，保证 padding 干净、每个成员都在生命周期与值域内；
（二）如果协议本身就是本地布局的镜像（同机器、同编译器、同编译选项），
先把目标对象零初始化，再 \Fn{memcpy} 进去，
或 \cpp{23} 下把存储清零后用 \Tp{std::start{\Bk}\_lifetime{\Bk}\_as} 取指针。
这条路只在“两端二进制格式已被明确约束”的前提下成立，
要固定宽度整数、显式填充，并禁用会改变布局的打包属性混用。
跨类型看单个字段时，用 \Tp{std::bit\_cast} 而不是联合体。

\clearpage

\chapter{指针的算术、比较与 provenance}\label{ch:ub05-pointer}

前面几章的注意力都在对象上：第~\ref{ch:ub04-lifetime} 章 问对象还在不在，第~\ref{ch:ub03-sanitizers} 章
问用什么手段能把它抓住，第~\ref{ch:ub02-inevitable} 章 问哪些坑本来就躲不开。这一章把注意力移到
指针值本身。一个指针值同时承担三件事：它是一句可以在对象内部平移的坐标，是一个可以参与比较
的值，也是一份携带来源信息的句柄。前两件事标准给了明确边界，第三件事只给了一半承诺，剩下的
那一半要由实现文档补齐。本章只谈语言契约；实现如何把指针放进寄存器、如何做区间检查、如何在
影子内存里重建来源信息，属于编译器与链接器那一册的主题，这里不展开。

\section{算术的合法域：对象之内，含尾后}\label{sec:ub05-domain}

\Tp{p + n}、\Tp{p - n}、\Tp{q - p} 这三类运算，标准允许的范围是\emph{同一个数组对象}的元素
序列，并且把尾后位置（one past the end）也算进这个范围。这里的数组按最小的包含数组对象
（smallest containment array）理解：多维数组按行优先展开后的一维视图参与运算，长度为
\Tp{N} 的数组合法下标是 \Tp{0} 到 \Tp{N}，其中 \Tp{N} 只能用来比较、做减法或往回退，
不能解引用。

非数组对象（一个单独的 \Tp{int}、一个类的标量成员）也享有同一条规则的空版本：合法范围只有
它自己和紧跟其后的那一个位置，等价于把它当作长度为 1 的数组。

\begin{clause}[title={标准怎么说：指针算术与差值}]
指针加减整数、指针相减，都以“两个指针落在同一个数组对象（含尾后）之内”为前提；落在范围之外
的算术结果是 \UB{}（\Std{expr.add}）。两个指针相减的结果类型是 \Tp{std::ptrdiff\_t}，
并且要求实际差值能被该类型表示，否则同样是 \UB{}。尾后指针可以参与算术与比较，
但对它解引用不在任何承诺之内。
\end{clause}

常见误判是把内存布局当成合法域。两个数组在链接顺序上恰好相邻，甚至地址上无缝相接，也不改变
“它们是两个对象”这件事；把 \Tp{a} 的尾后指针当作进入 \Tp{b} 的门票，是标准的 \UB{}，不是
\IDB{}。消毒器与静态分析在这一点上的判断通常与标准一致。

\begin{lstlisting}[style=cppstyle, caption={合法域内与合法域外的分界}]
#include <cstddef>

int sum_span(int* first, int* last) {
    int total = 0;
    // [first, last)：last 允许是尾后指针，用作哨兵完全良构
    for (int* p = first; p != last; ++p) { total += *p; }
    return total;
}

long diff_maybe_ok(int* a, int* b) {
    // 同一数组（含尾后）之内：良构；跨对象：UB，与地址是否连续无关
    // 差值还必须能放进 std::ptrdiff_t，放不进去同样是 UB
    long d = static_cast<long>(b - a);   // 窄化是另一处隐患，见 5.1.2 节
    return d;
}

int* walk_off_the_edge(int* a, int n) {
    // 尾后再往外一步：a + n + 1 的构造动作本身已经出事
    return a + n + 1;
}
\end{lstlisting}

\subsection{尾后指针能做什么、不能做什么}\label{sec:ub05-past}

“尾后指针可减不可解引用”这句话要拆细，否则容易写成两种相反的过度防御。精确的表述是四条：

\begin{itemize}
\item 尾后指针\emph{可以}构造，前提是算术本身没越界（\Tp{a + N}，\Tp{N} 为数组长度）；
\item 尾后指针\emph{可以}参与相等与关系比较、可以参与减法、可以从它往回退；
\item 尾后指针\emph{不可}解引用，这是无条件的一票否决；
\item 尾后指针\emph{不可}再往前加：\Tp{a + N + 1} 的构造动作本身就出事，不必等到解引用。
\end{itemize}

第四条最容易被忽略：出事的时刻是表达式求值，而不是后续那次解引用。以长度为前提写
\Tp{a + i}，只要 \Tp{i} 越界就已经进入 \UB{}；把这样一个值存下来，留到以后才解引用，
是同一处缺陷的两种表现，而消毒器报的是求值那一步。审查时若只顺着解引用往回找，
就会把根因留在中间那几次算术上。多维数组还有一处陷阱：用 \Tp{mat[3][0]} 这类以行指针
进入下一行的写法，当 \Tp{mat} 只有两行时等价于跨对象算术，在标准里站不住，编译器未必反对。

\subsection{\Tp{std::ptrdiff\_t} 的角色}\label{sec:ub05-diff}

指针差值的类型是有符号整型 \Tp{std::ptrdiff\_t}（\Hdr{cstddef}），它的取值范围由实现决定，
标准只要求它能表示同一对象内部两个指针的差。实践中要盯住两点。第一，把差值塞进
\Tp{int} 或 \Tp{size\_t} 之类的窄类型是常见缺陷：前者窄化溢出，后者把符号丢了，
于是 \Tp{p - q} 在 \Tp{p} 小于 \Tp{q} 时变成一个巨大的正数，越界判断反而放行。
第二，在 32 位地址空间里问题往往相反：差值本身能表示，但把它当作长度传给接受
\Tp{std::size\_t} 的接口时，符号被重新解释成大数，检查全部绕过。跨平台的正确写法是
一律用 \Tp{std::ptrdiff\_t} 承载差值，转成无符号之前先判断符号。

\begin{note}[title={区间检查的标准写法}]
先算偏移、再判断、最后才做算术：\Tp{off < 0 || off >= n} 用整数比较完成，
不要写成先比较 \Tp{p < first || p >= last} 再在外面继续对 \Tp{p} 加减。
后者的问题不在比较本身，而在为算边界所做的那些指针加减可能已经越界。
\end{note}

\section{比较：什么时候 \Tp{<} 才有意义}\label{sec:ub05-compare}

相等性比较与顺序比较在标准里不是一回事。任意两个同类型指针都可以问“相等吗”；顺序比较
（\Tp{<}、\Tp{<=}、\Tp{>}、\Tp{>=}）只在两者落在同一个数组对象（含尾后）时才有承诺。
\CPP{} 的措辞是不满足该前提时结果为未指定值（\USB{}），\CL{} 则把不相干对象之间的关系运算
列入未定义行为，附录 J 的清单里有它。这个语言差别值得在意：同一份代码在 \CL{} 里可能连
“随便给个结果”都算不上，在 \CPP{} 里至少实现可以给出稳定的排序。

\begin{longtable}{@{}>{\raggedright\arraybackslash}p{0.20\textwidth} >{\raggedright\arraybackslash}p{0.16\textwidth} >{\raggedright\arraybackslash}p{0.38\textwidth}@{}}
\caption{指针比较的可依赖情形与兜底手段}\label{tab:ub05-compare}\\
\toprule
情形 & 可用的运算 & 类别与说明 \\
\midrule
\endfirsthead
\toprule
情形 & 可用的运算 & 类别与说明 \\
\midrule
\endhead
\bottomrule
\endfoot
同一数组的两个元素，含尾后 & \Tp{== != < <= > >=} & 全部良定义，顺序即元素序（\Std{expr.rel}） \\
同一对象的两个视图 & \Tp{==} 为真 & 良定义；指向同一对象的不同视图相等，顺序仍按数组前提判断 \\
不同对象的两个指针 & 只有 \Tp{==} \Tp{!=} & 相等性良定义（不同对象必不相等）；顺序在 \CPP{} 是 \USB{}，在 \CL{} 是 \UB{} \\
空指针值参与的比较 & \Tp{==} \Tp{!=} & 与任何指向实际对象的指针都不相等；两个空指针相等；顺序同上 \\
兜底：\Tp{std::less} & 任意两个同类型指针 & \cpp{14} 起给出全序，且与 \Tp{==} 一致；\CPP{} 侧唯一可依赖的跨对象排序 \\
\end{longtable}

\Tp{std::less}（\Hdr{functional}）在 \cpp{14} 之后被明确要求对指针给出全序，即使被比较的两个
指针不属于同一个数组；它不承诺这个全序就是地址序，也不承诺不同数组之间的相对顺序有任何语义。
要用指针给容器排序、要判断两个地址区间是否重叠，就用它，别指望裸 \Tp{<}。三路比较运算符在
\cpp{20} 里对不相关指针返回 \Tp{unordered} 这一档，与 \Tp{std::less} 的全序正好对照：
前者是如实承认没有关系，后者是造一个关系出来，选哪个取决于你要的是正确性还是排序稳定性。

\section{provenance：指针转整数再转回来}\label{sec:ub05-prov}

这是本章唯一一处“标准给了整数化、没给回转”的地方，契约的边界要背清楚。

\Tp{std::uintptr\_t}（\Hdr{cstdint}）是可选类型：实现可以不提供它。如果提供，它是无符号整型，
且\emph{只有}在转换后的整数能表示原值时，往返才保证结果与原指针比较相等。请注意这句话的射程：
它保证的是相等性，不是“转回来之后可以像原来那样用”。\CL{} 侧更直白：指针转整数与整数转指针
都写在转换规则里（\Clause{6.3.2.3}），前者只在足以容纳原值的整数类型上不丢信息，后者被明确
标为实现定义；指向不完整类型的指针转换后还可能失去有效性。

为什么不能假定往返保留了优化所需的信息？因为指针值的来源不止地址一位。现实里至少有四类信息
挂在指针旁边，而不在指针的整数值里：

\begin{itemize}
\item 对齐位与低位标签：一些分配器与运行时把小整数塞进低位，整数值变了含义也变了；
\item 有效区间：实现需要知道这个指针有权访问哪一块内存，区间信息不随整数移动，
      边界检查与能力指针一类机制正是靠它判定合法性；
\item 对象类型与生命周期元数据：调试运行时与内存检查器按创建点登记这些信息；
\item 可访问性：跨模块、跨地址空间的指针在整数化后可能只剩下一个无用的数字。
\end{itemize}

于是有本册反复出现的那句判断：整数化是往实现文档里存钱。你写下的往返代码能不能跑，
取决于实现文档是否承认这笔存款，标准只承认到“比较相等”为止。

\subsection{\Tp{std::to\_address} 与 \Tp{std::pointer\_traits}}\label{sec:ub05-toaddr}

\cpp{20} 引入 \Tp{std::to\_address}（\Hdr{memory}），用途是把“类指针的东西”统一成裸指针：
对 \Tp{T*} 直接返回原值，对自定义智能指针走 \Tp{std::pointer{\Bk}\_traits} 的
\Tp{to\_address}。它是取用裸指针的正规手段，但它做的是视图转换，不是来源重建：
从一个已被判定无效的整数回转出的指针，用 \Tp{std::to\_address} 再走一遍并不会让它重新有效。
读别人代码时，看到 \Tp{to\_address} 就默认它只取裸指针，不要读出“这里做了合法性恢复”的意思。

\Tp{std::pointer{\Bk}\_traits}（\cpp{11} 起，同在 \Hdr{memory}）给出自定义指针类型的
接口清单：\Tp{element{\Bk}\_type}、\Tp{difference{\Bk}\_type}、\Tp{rebind}、
\Tp{pointer{\Bk}\_to} 与 \Tp{to{\Bk}\_address}。它是接口约定，
而不是把任意整数变成合法指针的通道。写自定义指针类型时把它当契约清单逐条实现；
用现成智能指针时，把 \Tp{to{\Bk}\_address} 当作唯一的裸指针出口。

\begin{lstlisting}[style=cstyle, caption={整数往返与低位标签的最小形态}]
#include <stdint.h>
#include <stdlib.h>

/* 指针 -> 整数 -> 指针：标准只保证“能表示则比较相等”，
   回转后能否继续解引用是实现文档的事（C 6.3.2.3：整数到指针为实现定义） */
void* round_trip(void* p) {
    uintptr_t bits = (uintptr_t)p;
    return (void*)bits;          /* 与 p 比较相等，但“有权访问什么”不在保证内 */
}

/* 低位标签：前提是实现真的只在 8 的倍数地址上分配 */
int   tag_of(void* p) { return (int)((uintptr_t)p % 8u); }
void* untag(void* p)  { return (void*)((uintptr_t)p - (uintptr_t)p % 8u); }

/* untag 的结果拿去解引用，就是那笔存进实现文档的钱：
   分配器换实现、对齐改变、或实现额外占用地址位时，这里先出事 */
long use_untagged(void* tagged) {
    long* q = (long*)untag(tagged);
    return *q;
}
\end{lstlisting}

\section{三种打包指针模式的取舍}\label{sec:ub05-packed}

指针要压缩、要打标签、要跨进程搬移，工程里有三种典型做法，它们的标准立场完全不同。
把三者混为一谈是本章最容易犯的错误。

\begin{longtable}{@{}>{\raggedright\arraybackslash}p{0.16\textwidth} >{\raggedright\arraybackslash}p{0.18\textwidth} >{\raggedright\arraybackslash}p{0.38\textwidth}@{}}
\caption{tagged pointer、压缩指针与 arena 偏移量在标准里的位置}\label{tab:ub05-packed}\\
\toprule
模式 & 标准立场 & 现实前提与代价 \\
\midrule
\endfirsthead
\toprule
模式 & 标准立场 & 现实前提与代价 \\
\midrule
\endhead
\bottomrule
\endfoot
tagged pointer & 站不住，纯实现文档范畴 & 依赖低位为空（即依赖对齐分配）与整数回转两步都被实现承认；对齐一变、分配器一换即失效，标签位还挤占地址位宽 \\
压缩指针（窄句柄） & 站不住，纯实现文档范畴 & 靠全局基址加偏移的句柄约定，每次访问做一次整数到指针的往返；跨模块、多基址、地址随机化都要额外规则兜住 \\
arena 偏移量 & 在标准里站得住 & 全程只在整数域算偏移，需要指针时对同一块缓冲区做数组下标，合法域清晰且不碰整数回转；代价是每次访问多一次基址加法 \\
\end{longtable}

结论值得记住：三者里只有 arena 偏移量能靠语言契约本身活下来，因为它把“变成指针”这件事推迟
到必须用的那一次，而那一次构造的合法域由一个明确的数组对象提供。tagged pointer 与压缩指针
不是不能用，而是用了就要接受“这段代码的可移植性等于你的实现文档质量”，并且要在评审清单里
显式记下依赖了哪些前提：分配粒度、对齐保证、地址空间模型、以及工具链是否会额外占用地址位。
带标签的写法还有一个隐蔽后果：把带标签的值原样交给按指针解释的库函数，比算术越界更难诊断。

\subsection{\Tp{std::atomic} 里存整数指针}\label{sec:ub05-atomic}

\Tp{std::atomic} 可以直接参数化为指针类型，也可以存 \Tp{std::uintptr\_t} 再自行解释。
前者是良构的：原子操作的对象是指针，读出的指针带着它的来源。后者多了一步整数到指针的转换，
于是本章前面那条边界原样适用：原子性只保证读出的是一个完整的值，并不保证读出的值是一个
有权访问的指针。比较并交换循环里按整数位模式比较预期指针与新指针，在带标签或有保留位的
实现上尤其容易得出反直觉结果：同一地址的两种整数表示可能并不相等。

并发还有第二层约束，不在本章射程内但必须留一句：整数化的指针被另一线程读出时，
它所指向的对象可能已经结束生命周期，此时出事的是访问本身而不是原子操作；用引用计数、
hazard pointer、读侧无锁结构等手法把“对象存活”与“指针可读”分开保证，
属于内存模型的内容，见 第~\ref{ch:ub10-memory-model} 章。

\section{空指针常量、\Tp{std::nullptr\_t} 与未初始化指针}\label{sec:ub05-null}

三个概念要分层，混用会带来一连串误判。

空指针常量是写代码时的字面量：\Tp{nullptr}、\Tp{0}、加圆括号的 \Tp{(void*)0}，以及 \Tp{NULL}。
空指针值是这些常量转换成指针类型后得到的值，它是否就是地址 \Tp{0} 属于实现定义；
两者在多数平台上重合，但不要靠这个重合写任何整数比较。空指针的解引用是 \UB{}，无条件。

\Tp{std::nullptr\_t}（\Hdr{cstddef}）是 \Tp{nullptr} 的类型：它是唯一对象的类型，可以转换成
任何指针类型（都得到空指针值），但不能转换成整型或枚举。这一点决定了重载行为：
\Tp{nullptr} 在 \Tp{void*} 与 \Tp{long} 两个重载之间只会选中前者，而 \Tp{0} 或 \Tp{NULL}
选中后者，于是产生完全不同的调用。\CL{} 允许 \Tp{NULL} 是实现定义的额外形式，
\CPP{} 则禁止把 \Tp{NULL} 定义成非整型类型（例如带 \Tp{(void*)} 的形态），
因此在 \CPP{} 代码里把 \Tp{NULL} 换成 \Tp{nullptr} 是安全的单向操作，反过来不成立。
标准明确规定 \Tp{sizeof(std::nullptr\_t)} 是 \Tp{1}，别按指针大小猜它。

\subsection{未初始化指针就用归哪一类}\label{sec:ub05-uninit}

这类问题很容易被归到“空指针解引用”，其实它是两处独立的 \UB{}，而且第一处更早。读一个
未初始化对象的值是 \UB{}（除 \Tp{unsigned char} 一类可复制字节序列的特例外），
指针类型不在例外里：\Tp{T* p;} 之后取值、\Tp{p + 1}、甚至把 \Tp{p} 拷给另一个指针，
都在取值那一刻出事。随后如果它恰好落在一个看似合理的地址上并被解引用，那是同一缺陷的
第二处表现。类别归属上，它属于 第~\ref{ch:ub04-lifetime} 章 那条“对象已存在但值不确定”的分支，
而不是本节的空指针分支；工具侧对应 第~\ref{ch:ub03-sanitizers} 章 里的未初始化内存检查，
以及编译器的 \Warn{uninitialized}、\Warn{maybe{\Bk}-uninitialized} 一族，两者报的是同一件事的两面。

实践判断：把指针的声明与初始化写成一句（\Tp{T* p = first;}），是消灭这类缺陷成本最低的做法；
确实必须分开时，给一个明确的空初值，并让区间检查使用相等比较，
而不是依赖“大概会先赋值”的路径顺序。

\begin{warning}[title={\Fn{realloc} 之后旧指针与新指针混用}]
\Fn{realloc} 的返回指针是唯一有效的那个入口。若发生搬移，旧对象已被释放，
再用旧指针读写就是生命周期结束后的解引用，属 \UB{}：\CL{} 侧由 \Fn{realloc} 的规定直接写明
旧空间被释放，\CPP{} 侧由 \Std{basic.stc.dynamic.deallocation} 认定旧指针值失效。
即使原地扩容、新旧指针数值相等，旧对象的有效范围也已按新大小重新划定，
用旧指针访问超出原对象的部分同样越界。最隐蔽的是失败路径：\Fn{realloc} 返回空指针时
原对象仍然有效，此时把返回值直接赋回原变量，就把有效指针换成了空指针，
后续解引用出事，而新旧混用在这一刻反而看不出异常。正确写法是用临时变量接收返回值，
判断成功后再覆盖原变量，并从覆盖那一刻起只认这一个入口。
\end{warning}

\begin{bestpractice}[title={需要跨对象算术时改用整数偏移量的接口设计}]
接口层不要用“一对指针”表达跨越多个对象的区间，改接受“基址、偏移、长度”，
偏移与长度用 \Tp{std::ptrdiff\_t} 或恰当的无符号类型承载。这样做有三个好处：
合法域可以写成一行整数比较；区间可以被序列化、被搬移、被并行复制而不动指针；
换实现（换分配器、换地址空间布局、从 32 位迁到 64 位）时，失效的是那段“基址到指针”
的局部代码，而不是所有做过差值的调用点。需要跨对象建立顺序时，
把顺序明确交给 \Tp{std::less} 或自定的序号表，不要在算术里藏假设。
\end{bestpractice}

\begin{guideline}[title={指针契约的一句话版本}]
算术只在对象内部做，顺序只在同数组内比；跨对象排序交给 \Tp{std::less}；
整数化当作往实现文档里存钱；指针能不能用由“对象还在不在”决定，而不是由“地址对不对”决定。
\end{guideline}

\section{可复现案例：内存池用指针差做索引}\label{sec:ub05-case}

\begin{casebox}[title={池换实现后，指针差索引整体失效}]
一个内存池把槽位编号算成“当前槽指针减去池基址”，用它把句柄映射回池内部。老实现把所有槽
一次性分配在同一块连续缓冲区里，差值恰好等于槽号，代码跑了多年。后来池改成按页分段的
free-list 实现，各段之间不再连续：\Tp{p - base} 落在两个数组之外，算术本身成为 \UB{}，
症状表现为句柄编号偶尔巨大、偶尔为负，而检查逻辑反而把越界值当合法槽号放行。
同时，段化之后每段大小也不再固定，即使差值可表示也不再对应槽号。
\end{casebox}

\begin{lstlisting}[style=cppstyle, caption={从指针差索引改写为显式基址加偏移}]
#include <cstddef>
#include <new>

struct Slot { void* link; char payload[24]; };

// 旧写法：依赖“所有槽落在同一数组里”，段化实现后算术越界
struct BadPool {
    Slot* base;
    std::size_t index_of(Slot* p) const {
        return static_cast<std::size_t>(p - base);      // 跨对象差值：UB
    }
};

// 新写法：句柄是偏移量，指针只在真正需要访问时从缓冲区基址构造
struct Pool {
    char*        region;              // 一整块缓冲区，合法域由它给出
    std::size_t  bytes;
    static constexpr std::size_t kSlotBytes = sizeof(Slot);  // 偏移按它对齐

    bool in_range(std::ptrdiff_t off) const {
        return off >= 0
            && static_cast<std::size_t>(off) % kSlotBytes == 0
            && static_cast<std::size_t>(off) + kSlotBytes <= bytes;
    }
    Slot* slot_at(std::ptrdiff_t off) {
        if (!in_range(off)) { return nullptr; }        // 先整数判断，失败给空指针
        void* where = region + off;                    // 缓冲区内的数组算术，合法
        return static_cast<Slot*>(::operator new(sizeof(Slot), where));
    }
};
\end{lstlisting}

放置形式的 \Tp{operator new} 只创建对象、不做分配；若槽位已在池初始化时构造好，
这里应改为直接使用已存在的对象，或用 \Tp{std::start{\Bk}\_lifetime{\Bk}\_as}（\cpp{23}）
让生命周期从这块存储开始，否则剩下的风险是生命周期而不是算术。
两种写法还都要求偏移落点的地址满足 \Tp{Slot} 的对齐，这一点由 \Tp{kSlotBytes}
与缓冲区基址的对齐共同保证，见 第~\ref{ch:ub07-alignment} 章。

改写要点有四条，逐条对应一处契约：句柄用 \Tp{std::ptrdiff\_t} 表达，把跨对象的差值
换成同一缓冲区内的整数加法；合法性判断改成整数区间与整数取模比较，不再依赖指针关系运算；
指针只在真正需要访问的那一层构造，构造后立即使用，不保存、不序列化；
对象创建与算术分开处理，避免把生命周期缺陷伪装成算术缺陷。
附录 \ref{ap:ubC-cases} 收录了这个案例的最小复现形式，类别标签见附录 \ref{ap:ubA-kinds}，
涉及的开关集中在附录 \ref{ap:ubB-switches}。

\begin{note}[title={与相邻章节的接口}]
本章只讨论指针值的算术、比较与整数化。指针所指对象的存在性见 第~\ref{ch:ub04-lifetime} 章；
用错误类型解释这块存储见 第~\ref{ch:ub06-aliasing} 章；缓冲区里按错误宽度读写见
第~\ref{ch:ub07-alignment} 章；指针参与求值时的副作用顺序见 第~\ref{ch:ub08-evaluation} 章。
\end{note}

\chapter{严格别名与类型双关}\label{ch:ub06-aliasing}

上一章问的是指针能不能碰这块内存，这一章问的是用什么类型去碰。两件事经常被混作一谈，
因此症状也经常被误判：越界访问有确定的地址线索，别名违规却没有——它表现为“同一份代码在
某个档位下算错”，于是被当成编译器的 bug 上报。本章给出别名规则的名单形式、四条合法的
类型双关路径、编译器依赖别名信息的理由，以及关掉它要付出什么。
实现内部如何构造别名图放在末尾单独划界，不混进语言契约里讲。

\section{别名规则的枚举形式}\label{sec:ub06-rule}

别名规则（type-based alias analysis 的语言依据）长得像一条禁令，实际上是一份\emph{白名单}：
它列出用哪几种类型的 glvalue 去碰同一块存储是允许的，落在名单之外就是 \UB{}。
把它按名单来背，比按“禁止强转”来背要准确得多。

\CPP{} 侧的名单（\Std{basic.lval}）大致是这几族。对某个存活的对象，可以通过下列类型的
glvalue 访问它：

\begin{itemize}
\item 与该对象类型相同的类型，允许差一层 cv 限定；如果最派生对象是多态类型，
      判据取它的动态类型而不是静态类型；
\item cv \Tp{void}，或指向 cv \Tp{void}（允许差一层 cv）的指针类型——这是 \Tp{void*}
      通用指针的语言依据；
\item 该对象类型的基类类型（允许差一层 cv）——这是“用基类引用碰派生对象里的基类子对象”的依据；
      反方向不在名单里：把一块只装着基类对象的存储当派生类用，是名单外的访问；
\item \Tp{std::byte}、\Tp{char}、\Tp{signed char}、\Tp{unsigned char}——本册称为字符豁免，
      它是唯一不以类型相符为前提的一族；
\item 与上述类型相差一层引用的类型。
\end{itemize}

\CL{} 侧的写法在 \Clause{6.5} 的第 7 段：允许通过相容类型的 lvalue、允许指针与指向
\Tp{void} 的指针互换、允许聚合或联合类型里存在公共初始序列的成员、允许字符类型碰任何对象。
两套名单的家族结构相同，细节（引用族、成员族）不同；后果相同：名单外的访问是 \UB{}，
不是 \USB{}，也不是 \IDB{}。这个归类很关键：编译器不需要为它保留任何行为，
因此“我关掉优化它就不坏了”不属于修复，只属于掩盖。

另外两套相邻机制要划清界限，它们常被混作一谈。\CL{} 的 \Tp{restrict}（\Clause{6.7.3}）是
程序员向编译器作出的额外承诺，违反它是 \UB{}，但它与别名名单互相独立：
加了 \Tp{restrict} 不代表类型可以乱来，反过来也一样。\Tp{volatile} 的访问保证与别名规则叠加时
更严格——通过类型不符的 glvalue 访问 \Tp{volatile} 对象，既破坏别名规则也绕过 volatile 语义，
见 第~\ref{ch:ub08-evaluation} 章。

\begin{clause}[title={标准怎么说：白名单与它的两处留白}]
标准不承诺“地址相同就意味着可以互相解释”。它承诺的是：只有名单内类型的 glvalue 才有资格
访问那个对象。留白有两处。其一，字符族给出的豁免只覆盖读写字节表示，
不覆盖“把这段字节继续当成某个 \Tp{T} 使用”；其二，对象表示里的填充位（padding bits）
在双关时取什么值是未指定的，别指望它等于源对象的填充位。
\end{clause}

\section{字符豁免与 \Tp{std::byte}：解决了什么，没解决什么}\label{sec:ub06-byte}

\Tp{std::byte}（\cpp{17}，\Hdr{cstddef}）的定义是 \Tp{enum class byte : unsigned char}
加一组位运算。它进入名单的意义有两个：一个语义上的，一个语法上的。语义上，它告诉读者
“这里操作的是字节，不是字符”，避免把内存内容误当文本；语法上，它和 \Tp{char} 一样
有资格碰任何对象，但不像 \Tp{char} 那样参与字符与数值的隐式转换，
于是 \Tp{std::byte} 数组不会被标准库的字符串接口意外接受。

豁免的边界必须说清楚，因为这里最容易写出双重缺陷。用 \Tp{char*} 或 \Tp{std::byte*}
读一块存储，解决的是“我用什么类型的 glvalue 去碰它”这一个问题；它没有解决
“这块存储里现在有没有一个 \Tp{T} 对象”这另一个问题。于是常见误判长这样：

\begin{lstlisting}[style=cppstyle, caption={四条合法双关路径与两条非法路径}]
#include <bit>
#include <cstdint>
#include <cstring>
#include <new>

std::uint32_t src_bits[1] = {0x3f800000u};   // 数组名即指针，示例避免取地址
float         dst_flt[1]  = {0.0f};

// 路径一：memcpy —— 全程不通过类型不符的 glvalue 访问，合法
void via_memcpy() {
    std::memcpy(dst_flt, src_bits, sizeof src_bits);
    std::memcpy(src_bits, dst_flt, sizeof dst_flt);      // 反向同理
}

// 路径二：std::bit_cast（C++20，<bit>）—— 语义即按表示拷贝，模板与 constexpr 友好
constexpr float bits_to_float(std::uint32_t b) { return std::bit_cast<float>(b); }
static_assert(bits_to_float(0x3f800000u) == 1.0f);       // 假定 float 为 IEEE 754（实现定义）

// 路径三：union —— C 侧是公认惯用法；C++ 侧读非活跃成员请按非法处理
union Pair { std::uint32_t bits; float as_float; };
float via_union(Pair p) { p.bits = 0x3f800000u; return p.as_float; }   // C++：UB

// 路径四：reinterpret_cast —— 只有实现文档明确支持时才可用，标准本身不给承诺
float via_cast(std::uint32_t* src) { return *reinterpret_cast<float*>(src); }  // 名单外：UB

struct Tag { std::uint32_t v; };
std::uint32_t bad_view(const unsigned char* buf) {
    return reinterpret_cast<const Tag*>(buf)->v;   // 别名不符 + 对象不存在：两处缺陷
}

// 正规做法：先让对象开始生命周期，再按它的类型访问
std::uint32_t good_view(void* buf) {
    Tag* t = static_cast<Tag*>(::operator new(sizeof(Tag), buf));  // 放置 new：创建对象
    std::memcpy(t, buf, sizeof(Tag));                              // 或用 start_lifetime_as
    return t->v;
}
\end{lstlisting}

上面 \Tp{bad\_view} 的两次出事是两回事：即便实现允许这种解释（比如关掉了严格别名），
缓冲区里也从未有过 \Tp{Tag} 对象，这是生命周期缺陷；而用 \Tp{Tag} 的 glvalue 去碰
\Tp{unsigned char} 数组是别名缺陷。修的时候只修一个，症状可能依旧。

\Tp{char} 与 \Tp{std::byte} 的差别在实践里只剩三条：\Tp{std::byte} 不能隐式参与算术
（必须先转成整数），因而写移位拼字时更不容易误用；它不会与文本语义混淆；它没有符号性问题。
别名豁免的范围则相同，名单把两者列在同一族里。至于 \Tp{signed char} 是否被 \CPP{}
各版本一并列入，措辞有过调整，跨版本代码按“\Tp{char}、\Tp{unsigned char}、
\Tp{std::byte} 三族必可，\Tp{signed char} 以你依循的草案为准”处理；\CL{} 的措辞以
字符类型为主，\Tp{signed char} 是否涵盖历来有争论，实现普遍同等对待，仍以各实现文档为准。

把字节缓冲变成对象的正规手段有两个：放置形式的 \Tp{operator new}（创建对象并可选初始化），
以及 \Tp{std::start{\Bk}\_lifetime{\Bk}\_as}（\cpp{23}，\Hdr{memory}，测试宏
\Feat{\_\_cpp{\Bk}\_lib{\Bk}\_start{\Bk}\_lifetime{\Bk}\_as}）。后者的定位正是
“我有一块对齐合适的存储，请让一个平凡可复制类型的对象从这里开始生命周期”，
它把本节的两个缺陷一次修掉；在支持不齐之前，放置 new 仍是可移植的那一个。

\section{四条合法路径的对照}\label{sec:ub06-paths}

\begin{compare}[title={同一件事的四种写法，代价各不同}]
\Tp{memcpy} 与 \Tp{std::bit{\Bk}\_cast} 都是拷贝表示，区别在类型检查的强度与可用性：
\Tp{bit{\Bk}\_cast} 要求两侧同宽、目标是可行类型，能在 \Tp{constexpr} 上下文里用，
读起来也直给意图；\Tp{memcpy} 更宽松（不必同宽，可以只拷一部分），
在常量表达式里则受限（\cpp{20} 起标准库版本可用于 \Tp{constexpr}，可用范围以你依循的版本为准）。
union 双关在 \CL{} 里是被普遍接受的惯用法，在 \CPP{} 里读非活跃成员按非法处理，
\CPP{} 只稳定保护公共初始序列这一小块。\Tp{reinterpret{\Bk}\_cast} 后直接解引用
完全落在实现文档里：它不是一种较弱的合法路径，而是一种需要外部担保的路径。
\end{compare}

\begin{table}[H]
\centering
\caption{四条双关路径的合法性来源与适用边界}\label{tab:ub06-paths}
\begin{tabularx}{0.78\textwidth}{>{\raggedright\arraybackslash}p{0.15\textwidth} X >{\raggedright\arraybackslash}p{0.20\textwidth}}
\toprule
路径 & 为什么合法 & 主要边界 \\
\midrule
\Tp{memcpy} & 从不通过类型不符的 glvalue 访问对象，别名名单根本没被触碰 & 源对象需可平凡复制；填充位内容不保证；小尺寸能否被折成一条移动是实现行为，不是承诺 \\
\Tp{std::bit{\Bk}\_cast} & \cpp{20} 明确赋予按对象表示解释的语义 & 两侧宽度必须相等；目标类型必须能表示该位模式；非法值表示要另行判断 \\
union 活跃成员 & \CL{} 惯用法，实现文档普遍接受；公共初始序列有明文保护 & \CPP{} 里读非活跃成员是 \UB{}；联合的大小与对齐由成员决定，末尾填充属 \IDB{} \\
\Tp{reinterpret{\Bk}\_cast} & 只有实现文档支持时才有立足点 & 标准层面既可能违反名单又可能缺对象；对齐要求另行成立（第~\ref{ch:ub07-alignment} 章） \\
\bottomrule
\end{tabularx}
\end{table}

一句话版本：能被标准担保的是前两条，第三条在 \CL{} 里靠惯例、在 \CPP{} 里靠禁用，
第四条从头到尾靠文档。需要“字节进、结构出”的场合，把 \Tp{memcpy} 或
\Tp{bit{\Bk}\_cast} 写成一层薄薄的函数，让所有类型解释集中在一处，
这既是可读性问题也是评审问题。

\section{编译器为什么依赖别名信息}\label{sec:ub06-why}

别名规则是标准发给优化器的许可证。有了“不同类型的指针不会指向同一对象”这个前提，
编译器才敢做这几件事：把一次加载提到循环外（它无法证明循环体里的写入会改它，但可以证明
类型不符的写入不会），把多个加载与存储重排以凑出向量化形态，把中间值留在寄存器里
而不写回内存，以及在两个类型不同的解引用之间省掉保守的顺序约束。反过来，代码违反别名规则时，
编译器并不是“优化出了 bug”，而是在执行一条标准允许它执行的推理：
既然你保证它们不重合，我就按不重合来排。

\begin{lstlisting}[style=shellstyle, caption={与严格别名相关的开关与诊断面}]
# 主流两家的文档都把严格别名挂在 O2 及更高级别；
# 低档位（-O1、-Os、-Oz）是否带上它，以各实现文档为准
gcc -O2 -fstrict-aliasing -Wstrict-aliasing=2 t.c
gcc -O2 -fno-strict-aliasing t.c              # 整段关闭：别名信息归零
clang -O2 -fno-strict-aliasing t.c
gcc -Wall t.c                                 # -Wall 含 -Wstrict-aliasing=3（GCC）

# 先用诊断把可疑点找出来，再决定是修代码还是整段关；
# 按函数关闭的属性看起来更精细，实际不可靠（见 6.5 节）
grep -n "reinterpret_cast" t.cpp
\end{lstlisting}

档位问题按文档处理：\Opt{fstrict-aliasing} 的主流约定是 \Opt{O2} 及以上带上它，
低档位是否包含以各实现文档为准；\Opt{fno{\Bk}-strict{\Bk}-aliasing} 的含义是放弃别名信息，
代价落在前面那几类优化上——被放弃的是整个翻译单元的循环与向量化质量，
不是一行代码的开销。工程上的两种做法由此区分：

\begin{itemize}
\item 局部修：把类型解释集中到少数函数，用 \Tp{memcpy} 或 \Tp{bit{\Bk}\_cast} 重写，
      保留其余代码的优化空间。这是首选。
\item 整段关：内核一类代码库会全项目加 \Opt{fno{\Bk}-strict{\Bk}-aliasing}，理由通常是
      历史代码里类型双关太多、且这些代码的优化重点本来不在别名上。
      这是可接受的策略，但要看清它换来的是什么：整个模块的别名信息一起没了，
      以后新写的代码也得不到这条规则的任何帮助，而且它并不能修对齐、生命周期
      与 \Tp{volatile} 那几个独立问题。
\end{itemize}

\begin{note}[title={别把 \Tp{restrict} 当别名开关}]
\Tp{restrict} 是加强承诺，不是放松承诺：它告诉编译器这个指针是访问该对象的唯一入口，
编译器据此可以更激进地重排。把它误当成允许别名的开关，方向和后果都是反的。
想让指针可以任意别名，\CPP{} 侧靠字符族与 \Tp{memcpy}，\CL{} 侧靠 \Tp{may\_alias} 属性，
都不靠 \Tp{restrict}。
\end{note}

诊断面要如实评价：\Warn{strict{\Bk}-aliasing=3}（GCC 由 \Warn{Wall} 打开）只覆盖模式明显的
用例，级别 1 与 2 会带上假阳性，跨函数与跨模块的双关它一般看不见；另一家实现是否提供同名
警告、语义是否一致，以其文档为准。静态警告不足以给结论，动态工具也不擅长这类缺陷
（它没有“错误的类型”这个概念），实践上靠代码审查与集中封装，工具面见 第~\ref{ch:ub03-sanitizers} 章。

\section{属性、内建兜底与字节缓冲区惯用法}\label{sec:ub06-attrs}

扩展属性用来在确认无法用标准手段表达的地方向实现申报意图。常见的是 \Tp{may\_alias}
（声明这个别名视图可以碰任何对象）、\Tp{aligned}（对齐另见 第~\ref{ch:ub07-alignment} 章）
与 \Tp{packed}。它们都是实现扩展，换编译器就要重新验证，写下来要连平台一起记进注释。

\Tp{optimize} 一类的按函数关优化的属性看似优雅，实际不可信：实现的文档明确提醒它不等于
命令行开关，可能只影响部分优化，也可能被后续内联抹掉。要“只关一个文件”就用构建系统
按源文件挂 \Opt{fno{\Bk}-strict{\Bk}-aliasing}，至少这个动作的效果可以核对；
要“只关一个函数”就不要指望属性，改成把该函数里的双关重写成标准路径。

\Tp{\_\_builtin{\Bk}\_memcpy} 是另一类兜底：实现把内建 \Tp{memcpy} 特殊对待，
它只拷贝对象表示，因而不会触发别名推理。写自定义字节拷贝函数（比如按字节翻转的序列化辅助）时，
尽量让编译器认出 \Tp{memcpy} 惯用法，比手写循环少一处别名歧义；但这是实现的识别能力，
不是标准承诺，识别失败就退化成普通循环，不影响正确性。

\begin{lstlisting}[style=cstyle, caption={\Hdr{cstdint} 与字节缓冲区：把整数写成字节再解析}]
#include <stdint.h>
#include <string.h>

/* 写：移位拼字节，不碰类型解释，天然避开别名规则 */
void put_u32be(uint8_t* out, uint32_t v) {
    out[0] = (uint8_t)(v >> 24);
    out[1] = (uint8_t)(v >> 16);
    out[2] = (uint8_t)(v >> 8);
    out[3] = (uint8_t)v;
}

uint32_t get_u32be(const uint8_t* in) {
    return ((uint32_t)in[0] << 24) | ((uint32_t)in[1] << 16) |
           ((uint32_t)in[2] << 8)  | (uint32_t)in[3];
}

/* 读整块结构：memcpy 进一个存活的对象，合法且可读 */
struct Wire { uint8_t op; uint8_t flags; uint16_t len; };
void parse_hdr(const uint8_t* buf, size_t n, struct Wire* dst) {
    if (n < sizeof *dst) { return; }
    memcpy(dst, buf, sizeof *dst);      /* 前提：dst 指向一个存活的 Wire */
}

/* 申报“这个视图可以碰任何对象”的实现扩展；可移植代码不要依赖它 */
typedef uint16_t __attribute__((may_alias)) any_u16;
uint16_t peek_u16(const uint8_t* p) { return *(const any_u16*)p; }  /* 仍受对齐约束 */
\end{lstlisting}

字节缓冲区的惯用法可以总结成一句：整数与字节之间用移位或 \Tp{memcpy} 换算，
结构体与字节之间用 \Tp{memcpy} 换算，任何情况下都不把缓冲区指针直接当结构体指针用。
最后一句里 \Tp{may\_alias} 那行给出的按 16 位读取只在缓冲区确实 2 对齐时安全；
把网络包按字节偏移解析时对齐往往不成立，那是另一类 \UB{}，与本章无关而属于
第~\ref{ch:ub07-alignment} 章。两个坑经常同时出现在同一段代码里，因此修的时候按顺序修：
先让类型解释合法，再让地址合法。

\begin{longtable}{@{}>{\raggedright\arraybackslash}p{0.24\textwidth} >{\raggedright\arraybackslash}p{0.14\textwidth} >{\raggedright\arraybackslash}p{0.34\textwidth}@{}}
\caption{常见写法与其类别归属}\label{tab:ub06-classes}\\
\toprule
写法 & 类别 & 触发点说明 \\
\midrule
\endfirsthead
\toprule
写法 & 类别 & 触发点说明 \\
\midrule
\endhead
\bottomrule
\endfoot
\Tp{memcpy} 在两个平凡可复制对象间搬表示 & 良定义 & 不通过类型不符的 glvalue 访问；填充位内容不计 \\
\Tp{std::bit{\Bk}\_cast}（\cpp{20}） & 良定义 & 同宽、目标可行的前提下按对象表示解释 \\
\Tp{char} 或 \Tp{std::byte} 数组读字节 & 良定义 & 豁免只给读写字节，继续按 \Tp{T} 用需另建对象 \\
把缓冲区指针强转成 \Tp{T*} 后解引用 & \UB{} & 别名不符，且对象多半未开始生命周期，两处独立缺陷 \\
\CPP{} 里读联合的非活跃成员 & \UB{} & \CL{} 按惯例接受，别把两边的结论互相搬 \\
\CL{} 里访问公共初始序列内的成员 & 良定义 & 两套语言都明文保护这一小块 \\
关掉 \Opt{fno{\Bk}-strict{\Bk}-aliasing} 后照常强转 & 仍是 \UB{} & 实现另给行为，但不是修复：换实现、换工具链后不复活 \\
\end{longtable}

\begin{warning}[title={字符豁免被用来“顺便”证明强转合法}]
看到标准允许 \Tp{char*} 访问任何对象，就推出那 \Tp{T*} 也行、编译器反正能优化，
是本章最常见的误读。豁免给字符族的特权来自一个特殊需求：程序必须能检查自己的内存表示。
它不给 \Tp{T*} 任何特权，也不改变“那块存储里有没有 \Tp{T} 对象”的判断。
连带一个更隐蔽的坑：\Tp{char} 的符号性是实现定义的（\IDB{}，别册的类型与转换主题里展开），
用 \Tp{char} 拼整数时高字节会被符号扩展，看起来像双关缺陷、其实是数值缺陷，
所以拼字节一律用 \Tp{unsigned char} 或 \Tp{std::byte}。
\end{warning}

\begin{bestpractice}[title={把类型解释压到一层薄薄的封装里}]
任何需要“字节进、类型出”的模块，先写两个内部函数（读取与写入），
内部只用 \Tp{memcpy} 或 \Tp{bit{\Bk}\_cast}，并让端序、宽度、对齐三件事在这里
一次性讲清；业务代码只调这两个函数。收益有三条：违反别名规则的写法没有藏身之处；
换平台时改动面收敛到两个函数；审查时“这一层是否只用合法路径”可以机械核对。
需要与只接受指针的 \CL{} 接口交互时，把缓冲区声明成有明确元素类型的数组，
让数组名充当指针，避免在接口边界上写类型解释。
\end{bestpractice}

\section{划界：语言契约与实现内部}\label{sec:ub06-boundary}

本章的结论全部来自语言标准：谁可以碰这块存储、碰到之后有没有对象、
编译器有权据此做什么推理。至于实现如何为一段代码构造别名图、别名图如何喂给指令重排与
向量化、内联与跨函数分析会不会让“只关一个文件”的意图失效——那是编译器与链接器分册的主题，
本章刻意不给实现内部细节，因为实现内部随时会变，而契约不会。
两册之间唯一要记住的接口是：契约告诉你这样写没有含义，
实现告诉你为什么它今天看起来有含义，两者都不能替代对方。

\section{可复现案例：协议头解析的三种写法}\label{sec:ub06-case}

\begin{casebox}[title={同一段解析代码的三种命运}]
一段从字节流里取 \Tp{uint32\_t} 长度字段的解析代码有三种写法：把缓冲区指针强转成
\Tp{uint32\_t*} 后解引用；用 \Tp{memcpy} 取四个字节；用 \Tp{std::bit{\Bk}\_cast}
配一个四字节数组。低档位、关掉严格别名时三者给出同样答案；开 \Opt{O2} 后第一种在
循环里被重排成“先读旧长度再写缓冲区”的形态，解析结果随输入顺序变化。
判定：第一种是 \UB{}（别名不符；缓冲区若只是 \Tp{char} 数组，\Tp{uint32\_t}
对象也从未存在），再加一次对未对齐地址的访问，能否被实现文档救回来只针对这一种写法；
第二、第三种各自合法，第三种还能进 \Tp{constexpr} 上下文而第一种不能。
症状最坑之处在于：第一种的错误结果在关掉内联后消失，于是被当成“优化引入的编译器问题”上报。
\end{casebox}

\begin{lstlisting}[style=cppstyle, caption={编译期可用的安全读取模板（只用合法路径）}]
#include <cstddef>
#include <cstdint>
#include <type_traits>

// 前提：bytes 指向 sizeof(T) 个来自大端流的字节；只走整数域，不做任何类型解释
template <class T>
    requires std::unsigned_integral<T>
constexpr T load_be(const std::byte* bytes) {
    T v = 0;
    for (std::size_t i = 0; i != sizeof T; ++i) {
        auto b = static_cast<unsigned char>(bytes[i]);   // 字符族豁免：读字节合法
        v = static_cast<T>(v << 8) | static_cast<T>(b);  // 无符号移位，不触有符号溢出
    }
    return v;
}

// std::byte 是 scoped enum，不接受整数：构造字节要显式 static_cast
constexpr std::byte msg[4] = {static_cast<std::byte>(0), static_cast<std::byte>(0),
                              static_cast<std::byte>(1), static_cast<std::byte>(0)};
static_assert(load_be<std::uint32_t>(msg) == 256u);      // 0x00000100，编译期即可核对
\end{lstlisting}

案例的结论有三条。第一，判定写法时要把别名缺陷、生命周期缺陷、对齐缺陷分开列，
它们往往同时出现但修法完全不同。第二，安全模板的特征是全程不做类型解释，
只在整数域拼值，因此既不受别名规则约束也不受对齐规则约束，还能进
\Tp{constexpr} 上下文。第三，选无符号类型不是风格问题：换成有符号类型就会把
第~\ref{ch:ub09-integers} 章 里的溢出问题请进来。附录 \ref{ap:ubA-kinds} 给出本案例的类别标签，
附录 \ref{ap:ubC-cases} 收录最小复现，涉及的开关集中在附录 \ref{ap:ubB-switches}。

\chapter{对齐：规定、代价与打包陷阱}\label{ch:ub07-alignment}

对齐是 C++ 里最少被讨论、却最容易在移植与重构时炸响的一条契约。它的麻烦在于：语言规定的只是要求，不是结果。标准告诉你每个类型有一个对齐要求，告诉你对象的地址必须满足这个要求，却把“这个要求是多少”“结构体到底长什么样”“未对齐访问会发生什么”全部留给实现。于是同一段代码，在一个平台上只是慢一点，在另一个平台上直接陷入，在第三个平台上悄悄产出错数据——三者都是合规的实现。

本章把对齐拆成两个方向来看：类型侧提出的要求，与对象侧实际拿到的地址。随后依次讨论过对齐、未对齐访问的平台差别、打包（压缩对齐）的三种写法及其共同风险、结构体布局与填充，以及对齐与并发之间那条“只是提示”的纽带。可以把本章当作第~\ref{ch:ub04-lifetime}、\ref{ch:ub05-pointer} 与 \ref{ch:ub06-aliasing} 三章的补充：那三章讲对象、指针与别名的约束，本章讲它们落在内存上时必须满足的位置约束。

\section{对齐的两个方向：类型的要求与对象的地址}\label{sec:ub07-two-directions}

\subsection{类型侧：\Lang{alignof} 与 \Tp{std::max\_align\_t}}\label{sec:ub07-type-side}

每个对象类型都带一个对齐要求，用 \Lang{alignof} 查询，单位是字节。它是编译期常量，含义是：这个类型的任何完整对象，其地址必须是该数值的整数倍。

\begin{lstlisting}[style=cppstyle,caption={类型侧的对齐要求是编译期常量，数值本身属于实现}]
#include <cstddef>
#include <cstdint>
#include <cstdio>

struct Pair { int a; int b; };

int main() {
    std::printf("%zu %zu %zu %zu\n",
        alignof(char), alignof(std::int32_t),
        alignof(std::int64_t), alignof(Pair));
    std::printf("baseline=%zu\n", alignof(std::max_align_t));
}
\end{lstlisting}

打印出来的具体数字属于实现，不属于语言：换一台机器、换一个数据模型，结果就可能不同。可移植的部分只有关系——数组的对齐要求与元素一致，聚合类型的对齐要求不小于其成员的最大值；至于档位是否总能排成 2 的幂，则由实现的 ABI 决定。

\Tp{std::max\_align\_t}（\Hdr{cstddef}）标记的是“实现不保证提供更宽对齐”的那条基线。一个类型的对齐要求若超过 \Lang{alignof}(\Tp{std::max\_align\_t})，它就带过对齐（extended alignment）；这类类型不能安全地出现在只承诺基本对齐的场合，除非实现明确支持——见 §\ref{sec:ub07-aligned-new}。

\subsection{对象侧：地址从哪里来}\label{sec:ub07-object-side}

类型只是提出条件，真正决定合规与否的是每一个对象的地址。地址的来源无非几处：静态存储区由编译器与链接器安排；栈上局部变量由编译器按类型要求摆放；堆内存由 \Fn{operator new}、\Fn{malloc} 一类接口给出，接口保证返回地址适合任何基本类型；剩下的、也是出事最多的来源，是指针算术与 \Lang{reinterpret\_cast}——人为把一个已有地址解释成对齐要求更高的类型。

\begin{lstlisting}[style=cppstyle,caption={地址由算术产生时，没有任何机制保证它满足目标类型的要求}]
#include <cstdint>
#include <vector>

std::int64_t read_at(const std::vector<unsigned char>& buf, std::size_t off) {
    const std::int64_t* p =
        reinterpret_cast<const std::int64_t*>(buf.data() + off);
    return *p;
}
\end{lstlisting}

这里要强调一个容易读漏的措辞：语言并不检查地址，语言规定的是“若不满足，该情形不在它的管辖范围内”。因此未对齐的对象在中间表示层面看起来完全正常，而优化器有权假定指针已经对齐，并在此前提上继续推导后续代码。

\subsection{一次未对齐读同时破坏的三件事}\label{sec:ub07-three-breaks}

很多人把未对齐访问当成性能问题。更准确的描述是：它同时动摇三条不同的保证。

第一条是可移植读取。硬件是否允许非对齐访问、允许到什么宽度、是否需要专门指令，都是实现层面的事。解引用一个未对齐指针属 \UB（未定义行为），因此“这次能跑通”不构成任何承诺，详见 §\ref{sec:ub07-platform}。

第二条是原子性保证。原子类型只在地址满足其对齐要求时，才提供无锁原子操作的语义。一个原本可移植原子（lock-free）的类型被搬到未对齐地址上，同步语义赖以成立的硬件前提就消失了；在容忍未对齐访问的平台上，这种破坏常表现为“单线程一切正常，并发压测下偶发读到撕裂值”。这与 第~\ref{ch:ub10-memory-model} 章 讨论的可见性与顺序问题是同一类故障的两个入口。

第三条是宽访存前提。编译器把标量循环改写为向量访存时，会假定基地址与步长满足宽寄存器要求的对齐。一旦这个假定来自人为伪造的地址，改写后的访存模式就与原代码不等价，越界、撕裂读以及更隐蔽的错误累加都会随之出现。

三条破坏的共同点是：它们都不在源代码层面报错，而且都不要求实现给出诊断（\NDR，无需诊断）。唯一可靠的判据是“这个地址是否真的满足类型要求”，而这恰恰是评审时最难验证的信息。

\section{过对齐：\Lang{alignas}、对齐化分配与 \Fn{std::assume\_aligned}}\label{sec:ub07-extended}

\subsection{\Lang{alignas} 只提高要求，不降低要求}\label{sec:ub07-alignas}

\Lang{alignas} 可以把一个类型或一个对象的对齐要求提高到实现支持的值上，它是过对齐的唯一语言层入口。

\begin{lstlisting}[style=cppstyle,caption={对类型加 alignas 与对单个变量加 alignas 是两件事}]
struct alignas(64) CacheLine {
    unsigned long long bits[8];
};

alignas(64) static unsigned char scratch[64];

CacheLine counter;
unsigned char raw[64];
\end{lstlisting}

最后两行的对比很关键：给类型加 \Lang{alignas}，该类型的每个对象都提出更高要求；给单个变量声明加 \Lang{alignas}，提高的只是这一个对象的要求，类型身份不变。这个区别决定了 §\ref{sec:ub07-aligned-new} 里“要不要走对齐化分配”的分岔。

\subsection{\cpp{17} 起按对齐请求分配}\label{sec:ub07-aligned-new}

\cpp{17} 起，\Lang{new} 表达式在为带过对齐的类型或数组分配存储时，会选择带对齐参数的分配函数重载；对应的语言入口是 \Hdr{new} 中的 \Tp{std::align\_val\_t} 与 \Fn{operator new}、\Fn{operator delete} 的对齐参数一族。在此之前，语言没有把对齐需求传给分配器的正式渠道，程序员只能依赖实现扩展或自行管理缓冲。部分实现在未启用 \cpp{17} 时还需要额外开关才提供这一族接口，开关名以实现文档为准。

\begin{lstlisting}[style=cppstyle,caption={过对齐类型经 new 分配时，对齐要求自动进入分配请求}]
#include <new>

struct alignas(128) State {
    double slot[16];
};

State* make() { return new State; }
State* make_n(int n) { return new State[n]; }

void destroy(State* s) noexcept { delete s; }
\end{lstlisting}

两条实践提醒。其一，替换全局 \Fn{operator new} 的实现必须同时提供对齐参数一族，否则过对齐类型会退回不保证对齐的路径，把问题推向运行时。其二，\cpp{17} 前后混编的翻译单元共享同一定义时，过对齐类型的分配路径不一致，会造成跨镜像的接口不对称；这类问题的完整讨论属于实现层与 ABI 话题，本章只标注它是风险点。

\subsection{\Fn{std::assume\_aligned}：编译器无法自行证明的断言}\label{sec:ub07-assume-aligned}

\cpp{20} 引入 \Fn{std::assume\_aligned}（\Hdr{memory}），它把一个指针声明为“其地址满足给定对齐”，让编译器据此优化后续访存。正确的定位是：这是一条编译器无法自行证明的断言。对齐值以模板实参形式给出，返回同类型指针，其可靠性完全依赖程序员的先验知识。

\begin{lstlisting}[style=cppstyle,caption={同一个断言，一个必然不成立、一个可能成立，差别只在地址来源}]
#include <memory>

double* misuse(unsigned char* buf, int i) {
    double* p = std::assume_aligned<8>(
        reinterpret_cast<double*>(buf + i));
    return p;
}

double* justified(double* p) {
    return std::assume_aligned<16>(p);
}
\end{lstlisting}

第一个函数里，偏移量可以是奇数，缓冲区起点也只按字节对齐，断言从一开始就是假的。误用 \Fn{std::assume\_aligned} 比误写一次 \Lang{reinterpret\_cast} 更危险：它把“未对齐”从一次可疑访问升级为编译器信任的前提，断言为假时被污染的不止这一次访问，而是所有基于该指针的后续推导。判据很简单——如果指不出“这个地址由谁、用什么机制保证对齐”，就不该加这条断言。

\section{未对齐访问的平台差别}\label{sec:ub07-platform}

同一份含未对齐访问的 \CPP 代码，在不同架构上会出现三类定性不同的结局。本章刻意不给周期数、不给具体机型，因为真正的差别在结局的种类，而不是慢多少。

\begin{table}[htbp]
\centering
\begin{tabular}{>{\raggedright\arraybackslash}p{0.18\textwidth}>{\raggedright\arraybackslash}p{0.34\textwidth}>{\raggedright\arraybackslash}p{0.28\textwidth}}
\toprule
结局种类 & 程序观察到的现象 & 对测试与移植的含义 \\
\midrule
透明承担 &
  访问由硬件拆成若干子操作或以额外微操作完成，功能正确，代价是访存变慢且宽度受限。 &
  本机测试全绿，问题只表现为性能回退或极端负载下的超时；移植到硬失败平台后，同一行代码立刻崩溃。 \\
陷入并由内核修复 &
  硬件产生对齐异常，操作系统在内核态用软件方式重放这次访问，再返回用户态继续执行。 &
  功能仍正确，但每次未对齐访问都付出异常级代价；这类回退容易被误读成锁竞争或分配器问题。 \\
直接失败 &
  异常不可修复，进程终止或返回错误状态，故障点与原因一一对应。 &
  这是最幸运的一类：可定位性最好，也最容易在测试阶段被抓住。 \\
\bottomrule
\end{tabular}
\caption{未对齐访问的三类定性结局}
\label{tab:ub07-platform-outcomes}
\end{table}

为什么同一份代码会走不同结局？因为语言对未对齐访问不做任何规定，三类结局全部是合法实现。指令集层面的历史选择决定了默认档位：有的架构为常规访存指令提供非对齐形式，有的只为特定向量加载提供，有的把非对齐列为保留异常；操作系统再在其上决定“能不能修、愿不愿修”。

结论只有一条：不要把任何一类结局当作“平台行为可预期”的证据。判定必须回到源码——只要存在“地址由算术或类型双关产生、又缺少对齐检查”的位置，就是缺陷，与本机是否崩溃无关。检测手段见 第~\ref{ch:ub03-sanitizers} 章：地址类动态检测器能在该路径被执行到的前提下报告未对齐访存，但“被执行到”本身是新加的前提。

\section{打包：三种写法与共同风险}\label{sec:ub07-packing}

打包指把结构体的对齐要求压缩到更小档位，从而消除填充，让布局贴近外部格式（协议头、文件格式、硬件寄存器表）。它在主流实现里都不是语言特性，而是扩展——这一点决定了它的全部风险。

\subsection{属性式：逐类型打包}\label{sec:ub07-packed-attr}

\begin{lstlisting}[style=cppstyle,caption={属性式打包只对声明了该属性的类型生效}]
struct __attribute__((packed)) Reg24 {
    unsigned char tag;
    unsigned long long value;
    unsigned short crc;
};

unsigned long long peek(const Reg24& r) { return r.value; }

const unsigned long long* addr(const Reg24& r) { return &r.value; }
\end{lstlisting}

同一个结构体，两个函数的风险等级完全不同。\Feat{peek} 读成员值：编译器知道该成员不对齐，会生成降级的访存序列。\Feat{addr} 取成员地址：它必须把一个“类型声称 8 对齐、实际只有 1 对齐”的指针交给调用方，问题被转移到接收端。属性自身的拼写与适用范围属于实现扩展，使用时以其文档为准，不要跨实现照抄。

\subsection{编译指示式：逐段打包}\label{sec:ub07-pragma-pack}

\begin{lstlisting}[style=cppstyle,caption={编译指示式的作用范围由成对的恢复指示决定}]
#pragma pack(push, 1)
struct Wire {
    int sequence;
    long long timestamp;
    short flags;
};
#pragma pack(pop)

static_assert(alignof(Wire) == 1, "packing applied");
\end{lstlisting}

编译指示式有两个固有风险：作用范围是“从指示点到恢复点之间的所有声明”，一旦它出现在公共头文件里，包含该头的翻译单元中位于区间内的其他类型也会被一并压缩；而压入与恢复不配对时，影响会溢出到后续声明。这两类错误都不报错，只静默改变布局。

\subsection{命令行式：全局档位}\label{sec:ub07-cli-pack}

\begin{lstlisting}[style=shellstyle,language=ShellUB,caption={全局档位把压缩规则施加给每一个类型}]
# global packing level: spelling and defaults vary by implementation
clang++ -fpack-struct=1 -c wire.cpp -o wire1.o
clang++ -fpack-struct -c wire.cpp -o wire2.o
cl /Zp1 /c wire.cpp
\end{lstlisting}

命令行档位是最应避免的一种用法：它不是给某个协议结构打补丁，而是把整个程序里所有类型的对齐要求都压低。于是标准库容器、第三方库对象、原子类型全部落在不合适的地址上，故障点通常离原因很远。上面出现的开关（\Feat{-fpack-struct}、\Feat{/Zp}）在不同实现上的拼写、默认档位与支持程度并不一致，务必以实现文档为准。

\subsection{为什么“取打包结构体成员地址”在语言层有问题}\label{sec:ub07-packed-member-address}

指针类型自带对齐要求。\Tp{T*} 这个类型表达的不只是“这里有一个地址”，还包括“这个地址满足 \Lang{alignof}(T)”这一承诺。把一个指向不对齐位置的指针交给 \Tp{T*}，等价于交出一个违反其类型不变量的值：

\begin{itemize}
\item 该指针一旦被解引用，即落入 \UB（未定义行为），即使本次硬件容忍了访问；
\item 该指针一旦进入模板库代码，库会按“\Tp{T} 已对齐”的前提生成宽访存或原子访存；
\item 该指针一旦被用于原子类型、锁，或用指针低位存标记的技巧，后果就从“错误数据”升级为“结构损坏”。
\end{itemize}

正因为这个错误足够隐蔽，某些实现提供专门的诊断，针对“取打包结构体成员地址并交给更宽类型”这一模式发出警告；具体诊断名与开关在不同实现上差别很大，以实现文档为准，也不要把它当作可靠防线。

\begin{lstlisting}[style=outstyle,caption={此类诊断的典型形态（示意措辞，实际文本以实现为准）}]
warning: taking the address of a packed member whose alignment
         requirement is not satisfied by its offset
note: the result is undefined behavior once that pointer is used
\end{lstlisting}

修复顺序也在这里给出：把打包结构体降级为纯粹的布局描述符，只读它的偏移与大小；需要取值时逐字段拷贝；需要把外部字节转成对象时按 §\ref{sec:ub07-raw-bytes} 的写法；绝不让打包结构体的成员地址离开当前函数。

\section{结构体布局与填充}\label{sec:ub07-layout}

填充出现在两个位置，作用不同。

成员之间：为了把某个成员的起始地址抬到它自己要求的位置，编译器在前一个成员之后插入无名字节。把对齐要求高的成员放在前面通常能减少中间填充——“重排成员”这一手段的全部原理就在这里。

结构体尾部：为了让该类型的数组满足对齐要求，结构体大小被向上取整到对齐要求的整数倍。换句话说，尾部填充不是浪费，它保证首元素合法的同时让后续元素的起始地址同样合法。压缩尾部填充的唯一办法是压低对齐，而那会把问题转给 §\ref{sec:ub07-packed-member-address}。

\begin{lstlisting}[style=cppstyle,caption={同样的成员、不同的顺序，大小可能不同}]
#include <cstdint>

struct BadOrder  { unsigned char b; std::int64_t v; unsigned char c; };
struct GoodOrder { std::int64_t v; unsigned char b; unsigned char c; };

static_assert(sizeof(BadOrder) >= sizeof(GoodOrder), "ordering matters");
\end{lstlisting}

一句话结论：成员顺序是可移植性的一部分——它决定填充、决定数组是否对齐、决定跨镜像的布局能否匹配；但布局的具体数字不由语言规定，标准只约定成员按声明顺序排列及其相对位置关系，不承诺任何偏移值。需要精确布局时，只能依赖实现文档，或者干脆改用显式序列化。

填充字节本身不承载值，对其内容做任何假设都落在 \USB（未指定行为）之内。把含填充的结构体整体写入文件、送进哈希或用于比较，都会把不确定的位带进输出——这是 第~\ref{ch:ub04-lifetime} 章 与 第~\ref{ch:ub06-aliasing} 章 同样会撞上的交叉问题。

\begin{warning}[title={全局打包档位改变所有类型的对齐，进而改变跨镜像的 ABI 契约}]
把打包档位设成全局默认（命令行式），改变的不是某个协议结构的布局，而是本次编译中每一个类型的对齐要求与大小。后果分三层：其一，同一份头文件在两个模块里算出不同的 \Lang{sizeof} 与成员偏移，跨镜像传对象时读写位置错位；其二，标准库与第三方库的二进制接口按正常对齐编译，与你压缩后的视图不一致；其三，原子类型与需要宽对齐的类型落在不合适的地址上，布局问题被升级为访存问题。如果确有某个结构需要紧凑布局，用逐类型属性或范围尽量小的打包区间，并把它限制在私有头文件内。永远不要为了省事使用全局档位。
\end{warning}

\subsection{三种改变对齐的手段对照}\label{sec:ub07-compare}

\begin{table}[htbp]
\centering
\begin{tabular}{>{\raggedright\arraybackslash}p{0.14\textwidth}>{\raggedright\arraybackslash}p{0.16\textwidth}>{\raggedright\arraybackslash}p{0.18\textwidth}>{\raggedright\arraybackslash}p{0.16\textwidth}>{\raggedright\arraybackslash}p{0.16\textwidth}}
\toprule
手段 & 生效时机 & 对可移植读取 & 对原子性保证 & 对跨镜像 ABI \\
\midrule
重排成员 &
  编译期，随类型定义一起分发。 &
  中性：地址始终满足要求，只是减少填充。 &
  中性：不改变任何类型的对齐要求。 &
  有风险：布局由顺序决定，两侧定义必须逐字节一致。 \\
打包指示 &
  编译期，按属性区间或全局档位生效。 &
  破坏：成员可能落在不满足要求的偏移上，取地址即成隐患。 &
  破坏：被压缩的类型不应承载原子成员。 &
  高风险：全局档位会让同一类型在镜像两侧大小不同。 \\
过对齐（\Lang{alignas}） &
  编译期提出要求；\cpp{17} 起运行期由对齐参数一族兑现。 &
  强化：地址比基本要求更宽，宽访存的前提更容易成立。 &
  有利：便于把独立数据摊到不同缓存行。 &
  有风险：对齐要求进入类型身份，两侧档位不一致时偏移与分配路径都不同。 \\
\bottomrule
\end{tabular}
\caption{改变对齐的三种手段及其影响面}
\label{tab:ub07-means}
\end{table}

\subsection{按字节解释外部数据的正确顺序}\label{sec:ub07-raw-bytes}

\begin{bestpractice}[title={先读进无对齐要求的字节缓冲，再逐字段构造对象}]
不要把外部字节流强转成对象然后读成员。可移植的顺序是：先把数据读进一个无对齐要求的字节缓冲（\Tp{std::array}、\Tp{unsigned char} 数组等）；再逐字段取出——用 \Fn{memcpy} 把按协议整理好的字节搬进局部对象，或用显式移位与掩码逐字节累加；若两侧尺寸相等、确需整体类型双关，用 \cpp{20} 的 \Fn{std::bit\_cast}（\Hdr{bit}）。这条顺序同时解决三件事：对齐（不再产生违规指针）、字节序（转换点显式可见）、填充（不会把不确定的位读进对象）。若确实想用打包结构体描述格式，把它当偏移表使用：只读 \Fn{offsetof} 与 \Lang{sizeof}，永不解引用它的成员地址。
\end{bestpractice}

\section{对齐与并发：缓存行只是提示}\label{sec:ub07-concurrency}

\cpp{17} 在 \Hdr{memory} 中提供了干扰尺寸常量，给出“实现认为适合避免（或促进）误共享的字节数”这一建议值。常见用法是把各自被独立写入的对象用 \Lang{alignas} 提到该档位，或用填充把它撑到该行大小。两个常量分别朝两个方向命名：需要避免误共享时看 \Tp{std::hardware\_destructive\_interference\_size}，希望促进共享时看 \Tp{std::hardware\_constructive\_interference\_size}。

\begin{lstlisting}[style=cppstyle,caption={用干扰尺寸常量把独立写入的计数器摊开}]
#include <atomic>
#include <cstddef>
#include <memory>

struct Counter { std::atomic<long long> value{0}; };

alignas(std::hardware_destructive_interference_size)
Counter producers[4];
\end{lstlisting}

但必须把话说清楚：这是提示，不是保证。第一，常量由实现给出，实现只能给近似值——同一实现可能对所有目标返回同一个保守数字，也可能与真实缓存行不符。第二，对齐只决定对象的起始位置，既不保证对象不与无关数据同处一行，也不保证跨内存域时的表现。第三，按缓存行对齐只把“落到同一行的概率”降下来，而误共享的根源是多个线程写入同一可寻址行内的不同字节，那属于第~\ref{ch:ub10-memory-model} 章的数据竞争与顺序话题，不是对齐话题——对齐不能让本就存在数据竞争的代码变正确，也不能替代原子类型或锁。反过来，为摊开缓存行而引入过对齐，就要接受 §\ref{sec:ub07-aligned-new} 那条分配路径的变化，以及跨镜像 ABI 上的额外约束。

\begin{casebox}[title={把外部协议头直接强转成含 64 位成员的结构}]
下面的写法读入协议头后直接解释为对象，在某些架构上失败。

\begin{lstlisting}[style=cppstyle,caption={原写法：缓冲区起点加偏移后直接双关为宽类型对象}]
#include <cstddef>
#include <cstdint>

struct Header {
    std::uint32_t magic;
    std::uint64_t timestamp;
    std::uint32_t length;
};

long long get_stamp(unsigned char* buf, std::size_t off) {
    Header* h = reinterpret_cast<Header*>(buf + off);
    return (long long)h->timestamp;
}
\end{lstlisting}

失败原因有三层，彼此独立。对象层：缓冲起点加偏移后，其地址不保证满足 \Lang{alignof}(Header)，在直接拒绝未对齐访存的平台档位上，这一句就是终止点。语言层：解引用未对齐指针落入 \UB（未定义行为）；即使硬件本次修复了访存，编译器仍可按“已 8 对齐”的前提生成宽访存。格式层：成员 \Feat{timestamp} 在结构内的偏移由实现决定，外部协议规定的是字节偏移与字节序，两者之间没有任何机制保证一致。

改写一：逐字段拷贝，不产生任何宽对齐指针。

\begin{lstlisting}[style=cppstyle,caption={改写一：字节先落到局部对象，指针从不越出函数}]
#include <cstring>

long long get_stamp_v1(const unsigned char* buf, std::size_t off) {
    std::uint64_t stamp = 0;
    std::memcpy(&stamp, buf + off + 4, sizeof stamp);
    return (long long)stamp;
}
\end{lstlisting}

改写二：保留解释式写法，但先用整数形式校验对齐，不满足则退回逐字段路径。

\begin{lstlisting}[style=cppstyle,caption={改写二：校验在先，类型转换在后}]
#include <cstdint>
#include <cstring>

long long get_stamp_v2(const unsigned char* buf, std::size_t off) {
    const std::uintptr_t a = reinterpret_cast<std::uintptr_t>(buf + off);
    if (a % alignof(Header) == 0) {
        const Header* h = reinterpret_cast<const Header*>(buf + off);
        return (long long)h->timestamp;
    }
    std::uint64_t stamp = 0;
    std::memcpy(&stamp, buf + off + 4, sizeof stamp);
    return (long long)stamp;
}
\end{lstlisting}

改写二里的分支顺序是关键：模数校验的对象是整数形式的地址，指针转换只在校验成立之后发生。如果先把地址做成 \Tp{Header*} 再去检查它，“生成违规指针”这一步本身已经出问题，检查救不回什么。

行为归类与检测：原写法属 \UB（未定义行为），与 第~\ref{ch:ub05-pointer} 章 讨论的“指针有效范围”是同一根源的两种表现；\Feat{Header} 的偏移与大小可以用编译期检查约束（\Fn{offsetof} 配 \Lang{static\_assert}），但那一只约束布局、不约束地址；运行时检测依赖 第~\ref{ch:ub03-sanitizers} 章 的地址类检查器。CI 至少要覆盖“缓冲区起点为奇数偏移”的用例——只测对齐起点，等于什么都没测。
\end{casebox}

\begin{note}[title={本章在内存与对象四条契约线中的位置}]
生命周期（第~\ref{ch:ub04-lifetime} 章）规定这块内存此刻是否是一个对象；指针（第~\ref{ch:ub05-pointer} 章）规定这个值是否指向它声称的位置；别名（第~\ref{ch:ub06-aliasing} 章）规定哪些指针可以被当作同一块存储来看待；本章的对齐规定对象落在地址的哪一格，以及指针类型对地址做了何种承诺。四条线各自独立成立，也各自独立失效：一个技术上合法的指针，可能同时是过期对象、违规别名与未对齐地址。评审时按这四条线逐一提问，比依赖“本机跑通了”可靠得多。
\end{note}

\begin{guideline}[title={对齐的三条底线}]
其一，永远不要让由算术或类型双关得到的地址绕过校验就交给宽类型：校验用整数形式做，转换在校验之后。
其二，打包只作用于最小必要的类型，并且只用作布局描述符；取成员地址、把成员做成原子类型、把结构跨镜像传递，三者都不可与打包同用。
其三，\Lang{alignas} 与 \Fn{std::assume\_aligned} 性质不同：前者提出要求、由分配机制兑现；后者声称兑现、由你独自负责。
\end{guideline}

\part{求值顺序与整型}

\chapter{求值顺序与副作用}\label{ch:ub08-evaluation}

表达式里写了两次自增、一次赋值、三个函数调用，编译器到底先算哪一个？
这一章讨论的不是优化技巧，而是语言在标准层面承诺了什么。
求值顺序与副作用的冲突之所以危险，是因为它的源码看起来完全合法：
没有越界、没有悬垂、没有类型双关，只有几个自增号的位置不同。
\CPP 早期长期借用一套近似 \CL「序列点」的说法描述这件事，
\CPP{} 11 起改用「求值之间的先后关系」这套代数式表述，
\CPP{} 17 又借助提案 \Feat{P0145R3} 把一批经典写法从 \UB{} 降级为 \USB{}。
理解这三层差别，才能判断一段代码是「不确定但有界」还是「彻底没有意义」。

\begin{note}[title={本章的判定口径}]
本章把结论分成三档：\UB{}（标准撤回全部承诺）、
\USB{}（实现在若干种结果里自选其一，通常由优化级别与目标架构决定）、
以及「已明确定义」。
判定顺序永远是先查标准是否定义了先后关系，再查实现是否给出额外文档说明，
最后才参考具体编译器的实测表现。
凡本章没有给出条款依据的顺序断言，一律按「不确定」处理，不得当作可依赖的约定。
\end{note}

\section{三个术语的准确差别}\label{sec:ub08-terms}

\subsection{sequenced before：同一线程内的强制先后}\label{sec:ub08-seq-before}

\Std{intro.execution} 把 \emph{sequenced before} 定义为求值（evaluation）之间的
非对称、可传递关系。
若求值 A \emph{sequenced before} 求值 B，则执行 A 的任何部分都严格早于执行 B 的任何部分，
中间不会被插入其它求值。
这是标准里唯一一档「强到可以据此推理」的关系。

需要区分两种先后：

\begin{itemize}
\item 值计算的先后：A 的值计算 \emph{sequenced before} B 的值计算，只保证读取的时刻有序。
\item 副作用的先后：A 的副作用 \emph{sequenced before} B 的副作用，才保证写入的可见次序。
\end{itemize}

标准对每个运算符分别规定「值计算之间」与「副作用之间」两套关系，
这也是阅读条款时最容易看漏的地方。
条件运算符与复合赋值只把左侧的一部分排在前面，并不自动把左侧对象的写入提前。

\subsection{indeterminately sequenced：有先后但不定，且不能交错}\label{sec:ub08-indet}

\emph{indeterminately sequenced} 表示：对某一次具体的求值而言，A 与 B 必然一个完整在前、
另一个完整在后，但标准不规定是哪一个。
关键在于「不能交错」：两个求值的指令序列不会互相穿插。
因此这类情形只让程序的结果变成 \USB{}，不会把程序推入 \UB{}。

\CPP{} 17 之后最常见的来源是函数调用的实参求值。
\Std{intro.execution} 规定各个实参表达式之间的求值关系是 indeterminately sequenced
（\CL 则是 unsequenced，见 第~\ref{sec:ub08-vs-c} 节）。
于是 \Tp{f(g(), h());} 中 \Fn{g} 与 \Fn{h} 谁先被调用不确定，
但不会出现两者的语句相互穿插、得到一个两边都不符合的结果。

\subsection{unsequenced：可以交错}\label{sec:ub08-unsequenced}

\emph{unsequenced} 表示两个求值之间没有任何先后保证，标准的措辞允许它们的执行相互交错。
只有这一档才有和副作用规则冲突的资格。
\Std{intro.execution} 中那条核心限制可以复述成：

\begin{quote}
若对一个标量对象的某个副作用，相对于同一个标量对象上的另一个副作用、
或者相对于使用该标量对象值的某个值计算而言是 unsequenced，则行为 \UB{}。
\end{quote}

这条规则是本章所有反例的判据。
注意三点细节：

\begin{enumerate}
\item 触发条件是 unsequenced，不是「顺序未规定」。indeterminately sequenced 不触发它。
\item 约束对象是标量对象。对 \Tp{std::vector<int>} 这类类对象的两个不同元素做写入，
      不受此规则约束，但会落到别的话题上：元素别名见 第~\ref{ch:ub06-aliasing} 章，
      迭代器失效见生命周期一章的讨论口径。
\item 「使用该对象值的值计算」也参与判定，所以读写冲突与写写冲突同等致命。
\end{enumerate}

\subsection{三档关系一览}\label{sec:ub08-relation-table}

\begin{table}[htbp]
\centering
\begin{tabularx}{0.78\textwidth}{>{\raggedright\arraybackslash}p{0.17\textwidth} >{\raggedright\arraybackslash}X >{\raggedright\arraybackslash}p{0.19\textwidth} >{\raggedright\arraybackslash}X}
\toprule
关系 & 能否交错 & 是否可能触发 \UB{} & 典型来源 \\
\midrule
sequenced before & 否，严格在前 & 否 & 逗号运算符、短路与、条件运算符、全表达式结尾 \\
indeter{\Bk}minately sequenced & 否，但先后不定 & 否，结果为 \USB{} & \CPP{} 17 函数实参之间 \\
unsequenced & 可以交错 & 与副作用落在同一标量时 \UB{} & \CL 函数实参之间 \\
\bottomrule
\end{tabularx}
\caption{三档求值关系的差别}
\label{tab:ub08-relations}
\end{table}

表~\ref{tab:ub08-relations} 的最后一列只列本章涉及的代表性来源，不是完整清单；
判定具体写法时仍要回到条款原文。

\section{副作用从哪里来}\label{sec:ub08-sideeffects}

\Std{intro.execution} 对副作用（side effect）的定义落在几类操作上：
通过 \Tp{volatile} 左值存取对象、修改对象、执行 I/O。
把这句话说白，标准关心的是「对外可见的状态变化」，而不是「看起来会改东西的表达式」。
下列情形都会被计入副作用：

\begin{itemize}
\item 修改标量对象：赋值、复合赋值、前置与后置自增减。
\item 修改类类型对象：调用非 \Tp{const} 成员函数、构造与析构写入成员、
      容器 \Fn{push\_back} 之类的可变操作。
\item 访问 \Tp{volatile} 对象：\CL 与 \CPP 都把 \Tp{volatile} 读写本身算作副作用，
      这与 第~\ref{ch:ub10-memory-model} 章 里原子量的可见性规则是两套独立机制，不要混用。
\item I/O：向流写入、调用 \Tp{printf}、写文件描述符。
\item 被调用函数内部的修改：调用本身就把被调用方对全局量、静态量、
      引用参数与指针所指对象的写入一并引入本次求值。
\item 库函数：\Hdr{cstdio} 与 \Hdr{cwchar} 的宽窄字符转换函数会改转换状态对象与
      \Tp{std::locale} 相关状态，因此这类调用彼此之间也有先后要求。
\end{itemize}

\begin{warning}[title={「看不见的写入」同样参与判定}]
一个名字看起来只是读取的函数调用，可能同时带来多处副作用：
它可能改全局量、改按引用传入的对象、改指针所指对象，甚至改 \Tp{mutable} 成员。
因此 \Tp{v[i] = f(i);} 是否安全，取决于 \Fn{f} 是否会写入 \Tp{v} 或它的元素。
只靠表达式形状做静态判断是不够的；代码评审时应要求被调用方在接口约定里
写明「是否修改与实参别名的对象」。
\end{warning}

\begin{guideline}[title={把副作用隔离到独立全表达式}]
一条语句只做一件「改变世界」的事，是本章唯一值得长期遵守的纪律。
需要多个修改时拆成多行，让每个修改都处在各自全表达式的结尾——
标准自动保证一个全表达式的副作用在进入下一个全表达式之前全部完成。
\end{guideline}

\section{与 \CL 的「序列点」模型对照}\label{sec:ub08-vs-c}

\begin{compare}[title={序列点模型与先后关系模型}]
\CL 用 \emph{sequence point}（序列点）这个「绝对时刻」描述副作用何时必须完成：
\Clause{5.1.2.3} 说相邻两个序列点之间，一个对象的存值至多被修改一次，
且若确实被修改，则对该对象的其他访问只能用于确定要存入的值。
这套说法直觉简单，坏处是它只能描述「点」，
无法表达「同一表达式内部两个子表达式之间的相对关系」。
\CPP{} 11 起改用二元关系 \emph{sequenced before} 刻画求值次序，
把 \emph{indeterminately sequenced} 与 \emph{unsequenced} 纳入词汇表，
于是可以精确说明「先算完 \Tp{f()} 再算 \Tp{g()}，但谁先未知」这类事实。
\CPP{} 17 的 \Feat{P0145R3} 进一步把大批运算符的次序写成明确的条款规则，
使一批原本落入 \UB{} 的写法降级为 \USB{}。
关系模型能覆盖 \CPP 特有的表达式子集，是这次替换的根本动机。
\end{compare}

\Clause{5.1.2.3} 给出的序列点集合大致包括：
全表达式结尾；\Tp{\&\&}、\Tp{||}、\Tp{,}、\Tp{?:} 中的指定位置；
部分函数调用的实参求值之后；以及一些声明与定义之间的点。
\CPP 中「全表达式结尾」仍是硬保证，但「函数调用」那一处被换成了分层描述：
被调表达式的值计算先于各实参表达式，实参之间只是 indeterminately sequenced。

\subsection{为什么序列点这套说法不够用}\label{sec:ub08-why-sp-fails}

三个具体理由：

\begin{enumerate}
\item \emph{重载运算符本质上是函数调用}。用户写的 \Tp{operator+} 找不到「序列点」这个位置，
      「点之间」的叙事无法说明左侧对象、实参表达式与临时量销毁三者的次序；
      换成对求值集合施加偏序，就能逐条规定。
\item \emph{带一串初始化器的表达式}。\Tp{std::initializer\_list} 参数、花括号列表、
      数组的列表初始化、\Tp{new} 表达式里可能存在的数组长度求值，
      都需要在同一表达式内部规定「按书写顺序」。序列点只能说「这些点之间没有副作用」，
      说不了「第 2 个元素先于第 3 个元素」。
\item \emph{与内存模型的接缝}。\CPP{} 11 的线程与原子量语义建立在「求值」这一粒度上，
      \emph{happens-before} 需要一个同构的关系概念才能和本章的偏序对齐；
      若沿用绝对时刻，同线程求值次序与跨线程可见性次序就成了两套词汇。
      这一点在 第~\ref{ch:ub10-memory-model} 章 展开。
\end{enumerate}

\subsection{同一构造在两种说法下的对照}\label{sec:ub08-c-vs-cxx}

\begin{longtable}{>{\raggedright\arraybackslash}p{0.19\textwidth} >{\raggedright\arraybackslash}p{0.23\textwidth} >{\raggedright\arraybackslash}p{0.21\textwidth} >{\raggedright\arraybackslash}p{0.15\textwidth}}
\caption{常见构造在 \CL 与 \CPP 中的顺序结论}\label{tab:ub08-c-vs-cxx}\\
\toprule
构造 & \CL 的说法 & \CPP{} 17 的说法 & 结论 \\
\midrule
\endfirsthead
\toprule
构造 & \CL 的说法 & \CPP{} 17 的说法 & 结论 \\
\midrule
\endhead
\bottomrule
\endfoot
\Tp{f(g(), h());} & 实参之间、实参与被调表达式之间都不规定顺序 & 实参之间 indeterminately sequenced，值计算先于调用 & 两语言下都是 \USB{} \\
\Tp{f(i, i++);} & 实参求值 unsequenced，读写落在同一标量 & 实参之间不能交错 & \CL 为 \UB{}，\CPP{} 17 为 \USB{} \\
\Tp{v[i] = i++;} & 更新与左侧求值 unsequenced & 赋值右操作数先于左操作数 & \CL 为 \UB{}，\CPP{} 17 已定义 \\
\Tp{i = i++ + i++;} & 两次修改落在同一对序列点之间 & 两个自增的副作用互不规定 & 两语言下都是 \UB{} \\
\Tp{a, b, c;} & 逗号运算符给出强制序列点 & 左操作数先于右操作数 & 已定义 \\
\Tp{fx() \&\& fy();} & 第一操作数之后是序列点，右侧可能不计算 & 左侧先于右侧，短路语义保留 & 两语言下都已定义 \\
\Tp{cout << x << y;} & 不适用 & 重载移位运算符左操作数先于右操作数 & \CPP{} 17 起已定义 \\
\Tp{P p\{fx(), fy()\};} & 不适用（\CL 无花括号列表） & 列表元素自左向右 & 已定义 \\
\end{longtable}

表中「重载移位运算符左操作数先于右操作数」是 \CPP{} 17 为移位运算符专门增加的规则，
它让链接形式的流输出不再依赖实现的自由度。
至于 \CPP{} 11 与 \CPP{} 14 的实现把这一情形判为 \UB{} 还是 \USB{}，
取决于当时的条款措辞，应以对应版本标准文本为准；迁移老代码时不要把这条规则当作可用。

\section{经典雷区与逐条判定}\label{sec:ub08-cases}

\subsection{同一标量的多次修改}\label{sec:ub08-multi-modify}

\begin{lstlisting}[style=cppstyle,caption={写写冲突：仍然没有意义的写法}]
int i = 0;
i = i++ + i++;                 // 两个自增的副作用互不规定：UB
int j = (i += 1) + (i += 2);   // 同一对象的两次复合赋值 unsequenced：UB
a[i] = i;                      // 安全：没有副作用冲突
i = f(i);                      // 安全：读取在调用内发生，写入在赋值时发生
\end{lstlisting}

判据始终是那句「副作用相对于同一标量对象的另一副作用或使用其值的计算是否 unsequenced」。
像 \Tp{i = i++ + 1;} 这种只牵涉一次额外修改的写法，在 \CPP{} 17 已被关系模型解释清楚；
在 \CL 与更早的 \CPP{} 里则不一定成立。
跨版本迁移代码时，这类「看起来只差一个括号」的表达式应当一律重写，
不要靠记忆去区分哪一年放宽了哪一条。

\begin{clause}[title={重载运算符的次序看写法，不看函数}]
\Std{over.match.oper} 把 \Tp{a + b} 形式的重载运算符解释成一次函数调用，
但这次调用套用哪套定序规则，取决于写法而不是取决于被调的那个函数。
\CPP{} 17 的 \Feat{P0145R3} 起，\emph{按运算符写法}调用时，
重载运算符沿用它所重载的那个内置运算符的定序规则：
\Tp{a \&\& b}、\Tp{a || b}、\Tp{a, b} 中左侧先于右侧，
\Tp{a[b]} 与 \Tp{a << b} 同样左侧先于右侧，赋值则先算右侧再执行写入。
把同一个表达式改写成函数调用写法 \Tp{operator\&\&(a, b)}，上面这条保证立刻消失，
退回普通调用规则：被调表达式的值计算先于各实参，两个实参彼此 indeterminately sequenced。
因此「重载 \Tp{operator[]} 的两个参数都改同一个对象」这类写法要分两种结论：
\CPP{} 17 起按运算符写已定义，按函数调用写则与函数调用同判。
\end{clause}

\subsection{自增与下标的组合}\label{sec:ub08-index}

\begin{lstlisting}[style=cppstyle,caption={下标表达式里的读与写}]
std::vector<int> v{1, 2, 3, 4};
v[i] = i++;             // CPP17 起已定义；CL 下仍是 UB
v[i++] = i++;           // 两个自增互不规定：UB
v[i] = v[i] + 1;        // 安全
int x = v[i] + v[i++];  // 读 i 与改 i 的表达式在不同子表达式中：UB 或 US，按版本判定
\end{lstlisting}

第二行是本章最重要的一例：即使 \CPP{} 17 放宽了赋值两侧的次序，
两个 \Tp{i++} 的副作用仍然互不规定，结论仍是 \UB{}。
把「\CPP{} 17 让求值顺序变安全了」理解成「自增号可以随便写」是危险误读。

\subsection{短路、逗号与条件运算符}\label{sec:ub08-shortcircuit}

除花括号列表之外，「操作数顺序已指定」的运算符就是这一族。
\Std{intro.execution} 把 \Tp{\&\&}、\Tp{||} 与 \Tp{,} 的第一个操作数排在第二个操作数之前，
把 \Tp{?:} 的第一个操作数排在被选取的那个分支之前。
这条规则不是 \CPP{} 17 才收回的自由度：\CPP 最初版本就用序列点的说法给出同一结论，
\CPP{} 11 换成定序关系的措辞；\CL 在 \Clause{6.5.13} 到 \Clause{6.5.17} 之间
也把这四个运算符逐一写成序列点。
两种语言在这一点上口径一致，可以放心依赖运算符内在顺序，不需要额外插入临时变量。

\begin{lstlisting}[style=cppstyle,caption={有强制顺序的构造}]
if (p != nullptr && *p > 0) { /* ... */ }   // 左侧先于右侧：安全
while (i < n && a[i] != key) { ++i; }       // 安全
int r = (++i, i * 2);                        // 逗号左侧先完成：r 用递增后的 i
\end{lstlisting}

条件运算符要注意一点：它只保证条件表达式的值计算先于被选取的那个分支，
分支内部与外部其他求值的关系仍按各自的规则判定。

\begin{warning}[title={顺序已指定与短路求值是两件事}]
重载版本在 \CPP{} 17 之后拿到了操作数顺序，却始终拿不到短路语义：
\Tp{operator\&\&} 与 \Tp{operator||} 是函数，两个实参必然都完成求值，
右侧不会因为左侧已经定出结果而被跳过。
标准库只为 \Tp{std::valarray} 重载这两个运算符，实现上也做不出惰性求值。
需要「左边不成立就不算右边」时，仍然要写成嵌套的 \Tp{if} 或者拆成两条语句，
不要指望自定义类型的 \Tp{\&\&} 替你短路。
\end{warning}

\subsection{链接的输出运算符}\label{sec:ub08-chain}

\begin{lstlisting}[style=cppstyle,caption={流输出的链接顺序}]
std::cout << f() << g() << h();   // CPP17：按书写顺序调用
                                  // 更早的版本：顺序依赖实现，实践结论不一
\end{lstlisting}

即使顺序已经确定，也不要在 \Tp{operator<<} 的实现里修改与实参别名的对象。
顺序保证解决的是「谁先执行」，不解决「执行期间对象是否被合法改写」。
自定义流插入运算符应保持逻辑上的常量性；确实需要携带状态时，
把状态放进显式的格式化辅助对象，例如 \Fn{std::setw}、\Fn{std::setfill} 这一类。

\section{函数调用、实参绑定与初始化列表}\label{sec:ub08-call}

\subsection{调用的分层顺序}\label{sec:ub08-call-order}

\Std{expr.call} 与 \Std{intro.execution} 一起给出如下分层：

\begin{enumerate}
\item 表示被调函数的表达式（含成员调用的对象表达式）的值计算，先于各个实参表达式；
\item 各实参的求值彼此 indeterminately sequenced；
\item 全部实参的副作用在进入被调函数体之前完成；
\item 形参初始化、默认实参求值与临时量构造之间的相对次序随版本条款而变，
      重载解析涉及的用户定义转换还额外引入自己的求值。
\end{enumerate}

第 4 条是本章唯一一处「必须查版本」的分层：不同版本对形参初始化与实参求值的
接缝写法并不相同，涉及默认实参的代码在移植时应重新核对，而不是套用第 2 条的结论。

\begin{warning}[title={\CL 与 \CPP 在这一点上结论不同}]
\Clause{6.5.2.4} 把 \CL 的实参之间、以及「被调表达式求值 vs 实参求值」都留成 unsequenced。
同一份「跨语言共享」的表达式片段（宏、参考实现、双语言头文件里的公共算法）
在两种语言下的判定可能不同：\Tp{f(i, i++);} 在 \CL 里是 \UB{}，
在 \CPP{} 17 里只是 \USB{}。
维护跨语言链接边界（\CL 链接约定下的共享头文件与内联实现）时不要假设规则相同。
\end{warning}

\subsection{花括号列表是少数有保证的顺序}\label{sec:ub08-initlist}

\begin{lstlisting}[style=cppstyle,caption={列表初始化自左向右}]
void log_all(std::initializer_list<int> xs);
log_all({next_id(), next_id(), next_id()});   // 按书写顺序求值

struct P { int x, y; };
P p{fx(), fy()};                               // x 的初始化先于 y
int arr[2] = {g(), h()};                       // 按索引顺序初始化
\end{lstlisting}

\Std{dcl.init.list}（\CPP{} 17 的 \Clause{8.5.4}）规定花括号列表中各元素的求值自左向右，
数组与聚合的初始化器也按索引/成员顺序依次求值。
当某个表达式必须在多个带副作用的子表达式之间取得确定顺序时，
把参数改成 \Tp{std::initializer\_list}、或者改用花括号列表，
是比插入临时变量更贴近语义的做法。

\subsection{默认实参}\label{sec:ub08-default-arg}

默认实参在每次调用时求值，求值发生在被调函数的作用域里，
因此它能看见形参与被调函数的其他名字。
它的副作用与调用方表达式的求值不在同一个语法层，
读代码时无法只从调用点判断它什么时候生效。
结论不是「默认实参危险」，而是「默认实参里不要放会改共享状态的表达式」：
自增、全局量读取、带副作用的函数调用都属于应当移出的内容。

\section{诊断手段}\label{sec:ub08-diag}

\subsection{编译器警告}\label{sec:ub08-warn}

GCC 与 Clang 对未定次序提供了专门的警告开关，
其中 \Opt{Wsequence-point} 在 \Opt{Wall} 之下会被启用，
能够捕获 \Tp{v[i] = i++;}、\Tp{f(i, i++)} 这类形态的多数实例。
不同版本对这些开关的覆盖范围与误报率都有变化，
某个具体写法是否被报告，应以所用编译器版本的文档与实测为准。

\begin{lstlisting}[style=shellstyle,caption={打开顺序相关警告}]
g++ -std=c++17 -Wall -Wextra -Wsequence-point -c eval_order.cpp
clang++ -std=c++17 -Wall -Wextra -Wsequence-point -c eval_order.cpp
\end{lstlisting}

警告只覆盖语法模式匹配得到的形态。
跨函数的副作用冲突（被调用方写了调用方也读的对象）无法靠这一类警告发现，
只能靠 sanitizer 或代码审查。

\subsection{运行期检测}\label{sec:ub08-run}

现有 sanitizer 中没有一个以「求值顺序」为检查对象。
\Opt{fsanitize=undefined} 一族（见 第~\ref{ch:ub03-sanitizers} 章 的介绍）检查的是有符号溢出、
除零、越界等可判定事件，顺序本身不在它的判定集合里；
线程检测器 \Opt{fsanitize=thread} 只在跨线程访问冲突上报警，
同一线程内的 unsequenced 副作用不会触发它。

可行的定位办法是差分构建：同一份源码在 \Opt{O0} 与 \Opt{O2} 下分别编译运行，
必要时再叠加不同目标架构。结果随优化级别漂移而源码没有可见的语义差别时，
绝大多数原因就落在求值顺序或未指定行为上。

\begin{bestpractice}[title={顺序敏感代码的自查清单}]
\begin{enumerate}
\item 一个全表达式里对同一标量最多出现一次修改；出现两次就拆行。
\item 修改与读取落在同一对象时，先把读取结果存进独立的临时变量。
\item 实参尽量不带副作用；必须带时确认被调用方不会写入与实参别名的对象。
\item 需要确定的从左到右顺序时，改用 \Tp{std::initializer\_list} 或花括号列表，
      而不是依赖「参数似乎总是从左算起」的经验。
\item 与 \CL 共享的表达式片段单独复核一遍，两种语言的结论可能不同。
\end{enumerate}
\end{bestpractice}

\section{跨版本与跨实现的差异}\label{sec:ub08-versions}

\begin{table}[htbp]
\centering
\begin{tabularx}{0.78\textwidth}{>{\raggedright\arraybackslash}p{0.13\textwidth} >{\raggedright\arraybackslash}X >{\raggedright\arraybackslash}p{0.24\textwidth}}
\toprule
版本 & 与求值顺序相关的变化 & 迁移影响 \\
\midrule
\CPP{} 03 & 沿用近似 \CL 的序列点叙事，重载运算符与列表元素次序无保证 & 老代码里的自增组合多数不可移植 \\
\CPP{} 11/14 & 引入 \emph{sequenced before} 与 \emph{indeterminately sequenced} 两个概念，但运算符细则尚未补齐 & 同一写法在不同实现上判定不同 \\
\CPP{} 17 & \Feat{P0145R3} 明确大批运算符次序，赋值右先于左、重载移位左先于右，且重载运算符按运算符写法沿用内置顺序 & \Tp{v[i] = i++;} 与链接流输出转为已定义 \\
\CPP{} 20 及以后 & 在 \Tp{new} 表达式的求值时机、比较运算等构造上继续收紧 & 旧判定需重新核对，勿沿用 \UB{} 结论 \\
\CL（C11/C17） & 序列点模型保持稳定，实参之间 unsequenced & \Tp{f(i, i++);} 一类仍是 \UB{} \\
\bottomrule
\end{tabularx}
\caption{语言版本对求值顺序结论的影响}
\label{tab:ub08-versions}
\end{table}

表~\ref{tab:ub08-versions} 最后一行的 \CL 结论跨版本稳定；
\CPP 各版本之间的差异则以对应版本标准文本为准，
本章刻意不给 \CPP{} 20 之后逐条的条款号。
实现层面的差异更直接：同一份 \USB{} 代码在两个编译器上可能给出两种都合法的结果，
这类差异不属于缺陷报告的范围，只能靠代码本身消除。

\begin{casebox}[title={「只改了一个括号」的事故形态}]
有一类反复出现的问题形态：维护者把 \Tp{buf[n++] = read();}
改成 \Tp{buf[n++] = buf[n] + 1;}，理由是少一次函数调用。
新写法里左侧下标对 \Tp{n} 的读取，与 \Tp{n++} 的副作用互不规定，
在 \CL 与旧版 \CPP 里都落到 \UB{}。
优化级别一变，实际写入的下标就可能与预期偏移量差一，
症状是某个数据被写到相邻位置，而源码审查看不出任何越界。
这类缺陷的定位特征是：地址消毒器报告的是数值差异而不是越界，
且只在特定优化级别出现。
修复方式不是补括号，而是把自增提到语句之外：先算值，再写入，再单独递增。
\end{casebox}

\section{本章要点}\label{sec:ub08-summary}

\begin{itemize}
\item \emph{sequenced before} 是强制先后，\emph{indeterminately sequenced} 是先后不定但不能交错，
      \emph{unsequenced} 才可能交错。只有第三种会和副作用规则冲突。
\item \UB{} 判据只有一句：同一标量对象上的副作用，相对于另一副作用或使用其值的计算，
      若为 unsequenced，行为未定义。
\item 副作用包括修改标量与类对象、\Tp{volatile} 访问、I/O、被调用方内部的修改，
      以及会改转换状态的库函数。
\item \CL 的序列点模型说「点」，\CPP 的关系模型说「边」；
      后者才能覆盖重载运算符、花括号列表与 \Tp{new} 表达式这些 \CPP 特有构造。
\item \CPP{} 17 把一批经典写法从 \UB{} 降级为 \USB{} 或已定义，
      但两次自增同时出现、实参与被调用方写入别名的组合仍然不安全。
\item 「操作数顺序已指定」的构造一直是那几种：\Tp{\&\&}、\Tp{||}、\Tp{,}、\Tp{?:}
      与花括号列表；\CPP{} 17 起这套顺序还覆盖按运算符写法调用的重载版本，
      而按函数调用写法退回普通调用，重载的 \Tp{\&\&}、\Tp{||} 也永远给不出短路。
\item 现有 sanitizer 没有针对求值顺序的检查器；
      差分构建与 \Opt{Wsequence-point} 是现实可用的主要手段。
\end{itemize}

求值顺序只在单线程范围内有意义。
一旦引入线程，同一对象上的冲突访问是否可观察，就不再由本章的偏序关系决定，
而由 \emph{happens-before} 与内存序决定，相关内容见 第~\ref{ch:ub10-memory-model} 章。
表达式里另一个高频冲突源是整型提升、转换与算术溢出，见 第~\ref{ch:ub09-integers} 章；
「全表达式结尾」这一保证与对象生命周期结束时点的差别，可对照 第~\ref{ch:ub04-lifetime} 章；
而本章反复用到的「已定义 / \USB{} / \UB{}」三档划分本身，
在 第~\ref{ch:ub01-taxonomy} 章 给出完整的类别体系。

\chapter{整型：有符号溢出、转换与整数模型}\label{ch:ub09-integers}

整型是本册里代价最低、也最容易被低估的一类未定义行为。它不需要越界的指针，
不需要悬垂的引用，一行算术就够；但它出现的场合恰恰是边界检查、长度换算、
计数循环、金额与时间戳——一旦出问题，最先被破坏的不是结果，而是程序自己写下的
那层防御。第~\ref{ch:ub01-taxonomy} 章给过的判据在这里格外好用：先确定这段话属于
哪个类别，再决定要不要修代码。整型这一章的多数条目，答案都是“修”。

\section{对称性差别：溢出未定义，回绕良定义}\label{sec:ub09-symmetry}

标准对整型算术的两句承诺是不对称的。有符号一侧，措辞落在 \Std{expr}：
当运算结果不能被结果类型表示时，行为未定义；这条不只覆盖加减乘，也覆盖一元取负、
自增自减以及任何产生超范围值的表达式。无符号一侧，措辞落在 \Std{basic.fundamental}：
无符号类型的运算按 $2^N$ 取模进行，$N$ 是结果类型的位宽——这是明确写死的承诺，
不是编译器“恰好”如此实现。

\begin{lstlisting}[style=cppstyle, caption={同一行代码的两种命运：有符号溢出与无符号回绕}]
int f(int x) {
    int y = x + 1;              // 有符号：x == INT_MAX 时整个表达式是 UB
    unsigned u = x + 1u;        // 无符号：模 2^N 回绕，标准明文保证
    return (y > x) + (int)(u > 0u);
}
\end{lstlisting}

上面这段还顺带暴露了第二个陷阱：\Tp{u > 0u} 对无符号量而言恒等于
\Tp{u != 0}，它不是“正数”的意思。有符号溢出是 \UB{}，无符号回绕是良好定义，
两者在语义上并不对等；把“回绕”当成“溢出”的同义词写进设计文档，是本类缺陷的常见起点。

\begin{clause}[title={标准怎么说：一句关于结果，一句关于类型}]
有符号一侧关心的是“结果能不能被结果类型表示”，它不给任何值；
无符号一侧关心的是“结果类型有哪些值”，它按模取一个。因此同一份源码在两个类型上的
正确性论证完全不同：前者要证明上界，后者要证明回绕后的值正是想要的值。
\end{clause}

\section{优化器为什么能据“溢出不可能发生”反推前提}\label{sec:ub09-optimizer}

\UB{} 的杀伤力从来不在于“结果不对”，而在于“标准撤回全部承诺之后，编译器有权把
这条表达式当成不可能到达的路径处理”。有符号溢出是这条推理链最典型的入口：
既然溢出是未定义，编译器就可以在任何优化里假定它不发生，再用这个假定反向缩小
输入范围。常见后果有三类：边界检查被判定为恒假而被删掉、循环上界被推断成
“必然在某个值之前退出”从而向量化、以溢出为终止条件的搜索被直接折叠成常量。

\begin{warning}[title={“它现在能跑”不是证据}]
未优化构建（\Tp{-O0}）下溢出的实际表现通常是回绕，因为那条路径根本没被优化过。
同一份代码在 \Tp{-O2} 之后可能整个分支消失。第~\ref{ch:ub01-taxonomy} 章的结论在这里
仍然成立：观察一次执行不能证明语义，只能证明那一次没有暴露。
\end{warning}

实现确实提供了改变这一前提的开关，但要知道它们的边界。GCC 与 Clang 都有
\Tp{-fwrapv}（把有符号溢出按补码回绕对待）与 \Tp{-fno-{\Bk}strict-{\Bk}overflow}
（禁止依赖“溢出不发生”的推断），它们是扩展开关而非标准设施；一旦开启，
程序里“靠溢出做终止条件”的写法在优化层面就再也得不到任何加速。
MSVC 没有对应开关，实践中有符号溢出通常表现为回绕，但这是实现行为（\IDB{}），
不是标准承诺，不能写进跨平台假设。消毒器一侧，第~\ref{ch:ub03-sanitizers} 章的
\Tp{-fsanitize{\Bk}=undefined} 已经覆盖 \Tp{signed-{\Bk}integer-{\Bk}overflow} 检查，报告文本一般形如
\Tp{signed integer overflow} 加表达式与目标类型；它的覆盖范围与 CI 分档结论同样适用本章。

\begin{bestpractice}[title={优先证界，不靠开关}]
把“绝不溢出”变成可检查的前提：用足够宽的中间类型算完再判界，或者用下一节的
显式检查函数。依赖 \Tp{-fwrapv} 只是给整个翻译单元贴了一张创可贴，
它既不能替你论证边界，也不能在别的编译器上生效。
\end{bestpractice}

\section{安全算术的可移植出口}\label{sec:ub09-safe}

标准本身直到最近都没有提供“带溢出标志的加法”这一设施，现实中的出口有三条，
按可移植性从高到低排：

\begin{enumerate}
  \item 显式构造不会溢出的比较式：把 \Tp{a + b} 的判界拆成两条单边检查——
        \Tp{b > 0} 时查 \Tp{a > max - b}，否则查 \Tp{a < min - b}，
        两条式子里减去的符号都保证结果落在范围内。这是纯标准的写法。
  \item 用更宽的类型算完再判界，例如升到 \Tp{int64\_t} 或用 \Tp{long long}
        承载中间量。前提是宽类型真的更宽（见第~\ref{sec:ub09-models} 节）。
  \item 用编译器扩展：\Tp{\_\_builtin{\Bk}\_add{\Bk}\_overflow}、\Tp{\_\_builtin{\Bk}\_sub{\Bk}\_overflow}、
        \Tp{\_\_builtin{\Bk}\_mul{\Bk}\_overflow} 一族（GCC 与 Clang 提供，返回是否为溢出，
        结果通过出参写出）。MSVC 走 x86/x64 的 intrinsic 路线，
        例如 \Fn{\_mulx\_u64}、\Fn{\_addcarry\_u64} 一类；不同体系结构上可用的
        intrinsic 集合差别很大，具体名字与可用性请以 MSVC 的 intrinsic 文档为准。
\end{enumerate}

\begin{lstlisting}[style=cppstyle, caption={三条出口：标准判界、宽类型、编译器扩展}]
#include <limits>

bool add_std(int a, int b, int *out) {           // 出口 1：只用不会溢出的表达式
    const int lo = std::numeric_limits<int>::min();
    const int hi = std::numeric_limits<int>::max();
    if (b > 0 ? a > hi - b : a < lo - b) return false;
    *out = a + b;
    return true;
}

bool add_wide(int a, int b, int *out) {          // 出口 2：宽类型承载中间量
    const long long t = static_cast<long long>(a) + static_cast<long long>(b);
    const long long lo = std::numeric_limits<int>::lowest();
    if (t < lo || t > std::numeric_limits<int>::max()) return false;
    *out = static_cast<int>(t);
    return true;
}

bool add_gcc(int a, int b, int *out) {           // 出口 3：编译器扩展，非标准
    return !__builtin_add_overflow(a, b, out);
}
\end{lstlisting}

\Tp{std::numeric{\Bk}\_limits} 需要提供三个边界：\Fn{max}、\Fn{min} 与 \Fn{lowest}。
\Fn{lowest} 不是 \Fn{min} 的别名——在无符号类型上 \Fn{min} 是零而 \Fn{lowest} 也是零，
在有符号类型上 \Fn{lowest} 才是 \Tp{INT\_MIN}；把 \Tp{-numeric{\Bk}\_limits<int>::min()}
当作下界，本身就会在最小值附近踩回第一节的坑。

\begin{note}[title={饱和算术尚未标准化}]
把溢出钳到边界值的“饱和运算”（saturating arithmetic）作为库设施仍停留在提案阶段，
截至写作时未标准化，别把它当成可移植设施来写；需要它时就用上面的判界函数手写钳位。
\end{note}

\section{整数提升与常用算术转换}\label{sec:ub09-promotions}

比大小之前先统一类型，这套规则叫整数提升与常用算术转换（usual arithmetic
conversions）。两条最容易咬人：

\begin{itemize}
  \item 比 \Tp{int} 位宽小的类型（\Tp{char}、\Tp{short}、枚举）先提升到 \Tp{int}
        或 \Tp{unsigned}，于是对窄无符号类型做算术常常并没有你预期的“窄”。
  \item \Tp{int} 与 \Tp{unsigned} 混合运算时，若 \Tp{unsigned} 的秩不低于
        \Tp{int}，有符号侧被整体转换为无符号。一个负的 \Tp{int} 会变成一个巨大的
        无符号数，比较结果与直觉完全相反。
\end{itemize}

\begin{lstlisting}[style=cppstyle, caption={反例与修法：size\_t 循环、有符号与无符号比较}]
#include <cstddef>
#include <vector>

std::size_t bad_forward(const std::vector<int> &v) {
    std::size_t n = 0;
    for (std::size_t i = 0; i <= v.size() - 1; ++i) {   // v 为空：size()-1 回绕成极大值
        if (v[i] > 0) ++n;                              // 下标一路越界
    }
    return n;
}

bool good_at(std::size_t i, const std::vector<int> &v) {
    return i < v.size() && v[i] > 0;                    // 先判空，比较不做差值
}

bool bad_cmp(int x, unsigned u) {
    return x < u;                       // x 被转成无符号：负数变成极大值
}

bool good_cmp(int x, unsigned u) {
    return x >= 0 && static_cast<unsigned>(x) < u;      // 先判符号，再显式转换
}

int bad_reverse(const std::vector<int> &v) {
    for (int i = v.size() - 1; i >= 0; --i) {           // size_t 窄化到 int，随后
        if (v[i] < 0) return v[i];                      // i>=0 靠“有符号”才成立
    }
    return 0;
}
\end{lstlisting}

最后一行值得多说一句：如果那个 \Tp{i} 声明成 \Tp{std::size\_t}，
\Tp{i >= 0} 就恒真，循环只能靠越界访问退出——那是另一种 \UB{}
（见第~\ref{ch:ub05-pointer} 章）。倒着遍历带无符号长度的容器，正确写法是把条件写成
“先减一再比较”的 \Tp{i-- > 0} 惯用法，或者改用 \Tp{std::vector} 的反向迭代器、
\Tp{std::span} 视图，让长度换算根本不出现。

\begin{table}
\centering
\begin{tabularx}{0.78\textwidth}{@{}>{\raggedright\arraybackslash}p{0.22\textwidth} >{\raggedright\arraybackslash}p{0.22\textwidth} >{\raggedright\arraybackslash}X@{}}
\toprule
表达式 & 谁被转换 & 建议的显式化写法 \\
\midrule
\Tp{int < unsigned} & \Tp{int} 转无符号 & 先判符号，再 \Fn{static\_cast} \\
\Tp{size\_t < int} & \Tp{int} 转 \Tp{size\_t} & 把长度先取成 \Tp{size\_t} 局部量 \\
\Tp{char + char} & 双方提升到 \Tp{int} & 需要窄结果就显式转回并判界 \\
\Tp{small u8 * 3} & 提升到 \Tp{int} & 无符号“回绕保证”此时并不适用 \\
\bottomrule
\end{tabularx}
\caption{混合类型比较与算术的转换归属一览}
\label{tab:ub09-convert}
\end{table}

\begin{guideline}[title={转换要么隐式无害，要么显式写清}]
凡是涉及有符号与无符号混合、或者宽度变化的比较，把它写成显式的
\Fn{static\_cast}，让读者在源码里就能看到“这里有一次转换、这里有判界”。
开启 \Tp{-Wsign-{\Bk}compare} 与 \Tp{-Wconversion} 这一类警告，并把它们
当作错误，比任何事后检查都便宜。
\end{guideline}

\section{窄化转换}\label{sec:ub09-narrow}

把宽类型写进窄类型，语义是按目标类型取值：无符号一侧按模 $2^N$ 截断（良定义），
有符号一侧超范围时归 \IDB{}；浮点到整型、且值不能被目标类型表示时是 \UB{}，
这一点常被误当成“截断”，实际它比截断更糟。

\begin{compare}[title={拷贝初始化与列表初始化的分工}]
\Tp{int i = 300u;} 是一次合法的显式转换链，编译器最多给警告；
\Tp{int i\{300u\};} 在 \cpp{11} 起被判为 \ILLF{}，标准库用“窄化转换”一词定义了这条禁令，
要求诊断。列表初始化不是万能防线：常量表达式在不损失精度时允许（例如
\Tp{int8\_t c\{127\};}），而运行期的值只要落在范围内也一样能通过。
\end{compare}

编译期防线还有 \Hdr{type\_traits} 的萃取一族：\Tp{is{\Bk}\_convertible{\Bk}\_v}
与 \Tp{is{\Bk}\_same{\Bk}\_v} 能把“能不能隐式转”变成编译期断言，
\Tp{common{\Bk}\_type{\Bk}\_t} 替你把“两个类型一起运算时落到哪个类型”变成可打印的结论，
\Tp{make{\Bk}\_unsigned{\Bk}\_t} 则把符号性直接写进类型名。
需要按位搬运而非数值转换时用 \Tp{std::bit{\Bk}\_cast}（\cpp{20}，\Hdr{bit}），
它给的是重新解释而不是数值转换，别把它当窄化工具用。

\section{除法与移位}\label{sec:ub09-div-shift}

整型除法的两条 \UB{}：除数为零；以及商不能被结果类型表示。第二条的经典实例是
\Tp{INT\_MIN / -1}——它不是因为“除以负数”而出事，而是因为商
$2^{31}$ 恰好落在 \Tp{int} 可表示范围之外。取余运算同样受“除数不得为零”的约束；
\CL{} 一侧把 \Tp{INT\_MIN} 对 \Tp{-1} 取余历史上列为未定义，\CPP{} 一侧则把商与
余数绑在同一个关系式上：被除数、除数、商与余数要满足
“商乘除数加余数等于被除数”，商一旦未定义，余数也就无从谈起。
需要“绝不在除法上翻车”的场合，把判界写在除法之前，别指望平台会给出陷阱。

移位这一侧的类别归属更细，而且版本敏感：

\begin{itemize}
  \item 移位量为负、或移位量不小于提升后左操作数的位宽：\UB{}。
        某些版本还把负移位量进一步收紧为 \ILLF{NDR}，请按所用版本的
        \Std{expr.shift} 原文核对，本册不代下断言。
  \item 左移有符号负值：\cpp{20} 之前是 \UB{}；\cpp{20} 起标准改为按结果类型
        位宽取模的良定义值。跨版本移植的代码里，不要靠“左移负数”表达任何语义。
  \item 右移有符号负值：\IDB{}（多数实现是算术右移，即复制符号位），
        不是 \UB{}。想要确定行为就显式转成无符号再右移，或者用带符号的除法。
\end{itemize}

\begin{warning}[title={别把“除二向下取整”写成“右移一位”}]
对负数，\Tp{x >> 1} 与 \Tp{x / 2} 在多数实现上不等价：前者向负无穷方向取整，
后者向零取整。两者都有定义时，它们是两种不同答案；这就是本册反复强调的
“\IDB{} 不等于 \UB{}，但你仍然不能假设跨平台一致”。
\end{warning}

\section{枚举的底层类型}\label{sec:ub09-enum}

无固定底层类型的枚举，其底层类型是 \IDB{}：实现要选一个足以容纳所有枚举数的
整型，选 \Tp{int} 还是更大的类型、有没有可能落到 \Tp{unsigned}，标准只写
“实现定义”。这意味着 \Tp{sizeof(Enum)} 与枚举能否塞进某个字段的结论，
在 \Tp{-m32} 与 \Tp{-m64} 下都可以不同。想要确定宽度，就写固定底层类型：
\Tp{enum class Kind : std::uint8{\Bk}\_t}（\cpp{11} 起）。

向枚举赋一个不在枚举数集合内的值，判定要看版本与是否有固定底层类型：
早期文本用“值的范围（range of an enumeration）”这一概念给出可容纳边界，
\cpp{20} 起相关措辞被调整。实践结论只有一条：不要靠“把整数转成枚举再读”
来携带数据，要么用 \Tp{std::to{\Bk}\_underlying}（\cpp{23}，\Hdr{utility}）明确地拿到整数值，
要么换成 \Tp{std::underlying{\Bk}\_type{\Bk}\_t} 加显式转换，把转换点写在源码上。

\begin{note}[title={作用域枚举救不了溢出}]
\Tp{enum class} 关掉的是隐式提升与隐式转换，它不改变算术溢出的类别归属。
对 \Tp{enum class} 的底层值做加法照样可能触发第一节的 \UB{}。
\end{note}

\section{数据模型与整数宽度}\label{sec:ub09-models}

同一份 C/C++ 源码在不同平台上“\Tp{int} 多宽、\Tp{long} 多宽”的答案不一样，
业界用三个缩写概括三种约定。机制细节（ABI、链接、序列化格式）属于
编译器与链接器那一册，这里只留语言侧的必要结论。

\begin{longtable}{@{}>{\raggedright\arraybackslash}p{0.14\textwidth} >{\raggedright\arraybackslash}p{0.14\textwidth} >{\raggedright\arraybackslash}p{0.16\textwidth} >{\raggedright\arraybackslash}p{0.28\textwidth}@{}}
\caption{整数数据模型与常用类型的宽度（典型值，具体实现以文档为准）}\label{tab:ub09-widths}\\
\toprule
模型 & \Tp{int} & \Tp{long} & 语言侧要记住的差别 \\
\midrule
\endfirsthead
\toprule
模型 & \Tp{int} & \Tp{long} & 语言侧要记住的差别 \\
\midrule
\endhead
\bottomrule
\endfoot
ILP32 & 32 & 32 & 32 位环境；\Tp{size\_t} 与指针同为 32 位 \\
LP64 & 32 & 64 & \Tp{long} 变宽，\Tp{int} 与 \Tp{long} 的隐式窄化在两个模型间不对称 \\
LLP64 & 32 & 32 & \Tp{long} 保持 32 位，需要 64 位必须写 \Tp{long\_long} 或 \Tp{int64\_t} \\
\end{longtable}

由这张表引出三条硬约束。第一，需要精确宽度时用 \Hdr{cstdint} 的
\Tp{int32\_t}、\Tp{int64\_t}；它们是可选设施，只有实现提供时才存在，
探测不到时标准还给了 \Tp{int\_least64{\Bk}\_t} 一族兜底。第二，需要“能装下指针”的整数时，
\Tp{intptr\_t} 同样是可选的，而 \Tp{std::bit{\Bk}\_cast} 更合适。第三，跨 LP64 与 LLP64
共享的持久化格式，绝对不要把 \Tp{long} 当作 8 字节来写；这是本册反复出现的
“把 \IDB{} 当契约”的典型翻车现场。

\section{案例：无符号差值把循环变成近无限循环}\label{sec:ub09-case}

下面这个骨架来自真实的调试经验，但它被剥离成最小复现形式，因此不含任何实测数据。
发送进度用 \Tp{unsigned} 表示，循环条件写成“还差多少个”：

\begin{casebox}[title={回绕的差值与带 \Tp{ptrdiff\_t} 的改写}]
\begin{lstlisting}[style=cppstyle, caption={反例：用无符号差值当循环条件}]
unsigned chunk();                       // 每轮实际发送量

void pump(unsigned sent, unsigned total) {
    while (total - sent > 0) {          // (1) 无符号：> 0 等价于 != 0
        sent += chunk();                // (2) 一旦 sent 越过 total，
        record(sent);                   //     (1) 的差值回绕成巨大值，循环近无限
    }
}
\end{lstlisting}

问题出在两处：\Tp{unsigned} 上的 \Tp{> 0} 没有“还剩一些”的含义，只有“非零”；
而 \Tp{total - sent} 在被减数更小时按 $2^N$ 回绕，标准保证它给出一个巨大值——
这是良好定义造成的死循环，不是 \UB{}。改写要做三件事：先把两个长度都放进
有符号且有界的类型，再显式判空/判倒退，最后用有符号差值驱动循环。

\begin{lstlisting}[style=cppstyle, caption={改写：显式判界与 \Tp{ptrdiff\_t} 差值}]
void pump_safe(std::size_t sent, std::size_t total, std::size_t step) {
    if (sent >= total) return;                       // (1) 显式判空与判倒退
    const std::ptrdiff_t rem =
        static_cast<std::ptrdiff_t>(total) -
        static_cast<std::ptrdiff_t>(sent);           // (2) 两者必须先在可表示范围内
    for (std::ptrdiff_t i = 0; i < rem; i += static_cast<std::ptrdiff_t>(step)) {
        record(sent + static_cast<std::size_t>(i));  // (3) 回到无符号只做加法，不做起点相减
    }
}
\end{lstlisting}

步骤 (2) 是整段改写的真正前提：\Tp{std::size\_t} 到
\Tp{std::ptrdiff\_t} 的转换在值超出有符号范围时并不安全，
\cpp{20} 之前写作实现定义，\cpp{20} 起随补码表示的统一而收紧口径，
具体结论请按所用版本的 \Std{conv} 核对。工程上的稳妥做法是：
凡是要用有符号类型表达“数量差”，先在接口上用 \Tp{assert} 或前置条件把长度
限制在 \Fn{std::numeric{\Bk}\_limits} \Tp{<std::ptrdiff{\Bk}\_t>::max()} 之内，
把宽度问题从求值阶段搬到契约阶段。
\end{casebox}

\subsection{本章检查清单}\label{sec:ub09-checklist}

\begin{itemize}
  \item 有符号算术是否都能给出上界证明？不能的话，走
        第~\ref{sec:ub09-safe} 节的三条出口之一，别用开关代替论证。
  \item 所有与 \Tp{size\_t} 交互的循环变量，类型是否统一、比较是否不做差值？
  \item 窄化点是否显式写了 \Fn{static\_cast}，编译警告是否已升级为错误？
  \item 除法之前是否判过零与商的范围？移位量是否保证严格小于位宽？
  \item 依赖枚举宽度、\Tp{long} 宽度的持久化格式，是否改用定宽类型？
\end{itemize}

\part{并发与库契约}
\chapter{内存模型与 data race}\label{ch:ub10-memory-model}

这一章单独成章，理由只有一个：在多线程里，“未定义行为”这个词失去了它熟悉的形状。
单线程的 \UB{} 通常能指着某一行说“这一行算错了”；而两个线程各自写下的每一行单看都合法，
出错的是它们之间的相对顺序——标准必须先把“顺序”定义成一套数学对象，
才谈得上说某个执行是否被允许。第~\ref{ch:ub01-taxonomy} 章的五分类在这里拿到了它最
重要的用武之地：竞态不是“读到旧值”，而是标准撤回承诺。

\section{\cpp{11} 把并发写进语言定义}\label{sec:ub10-terms}

\cpp{03} 及以前的 C++ 对多线程几乎一无所说是事实：没有内存模型，没有原子类型，
标准库里连一个线程类型都没有，程序能否并发完全取决于实现愿不愿意。
\cpp{11} 把这套契约补上了，代价是读者必须先掌握四个术语。

\begin{itemize}
  \item \Mx{modification order}：对同一个对象，所有写操作被要求排成一个全序。
        “谁最后写入”在标准里从此是一个有定义的问题。
  \item \Mx{happens-before}：跨线程的执行约束关系。若 A 早于 B，B 就不能被搬到 A 之前，
        读者也就能看到 A 的写入。日常讲法沿用这个词；现行文本把它拆成了
        同线程顺序、线程间同步边与依赖有序边等若干条边的组合，
        精确措辞请以 \Std{intro.races} 与 \Std{atomics} 的当前版本为准。
  \item \Mx{synchronizes-with}：建立 happens-before 的主要手段——一次原子写与一次
        原子读之间的配对关系，由内存序决定它是否成立。
  \item \Mx{release sequence}：解释“同一个原子量上，别的线程夹在中间的写入会不会
        断开同步链”的概念。它是并发论证里最容易被忽略的一环，
        日常结论可以记成：读者与同一线程的读改写不会打断链；一旦同一个原子量上
        掺入别的线程的普通 store 或 \Tp{consume} 读，就要回到 \Std{atomics.order}
        重新论证这条链是否延续。
\end{itemize}

\begin{clause}[title={标准怎么说：竞态的三个条件}]
标准用“冲突（conflict）”描述同一内存位置上的两个操作：至少一个是写。
当两个冲突操作分属两个线程、其中至少一个不是原子操作、并且二者之间没有
happens-before 关系时，程序含有 data race，行为未定义。
\end{clause}

\section{data race 的精确定义与它的后果}\label{sec:ub10-race-def}

定义里的三个条件值得逐个抠。第一个是“冲突”：两个操作访问同一个内存位置，
至少一个是写；两个线程并发读同一个对象不冲突，因此不是竞态。
第二个是“至少一个非原子”：两边都是原子操作时，标准把它们纳入原子语义，
不落入竞态条款——这就是为什么“把变量换成 \Tp{std::atomic}”能治愈竞态而不只是缓解。
第三个是“没有 happens-before 关系”：这是唯一真正决定对错的条款。
只要两个冲突操作之间有同步边（锁、原子配对、线程的建立与加入），
哪怕时间上看起来同时在跑，程序仍然合法。

\begin{lstlisting}[style=cppstyle, caption={最小竞态：两个线程对同一个普通对象做读改写}]
#include <thread>
unsigned long g_counter = 0;              // 普通对象，不是原子

void bump() {
    for (int i = 0; i < 100000; ++i) {
        ++g_counter;                      // 读、加、写三步，两个线程各做十万次
    }
}

int main() {
    std::thread a(bump), b(bump);
    a.join();
    b.join();
    sink(g_counter);                      // 标准没有承诺这里能看到什么
}
\end{lstlisting}

\Tp{++g\_counter} 至少含一个写，两侧都不是原子操作，两条线程的对应操作之间
也没有任何 happens-before 边——三个条件齐了。此处最容易误解的是后果：
它不是“少加了几次”，而是整个程序进入未定义状态。编译器有权假设你的程序里没有竞态
（否则无从优化），于是可以把这个变量缓存在寄存器、把读提到循环外、
甚至把看起来多余的检查删掉。第~\ref{ch:ub09-integers} 章讲有符号溢出时是同一条逻辑：
标准撤回承诺，优化器拿走前提。

\begin{warning}[title={“结果只是旧值”是最危险的自我安慰}]
把竞态当作“可能读到稍微旧一点的值”，等于承认“这段代码偶尔不对但不会更糟”。
真实案例里竞态会表现为读到未初始化字节、指针撕裂（一半新地址一半旧地址）、
循环永不退出、以及只在发布构建复现的崩溃。若一个变量可能被另一个线程改写，
它需要的不是运气，而是一条同步边。
\end{warning}

\section{\Tp{std::atomic} 与六种内存序}\label{sec:ub10-order}

\Hdr{atomic} 提供 \Tp{std::atomic} 主模板与定宽整型的若干特化，
每个操作都带一个默认为 \Tp{std::memory\_order\_seq\_cst} 的内存序参数。
枚举里有六个值：\Tp{relaxed}、\Tp{consume}、\Tp{acquire}、\Tp{release}、
\Tp{acq\_rel}、\Tp{seq\_cst}。读改写操作（\Tp{fetch\_add}、\Tp{compare\_exchange}
一族）取的是读写合一的档位，\Tp{acq\_rel} 正是为它们准备的。

\begin{longtable}{@{}>{\raggedright\arraybackslash}p{0.16\textwidth} >{\raggedright\arraybackslash}p{0.26\textwidth} >{\raggedright\arraybackslash}p{0.22\textwidth} >{\raggedright\arraybackslash}p{0.14\textwidth}@{}}
\caption{六种 memory order 各保证什么、不保证什么}\label{tab:ub10-orders}\\
\toprule
档位 & 保证 & 不保证 & 典型用途 \\
\midrule
\endfirsthead
\toprule
档位 & 保证 & 不保证 & 典型用途 \\
\midrule
\endhead
\bottomrule
\endfoot
\Tp{relaxed} & 单个操作的原子性，以及同一原子量的修改顺序一致 & 不建立同步边，周围读写可任意跨越 & 计数器、引用计数、位图标志 \\
\Tp{acquire} & 读到某次 release 写之后，该写之前的全部写入都可见 & 不阻止本线程后续写入被别人提前看到 & 读标志、取已发布的数据 \\
\Tp{release} & 本操作之前的全部写入，对配对的 acquire 读者可见 & 单独不构成边，必须有人 acquire & 写标志、发布对象 \\
\Tp{acq\_rel} & 用于读改写：读侧 acquire 与写侧 release 同时成立 & 只是同步边，不是全局顺序 & CAS 循环、无锁栈弹出 \\
\Tp{seq\_cst} & 以上之外，所有顺序一致操作共享一个全局修改顺序 & 不保证实时性，也不免锁 & 默认档；跨多个原子量的推理 \\
\Tp{consume} & 名义上只在依赖链上可见 & 至今没有可用的宽松语义，实际按 acquire 处理 & 不要用 \\
\end{longtable}

日常并发的默认选择应当是：需要跨对象传递信息时保持默认（\Tp{seq\_cst}），
只有被证明是热点、并且能写出同步边时才降级到 acquire/release。
“为了性能先写 relaxed”是本章反对的第一种写法：它省下的指令很少，
付出的却是整份正确性论证。多原子量之间的全局一致顺序是 \Tp{seq\_cst} 独有的承诺，
一旦降级就再也补不回来。

\section{同步边是怎么建立的}\label{sec:ub10-publication}

发布一个对象，靠的是“先写数据，再用一次 release 写发布指针；读者先 acquire，再碰数据”。

\begin{lstlisting}[style=cppstyle, caption={acquire/release 配对：普通数据的合法发布}]
#include <atomic>
#include <thread>
int g_data = 0;                             // 普通对象
std::atomic<int> g_flag{0};

void producer() {
    g_data = 42;                                       // (1) 普通写，尚无同步
    g_flag.store(1, std::memory_order_release);        // (2) release 写
}

void consumer() {
    if (g_flag.load(std::memory_order_acquire) == 1) { // (3) acquire 读，与 (2) 配对
        use(g_data);                                   // (4) 合法：(1) happens-before (4)
    }
}

int main() {
    std::thread t1(producer), t2(consumer);
    t1.join();                            // (5) t1 完成 happens-before join 返回
    t2.join();
}
\end{lstlisting}

如果 (2) 换成 \Tp{relaxed}，(3) 就与它构不成同步边，(4) 读 \Tp{g\_data} 与 (1) 写
\Tp{g\_data} 立刻成为一对无 happens-before 的冲突操作：竞态回来了，而且这一回
连“平台会替我保持顺序”的借口都没有——x86 上写与写的硬件顺序确实存在，
但编译器重排不受硬件约束。这就是“靠平台特性写并发”的标准教训。

\section{线程的建立、加入与分离}\label{sec:ub10-thread}

线程对象自己也是同步边的来源，这一条常被漏掉：

\begin{itemize}
  \item \Tp{std::thread} 构造：构造调用 happens-before 新线程执行其入口函数。
        因此在创建线程之前写好的共享数据，新线程读它是合法的。
  \item \Fn{join}：被 join 的线程完成执行 happens-before \Fn{join} 返回。
        这是一条完整的收口，也是绝大多数一次性并行任务真正需要的同步。
  \item \Fn{detach}：不留任何同步边。分离后的线程仍在跑，而它捕获的栈对象随时可能
        先它一步销毁——那是第~\ref{ch:ub04-lifetime} 章的生命周期问题，
        只是这里没有别的线程替你把它抓出来。
  \item \Tp{std::jthread}（\cpp{20}，\Hdr{thread}）：析构时自动 \Fn{join}，并附带一个
        停止令牌用于协作取消。它还顺手治好了“可 join 的 \Tp{std::thread} 析构会
        触发 \Fn{terminate}”这条老路（见第~\ref{ch:ub11-terminate} 章）。
\end{itemize}

\begin{guideline}[title={四条不需要内存序知识的并发写法}]
第一，能不共享就不共享：每个线程写自己的结果，\Fn{join} 之后主线程读。
第二，必须共享就加锁：同步边由互斥量免费提供，见下一节。
第三，只传递一个“就绪”信号时，用原子标志配 acquire/release。
第四，绝不在 \Fn{detach} 的线程里按引用捕获栈对象。
\end{guideline}

\section{非原子共享状态的合法处理}\label{sec:ub10-mutex}

把共享数据留在锁里，是最省推理的做法：\Tp{std::mutex} 的 \Fn{lock} 与 \Fn{unlock}
自带同步语义——前一次持有该锁期间对共享数据的写入，happens-before 后一次成功获得
同一把锁之后的访问。也就是说，只要所有访问都被同一把锁保护，
竞态的三个条件就凑不齐。

\begin{itemize}
  \item \Fn{std::lock\_guard}、\Fn{std::unique\_lock}、\Fn{std::scoped\_lock} 与
        \Tp{std::shared\_mutex} 都是同一对边的包装，别手写 \Fn{lock} 与 \Fn{unlock} 配对。
  \item \Tp{std::call\_once} 与 \Tp{std::once\_flag}（\Hdr{mutex}）：若干线程调用同一个
        \Fn{call\_once}，被选中那一次的副作用 happens-before 其余各次返回。
        它给的是“初始化一次”的同步边，不是对后续写入的保护。
  \item 局部静态对象：\cpp{11} 起线程安全的初始化写进了语言（俗称 magic statics），
        初始化只发生一次并且对后来者可见。用它做单例，比任何手写双检锁都短。
  \item 只读后发布：启动阶段填好、之后只读的数据，配合一条同步边
        （线程创建、\Fn{join}、release 写）发布，就不需要锁。
\end{itemize}

无锁结构不在这张豁免名单上。它只是把锁换成手工构造的同步边，
因此必须先证明每一次访问都在某条 happens-before 边之后：先证明无竞态，再谈优化。
跳过这一步的无锁代码，正确性只能靠运气维持。
标准库这一侧同样有契约：库函数不得造成数据竞态，也不得去碰没有通过参数交给它的对象，
展开的约定见第~\ref{ch:ub12-library-pre} 章。

\begin{compare}[title={原子、互斥与 volatile 的分工}]
原子量解决“单个位置上的读改写不撕裂，并按选定档位建立同步边”；
互斥量解决“一段临界区里多个位置的更新一起可见”；
\Tp{volatile} 解决“编译期不要优化掉这一次访问”。三者不可互换。
\end{compare}

\section{\Tp{volatile} 不解决线程可见性}\label{sec:ub10-volatile}

\Tp{volatile} 的语义是：对该对象的访问不得被省略，也不得在两次访问之间被随意合并或重排，
目的是让代码配合“编译器看不见的第三方修改者”工作。它的合法用户是两类：
映射到设备寄存器的内存；以及在异步上下文中被改写的对象——例如信号处理函数里写、
主逻辑读的 \Tp{sig\_atomic\_t}。\Hdr{csetjmp} 的 \Fn{setjmp} 语境也是这条用途线的一部分：
被 \Fn{longjmp} 回退穿越的非 \Tp{volatile} 局部对象，其值不可靠，
这正是标准把 \Tp{volatile} 留在这两条路上的原因。

\begin{warning}[title={不要把 volatile 当同步}]
\Tp{volatile} 不建立 happens-before，不提供原子性，也不保证内存顺序。
它甚至不是写给多线程的：两个线程访问同一个 \Tp{volatile} 非原子对象，
仍然是本章第二节那三个条件齐备，仍然是 \UB{}。
“用 \Tp{volatile} 修饰共享标志”是上一个时代某些语言的写法，在 \CPP{} 里它是错误答案。
对 \Tp{volatile} 限定类型做原子操作在 \cpp{20} 起被标为不推荐，后续版本口径更紧，
不要把两件事混起来用。
\end{warning}

\section{工具口径：TSan 抓什么、抓不到什么}\label{sec:ub10-tsan}

竞态检测靠插桩：编译器把每一处内存访问改写成可报告的调用，
运行库维护逻辑时钟，据此判断两个冲突访问之间有没有 happens-before 边。

\begin{lstlisting}[style=shellstyle, caption={ThreadSanitizer 的最小可用命令}]
# 整程序都要插桩：只要有一条链路没带 -fsanitize=thread，报告面就不完整
clang++ -std=c++20 -g -O1 -fsanitize=thread -fno-omit-frame-pointer race.cpp -o race
gcc -std=c++20 -g -O1 -fsanitize=thread race.cpp -o race
TSAN_OPTIONS="halt_on_error=1 second_deadlock_stack=1" ./race
# 报告形状（示意，非本仓库实测）：
# WARNING: ThreadSanitizer: data race on address 0x...
#   Read of size 4 by thread T2 / Previous write by thread T1
\end{lstlisting}

TSan 抓得到的是：这次运行里真实发生、且两条路径都被插桩覆盖的冲突访问对。
抓不到的至少包括：依赖具体数据才会走到的分支上的竞态（覆盖率决定成败）、
未被插桩的第三方库或汇编里的访问、“没有竞态但顺序错了”的原子性缺陷
（两个字段分开更新被读者看到中间状态，这完全合法，却是真实缺陷）、
以及需要特定交错才出现的窗口——工具只观察它跑到的那几种交错。
报告还会被内存序选择影响：代码大量降级到 \Tp{relaxed} 时，
工具不再能推断你本来想要的同步边，看起来反而更“干净”。

\begin{note}[title={TSan 全绿不等于没有竞态}]
第~\ref{ch:ub03-sanitizers} 章的 CI 分档结论在这里同样适用：消毒器证明的是
“这一轮执行没有暴露”，不是“所有执行都安全”。并发代码要额外做三件事：
在压力与随机调度下反复跑、改变并行度（例如用 \Tp{taskset} 限制核数）、
以及把每个共享变量“由谁写、在哪条同步边之后读”写进注释或类型设计里。
\end{note}

\section{划界}\label{sec:ub10-scope}

\begin{itemize}
  \item 并行算法册只讲执行策略（execution policy）的用法与适用条件；
        本章只关心它们背后“同一位置能否被并发访问”的语言契约。
  \item 原子操作在某个目标上如何落地（用哪条指令、是否无锁、
        与 \Tp{std::atomic\_ref} 的配合）属于编译器与目标代码册。
  \item 本章只讲语言契约：哪些执行被标准允许；“允许”意味着每个实现都必须给出
        允许集合中的一个，而不是任意的一个。
\end{itemize}

\section{案例：double-checked locking}\label{sec:ub10-dcl}

双检锁是并发章节绕不过去的例子，因为它同时踩中本章三个要点：
非原子的读与写、缺失的同步边、以及“加了锁就以为自己安全”的错觉。

\begin{casebox}[title={双检锁：裸指针版本与原子版本}]
\begin{lstlisting}[style=cppstyle, caption={两个版本：(1) 竞态版，(2) 正确版}]
struct Config { int port; };
std::mutex g_mtx;

// 版本一：教科书级的 data race
Config *g_cfg = nullptr;
Config *bad_dcl() {
    if (g_cfg == nullptr) {                      // (1) 非原子读
        std::lock_guard<std::mutex> lk(g_mtx);   // (2) 加锁
        if (g_cfg == nullptr) {                  // (3) 非原子读，仍在锁内
            Config *p = new Config();            // (4) 普通写：初始化对象
            g_cfg = p;                           // (5) 非原子写：发布指针
        }
    }
    return g_cfg;                                // (6) 非原子读，在锁外
}

// 版本二：原子指针 + release/acquire
std::atomic<Config *> g_ok{nullptr};
Config *good_dcl() {
    Config *p = g_ok.load(std::memory_order_acquire);   // (1') 原子读
    if (p == nullptr) {
        std::lock_guard<std::mutex> lk(g_mtx);          // (2') 加锁
        p = g_ok.load(std::memory_order_relaxed);       // (3') 原子读，锁已给边
        if (p == nullptr) {
            p = new Config();                           // (4') 普通写，尚未发布
            g_ok.store(p, std::memory_order_release);   // (5') 原子写：发布
        }
    }
    return p;                                  // (6') 返回本地副本
}
\end{lstlisting}

\textbf{版本一的逐行判定。}(1) 与 (6) 是非原子的读，(5) 是非原子的写，它们访问同一个
内存位置且至少一个是写——冲突成立。关键在于有没有 happens-before：
(2) 的锁只保护走到慢路径的那几个线程，(1) 与 (6) 都在锁外执行，与 (5) 之间没有边。
三条件齐备：程序含 data race，后果是 \UB{}，而不是“偶尔读到 \Tp{nullptr}
于是多进一次锁”。即便硬件把指针写做成不可分割，(4) 与 (5) 之间仍然可能被编译器重排，
读者完全可能拿到一个非空但成员还没写完的指针。

\textbf{版本二的逐行判定。}(1') 是原子读，(5') 是原子写，二者在同一个原子对象上：
只要 (1') 观察到 (5') 写入的值，就构成 release/acquire 配对，(4') 对 \Tp{Config}
成员的普通写入因此 happens-before (1') 之后的任何使用。慢路径里的 (2') 与 (3')
由互斥量的解锁/加锁边兜住，那里用 \Tp{relaxed} 也够。(6') 返回本地副本，
不再触碰共享量。两个细节必须盯住：第一，返回裸指针要求对象永不销毁，
否则读侧还得配合生命周期规则；第二，(5') 若写成 \Tp{relaxed}，
(1') 就再也配不成边，版本二立刻退回版本一的处境。

\textbf{更短的第三条路。}把单例写成函数内的局部静态对象，
用语言的线程安全初始化替代整套手写双检锁；或者在启动阶段构造好，
用一条同步边发布后只读。正确性论证的长度，本身就是本章的取舍标准。
\end{casebox}

\subsection{本章检查清单}\label{sec:ub10-checklist}

\begin{itemize}
  \item 每个共享变量能否回答三个问题：谁能写、谁能读、读写之间靠哪条同步边分隔？
  \item 所有非原子共享状态是否都在同一把锁的保护下？锁外有没有读或写？
  \item 内存序降级是否有明确论证与被测量过的收益？默认档有没有被顺手改掉？
  \item \Fn{detach} 的线程是否按引用捕获了栈对象？每个 \Tp{std::thread} 是否都 \Fn{join} 过？
  \item 除了 TSan 报告，是否补了压力运行与不同并行度下的复跑？
\end{itemize}

\chapter{自我移动、\Tp{noexcept} 与 terminate 的边界}\label{ch:ub11-terminate}

前两章把镜头对准了求值与内存模型：第~\ref{ch:ub09-integers} 章关心一个算术表达式的边界，第~\ref{ch:ub10-memory-model} 章关心两个线程
之间的可见性。本章把镜头推到程序的两个终局上：一个是\textbf{对象自己移动自己}，另一个是
\textbf{进程被终止函数收场}。这两件事的共同点在于，标准给的都不是提示而是承诺：
承诺被兑现时程序有意义，承诺被撤回时程序没有下一步。

本章的判定口径始终只用五个类别宏：\UB{}、\IDB{}、\USB{}、\NDR{}、\ILLF{}。凡涉及具体实现
的行为（虚表桩、\Fn{abort}、报告格式），都会显式标注为实现行为而非标准承诺。

\section{自我移动：标准到底承诺了什么}\label{sec:ub11-selfmove}

标准对库类型的移动赋值通常写成这样的效果描述：把源对象的资源转移给目标对象，源对象进入
\textbf{有效但未指定}的状态。请注意这句话的主语：它假设了源与目标是两个对象。当
\Tp{x = std::move(x)} 出现时，这句话的自变量落在自己身上，标准文本并没有为它准备单独的
承诺。于是形成三种口径并存的局面：

\begin{itemize}
  \item 多数容器与字符串：移动后处于 \USB{} 的有效状态；但自移动本身不在该句的覆盖范围内，
        实现可以额外保证，也可以不保证。
  \item 少数类型与少数重载：标准显式补了自移动与自赋值的规定，或显式给出前置条件把重叠排除掉。
        典型的是字符串的裸指针 \Fn{assign} 重载要求源区间不与目标缓冲区重叠，
        违反即越出承诺范围。
  \item 只移动语义的类型（\Tp{unique\_ptr} 一类）：自移动在多数实现里表现为先释放再接管，
        结果是对象自毁。这未必被标准判为 \UB{}，但等价于一个确定的资源错误，属于实现行为。
\end{itemize}

因为口径不统一，靠记忆去判断某个类型的自移动是否安全是不可行的。可行的做法是用一条与类型
无关的规则把它替代掉。

\subsection{可机械执行的三条规则}\label{sec:ub11-selfmove-rules}

把 \Tp{x = std::move(x)} 之后、对 \Tp{x} 的每一次使用都按下列三条之一归类，归类失败的就不写：

\begin{enumerate}
  \item \textbf{重新赋值}：任何形式的赋值（含 \Fn{assign}、\Fn{emplace}、\Fn{reset}）之后，
        \Tp{x} 重新回到自己有文档的状态，此后按普通对象对待。
  \item \textbf{析构}：局部对象的正常析构、容器元素的析构都在承诺范围内。
  \item \textbf{有复杂度或状态保证的成员}：\Fn{clear}、\Fn{swap}、\Fn{empty}、\Fn{size}
        这一类标准给出固定复杂度（常数或对数之外另有说明）的成员可以调用。
        但读到的值是 \USB{} 的，不得假定它等于零、等于空或等于移动前的值。
\end{enumerate}

其余一切（读旧值、比较旧值、把旧内部缓冲当作自己的资源再次转移）都不在规则内。

\begin{lstlisting}[style=cppstyle, caption={自移动的三种安全写法：自检查、先取副本、或压根不出现}]
struct Buf {
  std::vector<int> data;

  // 写法一：显式自检查。移动赋值的标准骨架本就该带这一句。
  Buf& operator=(Buf&& other) noexcept {
    if (this == &other) { return *this; }      // 自移动：什么都不做，值保持
    data = std::move(other.data);
    return *this;
  }
};

// 写法二：需要“可能的自我赋值”时先取副本，避免把源与目标当两个对象。
void reseed(std::vector<int>& v, std::vector<int> src) {
  v = std::move(src);        // src 是值形参，绝不与 v 同一对象
}

// 写法三：只读场景不要用移动表达意图。
void log_size(const std::vector<int>& v) { /* 接收 const 引用，不需要移动 */ }
\end{lstlisting}

\begin{guideline}[title={用规则替代记忆}]
不要在脑子里维护一张“哪些类型自移动安全”的表。标准没给这张表，实现也不会替你补。
把 \Tp{x = std::move(x)} 之后对 \Tp{x} 的使用限制在再赋值、析构与有保证的成员调用三件事上，
规则与类型无关，因而可以机械化检查。
\end{guideline}

\begin{table}
\centering
\begin{tabularx}{0.78\textwidth}{>{\raggedright\arraybackslash}p{0.22\textwidth} X >{\raggedright\arraybackslash}p{0.22\textwidth}}
\toprule
操作 & 移动后是否合法 & 说明 \\
\midrule
\Fn{clear} 与 \Fn{swap} & 合法 & 标准给出状态保证，调用后对象回到可文档化的状态 \\
\Fn{assign} 或 \Fn{emplace} & 合法 & 这是把对象拉回已知状态的唯一途径 \\
析构 & 合法 & 实现的析构路径不得依赖元素值 \\
\Fn{empty} 与 \Fn{size} & 合法，但值是 \USB{} & 返回值不由你决定，也不保证等于零 \\
\Fn{begin} 到 \Fn{end} & 合法，但区间内容是 \USB{} & 区间必须自洽、可遍历，长度不由你假定 \\
假定它等于空 & 不合法 & 标准写明的是有效而非空；实现可以在移动后保留缓冲 \\
读旧元素、比较旧值 & 不保证 & 一旦用于业务判断就成了跨实现漂移的根源 \\
把旧内部指针交给别人 & 危险 & 相当于把同一份资源转移两次，落在所有权错误一侧 \\
\bottomrule
\end{tabularx}
\caption{被移动对象（含自我移动后）的合法操作集合}\label{tab:ub11-moved-from}
\end{table}

\begin{warning}[title={自交换比自移动更容易被漏掉}]
\Fn{std::swap} 的效果被描述成“用移动构造一个临时量、再两次移动赋值”。当两个实参是同一个对象时，
这条算法中间就会出现一次自移动。因此自交换能否保住原值，实际上依赖移动赋值对自移动的态度。
标准对 \Fn{swap} 本身要求可交换，但对自定义类型的移动赋值没有替你补这一条保证，
自写移动赋值时把自检查当成骨架的一部分，比事后论证更安全。
\end{warning}

\section{\Tp{noexcept} 是承诺，不是提示}\label{sec:ub11-noexcept}

\Tp{noexcept} 出现在函数类型里，出现在调用点的编译期判定里，也出现在优化器的假设里。
它与 \Tp{const}、引用限定符一样是类型的一部分，这一点决定了它不能被当成“性能提示”随手加。

\subsection{违反承诺的几条路径}\label{sec:ub11-noexcept-paths}

被标记为 \Tp{noexcept} 的函数，只要有一次控制流越出它的边界，标准规定的结果就是调用终止函数。
触发点主要有：

\begin{itemize}
  \item 函数体内抛出的异常未被本地 \Fn{catch} 消化，沿边界外逃；
  \item \Tp{noexcept} 的移动构造或移动赋值被容器用于扩容，而它真的抛了；
  \item 析构函数在栈展开过程中再次抛出（见后面单独一节）；
  \item 直接调用终止函数，或调用一个显式 \Tp{noexcept} 却内含 \Fn{throw} 的回调；
  \item 未捕获异常一路走到线程与 \Fn{main} 的边界。
\end{itemize}

代价的定性描述只有一句：承诺换来的是编译器可以省掉一条展开路径，一旦违约，恢复点直接消失，
程序没有“捕获后再处理”的机会。具体省下多少代码量与运行时开销取决于实现与优化档位，
本册不给数字。

\begin{clause}[title={终止与未定义行为的两套规则}]
标准把终止的处理写成一组独立条文：未捕获异常、越出 \Tp{noexcept} 边界、栈展开期间再次抛出，
这三类都会走到终止函数。处理程序的安装接口是 \Fn{std::set{\Bk}\_terminate}，
读回接口是 \Fn{std::get{\Bk}\_terminate}，安装动作的返回值就是被替换掉的前一个处理程序。
主题名（\Std{exception.terminate} 一类）在不同草案里编号有变动，引用时以当前草案为准。
关键差别是：这一组条文给出的是\emph{有终局}的行为，而 \UB{} 那一组什么也没给。
\end{clause}

\begin{compare}[title={terminate 与 \UB{} 的三点差别}]
第一，结局确定性：终止会先走处理程序，实现默认再走 \Fn{abort} 一类收场，进程状态可预期；
\UB{} 之后程序的输出没有任何约束。第二，可观测性：终止可以挂处理程序，可以写日志、
可以触发转储；\UB{} 没有可挂的位置，只能靠前置检测（消毒器、断言）。
第三，可移植性：同一段违约代码在不同编译器上的终止路径一致，而 \UB{} 的表现随优化档位漂移。
因此“能被终止”反而比“落到 \UB{}”更好，这正是本章末尾三条写法的动机。
\end{compare}

\section{终止处理程序的用途与局限}\label{sec:ub11-terminate-handler}

终止处理程序相关接口位于异常支持设施中，签名固定为无参、返回 \Tp{void} 的函数指针。
它的合理用途只有一个：在进程即将结束之前做一次单向的取证输出。安装时的返回值给出
前一个处理程序，这一点让库可以临时接管、事后交还。

\begin{lstlisting}[style=cppstyle, caption={终止处理程序：只能报告，不能恢复}]
#include <exception>
#include <cstdlib>

namespace {
std::terminate_handler previous = nullptr;

[[noreturn]] void on_terminate() {
  // 这里能做的：写一条最后的日志、落一个标记文件、调用 abort。
  // 这里不能做的：抛异常（会递归触发终止）、试图继续执行业务流程、
  // 依赖任何共享状态（发生终止的原因也许正是一次数据竞争）。
  if (previous != nullptr) { previous(); }   // 交还实现原本的默认处理
  std::abort();
}
} // namespace

void install_terminate_probe() {
  previous = std::set_terminate(on_terminate);   // 返回安装前的处理程序
  // 局限一：处理程序是进程级全局，多线程下安装与卸载要自己同步。
  // 局限二：它不告诉你为什么终止；当前异常是否可得由实现决定，不可依赖。
  // 局限三：get_terminate 只用于读回，不做安装；别把它当查询机制滥用。
}
\end{lstlisting}

\begin{note}[title={别把终止处理程序当作兜底}]
它不参与恢复，也不参与资源回收：静态析构、智能指针的连锁释放、RAII 的清理都可能已经在
展开路径上跑过一半。真正需要“最后一次机会”的是取证，不是业务；业务需要的仍然是
不把控制流推到终止边界上。并发调用 \Fn{set{\Bk}\_terminate} 是在写进程级全局，
按第~\ref{ch:ub10-memory-model} 章的口径它需要外部同步，否则本身就是一次数据竞争。
\end{note}

\begin{lstlisting}[style=outstyle, caption={一次终止的报告形状（示意排版，非本仓库实测输出）}]
# 示意：越出 noexcept 承诺后的典型形状
terminate called after throwing an instance of 'std::bad_alloc'
  what():  std::bad_alloc
Aborted (core dumped)
# 对照：未定义行为往往连一行这样的报告都没有，程序只是给出了一个错的答案。
\end{lstlisting}

\section{析构函数的隐式承诺}\label{sec:ub11-dtor-noexcept}

从 \cpp{11} 起，析构函数默认就是 \Tp{noexcept}。更准确地说：用户未显式书写的析构函数是
隐式声明的，它的 \Tp{noexcept} 说明符由基类子对象与成员的析构所允许抛出的范围共同推导。
于是三条后果同时成立：

\begin{itemize}
  \item 析构里抛出异常，控制流越出 \Tp{noexcept} 边界，直接走终止函数；
  \item 栈展开过程中（异常正在寻找捕获点）析构再抛出，同样终止，这是双重异常的标准形态；
  \item 隐式移动构造与隐式移动赋值的 \Tp{noexcept} 同样由成员推导：
        一个成员用了不能不抛的移动，整个类型的承诺就没了。
\end{itemize}

\Tp{noexcept(false)} 可以为析构函数显式取消承诺，但取消之后调用方（尤其是容器与
\Fn{std::unique{\Bk}\_ptr} 的释放路径）必须能容忍展开，代价通常大于收益。

\section{\Tp{noexcept} 与容器扩容的移动要求}\label{sec:ub11-vector-growth}

\Tp{std::vector} 扩容时要搬运旧元素，标准给的算法选择条件是：当元素类型的不抛移动构造成立，
或者它压根不可复制时，用移动；否则退回拷贝。判定用的类型特征写作
\Tp{is{\Bk}\_nothrow{\Bk}\_move{\Bk}\_constructible}，实现侧对应的工具是
\Fn{move{\Bk}\_if{\Bk}\_noexcept}。这个条件解释了为什么一次忘记写 \Tp{noexcept}
的移动构造会让扩容悄悄退化成整段拷贝。

\begin{lstlisting}[style=cppstyle, caption={移动承诺决定扩容用搬还是拷；隐式承诺由成员推导}]
struct Old {
  std::vector<int> v;
  Old(Old&& o) /* 忘了 noexcept */ : v(std::move(o.v)) {}
  Old& operator=(Old&&) = default;
  Old(const Old&) = default;
  Old& operator=(const Old&) = default;
};                       // 有可复制性、又没有不抛移动 -> 扩容走拷贝

struct New {
  std::vector<int> v;
  New(New&& o) noexcept : v(std::move(o.v)) {}
  New& operator=(New&&) = default;
  New(const New&) = default;
  New& operator=(const New&) = default;
};

static_assert(std::is_nothrow_move_constructible_v<Old> == false);
static_assert(std::is_nothrow_move_constructible_v<New> == true);

// 自己的算法里表达同样的取舍：移动不安全就干脆别移动。
template <class T>
T pick(T&& t) { return std::move_if_noexcept(t); }

struct HasThrowingDtor {
  std::vector<int> v;
  ~HasThrowingDtor();               // 析构可能抛：类型整体不再默认不抛
};
static_assert(std::is_nothrow_destructible_v<New> == true);
\end{lstlisting}

\begin{warning}[title={三个会静默拆掉承诺的动作}]
其一，用户声明析构函数会抑制隐式移动的生成，于是容器扩容退回拷贝路径。
其二，把移动构造写成 \Tp{noexcept} 却在其中调用可能抛出的分配或第三方接口，
违约的后果是终止，不是回滚。其三，成员里混进一个不可复制、又不满足不抛移动的类型，
扩容会失去强异常保证的立足点。三者都不产生诊断，只在性能或终局上暴露。
\end{warning}

\section{其它终止与错误路径的归类}\label{sec:ub11-other-paths}

终止路径不止 \Tp{noexcept} 违约一条，而且不同路径在标准里的类别并不相同。把它们混为
“都会崩”是本章最常见的误读。

\begin{itemize}
  \item \textbf{未捕获异常}：沿调用链找不到捕获点，走到 \Fn{main} 或线程入口之外，
        标准规定调用终止函数。这是有终局的行为，不是 \UB{}。
  \item \textbf{纯虚调用}：构造或析构期间对象的多态身份停留在当前层，若这一层的目标函数是
        纯虚函数，标准把这种调用判为 \UB{}。实现侧常见做法是把虚表槽指向一个纯虚桩，
        桩内调用终止，于是“\UB{} 的实现行为恰好是终止”。别把后者当成前者。
  \item \textbf{构造期与析构期的向下转型}：这不是 \UB{}。标准明确规定此时操作数
        被视为当前类类型的对象，因此 \Tp{dynamic{\Bk}\_cast} 只认当前子对象层次：
        指针形式得到空指针，引用形式抛出 \Tp{bad{\Bk}\_cast}。
        抛出若发生在析构中，还可能升级为终止。
  \item \textbf{按下标取元素的两条路}：\Tp{std::tuple} 与 \Tp{std::array} 的下标是编译期常量，
        越界属于 \ILLF{}，压根编译不过去。\Tp{std::variant} 的类型或下标与当前持有的类型
        不匹配时抛出 \Tp{bad{\Bk}\_variant{\Bk}\_access}，
        这是定义良好的失败，不是 \UB{}。
  \item \textbf{\Fn{visit} 与无值 variant}：对处于“异常后无值”状态的 variant 调用 \Fn{visit}，
        标准规定抛出同一个异常类型。而默认构造的 variant 持有第一个类型的默认值，
        它并不空，\Fn{visit} 会正常分派。“空”与“无值”是两件事，
        别把它们写成同一条注释。
  \item \textbf{\Tp{std::optional}}：\Tp{value()} 在空实例上抛出
        \Tp{bad{\Bk}\_optional{\Bk}\_access}；\Fn{operator*} 与 \Fn{operator->}
        的前置条件要求实例有值，违约即越出承诺范围，落到 \UB{} 一侧。
\end{itemize}

\begin{compare}[title={抛出、返回错误与终止的取舍}]
抛出有类型化的错误信息，代价是展开路径与不可预测的耗时，且必须和 \Tp{noexcept} 边界商量好；
返回 \Tp{std::optional} 或错误码没有展开代价，能被 \Tp{noexcept} 函数携带，
代价是调用方可以忘记检查；终止不留恢复点，只适合“继续跑比停下更危险”的情形。
库边界上优先选中间那条路。
\end{compare}

\section{库边界上把 \UB{} 变成可诊断的失败}\label{sec:ub11-boundary}

库的对外接口是 \UB{} 唯一可以被拦住的地方。拦住它靠三种写法，它们互补而非可替换。

\begin{lstlisting}[style=cppstyle, caption={库边界的三条拦截写法}]
#include <cassert>
#include <cstddef>
#include <optional>
#include <stdexcept>
#include <vector>

// 一、前置断言：把“调用方违约”变成开发期可定位的诊断。
//     注意它只在开启断言的构建里存在，不能作为唯一的防线。
int& slot_at(std::vector<int>& v, std::size_t i) noexcept {
  assert(i < v.size() && "slot_at: index out of range");
  return v[i];
}

// 二、at 类接口：越界是抛出，不是终止也不是未定义。
//     它让 noexcept 之外的调用方拿到一个可捕获的确定失败。
const int& slot_checked(const std::vector<int>& v, std::size_t i) {
  return v.at(i);            // 越界抛出 std::out_of_range
}

// 三、返回 std::optional：让“没有结果”成为类型的一部分，
//     调用方忘记检查时编译器与静态分析就有抓手。
std::optional<std::size_t> find_pos(const std::vector<int>& v, int needle) {
  for (std::size_t i = 0; i < v.size(); ++i) {
    if (v[i] == needle) { return i; }
  }
  return std::nullopt;
}
\end{lstlisting}

\begin{bestpractice}[title={三条拦截写法的分工}]
其一，内部热路径用前置断言加上非抛接口：断言负责开发期，\Tp{noexcept} 负责编译期。
其二，对外边界一律提供 \Tp{at} 一类带检查的重载，把越界映射为
\Tp{out{\Bk}\_of{\Bk}\_range} 或库自己的错误类型，让违约的调用方拿到可捕获的失败而不是终止。
其三，可能无结果的查询返回 \Tp{std::optional}，把“忘记检查”从运行期问题变成编译期问题；
需要区分多种失败原因时返回类型化错误码，并把它写进接口文档的前置条件一节。
\end{bestpractice}

\begin{casebox}[title={自写移动赋值漏掉自检查导致的双重释放}]
某类型把内部缓冲的所有权保存在裸指针成员 \Tp{p} 上，移动赋值写成：先释放自己的旧缓冲，
再接管源的指针，也就是先对自身的 \Tp{p} 执行释放，再令 \Tp{p} 等于 \Tp{o.p}，
最后把 \Tp{o.p} 置空。当调用点写成 \Tp{x = std::move(x)}
（常见于通用模板里源与目标恰好别名）时，第一步就把自己正在托管的缓冲释放掉，
第二步把同一个地址交回给自己，于是 \Tp{x} 里存着一个已释放的指针；
析构时第二次释放，堆检查器报出双释放或读已释放内存。

类别判定：这条路径的根源是调用方越出了移动赋值对“源与目标不同一”的隐含假设。
对标准库类型而言它多半落在“标准未承诺”的一侧；对自写类型而言，
它是一次确定的资源管理错误，最终表现为 \UB{}（访问已释放存储并二次释放）。

检测面：AddressSanitizer 能看到双释放与释放后使用，并报出分配点，
因为这条路径上错误已经真正执行到了释放动作，属于可诊断的一类。
UndefinedBehaviorSanitizer 看不到它，因为这不是算术或类型层面的未定义模式。
更省事的修法不是加检测器，而是把自检查写回移动赋值的第一行，
或者改用 \Tp{std::unique{\Bk}\_ptr} 成员，让默认的移动赋值替你完成顺序正确的交换。
\end{casebox}

\section{本章判定清单}\label{sec:ub11-checklist}

\begin{itemize}
  \item 看到 \Tp{x = std::move(x)}：当作移动后状态处理，
        只做再赋值、析构与有保证的成员调用；
  \item 看到自写移动赋值：先确认有自检查，再看是否 \Tp{noexcept}，
        最后看成员是否全部不抛移动；
  \item 看到 \Tp{noexcept}：确认函数体内每条抛出路径都被本地消化，
        否则它不是提示而是一张期票；
  \item 看到析构函数里有 \Fn{throw}：默认承诺已经生效，抛出的终局是终止，
        除非你显式取消并承担后果；
  \item 看到类型特征判定为假但类型可复制：扩容会走拷贝，
        这是性能上的实现选择，不是错误；
  \item 看到纯虚调用或构造期的向下转型：分别按 \UB{} 与定义良好的失败处理，
        不要写进同一段注释；
  \item 安装终止处理程序时：只做单向取证输出，绝不恢复业务流程，
        并记得它接管的是进程级全局。
\end{itemize}

\begin{guideline}[title={本章的一句话结论}]
自我移动与 \Tp{noexcept} 的共同教训是：标准库里的“有效但未指定”只覆盖它写到的一半，
另一半要靠调用方不去越出前置条件。第~\ref{ch:ub12-library-pre} 章会把这条界线推广到整个标准库的前置条件条文上。
\end{guideline}

\chapter{标准库的前置条件}\label{ch:ub12-library-pre}

第~\ref{ch:ub11-terminate} 章在两个具体点上碰到了“调用方越出承诺”的形状：自移动与 \Tp{noexcept} 违约。
本章把这件事一般化。标准库的每一个函数条目都不只是一句效果描述，它是一组带名分的要素，
其中一部分违约的后果直接就是 \UB{}。读库条文的能力，本质上就是分辨哪一句是承诺、
哪一句只是说明的能力。

\section{库条文的要素结构}\label{sec:ub12-elements}

标准库函数条目按固定名分组织（主题上对应 \Std{structure.specifications}）：
Constraints（旧写作 Requires，位于模板层）、Mandates、Preconditions、Postconditions、
Remarks、Effects、Returns、Throws、Complexity、Errors 与 Error conditions。
每个名分承担不同的义务，违约的类别也不同：

\begin{itemize}
  \item \textbf{Constraints 与 Mandates}：属于编译期契约。不满足时程序 \ILLF{}，
        编译器有义务拒绝，Mandates 不满足同样是不合法程序。
  \item \textbf{Preconditions 与 Requires（函数层）}：属于运行期契约。违约的后果
        由 \Std{res.on.required} 一句话定死：行为未定义。
        这就是“前置条件违约通常直接落到 \UB{}”的规范来源。
  \item \textbf{Effects、Returns 与 Postconditions}：实现对调用方的承诺，违反它是实现缺陷。
  \item \textbf{Throws 与 Errors}：错误通道。走这条通道的失败是定义良好的，
        不该被算进 \UB{}。
  \item \textbf{Complexity}：算法代价的合同。达不到属于实现质量问题，不是程序非法。
\end{itemize}

\begin{clause}[title={两个典型措辞样例}]
容器与字符串条目里最常见的前置条件形状有这两种（主题化转述，具体编号随草案变动）：
下标访问写作 Requires 位置小于容器当前元素个数，越界即越出承诺；
字符串的裸指针赋值与构造重载写作 Requires 源区间有效且不与目标所控制的字符序列重叠。
后者特别值得记住：它排掉的正是“从自己的一段子串重新赋值给自己”这类看似无害的写法。
违约的落点都是 \UB{}，而不是某个诊断。
\end{clause}

\begin{table}
\centering
\begin{tabularx}{0.78\textwidth}{>{\raggedright\arraybackslash}p{0.24\textwidth} X}
\toprule
要素名分 & 违约之后标准给什么，以及由谁来发现 \\
\midrule
Constraints、Mandates & \ILLF{}：程序不合法，编译器必然拒绝，不需要任何外部检测器 \\
Preconditions、Requires & \UB{}：不保留任何承诺，只有前置检测（断言、消毒器、库调试档）能发现 \\
Effects 与 Postconditions & 实现缺陷：调用方无从依赖它，靠测试与断言暴露出来 \\
Throws 未覆盖的异常 & 视边界而定：越出 \Tp{noexcept} 就走终止路径，参见第~\ref{ch:ub11-terminate} 章的边界讨论 \\
Complexity & 实现质量问题：程序仍然合法，靠基准测试与实现文档来核对是否达标 \\
\bottomrule
\end{tabularx}
\caption{要素名分与违约类别}\label{tab:ub12-elements}
\end{table}

\begin{lstlisting}[style=cppstyle, caption={三种名分在同一段代码里的分工}]
#include <stdexcept>
#include <string>
#include <vector>

std::size_t demo(const std::vector<int>& v, const char* src, std::size_t n) {
  // Preconditions 违约 -> UB：越界下标不会带来任何诊断
  if (n >= v.size()) {
    throw std::out_of_range("demo: n out of range");  // 自己补的 Throws 通道
  }
  // Requires: [src, src+n) 有效，且不与 s 自己控制的序列重叠
  std::string s(src, n);
  return v[n] + s.size();
}
\end{lstlisting}

\section{迭代器失效总表}\label{sec:ub12-iterator-invalidation}

迭代器失效是标准库最大的一片前置条件区。判定规则可以压缩成两条：容器是否发生了
存储重排（重分配、桶重排、区间搬移），以及标准是否对该操作逐条列出失效范围。
下表按容器族整理，列出的都是标准承诺的范围；表外的操作请以“定点操作只影响被指元素”
为默认预期，并在库文档核对。

\begin{longtable}{@{}>{\raggedright\arraybackslash}p{0.19\textwidth} >{\raggedright\arraybackslash}p{0.20\textwidth} >{\raggedright\arraybackslash}p{0.29\textwidth} >{\raggedright\arraybackslash}p{0.12\textwidth}@{}}
\caption{常用容器的迭代器与引用失效范围}\label{tab:ub12-invalidation}\\
\toprule
容器 & 操作 & 失效的迭代器与引用 & 有无保证 \\
\midrule
\endfirsthead
\toprule
容器 & 操作 & 失效的迭代器与引用 & 有无保证 \\
\midrule
\endhead
\bottomrule
\endfoot
\Tp{vector} & 触发重分配的插入，\Fn{reserve}，\Fn{resize} 增大 & 全部迭代器与全部引用失效 & 有 \\
\Tp{vector} & 容量足够的尾部追加 & 仅尾后迭代器作废；元素引用不动 & 有 \\
\Tp{vector} & 中间插入与中间擦除 & 起始点及其后的迭代器与引用 & 有 \\
\Tp{deque} & 两端插入或两端擦除 & 迭代器全部失效，元素引用与指针仍有效 & 有 \\
\Tp{deque} & 中间插入 & 迭代器全部失效，引用不变 & 有 \\
\Tp{deque} & 中间擦除 & 迭代器全部失效，仅被擦元素的引用失效 & 有 \\
\Tp{list} & 任意插入、\Fn{splice}、\Fn{merge} & 不使任何迭代器与引用失效 & 有 \\
\Tp{forward{\Bk}\_list} & 任意插入与 \Fn{splice\_after} & 不使任何迭代器与引用失效 & 有 \\
\Tp{list} 与 \Tp{forward{\Bk}\_list} & 擦除与 \Fn{remove} 系列 & 仅指向被擦元素的迭代器与引用 & 有 \\
\Tp{map} 与 \Tp{set} & 插入、\Fn{emplace}、\Fn{extract} & 不使任何迭代器与引用失效 & 有 \\
\Tp{map} 与 \Tp{set} & 擦除 & 仅被擦元素的迭代器与引用 & 有 \\
\Tp{unordered{\Bk}\_*} & 插入、\Fn{reserve}、\Fn{rehash} 触发桶重排 & 迭代器失效，\textbf{元素引用保持有效} & 有 \\
\Tp{unordered{\Bk}\_*} & 擦除 & 仅被擦元素的迭代器与引用 & 有 \\
\Tp{string} & 追加、插入、擦除、重分配 & 规则同 \Tp{vector}；跨扩容时旧缓冲必失效 & 部分 \\
\Tp{array} & 任何成员操作 & 无失效概念：值语义、不重排 & 有 \\
适配器 \Fn{pop} 系列 & 对空容器调用 & 前置条件违约，落 \UB{} & 无 \\
\end{longtable}

\begin{guideline}[title={把长表压成两句话}]
节点式容器（\Tp{list}、\Tp{forward{\Bk}\_list}、\Tp{map}、\Tp{set}）只在“元素真的消失”时
使迭代器失效；连续式容器（\Tp{vector}、\Tp{string}）在重排时使一切失效。
\Tp{deque} 与 \Tp{unordered{\Bk}\_*} 是两种折中，它们把迭代器与引用分开承诺：
迭代器可能全废，但元素引用仍在。写循环时先问自己在追踪哪一层。
\end{guideline}

\begin{warning}[title={引用稳定不等于迭代器稳定}]
\Tp{unordered{\Bk}\_map} 重排桶之后，指向元素的引用与指针依旧有效，但所有迭代器都作废。
于是“拿着迭代器跨一次插入”是 \UB{}，而“拿着引用跨一次插入”是定义良好的。
反过来，\Tp{deque} 两端插入使迭代器全废但引用稳定；\Tp{vector} 扩容则两者同时作废。
把这三行记成一张对照表，比记十个函数签名有用。
\end{warning}

\section{视图类不持有所有权}\label{sec:ub12-views}

\Tp{std::string{\Bk}\_view}（\cpp{17}）与 \Tp{std::span}（\cpp{20}）
把“一段别人的存储加上长度”打包成可拷贝的值。它们不延长任何东西的生命周期，
因此悬垂风险与迭代器失效同族但更隐蔽：视图本身合法存在时，它指向的字节可能已经不在了。

\begin{lstlisting}[style=cppstyle, caption={视图悬垂的三个来源与对应的可移植写法}]
#include <string>
#include <string_view>
#include <vector>

// 悬垂一：返回值是视图，实体是局部字符串。函数返回即失去所指。
std::string_view bad_greeting() {
  std::string s = "hello";
  return s;                    // 视图指向即将析构的缓冲区
}

// 悬垂二：默认实参里的临时量。临时量在完整表达式末尾销毁。
void bad_log(std::string_view tag = std::string("default")) { /* tag 可能已悬垂 */ }

// 悬垂三：视图跨过一个会重排的写操作。
std::string_view bad_first(std::vector<std::string>& v) {
  std::string_view sv = v.front();
  v.emplace_back("more");      // 可能触发重分配，旧缓冲的地址随之改变
  return sv;
}

// 可移植一：需要跨边界存活就返回拥有者本身。
std::string good_greeting() { return "hello"; }

// 可移植二：视图只做参数，并在文档里写清它的有效区间。
void good_log(std::string_view tag) { (void)tag; }

// 可移植三：把长度与所有权一起交出去。
std::vector<char> good_bytes(const std::vector<char>& src) { return src; }
\end{lstlisting}

\begin{note}[title={span 接初始化列表的生命周期口径}]
\Fn{std::initializer{\Bk}\_list} 的底层数组活在该初始化列表对象的生命周期内。用
\Fn{std::span} 接一个临时初始化列表时，可用区间取决于所在完整表达式与被绑定实参的生命周期。
后续草案对这个构造的调整口径有过变动，跨版本写法请以目标实现的标准库文档
与特性测试宏为准，不要凭记忆断言它在你的编译器上一定有效或一定失效。
\end{note}

\begin{bestpractice}[title={视图入参的三条纪律}]
其一，视图作为参数时，约定“调用方保证所指存储在本次调用内有效”，
并把它写进函数文档的前置条件一栏；其二，绝不返回指向本地或临时存储的视图，
需要跨界就返回拥有者（\Tp{std::string}、\Tp{std::vector}）或由调用方传入缓冲区；
其三，视图跨越任何可能重排或扩容的写操作时，重新取一次视图，别信任旧的那个。
共享所有权是最后手段：它确实解决悬垂，但把生命周期问题换成了一次引用计数开销
与释放顺序的不确定。
\end{bestpractice}

\section{边界访问的两条路径}\label{sec:ub12-bounds}

同一个越界动作，走不同的接口会得到完全不同的类别。

\begin{itemize}
  \item \Fn{operator[]}：带前置条件（下标小于元素个数），违约即 \UB{}；
        它被设计为可与无检查指针访问同等代价，因此标准不要求任何诊断。
  \item \Tp{at()}：同样的输入越界时抛出 \Tp{std::out{\Bk}\_of{\Bk}\_range}，
        是定义良好的失败。\Tp{map} 与 \Tp{unordered{\Bk}\_map} 的 \Fn{at}
        还额外规定“键不存在也抛该异常”，而 \Fn{operator[]} 在键不存在时
        插入一个值初始化元素，两者语义并不对称。
  \item \Fn{front} 与 \Fn{back}：前置条件是容器非空。空容器上调用属于 \UB{}，
        标准库没有为它们提供抛出变体，须自己写 \Fn{empty} 检查，
        或用 \Tp{std::optional} 包装一层。
  \item 适配器的弹出系列：\Tp{queue}、\Tp{priority{\Bk}\_queue}、\Tp{stack} 的
        \Fn{pop} 与 \Fn{pop{\Bk}\_front} 同样带非空前置条件，空容器上调用是 \UB{}。
        \Tp{queue} 的 \Fn{front} 与 \Tp{priority{\Bk}\_queue} 的 \Fn{top} 是同一类契约。
  \item 区间参数：\Fn{erase} 的双迭代器形式、算法里的 \Tp{sort}(first, last) 一类
        要求区间由同一容器取得且起点不越过终点；把两个不同容器的迭代器配成区间
        属于越出承诺。
  \item 字符串的尾后约定：\Tp{std::string} 允许把尾后位置用作读取空字符的位置，
        但这个便利不适用于修改，也不适用于 \Tp{vector}；别把两者的边界规则混起来记。
\end{itemize}

\begin{compare}[title={检查与不检查的契约对比}]
两条路径的差别可以逐条列出来。失败形态：\Fn{operator[]} 落到 \UB{}，
\Fn{at} 落到一个可捕获的异常类型。代价：\Fn{at} 多一次比较与分支，
\Fn{operator[]} 没有额外要求，这也是标准不替它加检查的原因。
能否越过 \Tp{noexcept} 边界：\Fn{at} 需要展开路径，通常不能直接放在 \Tp{noexcept} 函数里；
\Fn{operator[]} 可以，但代价是违约即无承诺。适用位置：外部输入与库边界用 \Fn{at}
或显式检查，已在上游证明安全的内层热循环才用 \Fn{operator[]}。
\end{compare}

\section{内存与所有权的类别划分}\label{sec:ub12-memory}

这一节的四条经常被写进同一份“禁忌清单”，但它们在标准里的类别并不相同，
而划清类别决定了该用哪种检测手段。

\begin{itemize}
  \item \textbf{\Fn{delete} 与 \Fn{delete[]} 不匹配}：标准在 \Std{expr.delete} 里规定，
        释放单个对象的表达式必须作用于由非数组 \Fn{new} 创建的对象，
        数组形式必须作用于元素类型与操作数类型相同的数组对象。
        不匹配是 \UB{}：轻则析构次数不对，重则把单对象按数组元数据解析。
        这一条没有任何实现能替你兜住。
  \item \textbf{同一裸指针构造两个 \Tp{unique{\Bk}\_ptr}}：两个所有者各自释放一次，
        后果是双释放与释放后使用，属于 \UB{}。类型本身没有“所有权唯一”的运行期检查，
        唯一性只是文档层面的约定。需要共享就换 \Tp{shared{\Bk}\_ptr}，
        需要借用就传引用或裸指针参数。
  \item \textbf{\Tp{shared{\Bk}\_ptr} 的循环引用}：这是泄漏，不是 \UB{}。
        程序行为完全符合标准，只是控制块与对象的引用计数永不归零，析构不会发生。
        划清这条界线很重要：它意味着修复手段是弱指针或所有权重构，
        而不是任何“未定义行为检测器”。另外，循环里各对象的析构顺序
        与析构发生的线程，标准并不给出保证。
  \item \textbf{自定义删除器与分配器}：删除器必须与获取内存的方式配对；
        对超对齐类型，释放时必须把对齐信息传回分配侧，否则释放路径与实际分配路径不一致，
        后果落在 \UB{} 一侧（对齐专题见第~\ref{ch:ub07-alignment} 章）。
        同理，把一个容器的元素交给另一个容器的分配器去释放，
        或让不等价的分配器互相交换元素，也越出了分配器契约。
\end{itemize}

\begin{note}[title={删除器抛出与第~\ref{ch:ub11-terminate} 章的接缝}]
\Fn{std::unique{\Bk}\_ptr} 的析构是 \Tp{noexcept} 推导的：删除器可能抛出时，
释放路径一旦抛出就走到终止函数。也就是说“删除器不该抛”不是风格建议，
而是由隐式承诺撑住的硬约束。同理，容器元素的析构若经由会抛的删除器或用户析构，
边界处的终局同样是终止而不是回滚。
\end{note}

\section{库自带的检查模式}\label{sec:ub12-check-modes}

标准不要求库检查前置条件，但主流实现都提供可选的强化档位。它们的共同特点是：
改变 ABI 或改变模板实例，因此必须整个程序一致地开启；它们的覆盖面永远只是
前置条件的一部分，不是全部。

\begin{lstlisting}[style=shellstyle, caption={三种实现的标准开启方式（宏名与覆盖面以实现文档为准）}]
# libstdc++：容器调试版，覆盖面广但 ABI 不兼容，必须整档重编
g++ -std=c++23 -O1 -D_GLIBCXX_DEBUG app.cpp
# libstdc++：轻量断言档，只加检查不改类型，适合常开
g++ -std=c++23 -O2 -D_GLIBCXX_ASSERTIONS app.cpp
# libc++：强化模式，档位名以所用版本的库文档为准
clang++ -std=c++23 -O2 \
  -D_LIBCPP_HARDENING_MODE=_LIBCPP_HARDENING_MODE_DEBUG app.cpp
# MSVC STL：迭代器调试级别 2，须与运行时库档位一致
cl /std:c++latest /MDd /D_ITERATOR_DEBUG_LEVEL=2 app.cpp
\end{lstlisting}

\begin{warning}[title={库检查模式不是消毒器}]
其一，覆盖面：调试模式看的是迭代器是否属于所操作的容器、区间是否合法、容器是否空；
它看不见有符号溢出、对齐错误、别名违约。其二，一致性：调试容器与普通容器在多数实现里
是不同类型或不同布局，混用会把一次越界换成一次链接错误或更隐蔽的内存破坏。
其三，性能：轻量断言档与全量调试档的代价差别很大，具体倍数量级随实现与工作负载变化，
本册不给数字，请以实测为准。宏名与开关覆盖面各实现版本都可能调整，
写进构建脚本前请在库文档核对（构建口径见第~\ref{ch:ub14-conventions} 章）。
\end{warning}

\begin{lstlisting}[style=cmakestyle, caption={让检查档位与编译器绑定的最小门禁}]
add_compile_definitions(
  $<$<CONFIG:Debug>:_GLIBCXX_ASSERTIONS>)
target_compile_options(app PRIVATE
  $<$<CXX_COMPILER_ID:GNU>:-Wall>
  $<$<AND:$<CXX_COMPILER_ID:GNU>,$<CONFIG:Debug>>:-fsanitize=undefined>)
# 注意：启用调试容器要求全部目标文件同档重编，勿与预编译库混用
\end{lstlisting}

\begin{casebox}[title={erase 循环里的失效迭代器，以及它何时才暴露}]
错误形态：在 \Tp{vector} 上边遍历边按条件擦除，循环头里写了自增，
循环体里却调用 \Fn{erase}。擦除会使被擦位置及其后的迭代器失效，
于是循环末尾那次自增是在一个已作废的迭代器上运算；连续式容器上这通常表现为
跳过元素、越出尾部，或直接读到失效区间。

正确形态是让擦除返回下一个有效位置：循环头只写初始化与比较，递增完全交给分支决定；
擦除分支写成 \Tp{it = c.erase(it)}，非擦除分支才执行自增，两条分支互斥；
尾后迭代器只参与比较，绝不自增。

类别判定：在连续式容器上这是 \UB{}（对失效迭代器运算并解引用）；
在节点式容器上写同样的循环只是逻辑错误（漏检元素），因为其它迭代器并未失效。
同一个错误代码在两类容器上落在不同类别里，这是本章最想留下的结论。

检测面：libstdc++ 的 \Opt{D{\Bk}\_GLIBCXX{\Bk}\_DEBUG} 与 libc++ 的调试、强化档位
会在迭代器作废后立刻发现使用并报错；MSVC STL 的 \Opt{D{\Bk}\_ITERATOR{\Bk}\_DEBUG{\Bk}\_LEVEL}
设为 2 时同样报告迭代器不可解引用。而优化档的普通构建里，这段代码可能安静地
给出错误答案。UndefinedBehaviorSanitizer 一般看不到它，因为失效与否是库层语义，
不是语言层的类型或算术模式；AddressSanitizer 只在失效恰好踩到已释放堆存储时才可能报出。
\end{casebox}

\section{本章判定清单}\label{sec:ub12-checklist}

\begin{itemize}
  \item 读库条目先找名分：编译期契约交给编译器，运行期前置条件违约一定是 \UB{}，
        只有走 Throws 与 Errors 通道的失败才是可捕获的；
  \item 写涉及迭代器的循环前，先在本章长表里查该容器该操作的失效范围，
        再决定是持有迭代器还是持有引用；
  \item 视图（\Tp{string{\Bk}\_view}、\Tp{span}）永远不延长生命周期：
        作为参数可以，作为返回值要三思；
  \item 边界访问：外部输入走 \Fn{at} 或显式检查，内部热路径才用 \Fn{operator[]}，
        并且永远给 \Fn{front}、\Fn{back} 与弹出系列配一次非空检查；
  \item 泄漏、双释放、悬垂是三类问题：只有后两类是 \UB{}，
        检测工具与修复手段都不通用；
  \item 库检查模式必须整程序一致开启，且不能替代消毒器，
        也不能替代把前置条件写进接口文档。
\end{itemize}

\begin{guideline}[title={前置条件是接口的一部分}]
一个只写在注释里的前置条件等于没有。要么把它编码进类型（\Tp{std::optional}、
强类型索引、只接受切片而不接受容器的构造函数），要么编码进检查（\Fn{at}、断言、
库调试档位），要么编码进编译期契约（Constraints 与 Mandates）。
三条路都通往同一个目标：让违约的调用方在某处拿到确定的失败，
而不是让违约在优化后的代码里改写成别的形状。
\end{guideline}

\part{标准库约定与前向兼容}
\chapter{实现定义与未指定：可移植边界的写法}\label{ch:ub13-implementation}

前两章处理的是标准收回承诺的 \UB{} 与走到尽头的 \ILLF{}；本章处理另外
三种状态：\IDB{}、\USB{} 与 \NDR{}。它们同样是“同一份源码、不同答案”，
但性质完全不同：\UB{} 是没有人承诺过的悬崖，这三类则是标准明确允许
存在、并且把选择权交给实现的岔路口。工程上真正的问题不是“这里没有
唯一答案”，而是“我能不能在编译期或文档里把答案钉死”。

这一章给出三件事的清单与方法：各类的文档义务差别、常见的 \IDB{} 与
\USB{} 条目、把探测结果固化下来的写法。机制层面的细节——比如异常
展开表怎么生成、调用约定在机器上长什么样——属于另一册“编译器与
链接器”的范围；本章只讲语言把决定权交给实现这件事本身，以及据此
应该怎么写代码。

\section{三类合法不确定的文档义务}\label{sec:ub13-obligation}

三者共同点是都不报错：程序合法，任何答案都符合标准。区别在于
“谁能告诉你答案”，而这决定了工程手段完全不同。

\begin{clause}[title=标准怎么说：符合性条款里的三种承诺]
\Std{intro.compliance} 把行为分类时写明：\IDB{} 的每一处，实现都
必须在文档中说明它选择了哪个方案；\USB{} 则是标准给出若干合法解、
不要求实现说明理由；\NDR{} 属于不合法但不要求诊断的程序片段，编译器
保持沉默不构成违背承诺。换句话说，只有 \IDB{} 附带文档义务，
\USB{} 与 \NDR{} 想问答案，得去翻实现本身。
\end{clause}

把“谁能回答”列成一张表，遇到具体问题时先查这一层，再决定投入
多少防御代码：

\begin{table}
\centering
\begin{tabular}{@{}>{\raggedright\arraybackslash}p{0.14\textwidth} >{\raggedright\arraybackslash}p{0.30\textwidth} >{\raggedright\arraybackslash}p{0.32\textwidth}@{}}
\toprule
类别 & 谁能告诉你答案 & 常规工程手段 \\
\midrule
\IDB{} & 实现手册与头文件宏；多数可用编译期常量直接探测 &
探测写进 \Tp{static{\Bk}\_assert}，把选择变成配置项或实现特化 \\
\USB{} & 没人欠你答案；最多在实现的版本历史里找到当过的事实 &
写测试断言到两个实现上；把依赖改成显式的稳定算法 \\
\NDR{} & 标准明说可以不诊断；编译器版本更新随时改变诊断面 &
用更严格的静态检查与消毒器补位，不指望编译器报错 \\
\bottomrule
\end{tabular}
\caption{三类合法不确定的“可回答性”对照}\label{tab:ub13-who-answers}
\end{table}

\begin{warning}[title=把三类混为一谈的两种代价]
把 \IDB{} 当 \UB{} 防，会给每个查表操作套上运行时分支，代码变重却
没换来任何保证——\IDB{} 至少有两个可查的答案来源：文档和宏。反过来，
把 \USB{} 当 \IDB{} 去找手册，会浪费半天找不到“规定”，因为实现
根本没有说明义务；\Fn{std::sort} 对等价元素的相对顺序就是这样，
想知道只能问代码和版本历史。
\end{warning}

\subsection{NDR：三类里最安静的一种}\label{sec:ub13-ndr}

\NDR{} 与前两者的分界线不在“有没有唯一答案”，而在“程序还合不
合法”：\IDB{} 与 \USB{} 都是合法程序的多解，\NDR{} 则已经违反了
一条不要求诊断的规则，编译器沉默是标准授权的状态。它的工程含义
有两层。其一，编译器没报错不构成任何辩护，能用的外力只有消毒器
与更严格的静态检查，分类与识别方法见第~\ref{ch:ub01-taxonomy} 章。
其二，同一份含 \NDR{} 的代码在某个编译器版本上“看起来能用”，
只是那一次没有被暴露，版本更新随时把沉默改成报错、也可能反过来。
所以文档义务的梯度是三级递减：\IDB{} 有文档可查，\USB{} 无文档
但程序合法，\NDR{} 既不合法也无诊断承诺。本章其余部分聚焦前两级，
\NDR{} 的违例形态不再展开。

\section{常见实现定义行为清单}\label{sec:ub13-idb-catalogue}

下面这份清单收录语言层面最常咬人的 \IDB{} 条目。判定标准只有一条：
标准把选择权交给实现，同时要求实现写进文档。每一项都给出“影响面”
和“在代码里如何探测”，探测手段优先给编译期可得的。

\begin{longtable}{@{}>{\raggedright\arraybackslash}p{0.24\textwidth} >{\raggedright\arraybackslash}p{0.28\textwidth} >{\raggedright\arraybackslash}p{0.24\textwidth}@{}}
\caption{语言层常见 \IDB{} 项目、影响与探测手段}\label{tab:ub13-idb-list}\\
\toprule
项目 & 影响 & 探测手段 \\
\midrule
\endfirsthead
\toprule
项目 & 影响 & 探测手段 \\
\midrule
\endhead
\bottomrule
\endfoot
\Tp{char} 的符号性 &
位图与查表索引、十六进制解析、字符分类 &
\Feat{CHAR\_MIN}（等于 0 即无符号）、\Tp{std::is\_signed} \\
\addlinespace
无固定类型枚举的底层类型 &
枚举序列化位宽、与整型间的提升路径 &
\Tp{std::underlying{\Bk}\_type}（\cpp{11} 起可编译期读取） \\
\addlinespace
\Tp{wchar\_t} 的宽度 &
宽字符串跨平台读写、\Fn{wcs} 系族的行为 &
\Tp{sizeof}(wchar\_t) 写进编译期断言 \\
\addlinespace
窄字符与执行字符集的转换 &
多字节边界、\Fn{narrow} 遇到不可表示字符时的替换方案 &
以实现文档为准；可用转换测试覆盖两端 \\
\addlinespace
静态存储对象初始化与销毁次序的细节 &
跨翻译单元的全局依赖、进程退出阶段的析构时机 &
实现文档加两端实测；语言只保证主函数前的完成 \\
\addlinespace
异常展开机制 &
能否穿越语言连接说明为 \Tp{extern} 的 C 帧、\Fn{noexcept}
边界处的终止路径 &
实现文档；机制细节在编译器与链接器一册展开 \\
\addlinespace
默认对齐要求与最大基本对齐 &
SIMD 缓冲、自定义分配器、\Tp{std::max{\Bk}\_align{\Bk}\_t} 的取值 &
\Fn{alignof}、\Feat{\_\_STDCPP{\Bk}\_DEFAULT{\Bk}\_NEW{\Bk}\_ALIGNMENT{\Bk}\_}（\Hdr{new}） \\
\addlinespace
调用语言连接到实际调用约定的映射 &
跨语言边界的进出栈、名字修饰 &
ABI 与实现文档；机器面细节归编译器与链接器一册 \\
\end{longtable}

\begin{note}[title=清单的用法是查“谁来回答”，不是照抄结论]
表中每一项的正确姿势都是先找到实现给出的答案（文档或宏），再决定
是禁止依赖、配置化还是特化。清单本身也会随标准修订增删，某一项
在新版本里可能从 \IDB{} 升级为有默认值的规定，例如枚举固定底层类型
就是 \cpp{11} 补上的；以版本化文档为准。
\end{note}

\section{常见未指定行为}\label{sec:ub13-usb-catalogue}

\USB{} 的共同形态是：标准给出几个合法解，实现挑哪个都行，也不欠
文档。下面五个是最常在实践中被误当成有保证的条目：

\begin{itemize}
\item 函数实参的求值顺序。标准不规定各实参子表达式的相对次序，
    依赖它们互相影响的写法在第~\ref{ch:ub08-evaluation} 章作为
    高危模式专门处理；这里只强调它属于 \USB{} 而非 \UB{}——单看
    某一种顺序都是合法的，坏在代码假设了顺序。
\item \Fn{std::sort} 对等价元素的相对顺序。两个“相等”的元素谁先
    谁后是未指定的；需要保留输入次序就用 \Fn{std::stable\_sort}，
    这是行为差异里少数有现成替代品的。
\item 哈希容器的遍历顺序。桶的数量、插入引发的重组策略都留给
    实现，遍历序不能作为协议的一部分；要稳定顺序就排序后再输出。
\item \Fn{rand} 的具体算法。标准只承诺返回值落在 \Tp{0} 到
    \Feat{RAND\_MAX} 且 \Feat{RAND\_MAX} 不小于 32767，序列本身
    各实现自选；跨平台可复现的随机流请换显式引擎加固定分布。
\item \Fn{operator new} 在分配失败前调用 \Fn{new\_handler} 的
    重试轮数。失败前反复尝试的次数与终止条件由库实现推进，
    抛 \Tp{std::bad\_alloc} 之前的循环细节不要做数值假设。
\end{itemize}

对待 \USB{} 的通用姿势比 \IDB{} 更保守：因为连“查文档”这条路都
不存在，任何依赖都必须写成“两个实现上都验证过”的测试，或者干脆
换成承诺更强的接口（\Fn{stable\_sort}、显式引擎、排序后遍历）。

\subsection{\USB{} 的版本漂移}\label{sec:ub13-usb-drift}

\USB{} 还有一个 \IDB{} 没有的脾气：同一个实现内部的合法解也会换。
库作者升级排序实现、换哈希种子策略、调整分配器重试逻辑，都不算
破坏承诺——标准从没指定过那一个解。因此“我们一直用 GCC，行为
是稳定的”这类论证在 \USB{} 上不成立，成立的可能只有两种替代
写法：要么把需求换成有规定的接口（要稳定次序就 \Fn{stable\_sort}），
要么把需求写进测试（遍历序参与输出的，先把输出规范成集合再比较）。
判别方法也简单：翻实现文档找到“明文规定”之前，先假设它是
\USB{}；找到之后，再核对它属于哪一版标准、哪一版库。

\section{探测与固定}\label{sec:ub13-detect-pin}

把答案钉死的第一层是编译期探测。可用的工具按来源分四类：

\begin{itemize}
\item 数值特征：\Tp{std::numeric\_limits} 给出各整型的边界与
    数字位数，\Feat{CHAR\_BIT} 等 \Hdr{climits} 常量给出字节定义；
\item 布局特征：\Tp{alignof} 与 \Tp{sizeof} 回答对齐和宽度，
    \Tp{std::max\_align\_t} 的取值属于 \IDB{}；
\item 类型特征：\Tp{std::is\_signed}、\Tp{std::underlying\_type}
    一类在编译期读取实现的选择；
\item 预定义宏与特性宏：\Feat{\_\_SIZEOF\_WCHAR\_T\_} 这类尺寸宏
    由实现提供，标准库侧的 \Feat{\_\_cpp\_lib\_*} 体系见
    第~\ref{ch:ub14-conventions} 章。
\end{itemize}

四类工具各管一段，用错段会拿到假答案：\Tp{std::numeric\_limits}
回答“这个整型在本实现上能装下什么”，是换实现时最便宜的第二意见；
\Tp{alignof} 与 \Tp{sizeof} 回答布局，二者几乎都是编译期常量，
天然适合进断言；类型特征回答“实现替我选了哪个方案”，比如枚举
底层类型；宏体系回答“这个设施存在与否”，存在性问题和取值问题
要分开问。注意 \Tp{std::numeric\_limits} 里也有 \IDB{} 的镜像：
个别成员（如 \Tp{has\_infinity}）的取值本身就是实现给出的事实，
探测到了不等于理解了原因。

探测结果一旦决定代码结构，就该写进编译期断言而不是运行时分支：
运行时 \Tp{if} 只会在某一次执行里暴露问题，而 \Tp{static\_assert}
在移植的第一次编译就失败，失败点就在假设写下的位置。

\begin{lstlisting}[style=cppstyle, caption={把实现定义的选择写进编译期断言：换实现先断言，而不是上线后翻车}]
#include <climits>
#include <type_traits>

// 本模块按“char 无符号”设计：256 项查表直接用 char 做下标。
static_assert(CHAR_BIT == 8, "查表宽度按一字节 8 位设计");
static_assert(CHAR_MIN == 0, "本模块假设 char 是无符号的");

constexpr bool kCharSigned = std::is_signed_v<char>;  // 也允许反过来查
\end{lstlisting}

尺寸类宏可以一次性导出留档，作为移植检查表的一部分：

\begin{lstlisting}[style=shellstyle, caption={导出实现的预定义宏清单，两份实现相减即是差异面}]
# 只跑预处理、把全部宏导出来比对
g++ -std=c++20 -dM -E -x c++ /dev/null > macros_gcc.txt
clang++ -std=c++20 -dM -E -x c++ /dev/null > macros_clang.txt
diff macros_gcc.txt macros_clang.txt | grep -E 'SIZEOF|DEFAULT_NEW'
\end{lstlisting}

\begin{bestpractice}[title=探测进编译期，差异进配置，假设进断言]
一条准则：凡是被探测用来决定代码路径的取值，必须在同一位置写
\Tp{static{\Bk}\_assert} 或预处理检查；探测与断言分离的写法，等于把
“实现选了什么”变成“作者猜了什么”。断言消息里写清依赖的是哪一
类行为（\IDB{} 还是 \USB{}），以及在哪两个实现上验证过。
\end{bestpractice}

\section{面对实现定义行为的三种工程姿势}\label{sec:ub13-postures}

同一处 \IDB{}，在不同代码库里正确姿势不同。选姿势的依据是这个
假设的影响半径：

\begin{guideline}[title=禁止依赖、配置项、实现特化：按影响半径选姿势]
其一，禁止依赖并加诊断：库里没有跨平台差异需求时，直接规定唯一
用法，配静态检查与 \Tp{static\_assert} 让违反者编译失败。其二，
把它变成配置项：探测结果进入构建配置或类型别名，代码只消费别名
（例如统一用 \Tp{uint8\_t} 而不是 \Tp{char} 承载原始字节）。其三，
写成实现特化：两种选择都值得支持时，各给一份实现与一份测试，
公共代码通过显式分支选择，绝不允许隐式依赖。
\end{guideline}

三种姿势的取舍依据是假设的影响半径：只在单个函数内使用的假设，
就地断言即可（姿势其一）；跨模块传递的假设，升级为配置项让所有
消费者看到同一个事实（姿势其二）；被两个以上实现各自给出的假设，
才值得付出版本测试的成本做特化（姿势其三）。反过来用是常见的
过度工程：给局部断言配构建选项，或给一个从来不支持第二实现的
产品维护双份特化代码。

\section{案例：\Tp{char} 的符号性翻掉一张查表}\label{sec:ub13-char-case}

下面这段解析十六进制转义并置入位图的代码，在 \Tp{char} 无符号的
实现上工作正常，在 \Tp{char} 有符号的实现上对同一输入给出不同
结果，甚至直接越界。输入字节 \Tp{0x80} 以上的部分，两份实现里
\Tp{char} 的值分别是正半区与负半区：

\begin{lstlisting}[style=cppstyle, caption={问题写法：拿裸 char 当 256 项表的下标}]
bool bitmap[256];

bool accept(char c) { return bitmap[c]; }   // char 有符号时 c 可为负
\end{lstlisting}

\begin{casebox}[title=char 符号性：同一输入在两份实现上的两种结局]
无符号实现里 \Fn{accept} 的下标落在 \Tp{0} 到 \Tp{255}，查表正确；
有符号实现里带符号位字节的入参会把下标变成负数，读写翻到数组
之前的内存——这一步已经不是 \IDB{} 而是 \UB{}，越界的机理见
第~\ref{ch:ub05-pointer} 章。同一行代码，前半段差异是合法的多解，
后半段差异是越界的事故；这就是为什么 \IDB{} 必须在边界处固定。
正确写法先把字节显式化，再用编译期断言把假设写死在文档位置：
\end{casebox}

\begin{lstlisting}[style=cppstyle, caption={修正写法：入口即转 unsigned char，假设另有断言兜底}]
bool accept(unsigned char c) { return bitmap[c]; }   // 接口层消灭裸 char

inline bool accept_bytes(const char* s, std::size_t n) {
    for (std::size_t i = 0; i != n; ++i) {
        // 转换点唯一：显式经过 unsigned char，不经符号位
        if (!accept(static_cast<unsigned char>(s[i]))) return false;
    }
    return true;
}

// 若要声明本实现分支，探测与断言写在同一处，而不是散在逻辑里
// static_assert(std::is_signed_v<char>, "本分支只跑在有符号 char 的实现上");
\end{lstlisting}

若代码库坚持两种实现都要支持，就把 \Tp{std::is\_signed} 的探测结果
升格为配置项（姿势其二）：所有承载原始字节的接口统一用
\Tp{unsigned char} 或 \Tp{std::byte}，让符号性无处发挥。只在解析
十六进制字符的局部，则注意 \Tp{a} 到 \Tp{f} 的字母在两种符号性
下比较结果虽相同，但混入了高位字节的输入流时，负值会先撞上
范围比较——这类“看起来只差一点”的场景，正是上一节三种姿势里
禁止依赖加诊断的典型适用面。

\subsection{清单之外：新条目如何进入工程流程}\label{sec:ub13-workflow}

标准每修订一次，\IDB{} 与 \USB{} 的清单都会移动：某项 \IDB{}
可能升格为有默认值的规定（枚举固定底层类型就是 \cpp{11} 补上的
先例），某项今天无人抱怨的写法也可能在新版里被重新归类。因此
清单不是一次性背诵材料，而是流程入口。最小流程三步：

\begin{itemize}
\item 换编译器版本或语言模式时，重跑一次预定义宏导出并做双向
    \Fn{diff}，差异行里的每个宏都要归类：\IDB{} 的取值变了、
    特性宏换值，还是纯粹的噪声；
\item 归类结果落到三处之一：编译期断言（可静态判定的）、构建
    配置（需要按实现切换的）、双份特化加测试（运行期才见分晓的）；
\item 每次标准修订后，把代码里所有 \Tp{static\_assert} 的消息重读
    一遍——它们是整个流程里唯一成文的假设清单，过期的假设往往
    先是消息过时，然后才是断言失效。
\end{itemize}

这条线通向下一分册主题：本章把当下的实现差异固定住，
第~\ref{ch:ub15-forward} 章处理这些固定点在未来版本里的存活问题。
固定得越早、越靠近编译期，前向兼容的代价越低，这一判断在本册
第~\ref{ch:ub03-sanitizers} 章的消毒器投入产出讨论里同样成立。

\chapter{标准库约定：特性测试宏、调试模式与分配器契约}\label{ch:ub14-conventions}

上一章把语言交给实现的三类“合法不确定”列了清单；这一章进入标准库
自己的约定层。库实现比语言多出三个可依赖又容易踩空的面：特性探测
体系（用什么宏问库“你支持到哪一步”）、调试与加固档位（同一份代码
在开与不开检查的库里走向完全不同的失败形态）、以及分配与错误的
契约（谁保证前置条件、失败以什么形式浮现）。这些约定大多不是
\IDB{} 而是库文档层面的承诺，违约的后果要么在链接期、要么在运行期，
往往比语言层的坑更难定位。

\section{版本与特性探测体系}\label{sec:ub14-feature-detection}

问“这个库/编译器支持什么”有三个层次，各管一段，混用会出假阴性。

第一层是语言版本宏 \Feat{\_\_cplusplus}。它按 \cpp{11}/\cpp{14}/
\cpp{17}/\cpp{20} 给出形如 \Tp{201703L} 的日期常数，但存在一类
历史问题：某些实现默认值停旧、要专门的开关才更新到真实语言模式，
具体开关以实现文档为准。因此可靠规则是：\Feat{\_\_cplusplus}
只用来问语言核心特性（语法与标准内建设施），问库特性请走第二层。

第二层是库特性测试宏 \Feat{\_\_cpp\_lib\_*}。\cpp{20} 起标准规定
它们集中收录在 \Hdr{version}（\Std{version.syn}），值是“年年月月”
六位日期；\cpp{20} 之前的实现把它们散落在各自的功能头里。单个宏
按惯例包在 \Feat{\_\_has\_include} 保护下探测，缺头即缺特性：

\begin{lstlisting}[style=cppstyle, caption={三层探测的分工示例：语言宏、库特性宏、头文件存在性}]
#include <version>          // C++20 起集中提供 __cpp_lib_* 取值

#if defined(__has_include)
#  if __has_include(<span>)
#    include <span>
#  endif
#endif

#if defined(__cpp_lib_span)
constexpr std::span<const int> window(const int* p, std::size_t n) {
    return std::span<const int>(p, n);
}
#endif
\end{lstlisting}

第三层是编译器内建探测器 \Feat{\_\_has\_include} 与
\Feat{\_\_has\_cpp\_attribute} 一类：前者问头文件在不在，后者问
属性（如 \Tp{fallthrough}）认不认。它们早年是实现扩展，后来在
较新的标准修订中被收编为推荐的可选特性；早年在完全不支持这些
探测器的编译器上，直接调用它们会得到“未知标识符按零处理”还是
预处理错误，两种结局随实现而异，稳妥写法是像上面那样先
\Feat{defined} 再调用。

\begin{note}[title=三类探测器的分工边界]
\Feat{\_\_cplusplus} 问语言、\Feat{\_\_cpp\_lib\_*} 问库的某项
特性是否可用、\Feat{\_\_has\_include} 问头文件是否存在。三者答案
可以互相矛盾：有头不等于特性齐（可能被特性宏挡在门内），语言模式
开着不等于库跟上了（库版本落后于编译器）。工程上把三个答案都
导出留档比对，方法同第~\ref{ch:ub13-implementation} 章的预定义宏清单。
\end{note}

特性宏的消费只有一种安全写法：先 \Feat{defined} 再与期望的年月
值比较，而不是假设当前实现恰好带最新值。库对某项特性做后续修订
时特性宏的日期值会跟着升，绑定了具体值的代码分支要能同时容纳
新旧两档语义；值表由标准与库文档维护，移植时查 \Hdr{version}
或库文档，不要抄二手资料。这与上一章“探测与断言写在同一处”的
准则同源：假设要成文，出处要可查。

\section{调试与加固档位：三大实现的对照}\label{sec:ub14-debug-tiers}

主流标准库各自提供若干“开检查”的编译档位。宏名与档位数量随版本
变动，下表只收确证过的名称，细节一律以各库文档为准。档位改变的是
前置条件违约时的暴露方式：从静默 \UB{} 升级到断言、异常或固定
诊断后终止。

\begin{longtable}{@{}>{\raggedright\arraybackslash}p{0.13\textwidth} >{\raggedright\arraybackslash}p{0.26\textwidth} >{\raggedright\arraybackslash}p{0.17\textwidth} >{\raggedright\arraybackslash}p{0.20\textwidth}@{}}
\caption{三大标准库的调试/加固档位对照}\label{tab:ub14-tiers}\\
\toprule
库 & 档位与宏名 & 检查什么 & 失败形态 \\
\midrule
\endfirsthead
\toprule
库 & 档位与宏名 & 检查什么 & 失败形态 \\
\midrule
\endhead
\bottomrule
\endfoot
libstdc++ &
\Feat{\_GLIBCXX{\Bk}\_DEBUG} 整档切换；轻量断言档
\Feat{\_GLIBCXX{\Bk}\_ASSERTIONS} &
迭代器有效性、容器下标与区间、算法前置条件 &
诊断信息后终止（走报告与异常路径，默认进程中止） \\
\addlinespace
libc++ &
历史调试宏 \Feat{\_LIBCPP{\Bk}\_DEBUG}；现代拆为调试模式、断言与
加固（\Feat{\_LIBCPP{\Bk}\_ENABLE{\Bk}\_HARDENING} 等），名称以库文档为准 &
迭代器配对、容器越界、可硬化前置条件 &
加固失败打印固定消息后立即终止；断言失败走 \Tp{assert} 一路 \\
\addlinespace
MSVC STL &
\Feat{\_ITERATOR{\Bk}\_DEBUG{\Bk}\_LEVEL}（调试构建默认高于发布构建）&
迭代器与容器的配对关系、范围检查 &
依级别走无效参数处理器（默认报告后终止）或异常，以库文档为准 \\
\addlinespace
MSVC CRT &
调试堆检查接口（如 \Fn{\_CrtSet{\Bk}\_DbgFlag} 一族）&
释放后再用、分配块边界踩踏 &
调试报告与断言，仅在调试 CRT 下可用 \\
\end{longtable}

\begin{lstlisting}[style=shellstyle, caption={同一段代码在三种档位下的编译方式}]
# libstdc++：整档切换，注意会改变对象布局，必须全工程一致
g++ -std=c++20 -O0 -D_GLIBCXX_DEBUG -D_GLIBCXX_DEBUG_PEDANTIC t.cpp
# libc++：调试/加固宏的名称与取值组合以库文档为准
clang++ -std=c++20 -O0 -D_LIBCPP_HARDENING_MODE=_LIBCPP_HARDENING_MODE_FAST t.cpp
# MSVC：_ITERATOR_DEBUG_LEVEL 由运行库选项给默认值，也可显式覆盖
cl /EHsc /std:c++20 /MDd /D_ITERATOR_DEBUG_LEVEL=2 t.cpp
\end{lstlisting}

\begin{warning}[title=调试档位改变对象布局，混档链接是事故源]
整档切换类（libstdc++ 的调试档、MSVC 的不同 \Feat{\_ITERATOR\_\Bk\_DEBUG\_\Bk\_LEVEL}）
会改变容器与迭代器的内存布局或符号可见性：部分翻译单元开档、部分
不开，轻则链接报错，重则同一类型带两套定义在运行期互相踩踏——
那是内存层的 \UB{}，与第~\ref{ch:ub04-lifetime} 章的生命周期事故
同族。规则：档位是构建级配置，进入缓存、库与可执行的全家桶，
不允许某个目标单独开。
\end{warning}

\subsection{档位怎么选}\label{sec:ub14-tier-choice}

整档调试、轻量断言、加固三类不是同一维度的三个开关。整档改变
对象布局与失败语义，适合本地排障与专门的 CI 作业；轻量断言档以
常数成本换更多前置条件被检查，可以在测试档长期开启；加固档为
发布形态设计，检查范围收缩到不改变对象语义的子集，失败形态是
固定消息后立即终止，便于运维系统识别。各实现的具体名称与取值
以库文档为准；选择顺序是：先确认布局是否改变，再决定失败形态
是否被你的告警管道消化，最后才是检查覆盖率。

\begin{casebox}[title=同一行越界，三种档位的三种结局]
代码只有一行：\Tp{operator[]} 取越界下标（\Fn{at} 在所有档位都抛
\Tp{std::out\_of\_range}，这里故意用不检查的那一个）。关闭检查的
发布档：标准只要求 \Tp{operator[]} 的 \Tp{index} 有效，越界即
\UB{}，静默读到脏内存或崩溃；libstdc++ 调试档：前置检查失败，
输出诊断后进程终止；MSVC 调试默认档：进入无效参数处理器，默认
报告后终止。三档里没有任何一档“返回一个合法值”，差别只在坏消息
来得多快、多响。把差异固定下来的写法：其一，档位作为工程级定义
进构建系统并在 CI 矩阵里至少跑无档与一档有档；其二，公开接口层
禁用裸 \Tp{operator[]}，统一 \Fn{at} 或显式范围校验，让失败形态
收敛到自家文档承诺的那一种，不再看库档位的脸色。
\end{casebox}

\section{分配器契约：preconditions 与传播属性}\label{sec:ub14-allocator-contract}

用户自定义分配器进入容器时受一组契约约束（\Std{allocator.requirements}）。
核心三条：

\begin{itemize}
\item \Fn{allocate} 的前置条件是元素个数 \Tp{n} 不超过分配器可服务
    的上限，可选的 \Tp{hint} 参数只是局部性提示、不改变合法性；
    越界调用不在标准承诺内，库通常在进入前用 \Fn{max{\Bk}\_size}
    拦截并抛 \Tp{std::length\_error}，那是容器的行为而非契约本身。
\item \Fn{deallocate} 的两个参数都必须对得上：\Tp{p} 必须是
    \Fn{allocate} 返回且尚未释放的地址，\Tp{n} 必须与当初分配时
    给的个数一致。尺寸不符在带大小提示的重载里尤其致命：
    自定义池按 \Tp{n} 选回收区间，不一致的 \Tp{n} 直接破坏池结构。
\item 相等性谓词只回答一件事：内存是否可互换释放。
\end{itemize}

\begin{clause}[title=标准怎么说：分配器相等的含义]
\Std{allocator.requirements} 规定：两个分配器 \Tp{a1 == a2} 成立，
当且仅当由其中一个分配的存储可以由另一个释放。相等并不要求两个
对象状态相同，也不要求各自分配的块能被对方“重新分配”；它只是
内存等价性的声明。容器在复制、移动、交换时依赖这一声明决定是否
共用分配器。
\end{clause}

分配器携带的状态决定契约要写多严：无状态分配器（不持有会跨模块
悬空的指针与句柄）用默认的传播属性也安全，因为它的
\Tp{operator==} 恒真；一旦分配器带池指针，三条传播属性与相等性
就是同一个决定的四处表态——要么三处都声明不传播并让相等性
互斥，要么三处都声明传播并保证池的生命周期不早于容器。
容器侧的转发由 \Tp{std::allocator{\Bk}\_traits} 统一完成，
契约只能在分配器这一侧写好。

三个传播属性决定复制赋值、移动赋值、交换时分配器跟着不跟。
属性名是共同前缀 \Tp{propagate{\Bk}\_on{\Bk}\_container} 加后缀：
\Tp{\_copy{\Bk}\_assignment} 管复制赋值，
\Tp{\_move{\Bk}\_assignment} 管移动赋值，\Tp{\_swap} 管交换，
三者默认全为 \Tp{false{\Bk}\_type}。容器实现用
\Fn{std::select{\Bk}\_propagate{\Bk}\_allocator}（\cpp{17}）
按这三条挑出赋值后该用哪个分配器。契约里最容易被违约的
组合是：分配器带进程级状态（池指针、模块句柄），传播属性却保持
默认值——交换后一侧容器拿着另一侧模块的分配器继续释放内存，
就是所谓的跨模块释放。后果不是异常而是堆结构损坏，属于
第~\ref{ch:ub04-lifetime} 章之后最贵的一类静默 \UB{}；分配器状态要
跨模块时，三个属性必须逐条表态，谁不传播就把相等性也写成
\Tp{false}。ODR 与模块边界上的更细一层属于编译器与链接器一册，
本章只记契约面。

\section{全局 operator new 的替换约定}\label{sec:ub14-global-new}

全局分配函数是可以被程序替换的一组固定重载（\Std{new.delete.single}
及其相邻小节）：\Fn{operator new}、\Fn{operator new[]}、两者对应的
\Fn{operator delete}，\Tp{nothrow} 版本，带尺寸提示的
\Fn{operator delete}（\cpp{14}），以及带 \Tp{std::align\_val\_t}
参数的对齐版本（\cpp{17}）。替换必须是全局作用域、非内联的定义；
一条配对纪律：替换 \Fn{operator new} 就必须同时替换配套的
\Fn{operator delete}，因为释放路径要认得分配路径产出的块格式；
只换一半等于制造了一个不在册的分配器，跨路径释放即 \UB{}。
\Fn{new\_handler} 与抛出的 \Tp{std::bad\_alloc}（\Hdr{new}）负责
失败路径：抛出版在耗尽重试前反复调用 \Fn{new\_handler}，失败前
的重试轮数正是第~\ref{ch:ub13-implementation} 章列过的 \USB{}，
不要对它做数值假设；\Tp{nothrow} 版返回空指针。

\begin{itemize}
\item 对齐分配的坑：请求对齐超过 \Feat{\_\_STDCPP\_DEFAULT\_NEW\_\Bk\_ALIGNMENT\_}
    时走 \Tp{align\_val\_t} 重载，实现常在块首多留空间再回拨指针；
    替换函数若忽略对齐参数、用普通 \Fn{free} 归还这种块，
    归还的地址与真实块首不符，堆直接损坏。
\item 动态库边界的坑：替换定义只对自己可见的链接单元生效。
    主程序替换了、静态库用默认、第三方 DLL 又带一份自己的替换时，
    同一个指针可能在一侧分配、另一侧释放。标准把“程序级替换”的
    承诺限定得很清楚，多份不一致的替换不在承诺之内；工程上把
    替换函数收在唯一的基础库导出面里，其余模块只链接不调换。
\item \cpp{14} 的尺寸提示是优化通道而非义务：替换版若利用尺寸
    参数归池，必须接受库传来的值与自己分配时用的值一致这一
    前提，任何“猜尺寸”的实现都在制造新的 \UB{}。
\end{itemize}

\section{错误处理约定基线}\label{sec:ub14-error-baseline}

\cpp{17} 的 \Hdr{charconv} 给出了库内错误处理的现代模板：
\Fn{std::from\_chars} / \Fn{std::to\_chars} 不抛异常、不吃 locale、
不分配内存，失败以返回结构里的状态码表达（\Tp{std::errc} 一族，
如 \Tp{invalid\_argument} 与 \Tp{result\_out\_of\_range}）；浮点
支持在 \cpp{23} 加入。对照旧 \Fn{strtol} 系族（\Hdr{cstdlib}）：
成功值与错误值挤在同一个返回值里，越界靠 \Feat{errno} 侧面通报，
而“无转换时 \Feat{errno} 置什么”本身是实现定义的。

\begin{lstlisting}[style=cppstyle, caption={\Tp{from\_chars}：状态即返回值，没有全局副作用}]
#include <charconv>
#include <string_view>
#include <system_error>

std::errc parse_int(std::string_view s, int& out) {
    const char* b = s.data();
    auto [last, ec] = std::from_chars(b, b + s.size(), out);
    if (ec != std::errc{} || last != b + s.size()) {
        return ec == std::errc{} ? std::errc::invalid_argument : ec;
    }
    return std::errc{};                 // 完整消费才算成功
}
\end{lstlisting}

\Tp{std::error\_code}（\Hdr{system\_error}）则是另一端的约定：
把平台错误码与类别 \Tp{std::error\_category} 打包成可比较、可
格式化、不抛异常的小对象，设计意图是让“错误”作为值流转，而不是
作为控制流炸开。它的坑与分配器同构：类别通常以单例引用比较，
跨动态库各持一份静态单例时相等性可能失真，这类模块边界问题
与编译器与链接器一册的 ODR 话题同源。

旧式的 \Feat{errno} 并非一无是处：它不占用返回值通道，遗留代码
接入免费；但它的线程局部形态决定了“调用后立刻读”是唯一正确用法，
错误想跨过函数边界就必须复制，而复制经常被忘记。
\Tp{std::error\_code} 正是把这次复制正式化：带码与类别的小对象
可以存储、比较、放进返回值，错误的解释推迟到消费点再做。

\begin{compare}[title=契约与约定：谁检查、谁解释失败]
契约（contract）规定前置条件由调用方保证，违约即 \UB{}——检查
与否取决于档位；约定（convention）把失败编进返回值或状态对象，
任何档位都可见。分层原则：库内部热路径可以依赖契约并把检查交给
加固档，公开接口必须用约定把失败说清楚。三档对照如下。
\end{compare}

\begin{center}
\begin{tabular}{@{}>{\raggedright\arraybackslash}p{0.22\textwidth} >{\raggedright\arraybackslash}p{0.24\textwidth} >{\raggedright\arraybackslash}p{0.24\textwidth}@{}}
\toprule
检查责任 & 失败形态 & 可见性 \\
\midrule
调用方自证（契约） & 静默 \UB{}，或档位开启后的断言与终止 &
只在调试/加固档与消毒器里可见 \\
库替你检查（调试档位） & 诊断后终止或抛异常 &
开档即见，但依赖构建配置 \\
状态码约定（\Tp{errc} 或 \Tp{error{\Bk}\_code}） & 正常返回附带状态 &
每次调用都可见，跨层可传递 \\
\bottomrule
\end{tabular}
\end{center}

表里的可见性是选档依据：调试档只保护开档的那次运行，状态码在
每一次发布运行里都保护调用方。两者不是互相替代的关系，热路径靠
档位兜底、边界靠状态码说话。

\begin{bestpractice}[title=给项目定一条错误处理基线]
新接口优先选“返回值带状态”的约定族（\Fn{from\_chars} 式、
\Tp{std::error\_code} 式），异常留给确实要炸开控制流的层；
在项目头文件里把这条基线连同所选档位一起写死，调试/加固档的
宏进构建配置而不是散落的源文件。基线的作用不是消灭 \UB{}，
而是保证违约时整个项目用同一种坏消息说话。
\end{bestpractice}

回看本章三块内容与上一章的接口：特性宏是查库“文档义务”的通道，
档位是把检查责任移交给库的通道，分配器与错误约定是谁负责解释
失败的规则表。三类通道指向同一个目标——把库的行为变成成文契约，
再用断言与构建配置把它固定住；固定手法与
第~\ref{ch:ub13-implementation} 章完全同构，只是提问对象从编译器
换成了库。契约写进头文件与构建配置的那一刻，违约的代价就从
运行期事故降级为编译期失败，这正是本章全部约定的共同回报。

\chapter{前向兼容与缺陷报告回译}\label{ch:ub15-forward}

前面十四章把五种行为类别逐个摊开：哪些操作落在 \UB{}、哪些是 \IDB{}、
哪些只是 \USB{}，以及怎么用消毒器与版本探测把它们钉在可测的范围内。
本章收束全册，回答一个更隐蔽的问题：\textbf{判定本身会移动}。
标准的修订不只在追加新特性，核心语言工作组与库工作组还在持续处理缺陷报告：
一条原本含糊的措辞被改写、一处从未言明的边界被补全、一个 \NDR{} 的漏洞被关掉。
这些决议一旦通过，往往以回译的形式并入更早的标准档——
于是同一份源码，在两个出版年份不同的档里，会落入不同的类别。
读这一章时请随时回指前面的章节：本章只做口径的对照，不做机制的重述。

\section{回译：旧标准档不是化石}

\subsection{缺陷报告如何进入旧档}

国际标准化组织对语言的维护由两条工作组线并行推进：一条管核心语言，
一条管标准库（在 \CPP{} 一侧通常被称作核心组与库组）。它们的工作产物是
缺陷报告的决议：对标准文本含糊、矛盾或从未覆盖之处给出统一解释或修补措辞。
关键点在于——这些决议不只写进下一个新版本，而是\textbf{回译进此前已发布的版本}，
形成带增补的同一编号标准的最新文字。换句话说，\cpp{11} 档的“正式文本”
在你今年翻开时，和十年前印出的那一版并不逐字相同。

对行为类别而言，回译的后果有三层：

\begin{itemize}
  \item \UB{} 被改写：某操作原本没有规则可依，决议补出规则后，
        它可能变成良好定义，也可能被明写为 \ILLF{}；
  \item \IDB{} 被收窄：原本留给实现任选的空间被收小，旧代码依赖的
        那一侧选择可能不再是文档承诺之一；
  \item \NDR{} 被关掉：不要求诊断的不合法情形，被决议改写成要求诊断，
        甚至改写成一条硬性的程序格式约束。
\end{itemize}

\begin{clause}[title=标准的活体属性]
标准文档存在“档”与“出版年”的区分：同一档号在不同年份发布过带缺陷报告
增补的修订文字。编译器与库的实现者跟进的是最新文字，而不是某一版印刷稿。
因此“我按 \cpp{14} 写的”这句话，并没有说清判定所依据的措辞是哪一年的。
\end{clause}

\subsection{同一份源码，两个判定}

判定移动最典型的形状是：源码一字未动，只换了编译时所选的档标与实现，
结论就从“没有规则”变成“违反规则”，或从“违反规则”变成“被承诺的行为”。
第~\ref{ch:ub08-evaluation} 章讲过的求值顺序即是此类：\cpp{17} 之前的措辞对大量子表达式的先后
保持沉默，决议与修订之后，一部分先后关系被写成定序要求，另一部分
仍保持 \USB{}。依赖沉默时代的侥幸写法，在新档下不是“更坏了”，
而是\textbf{第一次被规则覆盖}——覆盖它的那条规则未必站在你这边。

\begin{warning}[title=判定移动不等于实现退化]
某段代码在新档、新实现上开始报错，通常不是编译器变差了，而是
标准文字追上了它。排查顺序应当是：先定位这段代码依赖哪一类约定，
再查该约定的措辞在两个档之间是否被缺陷报告改写过，最后才怀疑实现。
\end{warning}

\section{口径变化案例清单}

本节列出五组能确证的口径变化。每组按同一模板叙述：主题、早先口径、
现在的口径、对既有代码的影响，以及本册哪一章把机制讲透。
这里刻意\textbf{不写缺陷报告编号}：编号会随决议合并与回译而换载体，
复述它们不如复述口径本身稳定。

\subsection{求值顺序的收紧}

\textbf{主题}：表达式子项的求值先后与竞争判定。\textbf{回指}：详见第~\ref{ch:ub08-evaluation} 章。

\textbf{早先口径}：早期版本对绝大多数运算符的操作数先后不作承诺，
函数调用实参之间、赋值左右两侧之间谁先谁后全凭实现习惯；
竞争检测规则因此大片落空，许多交错读写只能靠实现的固定习惯来“稳”。
\textbf{现在的口径}：经由措辞重写与后续澄清，\cpp{17} 起一批规则被定序化：
链式输出按从左到右求值、标量赋值的右操作数先于左操作数完成求值排序；
而函数调用实参之间的顺序\textbf{仍然}是 \USB{}，并未因这次收紧被一并规定。
\textbf{对既有代码的影响}：把实参之间的实现习惯当承诺的写法，判定不变但
风险显式化了；在赋值两侧交错修改同一标量的写法，则从“没人管”进入
“有规则可依、且规则判你越界”的状态。回归测试要按档分组固定（见 第~\ref{sec:ub15-testing} 节）。

\subsection{移动后对象的保证范围}

\textbf{主题}：标准容器与字符串被移动之后，哪些操作仍被承诺。
\textbf{回指}：机制与析构路径见第~\ref{ch:ub11-terminate} 章。

\textbf{早先口径}：库的措辞只在个别条目上写了移动后对象处于
“有效但未指明的状态”，而“有效”到底覆盖哪些成员函数，长期含糊；
实践中只有一条最小共识：可赋值、可析构。
\textbf{现在的口径}：库组的一轮清理决议把这个短语统一为全库约定，
并明确它至少承诺：可析构、可赋值、可查询式的常规成员调用，
但每个查询返回什么值属于 \USB{}，标准从不承诺“移动后必为空”。
\textbf{对既有代码的影响}：移动之后立刻调用 \Tp{empty()} 或 \Tp{clear()}
在新口径下是被承诺的操作；而假设移动后 \Tp{size()} 归零的逻辑，
从“碰巧能用”被重新判定为对 \USB{} 的越权解读。

\subsection{视图类对象与悬垂}

\textbf{主题}：非拥有视图（字符串视图、span 类对象）与所指对象寿命的关系。
\textbf{回指}：库前置条件整体见第~\ref{ch:ub12-library-pre} 章。

\textbf{早先口径}：视图类刚引入时，标准只说明它不持有资源；
“通过悬垂视图访问”究竟算违反生命周期、算违反前置条件，还是算 \IDB{}，
文字上没有对齐；\Tp{data()} 是否保证以空字符终止，也一度存在读法分歧，
后来由库组决议明确：不承诺。
\textbf{现在的口径}：视图统一按“引用语义”归位——用它访问已终结的对象，
就是第~\ref{ch:ub04-lifetime} 章生命周期规则的直接违反，判定为 \UB{}，
而不是某个库专有条款问题；终止符承诺仅存在于实现扩展层面。
\textbf{对既有代码的影响}：以视图作参数做隐式转换、再把视图存进成员的
常见写法，风险位置没有变，但判定被提前到了调用点：函数接收视图时并
不为所指数据的寿命续期，悬垂的产生点往往在形参绑定那一刻。

\subsection{生命周期显式化工具的引入}

\textbf{主题}：从裸字节缓冲区安全地“造”出对象。\textbf{回指}：
生命周期全谱见第~\ref{ch:ub04-lifetime} 章。

\textbf{早先口径}：-placement \Tp{new}、把 \Hdr{cstring} 风格的字节缓冲
重新解释成对象指针，这类写法在措辞上长期是灰色地带：要么依仗
字符型访问的宽松条款，要么干脆没有条款可依，全靠实现的默契。
\textbf{现在的口径}：核心组的澄清把一部分“隐式启动生命周期”的情形写成了
规则（标准布局类型经字节存储访问等），库组则补上显式工具
\Tp{std::start{\Bk}\_lifetime{\Bk}\_as} 家族与 \Tp{std::launder}，
让这类转换从“默契”升格为“有名字、有前置条件的操作”。
\textbf{对既有代码的影响}：过去依赖默契的缓冲区复用代码，如今可以被逐处
判定：落在隐式规则内的是良好定义，落在显式工具前置条件内的是良好定义，
两处都不落的仍是 \UB{}——判定的粒度变细，代码没有变，是标准变细了。

\subsection{并发同步语义的补强}

\textbf{主题}：原子对象的构造、销毁与内存序承诺。\textbf{回指}：
内存模型全谱见第~\ref{ch:ub10-memory-model} 章。

\textbf{早先口径}：内存模型初立时，原子对象在构造期间与销毁之后能否
被其他线程安全引用，文字没有正面回答；引用计数场景里“最后一条释放
与析构的交界”只能按实现行为理解。
\textbf{现在的口径}：核心组围绕原子对象销毁语义与引用计数交接做了
专门澄清，把“销毁后仍被访问”明确归入数据竞争之外的另一条 \UB{} 路径；
库侧则增补了等待/唤醒类原子接口与按引用挂接的原子视图类型，
使过去只能自旋降级的写法有了被承诺的替代。
\textbf{对既有代码的影响}：自实现引用计数的库代码，若曾依赖“析构前
还有一次合法读取”的旧默契，新口径下这类读取被明确判为越界；
改用库的新同步设施需要库档支持，判定随实现漂移的问题见下一节。

\begin{compare}[title=五组口径变化的共同形状]
共同点只有一个：源码未动，规则后到。差异在方向——求值顺序是
“沉默改为规定”，移动后对象是“含糊改为承诺”，视图悬垂是
“分歧改为归位”，生命周期工具是“默契改为操作”，并发补强是
“空白改为封死”。五者都不以“新增诊断”为义务，除非实现自愿提供。
\end{compare}

\section{双维度版本：编译器、标准库与语言档}\label{sec:ub15-versions}

\subsection{三个版本号互不相等}

前向兼容排查里最常犯的错误，是把“编译器版本”当作“行为版本”。
实际上参与判定的是三件事：

\begin{enumerate}
  \item \textbf{语言档}：由 \Tp{-std=} 一类选项选定，决定核心语言用哪套
        措辞——求值顺序、生命周期规则的松紧都挂在这里；
  \item \textbf{编译器版本}：同一档下，不同前端对 \UB{} 的推理激进程度
        不同，消毒器的覆盖面也不同（第~\ref{ch:ub03-sanitizers} 章）；
  \item \textbf{标准库版本}：libstdc++、libc++ 与各厂商实现按自己的节奏
        跟进库组决议，与前端版本彼此独立发布。
\end{enumerate}

三者独立演进意味着：\textbf{只换库、不换编译器}，库档判定可以变——例如
同一前端同一 \cpp{17} 档下，升级库实现后，越界的 \Tp{operator[]} 从
默默读到相邻内存变为按断言契约中止；\textbf{只换编译器、不换库}，
核心语言判定可以变——新版本前端按最新回译文字做定序推理，优化改变
了旧代码侥幸成立的形状。探测这两个维度的规范写法在第~\ref{ch:ub14-conventions} 章：编译器三元组宏、\Tp{\_\_cplusplus} 取档、
库侧特性宏与库版本宏分列。

\subsection{把依赖旧行为的地方钉死}

对“本文件刻意依赖某旧口径”的代码，正确做法不是祈祷实现不升级，
而是用 \Tp{static\_assert} 与 \Feat{\_\_cpp\_lib\_*} 特性宏把假设
编译期钉死：库一旦回译到位，构建当场失败，逼人来改写。

\begin{lstlisting}[style=cppstyle, caption={cppstyle：用特性宏把旧行为依赖钉死在编译期}]
// 注意：C++20 起可由 <version> 头统一引入特性宏；
// 在更早的档里需先包含对应功能头文件再检测。
#ifdef __cpp_lib_bit_cast
static_assert(__cpp_lib_bit_cast >= 201806L,
              "bit_cast feature macro does not match the ratified wording");
#endif

#ifndef __cpp_lib_ranges
#error "this file assumes ratified range semantics; upgrade the library"
#else
static_assert(__cpp_lib_ranges >= 201911L, "ranges wording too early");
// 本翻译单元依赖视图在定稿档下的悬垂判定；若特性宏值对不上，宁可拒编。
#endif

// 钉死“旧库档才成立”的写法：库一旦升级就报错，而不是悄悄换语义。
#if defined(__cpp_lib_start_lifetime_as)
static_assert(false, "library now provides explicit lifetime tools; "
                     "rewrite this buffer cast instead of relying on old behavior");
#endif

// 编译器维度与库维度分开钉。
#if defined(__clang__)
static_assert(__clang_major__ >= 9, "sanitizer coverage too old for this file");
#elif defined(__GNUC__)
static_assert(__GNUC__ >= 9, "UBSan instrumentation incomplete here");
#endif
\end{lstlisting}

\begin{note}[title=特性宏只证实现跟进到哪一步]
\Feat{\_\_cpp\_lib\_ranges}、\Feat{\_\_cpp\_lib\_bit\_cast} 这类宏报告的是
\textbf{库实现采纳定稿措辞的程度}，不是语言档本身；而 \Tp{\_GLIBCXX\_ASSERTIONS}
之类库侧宏反映的是实现自愿加装的契约检查。用它们做回归开关时，
三者要各自记录，缺一角就是一次无法复现的比较。
\end{note}

\subsection{换库与换编译器的排查顺序}

同一 UB 判定在版本漂移中被改变时，先用一段脚本把三个维度固定成
一个可比较的元组，再谈“变没变”。

\begin{lstlisting}[style=shellstyle, caption={shellstyle：固定三元组并打印判定所依据的档位}]
# tuple = (语言档, 编译器版本, 库版本)；三者齐全才算同一次判定
g++ --version | head -n1
echo '__cplusplus = ' ; printf '#include <cstdio>\nint main(){printf("%ld",(long)__cplusplus);}\n' | g++ -std=c++17 -x c++ - -o /tmp/cxxlang && /tmp/cxxlang
g++ -std=c++17 -E -dM - </dev/null | grep -E '__GNUC__|__clang_major__' || true
# 库维度：让预处理器自己回答
printf '#ifdef _GLIBCXX_RELEASE\n_GLIBCXX_RELEASE\n#endif\n' | g++ -E -x c++ - | tail -n1
\end{lstlisting}

MSVC 一侧还要多钉一枚标志：\Tp{/Zc:\_\_cplusplus} 让 \Tp{\_\_cplusplus}
如实反映所选档，否则它长期报告旧值，库档探测会被误导——
这正是第~\ref{ch:ub13-implementation} 章所讲“\IDB{} 面”的实例。

\section{测试面：回归判定按实现固定}\label{sec:ub15-testing}

\subsection{为什么不能跨实现比较}

回译把判定的“依据文字”统一了，但没有把判定的“暴露方式”统一：
同一段违反生命周期的代码，A 实现给出消毒器报告，B 实现给出静默正确，
C 实现给出崩溃。三种表现对应同一条 \UB{} 判定。因此回归判定必须
\textbf{按实现固定}：期望结果与（档、编译器版本、库版本、检查开关）
四元组绑定，只允许同组内的前后比较。跨实现的“它也报了吗”不是回归
证据，只是又一次抽样。

\begin{warning}[title=以前能跑不是证据]
重述第~\ref{ch:ub02-inevitable} 章的论断：一次没有触雷的运行，
不构成“行为被承诺”的证据。旧版本上“能跑”只说明该次输入、该次优化
没有暴露规则；回译之后规则补上，暴露概率并不由你的历史运行决定。
把“历史通过”写进回归基线，等于把 \UB{} 洗成假想的 \IDB{}。
\end{warning}

\subsection{把一个 UB 缺陷写成可复现的消毒器用例}

与第~\ref{ch:ub03-sanitizers} 章的 CI 分档衔接，流程固定为六步：

\begin{enumerate}
  \item \textbf{最小化}：把触发代码削减到单一翻译单元，输入写死在源码里，
        不允许读取环境；
  \item \textbf{归类}：明确它落在五类的哪一类（第~\ref{ch:ub01-taxonomy} 章，
        第~\ref{ch:ub01-taxonomy} 章），并写下判定所依据的档与措辞主题；
  \item \textbf{选档}：按类别挑消毒器：整数与指针越界走未定义行为档，
        容器前置条件与生命周期走地址档，并发走线程档；
  \item \textbf{固定三元组}：在用例头部注释里钉死语言档、编译器版本、
        库版本与消毒器标志，缺一不予合并；
  \item \textbf{进分档}：按快速门禁、夜间全量、发布前取证三档摆放，
        期望是“报告必然出现”，而不是“退出码为某值”；
  \item \textbf{留档}：把报告形状（非实测数据）记为模式而非整行比对，
        避免工具措辞升级把门禁变成噪声源。
\end{enumerate}

\begin{lstlisting}[style=shellstyle, caption={shellstyle：六步流程落到 CI 三档的骨架}]
# 快速档：只挂未定义行为消毒器，要求出现报告即失败
timeout 120 g++ -std=c++17 -O1 -g -fsanitize=undefined,pointer-overflow \
    repro.cpp -o repro && ./repro
# 夜间档：地址档与线程档全量，收集报告模式供版本间比对
for f in tests/repro/*.cpp; do
  g++ -std=c++17 -O1 -g -fsanitize=address "$f" -o "/tmp/$(basename $f .cpp)" || exit 1
done
# 发布前档：固定三元组，把判定基线与工件一同封存
printf 'lang=17 cc=%s lib=%s flags=ubsan,asan\n' "$(cc -dumpversion)" "$(_libver)" >> build.lock
\end{lstlisting}

\begin{casebox}[title=同一越界依赖在两个标准档下的判定链]
源码片段：取“尾后一格的指针形成”\Tp{\&a[n]}（\Tp{n} 恰为数组长度），
随后\textbf{不解引用}，仅用于比较或作减法锚点。判定链：
\begin{itemize}
  \item 旧档口径：下标与指针算术的措辞把这类形成排除在外，属于 \UB{}；
        未挂消毒器、未开检查的旧实现通常“看起来能用”，因为它只是
        算了一个地址；
  \item 新档口径：数组下标语义被重写后，尾后指针的\textbf{形成}本身成为
        良好定义的操作——同一行代码从 \UB{} 判定升格为有规则可依；
        这是“从报错面变回能用面”的方向；
  \item 反向情形：对 \Tp{a[n]} 的\textbf{解引用}在两个档下都仍是 \UB{}，
        但新版未定义行为消毒器把它变成显式报告——从“能用”翻到“报错”；
  \item 结论：换档改变的是判定所在类别与诊断面，不是实现好坏；
        回归基线必须按 第~\ref{sec:ub15-versions} 节 的三元组分组。
\end{itemize}
\end{casebox}

消毒器报告的可复现形状大致如下（示意，非本仓库实测值）：

\begin{lstlisting}[style=outstyle, caption={outstyle：报告形状的模式化记录}]
# 只记录模式：文件与行号随重构移动，不作逐行比对基线
runtime error: index 5 out of bounds for type 'int [5]'
SUMMARY: UndefinedBehaviorSanitizer: undefined-behavior repro.cpp:8:6
\end{lstlisting}

\section{收束：本册的三条底线}\label{sec:ub15-closing}

十五章读到这里，方法论收敛成三句话。每句都是一个可执行的门禁，
而不是态度。

\begin{guideline}[title=底线一：契约可见]
每个操作都要能回答：它落在五类（良好定义、\UB{}、\IDB{}、\USB{}、
\ILLF{} 及其 \NDR{} 细分）中的哪一类，判定依据是哪一档的哪段措辞。
回答不出来的调用点，就是回译最可能偷袭的位置。
\end{guideline}

\begin{guideline}[title=底线二：检查可跑]
消毒器档位与库契约开关必须进 CI，而不是“本地可以试”：
未定义行为档、地址档、线程档各有覆盖面，库档（断言与调试容器）
另算一维。检查跑不到的分支，等于没有判定。
\end{guideline}

\begin{guideline}[title=底线三：行为可测]
回归判定按（语言档、编译器、库、开关）四元组固定并可复现：
六步流程产出的用例，任何人任何机器重跑都应落在同一判定类别上。
跨实现的比较、无基线的“以前能跑”，都不算测到。
\end{guideline}

最后把全册十五章各自回答的问题与附录的落点串成一张表：
正文讲机制，附录给清单；遇到判定移动，先回本表定位章与附录再取证。

\begin{longtable}{@{}>{\raggedright\arraybackslash}p{0.26\textwidth} >{\raggedright\arraybackslash}p{0.38\textwidth} >{\raggedright\arraybackslash}p{0.16\textwidth}@{}}
\caption{十五章与它回答的问题及对应附录总览}\label{tab:ub15-map}\\
\toprule
章 & 它回答的问题 & 对应附录 \\
\midrule
\endfirsthead
\toprule
章 & 它回答的问题 & 对应附录 \\
\midrule
\endhead
\bottomrule
\endfoot
第~\ref{ch:ub01-taxonomy} 章 分类学 & 一个操作落在五类中的哪一类 & \ref{ap:ubA-kinds} \\
第~\ref{ch:ub02-inevitable} 章 不可避免 & 哪些 UB 无法从实践中根除 & \ref{ap:ubC-cases} \\
第~\ref{ch:ub03-sanitizers} 章 消毒器 & 检查如何跑进 CI 分档 & \ref{ap:ubB-switches} \\
第~\ref{ch:ub04-lifetime} 章 生命周期 & 对象何时存在、何时终结 & \ref{ap:ubA-kinds} \\
第~\ref{ch:ub05-pointer} 章 指针 & 哪些指针值与算术有承诺 & \ref{ap:ubA-kinds} \\
第~\ref{ch:ub06-aliasing} 章 别名 & 严格别名与放松开关的边界 & \ref{ap:ubB-switches} \\
第~\ref{ch:ub07-alignment} 章 对齐 & 未对齐访问在哪些实现上会翻车 & \ref{ap:ubC-cases} \\
第~\ref{ch:ub08-evaluation} 章 求值顺序 & 定序收紧后谁在违反竞争规则 & \ref{ap:ubC-cases} \\
第~\ref{ch:ub09-integers} 章 整数 & 溢出、移位与除零的判定面 & \ref{ap:ubA-kinds} \\
第~\ref{ch:ub10-memory-model} 章 内存模型 & 数据竞争与内存序的承诺范围 & \ref{ap:ubC-cases} \\
第~\ref{ch:ub11-terminate} 章 终止 & 异常路径与移动后保证的边界 & \ref{ap:ubC-cases} \\
第~\ref{ch:ub12-library-pre} 章 库前置条件 & 视图与前置条件违反如何归位 & \ref{ap:ubC-cases} \\
第~\ref{ch:ub13-implementation} 章 实现定义面 & \IDB{} 去哪里查、怎么固化 & \ref{ap:ubB-switches} \\
第~\ref{ch:ub14-conventions} 章 约定与探测 & 版本与特性宏怎么写成门禁 & \ref{ap:ubB-switches} \\
第~\ref{ch:ub15-forward} 章 回译 & 判定本身如何随标准修订移动 & 三附录合参 \\
\end{longtable}

\begin{bestpractice}[title=把本章装进代码库]
在仓库里维护一份判定基线：每个消毒器用例绑定其四元组，每条
“依赖旧口径”的注释引用它等待回译的措辞主题。标准文字移动时，
失败应当首先出现在基线比对里，而不是出现在生产环境里。
\end{bestpractice}

\appendix
\chapter{五种行为类别速查}\label{ap:ubA-kinds}

前面各章按机制纵向讲未定义行为，查起来要先知道机制；附录~\ref{ap:ubA-kinds} 反过来，按语言构造横向查。
一行一个构造，五个字段依次是：构造本身、它落进哪一类行为、依据的条款主题、
能被什么工具抓到、以及不用它也能把活干完的写法。
凡标准里编号可能随版本移动的位置，本表一律写成“某小节，见正文章”，不写死条款号；
类别一列里出现“良定义”的行，表示它常被误当成 \UB{}，实际标准给的是确定结局，
多数是调用 \Fn{terminate}。

\begin{guideline}[title=这张表怎么用]
先判类别，再决定动作：\UB{} 与 \ILLF{} 必须改代码；\IDB{} 只要把依赖写进移植文档，
并至少在两个实现上验证；\USB{} 通常什么都不用改，只要别假设它是某一个固定值；
\NDR{} 最危险，因为它连诊断都不承诺。检测手段一列只给消毒器族与检查方向，
具体开关名、能否同时打开与失败策略见 附录~\ref{ap:ubB-switches}。
\end{guideline}

\begin{clause}[title=五个词的边界]
\UB{} 是标准撤回全部承诺；\IDB{} 是标准要求每个实现自选一个答案并写进文档；
\USB{} 同样自选，但不要求写进文档；\ILLF{} 是程序本身不合法且要求实现给出诊断；
\NDR{} 是不合法但不要求诊断。判定的难点从来不在背下这五句，
而在认出一段代码到底落进哪一句。同一句话在 \CL{} 与 \CPP{} 里覆盖的范围并不相同，
本表凡有分歧都会写进类别一列。
\end{clause}

\begin{longtable}{@{}>{\raggedright\arraybackslash}p{0.16\textwidth} >{\raggedright\arraybackslash}p{0.09\textwidth} >{\raggedright\arraybackslash}p{0.17\textwidth} >{\raggedright\arraybackslash}p{0.19\textwidth} >{\raggedright\arraybackslash}p{0.19\textwidth}@{}}
\caption{语言构造到行为类别的速查表}\label{tab:apA-kinds}\\
\toprule
语言构造 & 类别 & 依据条款主题 & 检测手段 & 可移植替代写法 \\
\midrule
\endfirsthead
\toprule
语言构造 & 类别 & 依据条款主题 & 检测手段 & 可移植替代写法 \\
\midrule
\endhead
\bottomrule
\endfoot
有符号加、减、乘溢出 & \UB{} & 算术结果无法用结果类型表示，见 \Std{expr{\Bk}.mul} 一类小节与第~\ref{ch:ub09-integers} 章 & \Mx{UBSan} 的 signed-integer-overflow 档；陷入式手段另有一档 & 先在无符号域判溢出再转回，或写显式饱和函数 \\
除法与取模的零除数 & \UB{} & 第二操作数为零，见 \Std{expr{\Bk}.mul} & \Mx{UBSan} 的 integer-divide-by-zero 档；浮点另有同类档 & 先判非零再除，把判断与运算分成两条语句 \\
\Tp{INT{\Bk}\_MIN} 除以 \Tp{-1} & \UB{} & 商超出可表示范围，仍属溢出 & \Mx{UBSan} 报成溢出还是报成除法随实现 & 先特判最小值，或转成更宽类型再除 \\
移位量不小于类型宽度 & \UB{} & \Std{expr{\Bk}.shift} 对右操作数的约束 & \Mx{UBSan} 的 shift-exponent 档 & 移位量先夹紧到零与宽度减一之间，或用 \Hdr{bit} \\
移位量为负 & \UB{} & \Std{expr{\Bk}.shift} & \Mx{UBSan} 的 shift 族 & 移位量取无符号并在入口处校验 \\
负值左移 & \UB{} & \Std{expr{\Bk}.shift}，与符号有关 & \Mx{UBSan} 的 shift-base 档 & 先转无符号做位移，再转回有符号 \\
读取未初始化标量 & \UB{} & 不确定值的读取，见 第~\ref{ch:ub04-lifetime} 章 & \Mx{MSan} 最直接；编译器告警与 \Mx{MSVC} 分析各覆盖一部分 & 声明即初始化，别信“紧跟着就会被赋值” \\
数组下标越界读写 & \UB{} & 解引用必须落在有效对象内，见 第~\ref{ch:ub05-pointer} 章 & \Mx{ASan} 的栈与全局越界报告；\Mx{UBSan} 的 bounds 档 & \Fn{at}、\Tp{std::{\Bk}span} 或带索引校验的访问封装 \\
解引用尾迭代器与越界迭代器 & \UB{} & 迭代器的有效区间要求，见 第~\ref{ch:ub12-library-pre} 章 & 容器调试模式与断言模式；\Mx{MSVC} 侧靠安全特性开关 & 改写为区间或 range-for，禁止手写指针比较到尾 \\
悬垂指针与悬垂引用 & \UB{} & \Std{basic{\Bk}.life}，见 第~\ref{ch:ub04-lifetime} 章 & \Mx{ASan} 的返回后使用；编译器的悬垂告警与 \Mx{MSVC} 分析 & 按值返回，或把生命期交给所有权对象 \\
析构之后继续访问 & \UB{} & \Std{basic{\Bk}.life} & \Mx{ASan} 的释放后使用报告 & 缩小作用域；所有权结束时让句柄立刻失能 \\
placement 构造之后沿用旧指针而未 launder & \UB{} & \Std{support{\Bk}.{\Bk}launder}，见 第~\ref{ch:ub04-lifetime} 章 & 无自动检测，只能人工审查指针来源 & 复用前调用 \Fn{std::{\Bk}launder}，或干脆重取指针 \\
\Tp{bool} 既非真也非假的左值使用 & \UB{} & 基本类型的取值域，见 第~\ref{ch:ub01-taxonomy} 章 & \Mx{UBSan} 的 bool 档 & 用 \Fn{memcpy} 把位模式规范成合法值后再用 \\
读取含陷阱表示的类型 & \UB{} & 值表示与陷阱表示，见 第~\ref{ch:ub01-taxonomy} 章 & 几乎无从稳定检测，只能靠约定 & 只在有明确定义的位模式之间转换，例如 \Fn{std::{\Bk}bit{\Bk}\_cast} \\
严格别名违规 & \UB{} & 对象类型与指针别名，见 \Std{basic{\Bk}.com{\Bk}pound} 与 第~\ref{ch:ub06-aliasing} 章 & 别名告警只在开启优化后给出部分形态 & \Fn{memcpy}、位模式转换设施，或走无符号字节的豁免通道 \\
对齐不足的引用与解引用 & \UB{} & 对齐要求，见 第~\ref{ch:ub07-alignment} 章 & \Mx{UBSan} 的 alignment 档；部分目标另有对齐告警 & 先拷贝到已对齐的局部量，再取引用或指针 \\
packed 结构成员取地址 & \UB{}，打包属性本身是实现扩展 & 成员地址须满足其对齐，见 第~\ref{ch:ub07-alignment} 章 & 编译期的 packed 成员告警，加 \Mx{UBSan} 的 alignment 档 & 把成员值取到局部变量，再对局部取地址 \\
跨对象指针算术 & \UB{} & 指针算术的有效区间，见 第~\ref{ch:ub05-pointer} 章 & 无常规手段；\Mx{Clang} 有指针溢出子档，形态随实现 & 用下标与区间表达，指针只做单对象访问 \\
无关指针比大小 & \UB{} & 关系运算符要求同一对象域内 & 无自动检测，评审才靠得住 & 用 \Fn{std::{\Bk}less}，或在 \Tp{std::{\Bk}uint{\Bk}ptr{\Bk}\_t} 上比较 \\
普通 \Tp{char} 的符号性 & \IDB{} & 基本类型表示的实现定义部分，见 第~\ref{ch:ub13-implementation} 章 & 无开关可查，只能靠跨平台测试矩阵 & 要位模式就写 \Tp{unsigned} \Tp{char}，要数值就先转 \Tp{int} \\
枚举底层类型与超范围取值 & \IDB{} 与 \UB{} 交界 & 底层类型可由实现选定，超范围读入受限 & \Mx{UBSan} 的 enum 档；严格枚举假设会收紧优化 & 用带作用域枚举，读入时先校验再转换 \\
实参与部分操作数的求值顺序 & \USB{} & 多数实参次序未指定，布尔、逗号与列表已排序，见 第~\ref{ch:ub08-evaluation} 章 & 无（每种答案都合法，只是不固定） & 把子表达式落到具名变量，一语句只做一件事 \\
同一表达式内对同一标量重复副作用 & \UB{} & \Std{intro{\Bk}.ex{\Bk}ecution} 的定序约束，见 第~\ref{ch:ub08-evaluation} 章 & 序列点告警只覆盖能静态看出的子集 & 拆成两条语句，读与写分离 \\
自移动，把对象移动到自身 & \UB{} & 移动赋值对源对象未作承诺，见 第~\ref{ch:ub12-library-pre} 章 & 无稳定检测，\Mx{ASan} 只是偶发命中 & 自赋值与自移动检查，或改用 swap 惯用法 \\
moved-from 之后读取内容 & \USB{} & 有效但未指定的状态，见 第~\ref{ch:ub12-library-pre} 章 & \Mx{clang-tidy} 的移动后使用类检查 & 移动后立刻赋新值，或限制在原作用域内用完 \\
\Tp{noex{\Bk}cept} 函数内抛出 & 良定义 & 标准给出确定结局：调用 \Fn{ter{\Bk}mi{\Bk}nate}，见 \Std{except{\Bk}.ter{\Bk}minate} 与 第~\ref{ch:ub11-terminate} 章 & 不需要消毒器，运行即终止 & 把可能抛出的语句移到边界之外 \\
构造或析构期调用纯虚函数 & \UB{} & 构造期虚分派表状态，见 第~\ref{ch:ub04-lifetime} 章 & 静态分析能覆盖一部分；运行期常直接崩 & 两阶段初始化，派生类就绪后再调用 \\
析构函数抛出且异常外溢 & 良定义 & 同样是调用 \Fn{ter{\Bk}mi{\Bk}nate}，见 \Std{except{\Bk}.ter{\Bk}minate} & 运行即终止，没有报告可取 & 析构内捕获全部异常，转为返回值或日志 \\
异常跨越 \Tp{noex{\Bk}cept} 边界 & 良定义 & 调用 \Fn{ter{\Bk}mi{\Bk}nate}，见 \Std{except{\Bk}.ter{\Bk}minate} 与 第~\ref{ch:ub11-terminate} 章 & 运行即终止 & 边界处设 catch-all，转成错误码后再返回 \\
违反库条目的 Requires 与前置条件 & \UB{} & 库契约，见 第~\ref{ch:ub12-library-pre} 章 & 调试模式与硬化模式；两侧标准库各有独立开关 & 调用前显式校验，或改用会抛的 \Fn{at} 一族 \\
\Tp{string{\Bk}\_view} 指向已析构缓冲 & \UB{} & \Std{basic{\Bk}.life}，见 第~\ref{ch:ub04-lifetime} 章 & \Mx{ASan}；\Mx{clang-tidy} 有视图悬垂类检查 & 返回拥有内容的 \Tp{std::{\Bk}string}，或保证宿主存活 \\
分配器与释放函数不匹配 & \UB{} & \Std{basic{\Bk}.stc{\Bk}.dyn{\Bk}amic{\Bk}.deal{\Bk}location} & \Mx{ASan} 的分配类型不匹配；编译器有配对告警 & 交给容器与智能指针，裸 \Fn{new} 只出现在工厂点 \\
\Fn{delete} 与 \Fn{del{\Bk}ete[]} 错配 & \UB{} & 同上，见 第~\ref{ch:ub12-library-pre} 章 & 同上；部分实现只在释放时报尺寸不符 & 用 \Tp{std::{\Bk}unique{\Bk}\_ptr} 的两个特化，或改用 \Tp{vector} \\
data race，非原子并发访问同一内存 & \UB{} & \Std{res.on{\Bk}.data{\Bk}.races}，见 第~\ref{ch:ub10-memory-model} 章 & \Mx{TSan}，必须全量重编译且与地址消毒器互斥 & 锁保护或原子变量，二选一并写清 \\
非法的 \Tp{memory{\Bk}\_order} 组合 & \UB{} & \Std{atom{\Bk}ics{\Bk}.order} 对各操作允许顺序的约束 & 无稳定检测，编译期只有少数形态能拒绝 & 从 \Tp{seq{\Bk}\_cst} 起步，确有证据再降级 \\
把 \Tp{vol{\Bk}atile} 当同步手段 & 仍是 \UB{} & \Std{vol{\Bk}atile} 只约束访问可见性，不提供原子与定序 & \Mx{TSan} 照报竞争 & 改用原子类型与显式栅栏 \\
\Fn{rand} 的序列与线程安全 & 序列是 \IDB{}，并发调用另有 \UB{} & 库随机数的实现定义范围，见 第~\ref{ch:ub13-implementation} 章 & \Mx{TSan} 报竞争；\Mx{clang-tidy} 有并发不安全函数检查 & 每线程一个显式引擎，配分布对象使用 \\
\Fn{qsort} 的比较函数抛出异常 & \UB{} & 经不抛出的函数指针类型传播异常，见 第~\ref{ch:ub11-terminate} 章 & \Mx{clang-tidy} 的异常逃逸类检查；编译期告警形态随实现 & 比较函数写成不抛出的静态函数，或改用 \Fn{std::{\Bk}sort} \\
区域设置与窄宽字符转换失败 & \IDB{} & 区域设置状态与可表示字符集合，见 第~\ref{ch:ub13-implementation} 章 & 无，只能靠固定测试向量 & 显式指定转换设施并检查失败返回值 \\
在动态库边界之外释放内存 & \UB{} 或 \IDB{} & 分配器作用域与运行时链接选项，见 第~\ref{ch:ub14-conventions} 章 & \Mx{ASan} 部分命中；多数表现为堆损坏 & 谁分配谁释放，跨边界只传值或走工厂 \\
本可编译期固定的不变式 & 不属 \UB{}：是 \ILLF{} & 常量求值与 \Tp{static{\Bk}\_assert}，见 第~\ref{ch:ub15-forward} 章 & 编译器直接诊断，无需运行 & 能固定的一律写成 \Tp{static{\Bk}\_assert} \\
\end{longtable}

\section{判定依据逐类说明}

这张表最容易起争执的是类别一列。下面按六组把判定理由摊开写清楚，
每一段对应表里若干行，读者可以把这几段当成表格的脚注来读。
读的时候请记住一个共同前提：类别不是对后果严重性的评级，
而是对“标准在这一句上还留没留承诺”的描述。

\subsection{内存类}

内存类里“越界”“悬垂”“生命周期之外”看着是三种毛病，其实是同一把尺子：
解引用一个指针或引用时，标准要求它指向一个\emph{当前存在的对象}，
且访问落在该对象的存储区间之内。数组下标越界之所以是 \UB{}，
不是因为会读到脏数据，而是因为越界那个位置根本没有被声明的对象；
解引用尾迭代器同理，尾按定义不对应任何元素，也就没有可供指向的对象。
这就是表里把这三行写成同一依据的原因，细节见第~\ref{ch:ub04-lifetime} 章与第~\ref{ch:ub05-pointer} 章。

\Fn{std::launder} 那一行最常被质疑：明明对象就在原来的存储地址上重建了，
为什么还算 \UB{}。判定的落点不在“对象在不在”，而在“旧指针是否还被允许当作新对象的指针”。
标准只在若干条件同时满足时允许沿用同一地址上的旧指针，其余情形撤回承诺。
把它当成“编译器可能仍然记得旧对象身份”的提醒来用就够，具体边界见第~\ref{ch:ub04-lifetime} 章。

\Tp{bool} 与陷阱表示两行属于同一类问题的两面：位模式落在类型域之外。
它们的危险在于代码通常不会自己造出坏值，而是从序列化缓冲区、共享内存或联合体里读出来，
所以检测手段那列特意把“有消毒器档可开”与“只能靠约定”分开写。
读这类数据时的正确姿势是先落到有明确定义的整数类型，校验之后再转成目标类型。

内存类这一组的检测手段最容易记错归属：边界越界与释放后使用属内存族，
未初始化属另一支内存族，语义层面的越界下标属语义族，三族报出来的形态互不替代。
把“越界”一律推给同一个开关，是很多团队跑了一年还漏掉悬垂引用的真正原因，
各族的分工见 附录~\ref{ap:ubB-switches}。

\subsection{求值类}

求值类要分两个动作：定序（sequenced-before）与求值顺序（evaluation order）。
前者管“同一标量的副作用与值计算之间必须有确定先后”，违反它就是 \UB{}，
于是自增与副作用重复那一行落在 \UB{}；后者管“先算左操作数还是右操作数”，
标准不指定，但两种答案都合规，于是实参求值顺序落在 \USB{}。
判定时问自己一句：换一种顺序，程序的结果是仅仅\emph{不同}，还是根本\emph{非法}。
前者是 \USB{}，后者是 \UB{}。详见 第~\ref{ch:ub08-evaluation} 章。

求值类还有一个反复出现的误判：把“今天所有编译器都自右向左”当作契约。
\USB{} 的定义就是实现可以自选且不要求写进文档，
所以这类代码在换编译器或换优化档位时给出的不是错误，而是另一种合法结果，
而下游断言往往只认识其中一种。要固定次序，唯一的手段是把中间结果落成具名对象。

\Tp{volatile} 那一行也常被归进求值类：它保证的是“每次求值都做真实访问”，
不保证原子性，也不提供跨线程定序。因此它与并发类同判，
不能因为“看起来像顺序问题”就降级成 \USB{}。

\subsection{整型类}

整型类的判定核心是一句话：溢出无处安放，标准不给有符号回绕的承诺，回绕只在无符号域成立。
表头四行（加减乘、零除、\Tp{INT\_MIN} 除以 \Tp{-1}、移位越界）都由此直接得出。
真正需要额外解释的是“为什么优化器敢用溢出反推条件”：既然溢出是 \UB{}，
优化器有权假设任何会溢出的位置永不到达，于是用“和是否比加数小”来判分支的代码
在高档位优化下可能被整段删掉。这类现象的机制在 第~\ref{ch:ub02-inevitable} 章 展开，
附录~\ref{ap:ubC-cases} 里给了可直接复现的最小例子。

\Tp{char} 符号性与枚举两行写成 \IDB{}，理由是实现必须自选一个答案并写进文档：
平台文档会说明普通 \Tp{char} 与哪一种显式类型同义，枚举的底层类型也能由实现选定。
一旦代码依赖某个具体选择，可移植性缺口与位图越界会同时出现，
所以替代写法那列强调的是“显式”，而不是“别依赖”。
无作用域枚举的取值范围还额外牵扯到底层类型选择，跨实现差异比 \Tp{char} 更隐蔽。

整型类还有一条实用的分界判据：同一处减法，操作数是无符号时下溢是良定义的回绕，
操作数是有符号时超出范围才是 \UB{}。很多“循环多跑一圈”的缺陷其实是无符号回绕
把越界读写放大出来的，此时判 \UB{} 的是越界那一步而不是回绕本身，
附录~\ref{ap:ubC-cases} 里给了一例。写判定说明时把这两步分开写，比笼统写“整型问题”更容易复查。

\subsection{并发类}

并发类只有一句主干：两个线程访问同一内存位置，至少一个是写，且都不是原子操作，
就构成 data race，而 data race 是 \UB{}，不是实现定义。
判成 \UB{} 而非 \IDB{} 的理由很实在——如果它只是实现定义，
那么每份实现文档都必须给出确定的竞争结果，而标准明确拒绝这一承诺。
详见 第~\ref{ch:ub10-memory-model} 章。

\Tp{memory\_order} 那一行需要单独说明：合法顺序的组合是规定好的，
非法组合落进撤回承诺的区域。因为这类形态极少且多数能静态看出，
表里写“无稳定检测”，别指望消毒器兜底。反过来，合法的顺序但语义用错
（例如该用获取的地方用了松弛）不是 \UB{}，而是普通的功能缺陷，
两者在报告里的表现完全不同，不要混在一条问题单里。

\Fn{rand} 那一行写成两类叠加：序列与周期属 \IDB{}，查文档可得；
多线程共用同一份全局状态则构成竞争，落到 \UB{}。
一行里出现两个类别时，判定动作按更严格的那个走，也就是先解决竞争再谈序列。
并发类还有一个常见误解值得写在这里：把 \Mx{TSan} 报出来的竞争当成“性能问题”或“偶发噪声”，
实际上标准在这里撤回了全部承诺，任何依赖竞争结果实现的缓存、旗标或无锁结构，
都可能在下一次编译器版本升级后给出完全不同的行为，具体手段的边界见 附录~\ref{ap:ubB-switches}。

\subsection{库契约类}

库契约类的判定标准是：条目自己的 Requires、调用前约束与调用后约束就是标准的一部分，
违约的后果直接是 \UB{}，而不是“未定义在实现文档里”。
这就是把库前置条件、迭代器有效区间与容器越界写成同一依据的原因，
也解释了检测手段那列为什么首选调试与硬化模式而非消毒器：
\Opt{D{\Bk}\_GLIBCXX{\Bk}\_DEBUG}、\Opt{D{\Bk}\_GLIBCXX{\Bk}\_ASSERTIONS}
与 \Mx{MSVC} 的安全特性开关补的正是这一层，开关的完整族名见 附录~\ref{ap:ubB-switches}。

moved-from 行判成 \USB{} 值得强调：移动之后对象仍然有效，标准只说取值未指定。
读它是 \USB{}；读它之后还假定它是原值，才是真正的问题。
表里把“读取”与“自移动”拆成两行，正因为后者让源对象的引用指向自身成员，
标准对这种形态没有任何承诺，于是判 \UB{}。

\Fn{qsort} 抛出、分配器不匹配、跨动态库边界释放三行的类别不完全相同，但动作完全一致：
把所有权与调用约定收进同一个模块。判定的细节在第~\ref{ch:ub14-conventions} 章，
库前置条件的展开在第~\ref{ch:ub12-library-pre} 章。
至于 \Tp{noexcept} 违约、异常越界与析构抛出这三行判成“良定义”，
是因为标准在这里给的不是撤回承诺，而是一条明确规则：调用 \Fn{terminate}。
它们的工程代价是崩溃而非误算，别把它记成 \UB{}。

\subsection{实现定义类}

\IDB{} 与 \USB{} 最容易抄错，区别只在“有没有文档”。
\Tp{char} 符号性、\Tp{wchar\_t} 宽度、\Fn{rand} 的序列、窄宽转换的可表示字符集合、
对齐上限与运行时链接选项都属 \IDB{}：标准承认答案由实现选，但要求实现把选择写出来。
函数实参求值顺序属 \USB{}：标准连文档都不要求。这条界线决定了动作差异——
\IDB{} 可以查、可以测两个实现然后写进移植文档；\USB{} 只能改代码把次序固定下来。
详见 第~\ref{ch:ub13-implementation} 章。

跨动态库边界释放那一行写成“\UB{} 或 \IDB{}”，是本册刻意保留的一处不确定：
它取决于两侧运行时与堆管理器是否同一、链接选项如何配置，
而这些细节在不同实现文档里给出的答案并不一致。本册不下死结论，
只把“谁分配谁释放”作为通用规则，具体到某工具链请以实现文档为准。
同理，本表所有“依据条款主题”一列都刻意不给条款编号：
标准各版本之间小节会挪动，编号写死反而更容易误导。

\begin{warning}[title=类别与严重性无关]
判成 \IDB{} 或 \USB{} 不等于安全。\Tp{char} 符号性只在换平台时咬人，
换平台那天所有位图索引一起坏；moved-from 读取只在换实现时才变，
换实现那天它给出的是完全没测过的值。类别说的是标准给不给承诺，不是后果多严重。
\end{warning}

\begin{table}[htbp]
\centering
\begin{tabularx}{0.80\textwidth}{@{}>{\raggedright\arraybackslash}p{0.20\textwidth} X@{}}
\toprule
常被混为一谈的一组 & 拆开的依据 \\
\midrule
\UB{} 与 \NDR{} & 前者是程序合法、但某处含义未被定义；后者是程序本身不合法，只是没人报错。
日常被咬到的是前者，做移植审计时难发现的是后者，因为它连诊断都不要求。 \\
越界与未初始化 & 都是“读到的不是你要的值”，但入口不同：越界看那块存储有没有对应对象，
未初始化看对象存在但其值不确定。检测手段因此分属 \Mx{ASan}、\Mx{UBSan} 与 \Mx{MSan} 三族，
见 附录~\ref{ap:ubB-switches}，最小复现见 附录~\ref{ap:ubC-cases}。 \\
良定义与“不会坏” & 调用 \Fn{terminate} 是确定的结局，不是未定义；但它在生产环境里同样是崩溃。
判成良定义只意味着不必再找消毒器，不意味着可以放着不管。 \\
\bottomrule
\end{tabularx}
\caption{三组最常被混判的边界}\label{tab:apA-border}
\end{table}

\section{判定流程}

\begin{enumerate}
\item 先定位这段代码触及的语言构造：是算术、解引用、生命周期、定序、并发还是库调用。
构造选错了，后面每一步都会跟着错。
\item 查本表的类别列，确认标准在这里是否还留承诺。留了但由实现选，是 \IDB{}；
留了但不写文档，是 \USB{}；没留，是 \UB{}。
\item 若判成 \UB{}，去 附录~\ref{ap:ubC-cases} 找最接近的案例，写一份最小复现，
别急着改原始代码。
\item 按 附录~\ref{ap:ubB-switches} 选一个能报出这个形态的开关族，把它跑进自动测试，
而不是手工跑一次。
\item 按本表最后一列改写代码，并确保改写后的版本在同一开关下静默。
\item 若判成 \IDB{} 或 \USB{}，把结论写进移植文档或代码注释，
注明在哪两个实现上验证过；这一步与改代码同等重要，否则下一个维护者会重新踩进来。
\end{enumerate}

\begin{bestpractice}[title=把判定写进评审清单]
评审一段可疑代码时按三步走：先指出它触及哪个语言构造，再从本表读出类别，
最后写下用哪个开关验证。三步都答不上来的位置，就是既没有契约也没有测试的位置，
应当优先补断言与 \Tp{static\_assert}，把结论固定到编译期。
\end{bestpractice}

\chapter{消毒器与诊断开关对照}\label{ap:ubB-switches}

附录~\ref{ap:ubB-switches} 把散在各章的“怎么抓”收进三张表：第一张按实现列开关族，
第二张按诊断信息反查行为类别与责任章节，第三张给出分档的持续集成建议。
三张表都不是开关手册，而是给评审与门禁用的索引：它们回答的是“该开什么”，
而不是“开了之后一定能抓到什么”。

\begin{guideline}[title=记法与适用范围]
表中开关按本册统一记法给出：\Opt{} 与 \Warn{} 这类宏一律排成双横线形式，
实际命令行里 \Tp{-fsanitize=...}、\Tp{-D{\Bk}\_GLIBCXX{\Bk}\_DEBUG} 这些是单横线，
\Mx{MSVC} 前端用斜杠形式；写构建脚本时以工具文档为准，别把本表的排版当拼写依据。
凡本册未确证到具体拼写的能力，表中不编开关名，只写“以对应工具文档为准”，
这类行不能作为配置的直接依据。
\end{guideline}

\section{实现与开关族}

\begin{guideline}[title=这张表怎么用]
先按你的构建矩阵选实现，再按最关心的行为类别选开关族，两列同时命中的行才值得进门禁。
覆盖类别一列写的是该族\emph{设计上}关注的类别，不是保证：动态检测只报告被执行到的那次错误；
静态分析与编译器告警只覆盖能静态看出的子集。两者叠加也推不出“全绿即无未定义行为”，
原因见 第~\ref{ch:ub02-inevitable} 章。表中“族”指同一开关下的一组子档：
把逐条子档名按版本列死更容易写错，宁可在构建脚本里固化再回来核对。
\end{guideline}

\begin{longtable}{@{}>{\raggedright\arraybackslash}p{0.12\textwidth} >{\raggedright\arraybackslash}p{0.32\textwidth} >{\raggedright\arraybackslash}p{0.19\textwidth} >{\raggedright\arraybackslash}p{0.17\textwidth}@{}}
\caption{实现、开关族与覆盖的行为类别}\label{tab:apB-switches}\\
\toprule
实现 & 开关族 & 覆盖的行为类别 & 代价与限制 \\
\midrule
\endfirsthead
\toprule
实现 & 开关族 & 覆盖的行为类别 & 代价与限制 \\
\midrule
\endhead
\bottomrule
\endfoot
\Mx{GCC} & \Opt{f{\Bk}sanitize{\Bk}=undefined} & \UB{} 的多数算术与解引用形态：有符号溢出、零除、移位、越界下标、对齐、\Tp{bool} 与枚举取值、空指针、越过函数尾的返回 & 开销与报告形态随版本；默认只报告不中止，需再开 \Opt{fno{\Bk}-sanitize{\Bk}-recover} \\
\Mx{GCC} & \Opt{f{\Bk}sanitize{\Bk}=signed-{\Bk}integer-{\Bk}overflow} & 有符号溢出，见 第~\ref{ch:ub09-integers} 章 & 只报告被执行到的那次；优化阶段已消失的分支无从上报 \\
\Mx{GCC} & \Opt{f{\Bk}sanitize{\Bk}=bounds-{\Bk}strict} & 数组下标与迭代器越界，见 第~\ref{ch:ub05-pointer} 章 & 比 \Opt{f{\Bk}sanitize{\Bk}=bounds} 覆盖更宽，代价是略高开销 \\
\Mx{GCC} & \Opt{f{\Bk}sanitize{\Bk}=alignment} 与 \Opt{f{\Bk}sanitize{\Bk}=object-{\Bk}size} & 对齐不足、访问超出对象尺寸，见 第~\ref{ch:ub07-alignment} 章 & 指针来源必须可信；由整数还原的指针查不出越界 \\
\Mx{GCC} & \Opt{ftrapv} & 有符号溢出的陷入式检测 & 性能很差且不带位置信息，只适合一次性复现，不进流水线 \\
\Mx{GCC} & \Opt{f{\Bk}sanitize{\Bk}=address} 配 \Opt{f{\Bk}sanitize-{\Bk}address-{\Bk}use-after-{\Bk}scope} & 堆栈与全局越界、释放后使用、离开作用域后使用、泄漏（部分平台），见 第~\ref{ch:ub04-lifetime} 章 & 与线程消毒器互斥；影子内存与分配元数据显著放大占用 \\
\Mx{GCC} & \Opt{f{\Bk}sanitize{\Bk}=thread} & data race 与锁次序问题，见 第~\ref{ch:ub10-memory-model} 章 & 数倍到十余倍开销，内存放大明显，必须全量重编译 \\
\Mx{GCC} & \Opt{f{\Bk}sanitize{\Bk}=leak} & 泄漏 & 只在部分平台可用；分配器不匹配那一类它报不出来 \\
\Mx{GCC} & \Opt{fanalyzer} & 路径敏感的释放后使用、双重释放、空指针解引用、未初始化读取、越界 & 编译时间大幅上升，误报需先基线化；分析深度有限会漏报 \\
\Mx{GCC} & \Tp{-Wuninitialized} 与 \Tp{-Wmaybe-{\Bk}uninitialized} & 未初始化读取的静态可见子集，见 第~\ref{ch:ub04-lifetime} 章 & 只在开启优化后才跑完整数据流，噪声随档位变化 \\
\Mx{GCC} & \Tp{-Warray-{\Bk}bounds=2}、\Tp{-Wdangling-{\Bk}pointer}、\Tp{-Wuse-{\Bk}after-{\Bk}free} & 越界、悬垂与释放后使用的静态形态 & 覆盖的是编译器能证明的部分，证明不了就不报 \\
\Mx{GCC} & \Tp{-Wfree-{\Bk}nonheap-{\Bk}object}、\Tp{-Wmismatched-{\Bk}new-{\Bk}delete} & 释放配对错误，见 第~\ref{ch:ub12-library-pre} 章 & 需要分配点与释放点在同一翻译单元内可见 \\
\Mx{GCC} & \Tp{-Wstrict-{\Bk}aliasing} & 严格别名违规的静态线索，见 第~\ref{ch:ub06-aliasing} 章 & 依赖优化器已经做出的假设，档位不同结果不同 \\
\Mx{GCC} & \Tp{-Wsequence-{\Bk}point} 一类定序告警 & 求值与副作用定序问题，见 第~\ref{ch:ub08-evaluation} 章 & 只覆盖能静态看出的子集，且对 \CPP{} 的覆盖比 \CL{} 更窄；报告形态随版本变动，逐项以工具文档为准 \\
\Mx{GCC} & \Opt{Wall}、\Opt{Wextra}、\Opt{Wpedantic}、\Opt{Werror} & 编译期可看出的 \ILLF{} 与部分 \NDR{} & \Opt{Werror} 对版本新增告警很脆，门禁要按编译器版本分档 \\
\Mx{GCC} 与标准库选型 & \Opt{D{\Bk}\_GLIBCXX{\Bk}\_DEBUG}、\Opt{D{\Bk}\_GLIBCXX{\Bk}\_ASSERTIONS} & 容器与算法的库契约、迭代器有效区间，见 第~\ref{ch:ub12-library-pre} 章 & 改变类型与二进制接口，不能与未一起重编译的库混链 \\
\Mx{GCC} & \Opt{fno-{\Bk}strict-{\Bk}aliasing} 与 \Opt{fwrapv} & 不是检测，是关掉承诺：别名假设与溢出回绕假设 & 掩盖问题且可能改变结果，只适合移植过渡期，要带退场日期 \\
\Mx{Clang} & \Opt{f{\Bk}sanitize{\Bk}=undefined} 的各子档 & 与 \Mx{GCC} 同族，档位名基本一致，报告文案不同 & 子集与默认开启项随版本变动，逐项以工具文档为准 \\
\Mx{Clang} & \Opt{f{\Bk}sanitize{\Bk}=memory} & 未初始化读取，见 第~\ref{ch:ub04-lifetime} 章 & \Mx{GCC} 没有这一族；所有依赖必须一起重编译，否则漏报与误报都多 \\
\Mx{Clang} & \Tp{-Wthread-{\Bk}safety} & 共享状态的加锁约定违约，见 第~\ref{ch:ub10-memory-model} 章 & 依赖源码里的线程安全标注，属性写法以对应编译器文档为准；没标注等于没检查 \\
\Mx{Clang} & \Mx{clang-tidy} 的 bugprone 检查族 & 移动后使用、悬垂句柄、有符号与无符号混用、可疑 \Tp{string{\Bk}\_view} 取数据、隐式窄化、异常逃逸 & 只读源码不看执行；检查名与可用性随版本增减，逐个以工具文档为准 \\
\Mx{Clang} & \Mx{clang-tidy} 的 clang-analyzer 检查族 & 路径敏感的释放后使用、泄漏、未初始化操作数、空指针解引用 & 编译命令要能被复原，否则静默什么也不报；适合放夜间全量档 \\
\Mx{Clang} & \Opt{f{\Bk}sanitize{\Bk}=fuzzer} & 不是类别检测，是路径生成；与语义族、地址族组合才产出类别报告 & 需要可重复的语料库与崩溃样本入库流程 \\
\Mx{Clang} & \Opt{f{\Bk}sanitize{\Bk}=local-{\Bk}bounds} & 由指针推断的局部数组越界，见 第~\ref{ch:ub05-pointer} 章 & \Mx{GCC} 无同名档；对动态长度数组覆盖有限，误报需甄别 \\
\Mx{Clang} & \Opt{f{\Bk}sanitize{\Bk}=pointer-{\Bk}overflow} & 跨对象指针算术与指针溢出，见 第~\ref{ch:ub05-pointer} 章 & 只覆盖能被插桩到的算术点，指针还原后失效 \\
\Mx{Clang} & \Opt{f{\Bk}sanitize{\Bk}=unsigned-{\Bk}integer-{\Bk}overflow} & 无符号回绕的可疑用法；回绕本身是良定义，不是 \UB{} & 属于代码风格档，用来暴露下面那类差值循环 \\
\Mx{Clang} 与 \Mx{GCC} & \Opt{f{\Bk}sanitize{\Bk}=vptr} 与 \Opt{f{\Bk}sanitize{\Bk}=return} 档 & 虚表类型不符与越过函数尾返回，见 第~\ref{ch:ub04-lifetime} 章 & 需要带类型信息的调试构建才给得出可读的定位 \\
\Mx{MSVC} & \Opt{analyze} & 代码分析告警族：初始化、越界、空指针、生命周期与所有权类问题 & 部分规则与语言标准选项互相制约；告警编号与新增项以文档为准 \\
\Mx{MSVC} & \Opt{f{\Bk}sanitize{\Bk}=address} & 与 \Mx{GCC} 同族的越界与释放后使用 & 支持的目标平台与运行时库组合有限；未初始化检查是否可用以对应版本文档为准 \\
\Mx{MSVC} & \Opt{sdl} 与迭代器调试级别 & 库契约与迭代器有效区间，见 第~\ref{ch:ub12-library-pre} 章 & 只覆盖标准库的断言点，不影响自定义容器与裸数组 \\
\Mx{MSVC} & \Opt{RTC1} 一类运行期检查 & 栈损坏、未初始化局部变量、构造期纯虚调用 & 仅调试构建有效，发布档没有等价物；子开关拼写以 \Mx{MSVC} 文档为准 \\
\Mx{MSVC} & \Opt{analyze} 与 \Opt{sdl} 叠用 & 静态线与库断言线互补，见 第~\ref{ch:ub12-library-pre} 章 & 两者都不是消毒器，执行路径上的错误仍要地址族 \\
跨实现 & \Opt{fstack-{\Bk}protector-{\Bk}strong} 与 \Opt{guard:cf} 一类防护开关 & 属于缓解而非检测：栈溢出与控制流劫持的代价 & 不报类别，也不改变 \UB{} 的判定，别把它当成检测结果计入门禁 \\
跨实现 & \Opt{D{\Bk}\_FORTIFY{\Bk}\_SOURCE=2} 与 \Opt{D{\Bk}\_FORTIFY{\Bk}\_SOURCE=3} & 已知尺寸对象上的越界读写，见 第~\ref{ch:ub05-pointer} 章 & 只在优化构建生效，且只有长度可静态推导的调用才被加固 \\
跨实现 & \Opt{fno-{\Bk}omit-{\Bk}frame-{\Bk}pointer} 与报告类环境变量选项 & 不新增检测，只决定报告能不能定位到调用栈 & 门禁要开；关掉它会显著削弱 \Mx{ASan} 与 \Mx{UBSan} 的可查性 \\
\end{longtable}

\begin{note}[title=三族开关的分工]
语义族管“这一步标准有没有承诺”，内存族管“这块存储能不能碰”，线程族管“两个线程之间有没有约定”。
三族互相替代不了：严格别名违规只有部分档位能报，\Tp{volatile} 误用只有线程消毒器会抱怨，
\Tp{char} 符号性三族都不管，只能靠跨平台测试矩阵与移植文档。
\end{note}

\begin{compare}[title=编译器告警与消毒器的差别]
编译器告警给的是“静态上可疑”，同一行代码在不同档位、不同版本下会跳出来或闭嘴；
消毒器给的是“这一条路径真的走错过一次”。前者的漏报靠换档位与换实现补，
后者的漏报只能靠测试覆盖补。两类噪声的处理策略完全不同，
混在同一道门禁里必然被整体关掉，这就是第三张表把它们分档的原因。
\end{compare}

把表里的行真正落到一次构建上时，还有四条组合规则要遵守。这四条不是风格偏好，
而是由检测原理决定的：前两条来自插桩方式互不相容，后两条来自报告质量的依赖项。

\begin{enumerate}
\item 地址族与线程族不能同开。两者都要接管全部内存访问，插桩方式互斥，
强行同开得到的不是双份报告，而是构建失败或不可信的报告。
\item 内存族要求把所有依赖一起重编译。只重编译自己的目标，
未插桩的库读写的内存会被当成合法值，漏报与误报同时上升。
\item 语义族默认是“报告但继续执行”。门禁必须显式改成失败，
否则进程会以零状态退出，流水线看起来一片绿。
\item 任何一族都要配 \Opt{fno-{\Bk}omit-{\Bk}frame-{\Bk}pointer} 与调试信息，
并打开报告栈的环境变量选项；没有栈的报告只能定位到文件，定位不到责任调用。
\end{enumerate}

还有一条与类别判定相关的规则：静态检查族的报告不能直接当作 \UB{} 的证据，
它给的是可疑形态；反过来，消毒器的一条报告就是真实执行过一次的证据，
可以按 附录~\ref{ap:ubA-kinds} 的类别列直接查表决定动作。

\section{诊断信息到类别与章节的反查}

\begin{guideline}[title=反查表怎么用]
下面每行的“诊断文本”是本册按常见报告形状写的\emph{示意措辞}，用来认方向，
不是任何版本的逐字输出；逐字文案、文件行号格式与错误码编号都以工具文档为准。
第一反应一列只给动作方向，最小复现与改写的完整写法见 附录~\ref{ap:ubC-cases}。
\end{guideline}

\begin{longtable}{@{}>{\raggedright\arraybackslash}p{0.30\textwidth} >{\raggedright\arraybackslash}p{0.10\textwidth} >{\raggedright\arraybackslash}p{0.13\textwidth} >{\raggedright\arraybackslash}p{0.27\textwidth}@{}}
\caption{诊断信息示意与责任章节}\label{tab:apB-diag}\\
\toprule
诊断文本（示意措辞） & 类别 & 责任章节 & 第一反应 \\
\midrule
\endfirsthead
\toprule
诊断文本（示意措辞） & 类别 & 责任章节 & 第一反应 \\
\midrule
\endhead
\bottomrule
\endfoot
语义类报告：有符号整数溢出，结果无法用该类型表示 & \UB{} & 第~\ref{ch:ub09-integers} 章 & 先分清这条溢出是算错还是被拿来判分支，再决定换无符号域还是夹紧 \\
语义类报告：除数为零 & \UB{} & 第~\ref{ch:ub09-integers} 章 & 判零与除法拆两条语句；零若是合法输入就改成显式错误返回 \\
语义类报告：移位量不小于类型宽度 & \UB{} & 第~\ref{ch:ub09-integers} 章 & 夹紧移位量；重点检查是否把字节数当成了位数 \\
语义类报告：移位量为负 & \UB{} & 第~\ref{ch:ub09-integers} 章 & 移位量改无符号，并在入口处校验来源 \\
语义类报告：载入值超出枚举可表示范围 & \UB{} & 第~\ref{ch:ub09-integers} 章 & 读入处校验；能改带作用域枚举就改 \\
语义类报告：\Tp{bool} 的左值既非真也非假 & \UB{} & 第~\ref{ch:ub01-taxonomy} 章 & 追数据来源，多半是缓冲区或联合体读出的原始位模式 \\
语义类报告：存储未按要求对齐 & \UB{} & 第~\ref{ch:ub07-alignment} 章 & 改成先拷贝到对齐的局部量，或调整布局与转换方式 \\
语义类报告：解引用空指针 & \UB{} & 第~\ref{ch:ub05-pointer} 章 & 判空的位置是否太晚；能否换成引用或非空视图 \\
语义类报告：使用了未初始化的成员指针一类值 & \UB{} & 第~\ref{ch:ub04-lifetime} 章 & 检查所有构造路径，尤其是条件初始化与被跳过的分支 \\
语义类报告：控制流到达函数尾且返回类型非 \Tp{void} & \UB{} & 第~\ref{ch:ub08-evaluation} 章 & 补 \Tp{return}；确不可达的分支改成显式终止 \\
语义类报告：虚表指针类型与预期不符 & \UB{} & 第~\ref{ch:ub04-lifetime} 章 & 检查向下转型的正确性与构造析构期的调用 \\
内存类报告：释放后使用堆内存 & \UB{} & 第~\ref{ch:ub04-lifetime} 章 & 先看释放点与使用点的先后，再谈缓存与引用计数 \\
内存类报告：栈缓冲区越界 & \UB{} & 第~\ref{ch:ub05-pointer} 章 & 复核长度参数的单位（元素数还是字节数）与复制来源长度 \\
内存类报告：分配类型不匹配 & \UB{} & 第~\ref{ch:ub12-library-pre} 章 & 回到分配点检查 \Fn{new} 与 \Fn{new[]} 是否配对 \\
内存类报告：使用了已返回函数的栈地址 & \UB{} & 第~\ref{ch:ub04-lifetime} 章 & 返回拥有值的对象，或把缓冲交给调用方持有 \\
线程类报告：同一内存位置上的 data race & \UB{} & 第~\ref{ch:ub10-memory-model} 章 & 先决定这块内存归锁还是归原子，别先去降级内存顺序 \\
内存类报告：读取了未初始化的值 & \UB{} & 第~\ref{ch:ub04-lifetime} 章 & 确认真的漏了初始化，而不是依赖了结构体的填充字节 \\
编译器告警：局部变量可能未初始化即被使用 & \UB{} & 第~\ref{ch:ub04-lifetime} 章 & 声明即初始化；不要用抑制告警代替修改 \\
编译器告警：可能返回局部量的地址 & \UB{} & 第~\ref{ch:ub04-lifetime} 章 & 改按值返回，或改成调用方提供的输出参数 \\
编译器告警：有符号与无符号比较，条件恒真或恒假 & \IDB{} 导致 \UB{} & 第~\ref{ch:ub09-integers} 章 & 统一索引类型；检查长度为零时下溢的那条路径 \\
编译器告警：违反严格别名规则 & \UB{} & 第~\ref{ch:ub06-aliasing} 章 & 换拷贝函数或位模式转换设施，不要长期挂着关闭别名优化的开关 \\
编译器告警：取了打包结构成员的地址 & \UB{} & 第~\ref{ch:ub07-alignment} 章 & 取值到局部再取地址；协议结构改成字节数组加显式解码 \\
静态检查：值被移动后仍被使用 & \USB{} & 第~\ref{ch:ub12-library-pre} 章 & 检查移动后是否立刻赋新值；确需读取就写明契约与测试 \\
静态检查：\Tp{string{\Bk}\_view} 的数据在宿主销毁后被使用 & \UB{} & 第~\ref{ch:ub04-lifetime} 章 & 让视图作用域不超过宿主，或返回拥有内容的类型 \\
静态检查：异常可能逃逸出禁止抛出的上下文 & 良定义 & 第~\ref{ch:ub11-terminate} 章 & 在边界加捕获全部异常的处理，或重新评估这个 \Tp{noexcept} 该不该写 \\
代码分析：使用了未初始化的内存 & \UB{} & 第~\ref{ch:ub04-lifetime} 章 & 与另一家编译器的同类告警一起处理，不要只压某一家 \\
代码分析：读取的数据大小无效 & \UB{} & 第~\ref{ch:ub05-pointer} 章 & 复核缓冲区尺寸与长度参数的一致性 \\
运行即终止且没有诊断报告（异常越过 \Tp{noexcept}） & 良定义 & 第~\ref{ch:ub11-terminate} 章 & 先装 \Fn{set{\Bk}\_ter{\Bk}minate} 取回栈，再定位抛出点 \\
得到错误结论但没有任何工具抱怨 & 可能是 \USB{}、\IDB{} 或 \NDR{} & 第~\ref{ch:ub01-taxonomy} 章 & 回到附录~\ref{ap:ubA-kinds} 先判类别；工具查不到不等于标准还给了承诺 \\
\end{longtable}

\begin{warning}[title=报告形态会消失]
同一处错误在 \Tp{-O0} 下报得清清楚楚，在 \Tp{-O2} 下一条也不报：优化器可能已经把
依赖溢出的分支整体删掉，或把越界访问折叠成常量地址。没报不等于没有，
只是那次运行没有走到标准撤回承诺的位置。这一现象的机制见 第~\ref{ch:ub02-inevitable} 章。
\end{warning}

一条报告落地时要同时记下三件事，缺一件都会在下一轮排查里重新花时间：
完整的命令行与编译器版本（换档位往往就复现不出来）、能稳定重放的输入或测试名、
以及按本表判出的类别。类别记下来还有一个实际用处：
同一处代码上的两条报告，如果一条判成 \UB{} 一条判成 \IDB{}，
它们的修复路径完全不同，前者改代码，后者补移植文档，不该混进同一张问题单。

\Mx{MSVC} 与 \Mx{Clang} 的分析类报告常按编号给出，编号在不同版本里会新增、会改含义，
本表因此只按能力描述反查而不列编号；要把某条编号固定成门禁规则，
请写进构建脚本与规则文件，并在注释里注明它对应本表哪一行。

\section{持续集成分档建议}

\begin{guideline}[title=分档表怎么用]
分档的前提是：一档只做一件事，失败策略要能直接落到流水线配置里。
全绿不等于无 \UB{}，理由有两条且都写在正文：动态检测只覆盖被执行路径（第~\ref{ch:ub02-inevitable} 章），
而标准允许实现换答案、实现也确实会换（第~\ref{ch:ub15-forward} 章）。
所以本表的“通过”只表示：这一档的检查项在当前测试集上没有报告，不是一句安全证明。
\end{guideline}

\begin{longtable}{@{}>{\raggedright\arraybackslash}p{0.12\textwidth} >{\raggedright\arraybackslash}p{0.36\textwidth} >{\raggedright\arraybackslash}p{0.18\textwidth} >{\raggedright\arraybackslash}p{0.14\textwidth}@{}}
\caption{持续集成的分档建议}\label{tab:apB-ci}\\
\toprule
档位 & 构建类型与打开的检查 & 失败策略 & 适用规模 \\
\midrule
\endfirsthead
\toprule
档位 & 构建类型与打开的检查 & 失败策略 & 适用规模 \\
\midrule
\endhead
\bottomrule
\endfoot
本机快速档 & \Tp{-O0} 或 \Tp{-O1} 加调试信息；\Opt{Wall}、\Opt{Wextra}、\Opt{Wpedantic}；语义族开关并配 \Opt{fno-{\Bk}sanitize-{\Bk}recover} & 任一报告即中止，不给重试；耗时控制在秒级到分钟级 & 任意规模，作为提交前的最低门槛 \\
门禁档 & \Tp{-O1} 或 \Tp{-O2} 加调试信息；上一档全部，再加地址族与其作用域后使用子档；标准库断言模式 & 新报告一律拒收；已有基线只在数量上升时拒收 & 中小型项目每个合并请求；大型项目按受影响目标跑 \\
并发档 & 构建与门禁档同源，但关掉地址族改开线程族；测试集只放多线程用例 & 竞争报告即拒收；与地址族分任务、分进程跑 & 有共享可变状态的所有目标 \\
夜间全量档 & 语义族全子档加地址族；有条件时另配内存族；再加静态分析族与两族静态检查 & 汇总成趋势报告并按周设阈值；单条新报告开问题而不阻塞合并 & 全仓库；数小时量级可接受 \\
模糊测试档 & 面向解析与边界入口的目标，路径生成开关叠语义族与地址族；语料与崩溃样本入库 & 任一最小复现用例转成回归测试；崩溃样本必须可重放 & 所有暴露给不可信输入的接口 \\
发布候选档 & 不带消毒器的发布构建（\Tp{-O2} 以上）；保留零依赖断言与加固检查；\Mx{MSVC} 侧开安全特性开关 & 断言触发即视为回归；本档通过不解除前面各档的门禁 & 全部对外发布物 \\
移植审计档 & 换实现的交叉构建：至少两个编译器族，外加一个 \Tp{char} 符号性相反的平台 & 差异一律记入移植文档；目标不是全绿，而是把差异写明 & 依赖 \IDB{} 与 \USB{} 行为的模块 \\
\end{longtable}

\begin{lstlisting}[style=shellstyle, caption={最小门禁脚本：互斥的两族分进程、分产物跑}]
#!/usr/bin/env bash
set -euo pipefail

SRC=main.cpp
COMMON='-fno-omit-frame-pointer -g -O1 -Wall -Wextra'
UB='-fsanitize=undefined -fno-sanitize-recover=all'
AS='-fsanitize=address -fsanitize-address-use-after-scope'
TS='-fsanitize=thread'

g++ -std=c++23 $COMMON $UB -o app_ubsan $SRC
UBSAN_OPTIONS=print_stacktrace=1 ./app_ubsan

g++ -std=c++23 $COMMON $AS -o app_asan $SRC
ASAN_OPTIONS=detect_stack_use_after_return=1 ./app_asan

g++ -std=c++23 $COMMON $TS -o app_tsan $SRC
TSAN_OPTIONS=halt_on_error=1 ./app_tsan
\end{lstlisting}

三个环境变量名是各家消毒器的运行时选项入口：
\Mx{UBSan} 的 \Tp{UBSAN{\Bk}\_OPTIONS}、\Mx{ASan} 的 \Tp{ASAN{\Bk}\_OPTIONS}、
\Mx{TSan} 的 \Tp{TSAN{\Bk}\_OPTIONS}。它们决定报告是否带栈、是否中止、是否统计泄漏，
与编译期开关是两套东西：编译期开关决定\emph{能不能}检测到，运行时选项决定\emph{报出来之后怎么办}。
门禁档要让它中止，本机复现档反而要让它继续跑完，以便一次看全所有报告点。

\begin{lstlisting}[style=cmakestyle, caption={把消毒器选项按编译器族落到目标上}]
target_compile_options(app PRIVATE
  $<$<CXX_COMPILER_ID:GNU>:-Wall -Wextra -Wpedantic>
  $<$<CXX_COMPILER_ID:Clang>:-Wall -Wextra -Wpedantic>
  -fno-omit-frame-pointer -g)
if(UBSAN_MODE)
  target_compile_options(app PRIVATE -fsanitize=undefined)
  target_link_options(app PRIVATE -fsanitize=undefined)
endif()
if(TSAN_MODE)
  target_compile_options(app PRIVATE -fsanitize=thread)
  target_link_options(app PRIVATE -fsanitize=thread)
endif()
\end{lstlisting}

\begin{bestpractice}[title=基线与噪声]
第一次给老项目开静态分析族或代码分析开关时，报告数量足以淹没评审。
做法是把当轮结果存成基线，门禁只拒收新增项，同时给基线设一条只减不增的曲线；
没有基线的静态分析会在两周内被整体关掉，反而比不开更糟。
消毒器不需要基线：它的每一条都对应一次真实执行，出现就改。
\end{bestpractice}

\begin{note}[title=为什么表里没有版本号与告警编号]
开关集合与告警编号在每个实现里都在动，写死版本或逐条列编号反而更容易误导。
本表的取舍是：只列本册确证存在的开关族与能力描述；具体拼写、可用组合与编号落在你
工具链上的形态，一律以实现文档与构建脚本为准。需要跨版本固定的，
请把确证后的开关名固化进构建脚本与文档化配置，而不是固化进一张附录表。
还有一条判据值得记住：同一份代码至少要在家用编译器与另一族编译器上各跑一遍门禁，
否则你验证的其实是某个实现，不是某段代码。
\end{note}

\chapter{常见未定义行为案例与最小复现清单}\label{ap:ubC-cases}

附录~\ref{ap:ubC-cases} 收十四个案例，每个都按同一个顺序写：最小代码、触发条件、类别判定、检测手段、
可移植改写。代码压到十行以内，为的是能直接贴进回归测试；判定与手段两栏回指
附录~\ref{ap:ubA-kinds} 与附录~\ref{ap:ubB-switches}，章末给一张案例与可检测性的汇总表。
代码里的注释用英文，与正文的排版规则分开。

\begin{guideline}[title=最小复现的写法要求]
一份能进问题单的最小复现要满足四条：单文件能编译、触发输入写死在代码里、
不依赖环境变量或偶然配置、同一份代码在关闭优化与高档优化下各跑一遍。
最后一条尤其重要：两边结果不一致，本身就是这里有 \UB{} 的证据，机制见 第~\ref{ch:ub02-inevitable} 章。
\end{guideline}

\section{有符号溢出被当作饱和判断}
\begin{lstlisting}[style=cppstyle, caption={反例：先算再判，判断本身已经站在溢出之后}]
int add_clamped(int a, int b) {
    int s = a + b;                  // UB already happened here
    if (s < a) return 2147483647;   // relies on wraparound
    return s;
}
\end{lstlisting}
\begin{itemize}
\item 触发条件：两个正数之和超过上限；高档优化最容易暴露。
\item 类别判定：\UB{}，见 第~\ref{ch:ub09-integers} 章。
\item 检测手段：语义族的溢出档，再用 \Tp{-O0} 与 \Tp{-O2} 各跑一次（附录~\ref{ap:ubB-switches}）。
\item 可移植改写：判断挪到运算之前，判断本身不溢出。
\end{itemize}
\begin{lstlisting}[style=cppstyle, caption={改写：用减法做上限判断}]
#include <limits>
int add_clamped(int a, int b) {
    using Lim = std::numeric_limits<int>;
    if (b > 0 && a > Lim::max() - b) return Lim::max();
    if (b < 0 && a < Lim::min() - b) return Lim::min();
    return a + b;                   // cannot overflow past this point
}
\end{lstlisting}

\section{同一表达式内重复副作用}
\begin{lstlisting}[style=cppstyle, caption={反例：一个表达式里对同一标量既读又写}]
int bump(int i) {
    return i = i++ + i;             // reads and writes of i are unsequenced
}
\end{lstlisting}
\begin{itemize}
\item 触发条件：任何输入都走这条路；不同实现各自给出一个稳定值，稳定的是未被承诺的结果。
\item 类别判定：\UB{}，依据 \Std{intro{\Bk}.ex{\Bk}ecution}，见 第~\ref{ch:ub08-evaluation} 章。
\item 检测手段：没有运行时报告可指望，只有告警与静态检查的部分覆盖（附录~\ref{ap:ubB-switches}）。
\item 可移植改写：读与写拆成两条语句，见下方。
\end{itemize}
\begin{lstlisting}[style=cppstyle, caption={改写：一语句一意图}]
int bump(int i) {
    int const old = i;
    i = old + 1;
    return old + i;                 // spells out the intended value
}
\end{lstlisting}

\section{发布版里消失的越界断言}
\begin{lstlisting}[style=cppstyle, caption={反例：边界检查写在断言里}]
#include <cassert>
#include <vector>
int at(std::vector<int> const& v, std::size_t i) {
    assert(i < v.size());
    return v[i];                    // with NDEBUG the check line is gone
}
\end{lstlisting}
\begin{itemize}
\item 触发条件：定义 \Tp{NDEBUG} 之后的发布构建；调试版一切正常容易误判成已修好。
\item 类别判定：\UB{}，见 第~\ref{ch:ub05-pointer} 章 与 第~\ref{ch:ub12-library-pre} 章。
\item 检测手段：容器调试与断言模式、边界档、地址族；必须在发布档位下跑（附录~\ref{ap:ubB-switches}）。
\item 可移植改写：边界检查写成不受 \Tp{NDEBUG} 影响的真实分支。
\end{itemize}
\begin{lstlisting}[style=cppstyle, caption={改写：越界是错误路径，不是断言}]
#include <stdexcept>
int at(std::vector<int> const& v, std::size_t i) {
    if (i >= v.size()) throw std::out_of_range("bad index");
    return v[i];                    // 或直接调用 v.at(i)
}
\end{lstlisting}

\section{\Tp{string\_view} 返回局部量的视图}
\begin{lstlisting}[style=cppstyle, caption={反例：视图短于它指向的缓冲}]
std::string_view banner() {
    std::string s = make_title();
    return s;                       // s dies, the view dangles at once
}
\end{lstlisting}
\begin{itemize}
\item 触发条件：调用方在下一次分配或调用之后才用到视图；可能读到垃圾，也可能看着完全正确。
\item 类别判定：\UB{}，依据 \Std{basic{\Bk}.life}，见 第~\ref{ch:ub04-lifetime} 章。
\item 检测手段：地址族配合返回后使用子档；静态侧有视图悬垂类检查（附录~\ref{ap:ubB-switches}）。
\item 可移植改写：返回拥有内容的类型，一行即可：\Tp{std::{\Bk}string} \Fn{banner} 按值返回。
\end{itemize}

\section{容器扩容之后仍用旧引用}
\begin{lstlisting}[style=cppstyle, caption={反例：引用跨过可能发生的重新分配}]
void append_all(std::vector<std::string>& v, std::string const& extra) {
    std::string& head = v[0];
    v.push_back(extra);             // may reallocate the buffer
    head += extra;                  // use after free
}
\end{lstlisting}
\begin{itemize}
\item 触发条件：插入真的触发扩容时；容量够用时一切正常，小测试集上永远不坏。
\item 类别判定：\UB{}，见 第~\ref{ch:ub04-lifetime} 章 与 第~\ref{ch:ub12-library-pre} 章。
\item 检测手段：地址族的释放后使用报告最直观；容器调试模式也会抱怨迭代器失效。
\item 可移植改写：先改变形状，再用下标定位，不缓存跨扩容的引用。
\end{itemize}
\begin{lstlisting}[style=cppstyle, caption={改写：扩容之后再取元素}]
void append_all(std::vector<std::string>& v, std::string const& extra) {
    v.push_back(extra);
    v[0] += extra;                  // no reference cached across growth
}
\end{lstlisting}

\section{联合体双关与非活跃成员读取}
\begin{lstlisting}[style=cppstyle, caption={反例：写入一个成员，读取另一个成员}]
union Bits { std::uint32_t u; float f; };
std::uint32_t to_bits(float x) {
    Bits b;
    b.f = x;                        // f is the active member
    return b.u;                     // reads the inactive one
}
\end{lstlisting}
\begin{itemize}
\item 触发条件：任何输入都走这条路；结果看着正确是因为多数平台按同一布局解释存储。
\item 类别判定：\CPP{} 里判 \UB{}；\CL{} 传统上按实现定义处理表示与填充，本册不给统一结论，
以实现文档为准（第~\ref{ch:ub01-taxonomy} 章、第~\ref{ch:ub13-implementation} 章）。
\item 检测手段：这一形态一般没有消毒器报告，只能靠改写与评审（附录~\ref{ap:ubB-switches}）。
\item 可移植改写：用标准给的位模式转换设施。
\end{itemize}
\begin{lstlisting}[style=cppstyle, caption={改写：显式且可移植的表示转换}]
#include <bit>
std::uint32_t to_bits(float x) { return std::bit_cast<std::uint32_t>(x); }
\end{lstlisting}

\section{结构体填充字节参与哈希}
\begin{lstlisting}[style=cppstyle, caption={反例：对含填充的对象逐字节哈希}]
struct Key { std::uint16_t a; std::uint32_t b; };
std::size_t raw_hash(Key const& k) {
    std::size_t h = 146959810393466ULL;
    auto const* p = reinterpret_cast<unsigned char const*>(&k);
    for (std::size_t i = 0; i < sizeof k; ++i) h = (h ^ p[i]) * 1099511628211ULL;
    return h;                       // padding holds indeterminate values
}
\end{lstlisting}
\begin{itemize}
\item 触发条件：同一键值在不同对象或不同实现下给出不同哈希；表现为等价对象被判不等。
\item 类别判定：读取不确定值属 \UB{}；填充的存在与位置属 \IDB{}。
\item 检测手段：内存族最直接，地址族与容器调试模式管不到这一层（附录~\ref{ap:ubB-switches}）。
\item 可移植改写：只对成员做哈希。
\end{itemize}
\begin{lstlisting}[style=cppstyle, caption={改写：逐字段混合，跳过对象表示}]
std::size_t hash_key(Key const& k) {
    std::size_t h = 146959810393466ULL;
    auto mix = [&h](std::uint64_t x) {
        for (int i = 0; i < 8; ++i) h = (h ^ ((x >> (8 * i)) & 0xFFu)) * 1099511628211ULL;
    };
    mix(k.a); mix(k.b); return h;
}
\end{lstlisting}

\section{取 packed 结构成员的地址}
\begin{lstlisting}[style=cstyle, caption={反例：把打包成员的地址当对齐指针返回}]
struct __attribute__((packed)) Header { uint16_t kind; uint32_t size; };
uint32_t* size_field(struct Header* h) {
    return &h->size;                // may be unaligned
}
\end{lstlisting}
\begin{itemize}
\item 触发条件：在不支持非对齐访问的体系结构上解引用返回的指针。
\item 类别判定：\UB{}，且打包属性是实现扩展而非标准设施，见 第~\ref{ch:ub07-alignment} 章。
\item 检测手段：编译期的打包成员告警，配合语义族的对齐档（附录~\ref{ap:ubB-switches}）。
\item 可移植改写：按字节读出再组装，拷贝本身就是允许的手段。
\end{itemize}
\begin{lstlisting}[style=cstyle, caption={改写：先拷贝到对齐的局部量}]
uint32_t size_of(struct Header const* h) {
    uint32_t tmp;
    memcpy(&tmp, &h->size, sizeof tmp);
    return tmp;
}
\end{lstlisting}

\section{\Tp{char} 符号性被用来做位图下标}
\begin{lstlisting}[style=cppstyle, caption={反例：依赖普通 char 是无符号}]
std::array<int, 256> table{};
void mark(char c) { table[c] = 1; } // negative index on signed-char platforms
\end{lstlisting}
\begin{itemize}
\item 触发条件：输入里出现最高位为一的字节，且该平台把普通 \Tp{char} 当有符号。
\item 类别判定：符号性是 \IDB{}；由它导致的越界写是 \UB{}（第~\ref{ch:ub13-implementation} 章）。
\item 检测手段：越界那一步可被地址族或边界档抓到，符号性前提抓不到，要靠构建矩阵。
\item 可移植改写：在 \Fn{mark} 里把下标显式转成无符号字节，一行即可，不改函数签名。
\end{itemize}

\section{无符号差值被当作循环上界}
\begin{lstlisting}[style=cppstyle, caption={反例：长度为零时上界回绕成极大值}]
std::size_t count_rises(std::vector<int> const& v) {
    std::size_t best = 0;
    for (std::size_t i = 0; i < v.size() - 1; ++i)
        if (v[i + 1] > v[i]) ++best;   // reads far past the end
    return best;
}
\end{lstlisting}
\begin{itemize}
\item 触发条件：空容器。回绕本身是无符号的良定义行为，坏在拿它当上界之后的越界读写。
\item 类别判定：越界那一步是 \UB{}，回绕本身不是（第~\ref{ch:ub09-integers} 章、第~\ref{ch:ub05-pointer} 章）。
\item 检测手段：边界档与地址族都能报；有符号与无符号比较的告警是最早的线索。
\item 可移植改写：上界用加法比较，空容器时一次都不进循环。
\end{itemize}
\begin{lstlisting}[style=cppstyle, caption={改写：用加法比较代替减法}]
    for (std::size_t i = 0; i + 1 < v.size(); ++i)
        if (v[i + 1] > v[i]) ++best;
\end{lstlisting}

\section{转发引用被二次转发}
\begin{lstlisting}[style=cppstyle, caption={反例：同一对象被转发两次}]
template <class T>
void record(T&& x) {
    store(std::forward<T>(x));      // may move out of x
    audit(std::forward<T>(x));      // sees a moved-from object
}
\end{lstlisting}
\begin{itemize}
\item 触发条件：实参是右值且第一次调用按值接收或移动；实参是左值时一切正常。
\item 类别判定：moved-from 之后读取属 \USB{}；若目标类型对被移空状态另有前提则升为 \UB{}。
\item 检测手段：静态检查里的移动后使用类检查是首选，消毒器只有偶发命中。
\item 可移植改写：消费点唯一，只保留一次转发。
\end{itemize}
\begin{lstlisting}[style=cppstyle, caption={改写：一次转发一处消费}]
template <class T>
void record(T&& x) { store(std::forward<T>(x)); }
\end{lstlisting}

\section{双重检查锁定}
\begin{lstlisting}[style=cppstyle, caption={反例：快路径上的读取没有同步}]
std::mutex m;
Config* g_conf = nullptr;           // plain pointer, not atomic
Config* conf() {
    if (!g_conf) {                  // unsynchronised read races with the write
        std::lock_guard<std::mutex> g(m);
        if (!g_conf) g_conf = new Config();
    }
    return g_conf;
}
\end{lstlisting}
\begin{itemize}
\item 触发条件：多线程且首次构造发生在别的线程；单线程测试永远通过。
\item 类别判定：\UB{}（data race，依据 \Std{res.on{\Bk}.data{\Bk}.races}，见 第~\ref{ch:ub10-memory-model} 章）。
\item 检测手段：线程族开关是唯一可靠手段；静态分析族都不覆盖并发定序（附录~\ref{ap:ubB-switches}）。
\item 可移植改写：让语言设施负责初始化的定序。
\end{itemize}
\begin{lstlisting}[style=cppstyle, caption={改写：函数内静态对象}]
Config* conf() {
    static Config instance;         // thread-safe initialisation
    return &instance;
}
\end{lstlisting}

\section{自写移动赋值导致双重释放}
\begin{lstlisting}[style=cppstyle, caption={反例：接走指针却没清空来源}]
class Buf {
    char* data = nullptr;
public:
    Buf& operator=(Buf&& other) {
        data = other.data;          // leaks old buffer, frees one block twice
        return *this;
    }
    ~Buf() { delete[] data; }
};
\end{lstlisting}
\begin{itemize}
\item 触发条件：源对象在移动之后仍然活着并正常析构；源对象立刻析构时泄漏那一半被隐藏。
\item 类别判定：\UB{}，见 第~\ref{ch:ub12-library-pre} 章。
\item 检测手段：地址族的双重释放报告最直接，泄漏那一半可由泄漏档补（附录~\ref{ap:ubB-switches}）。
\item 可移植改写：先释放自己的旧资源，接管之后把来源置空。
\end{itemize}
\begin{lstlisting}[style=cppstyle, caption={改写：接管与置空成对出现}]
    Buf& operator=(Buf&& other) noexcept {
        delete[] data;
        data = other.data;
        other.data = nullptr;
        return *this;
    }
\end{lstlisting}

\section{把 volatile 当原子用}
\begin{lstlisting}[style=cppstyle, caption={反例：用 volatile 变量做线程间旗标}]
volatile bool g_stop = false;       // only forces real accesses
void worker() { while (!g_stop) tick(); }
void shutdown() { g_stop = true; }
\end{lstlisting}
\begin{itemize}
\item 触发条件：两个线程读写同一旗标且至少一方是写；编译器可能把循环里的读提到循环外。
\item 类别判定：\UB{}，竞争不因 \Tp{vol{\Bk}atile} 而豁免（\Std{vol{\Bk}atile}，见 第~\ref{ch:ub10-memory-model} 章）。
\item 检测手段：线程族会照常报竞争；\Mx{MSVC} 侧只能靠代码分析与评审。
\item 可移植改写：旗标换成原子，并写清它需要的顺序。
\end{itemize}
\begin{lstlisting}[style=cppstyle, caption={改写：原子旗标加显式定序}]
std::atomic<bool> g_stop{false};
void worker() { while (!g_stop.load(std::memory_order_acquire)) tick(); }
void shutdown() { g_stop.store(true, std::memory_order_release); }
\end{lstlisting}

\section{案例汇总}

\begin{guideline}[title=汇总表怎么用]
“能自动检测”指存在一个开关族，在\emph{被执行到}的前提下给出报告；“不能”指该形态
要么只能静态看出，要么根本无从观察，必须靠改写与评审。类别定义见 附录~\ref{ap:ubA-kinds}，
完整开关族见 附录~\ref{ap:ubB-switches}。
\end{guideline}

\begin{longtable}{@{}>{\raggedright\arraybackslash}p{0.24\textwidth} >{\raggedright\arraybackslash}p{0.10\textwidth} >{\raggedright\arraybackslash}p{0.12\textwidth} >{\raggedright\arraybackslash}p{0.34\textwidth}@{}}
\caption{案例与类别、可检测性汇总}\label{tab:apC-summary}\\
\toprule
案例 & 类别 & 能否自动检测 & 首选手段与限制 \\
\midrule
\endfirsthead
\toprule
案例 & 类别 & 能否自动检测 & 首选手段与限制 \\
\midrule
\endhead
\bottomrule
\endfoot
有符号溢出被当作饱和判断 & \UB{} & 能 & 语义族的溢出档；与优化档位强相关，须两档都跑 \\
同一表达式内重复副作用 & \UB{} & 部分 & 告警与静态检查；没有运行时报告可指望 \\
发布版里消失的越界断言 & \UB{} & 能 & 容器调试与断言模式、边界档；须在发布档位下跑 \\
\Tp{string{\Bk}\_view} 返回局部量 & \UB{} & 能 & 地址族的返回后使用；栈未复用时可能静默 \\
容器扩容后仍用旧引用 & \UB{} & 能 & 地址族的释放后使用；最能体现覆盖依赖测试集 \\
联合体双关读非活跃成员 & \CPP{} 下 \UB{} & 不能 & 只能改写为位模式转换设施；消毒器一般不报 \\
填充字节参与哈希 & \UB{} 叠加 \IDB{} & 部分 & 内存族最有效，但要求全量重编译 \\
packed 成员取地址 & \UB{} & 部分 & 编译期告警为主；对齐档只在真的解引用时报告 \\
\Tp{char} 符号性做位图下标 & \IDB{} 导致 \UB{} & 部分 & 越界一步可抓；符号性前提只能靠跨平台矩阵 \\
无符号差值当循环上界 & \UB{} & 能 & 边界档与地址族；空容器用例是触发的唯一门槛 \\
转发引用二次转发 & \USB{} 或 \UB{} & 部分 & 静态检查的移动后使用类；动态手段命中偶发 \\
双重检查锁定 & \UB{} & 能 & 线程族独占，需全量重编译且与地址族互斥 \\
自写移动赋值双重释放 & \UB{} & 能 & 地址族最直接；泄漏一半可由泄漏档补 \\
\Tp{volatile} 当原子用 & \UB{} & 能 & 线程族；是否想拿它当同步，只能靠评审认定 \\
\end{longtable}

\begin{note}[title=从案例回到手段与契约]
表里能被消毒器稳定抓到的不到一半，其余要靠静态检查、跨平台矩阵与评审。
三条回指正文的结论：各种检测手段的能力边界见 第~\ref{ch:ub03-sanitizers} 章；
并发类案例的判据都落在同一处，见 第~\ref{ch:ub10-memory-model} 章；
容器、字符串视图与移动语义相关的案例，根子在库的前置条件与契约，见 第~\ref{ch:ub12-library-pre} 章。
最小复现的价值不止定位一次缺陷，它同时是给这段代码补上的回归测试。
\end{note}

\begin{bestpractice}[title=案例进代码库的三步]
第一步把改写后的版本作为唯一实现路径保留下来；第二步用表里的首选手段构造一条
会在旧代码上报告、在新代码上静默的测试；第三步把这条测试挂到 附录~\ref{ap:ubB-switches}
的门禁档而不是夜间档，因为只有门禁档会拦住下一次同样的写法。
\end{bestpractice}


\cleardoublepage


\end{document}
